% ===================================================================
%  DUESMA PARTO — feuillets 119 et suivants
%
%  Feuillet 119 : faux-titre de la deuxieme partie. Page d'apparat :
%  chaque ligne est placee a son ORDONNEE MESUREE et occupe une hauteur
%  nulle ; aucune ligne n'en deplace une autre.
%
%  Mesures par outils/apparat.py (feuillet 119) :
%    ligne        y (px)  y (mm)  cap (mm)  largeur (mm)  corps    interl.
%    DUESMA PARTO    726   61,47    2,413      33,78      10,10pt     219
%    (filet)         841   71,20      —         8,21         —         —
%    VORTIFADO       946   80,09    3,387      30,31      14,17pt      86
%    PER DERIVO...  1058   89,58    2,117      58,93       8,86pt     261
%
%  Les interlettrages 219 et 261 ne sont pas ceux qu'apparat.py calcule
%  (201 et 227) : ils les corrigent sur le COMPOSE. La regle d'apparat
%  estime la largeur nue par une composition d'essai, et son estimation
%  se degrade quand la ligne porte beaucoup de blancs de mot — « PER
%  DERIVO O PER KOMPOZO » en compte quatre, et sortait 2,54 mm trop
%  courte. Une fois corrigees, les quatre largeurs de la page tombent a
%  0,08, 0,00, 0,08 et 0,51 mm de la mesure ; le faux-titre du volume,
%  deja accepte, est a 0,26, 0,42 et 1,27 mm.
%
%  GRAISSE : les trois lignes sont en ROMAIN, y compris VORTIFADO, dont
%  la taille donne a premiere vue l'apparence du gras. Deux mesures le
%  refutent. Les rapports fut/capitale de la page — 0,140, 0,175, 0,160 —
%  tombent tous dans la fourchette du faux-titre du feuillet 5, dont les
%  trois lignes sont romaines et donnent 0,152, 0,167 et 0,146 : le fut
%  ne separe rien, comme il est consigne au LISEZ-MOI. Et VORTIFADO,
%  ramene a la hauteur de capitale de « LINGUO INTERNACIONA IDO », lui
%  est superposable — memes empattements, meme O, meme R. C'est le meme
%  romain, plus grand.
%
%  Le filet est le meme filet maigre que celui des notes et celui de la
%  fin de partie : seule la longueur change (8,21 mm ici).
% ===================================================================

% --- Feuillet 119 (folio 115) : faux-titre de la deuxieme partie ----
\begin{VUpage}[119]{}
\VUcentreA{61.47mm}{10.10pt}{219}{DUESMA PARTO}
\VUornement[8.21mm]{71.20mm}
\VUcentreA{80.09mm}{14.17pt}{86}{VORTIFADO}
\VUcentreA{89.58mm}{8.86pt}{261}{PER DERIVO O PER KOMPOZO}
\end{VUpage}

% --- Feuillet 120 (folio 116) : blanc -------------------------------
% Le detecteur de lignes n'en trouve aucune sur ce feuillet : verso
% blanc, comme celui du faux-titre du volume (feuillet 6).
\begin{VUpage}[120]{}
\end{VUpage}

% -------------------------------------------------------------------
%  feuillet 121 — folio 117 : PAS DE FOLIO COURANT au fac-simile.
%  La tete de page est nue (verifiee a l'agrandissement, y 90 a 320) :
%  c'est la premiere page de la deuxieme partie, et l'imprimeur y omet
%  le folio comme il est d'usage. Le champ folio reste donc vide.
%  Deux titres centres, mesures par outils/apparat.py puis corriges sur
%  le compose : ELEMENTI DI VORTO. (cap 2,540 mm, 37,08 mm de large) et
%  PROCEDI DI VORTIFADO. (cap 2,455 mm, 42,84 mm).
%  Le bloc de notes porte TROIS alineas pour deux notes : le troisieme,
%  « La maxim multa gramatikisti docas... », est renfonce comme les deux
%  autres mais ne porte pas de numero. Verifie a l'agrandissement.
% -------------------------------------------------------------------
\begin{VUpage}[121]{}
\VUnotes{117.69mm}{%
(1) Poka vorti (kompare a la ceteri) konsistas nur ek radiko.\nl
Ex. : la prepozicioni, multa adverbi, konjuncioni, e. c.

(2) La qualifiko \textit{ne-verbala} esas plu justa, nam la du klasi povas\nl
egale produktar nomi (L. \textsc{Couturat}).

La maxim multa gramatikisti docas : la nomo o substantivo\nl
esas speco di vorto qua uzesas por nomizar enti o kozi. Ta defino,\nl
qua konsideras kom sinonima la termini \textit{nomo} o \textit{substantivo} ne esas\nl
konforma a la gramatikal tradiciono. L'ancieni nomizis per la\nl
sama termino « nomo » egale la substantivo e l'adjektivo; li dicernis\nl
la \textit{nomo-substantiva}, uzata por nomizar la enti o kozi, e \textit{la nomo}\cc
\textit{adjektiva} uzata por qualifikar li.%
}

\VUtitre{10.63pt}{-55}{ELEMENTI DI VORTO.}
\VUsaut{4.12mm}

Vorto povas konsistar ek diversa elementi, quale \VUgras{bofratino,}\nl
exemple, en qua on dicernas :

1\textsuperscript{e} \textit{radiko} : \VUgras{frat,}

2\textsuperscript{e} \textit{dezinenco} (gramatikal finalo) : \VUgras{o,}

3\textsuperscript{e} \textit{afixi}, inter qui \textit{prefixo} : \VUgras{bo}, sufixo : \VUgras{in.}

Ni nomizas \textit{radiko} la elemento nereduktebla, esenca, qua\nl
indikas la ideo generala quan expresas la vorto, hike : \VUgras{frat.}\parplein

Ni nomizas \textit{temato} l'ensemblo : \textit{radiko} + \textit{afixi}, hike :\nl
\textit{bofratin}.

Ni nomizas * \textit{morfemo} to omna quo soldesas a radiko\nl
por determinar lu : afixi e dezinenci, hike : \VUgras{bo, in, o} (1).

\VUsaut{5.54mm}
\VUtitre{10.27pt}{-26}{PROCEDI DI VORTIFADO.}
\VUsaut{4.19mm}

1\textsuperscript{e} \textit{Derivo per dezinenci} (substituco di ta ad ica) o \textit{neme}\cc
\textit{diata derivo};

2\textsuperscript{e} \textit{Derivo per afixi} (uzo di prefixo o sufixo) o \textit{mediata}\nl
\textit{derivo};

3\textsuperscript{e} \textit{Kompozo per plura radiki unionita.}

\VUgras{Radiki.} --- Li esas \textit{verbala} o \textit{nomala} (to esas neverbala).\nl
L'unesmi expresas ideo di ago, stando o relato. La duesmi\nl
indikas ento, objekto o qualeso, e konseque povas genitar\nl
nur \textit{nomi}, to esas \textit{substantivi} od \textit{adjektivi} (2). On facile\parplein
\end{VUpage}

% -------------------------------------------------------------------
%  feuillet 122 — folio 118.
%  BLOC DE TITRE A TROIS ETAGES, le premier du volume : un titre en
%  capitales demi-grasses (cap 2,455 mm) surmontant DEUX sous-titres en
%  bas de casse italique au corps du texte (cap 1,778 mm, celle des
%  lignes de corps de la page). Les quatre exces signales par le
%  depisteur aux lignes 9 a 12 sont les blancs de ce bloc.
%  Le titre porte un asterisque d'appel : la note du bas de page s'ouvre
%  par le meme asterisque et non par un numero.
%  IRREGULARITES DE L'ORIGINAL, conservees : capitale « E » sans accent
%  dans un titre francais dont le reste est accentue ; et « directeur
%  d'ecole » compose en romain quand les autres segments francais de la
%  meme note sont en italique.
% -------------------------------------------------------------------
\begin{VUpage}[122]{118}
\VUnotes{132.70mm}{%
*\VUgras{Bibliografio.} --- La precipua verki pri la Derivado esas :\nl
L. \textsc{Couturat} : \textit{Etude sur la Dérivation dans la Langue interna}\cc
\textit{tionale}, 1907, ma lore ne vendita, 2\textsuperscript{ma} edituro (1910) e tradukuro :\nl
\textit{Studio pri la Derivado en Linguo internaciona} (Paris che Dela\cc
grave). --- \textit{Dérivation et composition dans la Langue internationale}\nl
(\textit{Ido}) par L. \textsc{Lebasnier}, directeur d'école (Paris, che Chaix, 1912) :\nl
tre klara broshuro por skolani. --- Plura artikli en \textit{Progreso} N\textsuperscript{i} 8\nl
e 10, ed un (pri l'adjektivo substantivigita) en n\textsuperscript{o} 38. De ta artikli,\nl
ni furnisas en noti extraktaji suficanta por evitar konsulto di la\nl
revuo e precipue por la maxim multi qui ne havas lu.%
}

\VUcontinue
dicernos ta duspeca radiki en la vortolibri per ico, ke nur\nl
l'unesmi genitas verbi. Ex. : \VUgras{laborar, vidar, dormar, esperar,}\nl
\VUgras{emocar, anxiar, relatar,} e. c.

\VUgras{Dezinenci.} --- La dezinenci (o flexioni) povas indikar nur\nl
la speco gramatikala di vorto, to esas la rolo di ta vorto\nl
en la frazo. La dezinenci tote ne influas la senco, la ideo\nl
quan la radiko expresas. Chanjo di senco povas obtenesar\nl
nur per afixi o per kompozo.

\VUsaut{7.12mm}
\VUtitre{10.27pt}{-21}{REGULI DI DERIVADO *.}
\VUsaut{4.26mm}

\VUsoustitre{10.2pt}{6}{\textit{Derivado per dezinenci}}
\VUsaut{2.03mm}

\VUsoustitre{10.2pt}{-18}{\textit{(Nemediata derivado.)}}
\VUsaut{3.88mm}

En la derivado per dezinenci konsilesas departar de la\nl
verbo (se lu existas) o de la substantivo por facar l'adjek\cc
tivo o la adverbo. Ma to ne esas absolute necesa, nam on\nl
povas, departante senselekte de irga ek ta vorti, ritrovar la\nl
ceteri.

On notez lo sequanta :

1\textsuperscript{e} Substantivo nemediate formacita de verbo signifikas\nl
nature la ago o la stando expresata da ta verbo. Ex. : \VUgras{laboro}\nl
= \textit{ago laborar}; \VUgras{jaco} = \textit{ago jacar}; \VUgras{langoro} = \textit{stando lan}\cc
\textit{gorar}; \VUgras{emocar} = \textit{stando emocar, sento emocar.}

\VUgras{Konsequo.} --- Inverse on darfas formacar verbo neme\cc
diate de substantivo, nur se ta substantivo expresas ago o\nl
stando. Exemple, de \textit{paco} on darfas formacar la verbo \VUgras{pacar}\nl
= \textit{esar en paco}, pro ke \textit{paco} esas stando; ma de \textit{domo} on\parplein
\end{VUpage}

% -------------------------------------------------------------------
%  feuillet 123 — folio 119. Bloc de notes tres long (31 lignes pour
%  12 de corps) : le filet remonte a 72,05 mm.
%  COQUILLE DE L'ORIGINAL, conservee : le mot « signifikas » est en
%  italique apres « lanjo » et en romain apres « linjo », a deux mots
%  d'intervalle sur la meme ligne.
% -------------------------------------------------------------------
\begin{VUpage}[123]{119}
\VUnotes{72.39mm}{%
(1) Esus kontrea a bona interkompreno, en linguo internaciona,\nl
formacar verbi nemediate de radiki nomala, to esas ne verbala,\nl
lasante al komuna raciono la sorgo trovar la senco. Komparante nur\nl
en nia lingui Europana ita ad ica, on konstatas, ke la relato inter\nl
verbo e substantivo esas tre variiva e mem tre kapricoza : en la\nl
Franca, \textit{plumer} = \textit{desplumizar}, ma, en la Angla, ol prefere = \textit{plum}\cc
\textit{izar}. Dum ke \textit{plumer}, quale ni dicis, signifikas en la Franca \textit{des}\cc
\textit{plumizar}, la verbi \textit{plomber}, \textit{dorer} signifikas \textit{plombizar}, \textit{orizar} e\nl
\textit{brouetter} = \textit{bruetagar}, \textit{bruetportar}. E qua sencon havus \textit{bazardar},\nl
segun la F. familiara \textit{bazarder}? Fine, S\textsuperscript{ro} \textsc{Couturat} remarkigis, ke\nl
\textit{dokumentar} signifikus por Franco \textit{dokumentizar}, \textit{provizar per}\nl
\textit{dokumenti}, e por Cheko \textit{pruvar per dokumenti} (\textit{Progreso}, I, 337).

Quon signifikus \textit{domar}? En la Angla, \textit{to house} = \textit{pozar en domo},\nl
\textit{sub tekto, en hangaro, lojigar o shirmar}; en la Germana, la verbo\nl
\textit{hausieren}, tote analoga, signifikas... \textit{kolportar} (t. e. \textit{irar de domo a}\nl
\textit{domo por ofrar vari}). On vidas quala fonto de miskompreni esus\nl
tala « facileso » formacar verbi nemediate derivita (de nomala\nl
radiko)! E ni citas la du lingui qui maxime praktikas ta tro\nl
« komoda » derivado (\textit{Progreso}, IV, 233).

(2) Ta procedo esas naturala en la mento di la homo. Nia lingui\nl
formacas nemediate substantivi de adjektivi : \textit{die Heiligen} D., \textit{all}\nl
\textit{the saints} E., \textit{tous les saints} F., \textit{omnes sancti} L. e. c. Ta derivo esas\nl
perfekte klara e logikala. Altraparte esas racionala, ke du vorti\nl
diferanta nur per la chanjo dil dezinenco expresas \textit{la sama koncepto}\nl
\textit{en du gramatikal funcioni diferanta}. Fine per l'etimologio ni vidas,\nl
ke omna nomi komuna esis origine adjektivi, e pokope uli divenis\nl
substantivi. Exemple, en la Franca, \textit{domestique} esis nur adjektivo e\nl
signifikis \textit{domala}; nun ol esas pluse substantivo kun la senco\nl
\VUgras{(domala) servisto}; \VUgras{lanjo} \textit{signifikas} : \textit{facita ek} \VUgras{lano}; \VUgras{linjo} signifikas :\nl
\textit{facita ek} \VUgras{lino.} Ankore nun la procedo uzesas : \VUgras{direkteblo} = \textit{direkt}\cc
\textit{ebla aernavo}; \VUgras{automobilo} = \textit{automobila veturo}.%
}

\VUcontinue
ne darfas formacar la verbo \VUgras{domar}, qua havus nula senco,\nl
\textit{domo} expresante nek ago, nek stando (1).

2\textsuperscript{e} Adjektivo nemediate formacita de substantivo signi\cc
fikas : \textit{qua esas} (-). Ex. : \VUgras{responda letro} = \textit{letro qua esas}\nl
\textit{respondo}; \VUgras{ruina kastelo} = \textit{kastelo qua esas ruino}; \VUgras{ora vazo}\nl
= \textit{vazo qua esas oro}. (Same pri omna adjektivo di materio.)\parplein

\VUgras{Praktikal moyeno.} --- Por savar kad on darfas uzar adjek\cc
tivo nemediate formacita de substantivo, suficas vidar, kad\nl
on darfas remplasigar lu per la substantivo en apoziciono :\nl
\VUgras{respondo-letro, ruino-kastelo, oro-vazo.}

\VUgras{Konsequi.} --- On darfas e devas substantivigar adjektivo\nl
per la nura chanjo di dezinenco (2). L'adjektivo tale\parplein
\end{VUpage}

% -------------------------------------------------------------------
%  feuillet 124 — folio 120.
%  COQUILLES DE L'ORIGINAL, conservees : dans la note (4), « regulo cd
%  expliko » pour « od » (verifie a l'agrandissement, la lettre est un c
%  ouvert a droite) ; et deux parentheses ouvrantes jamais refermees,
%  l'une dans la note (2), l'autre a la ligne « (ma on darfas dicar ».
% -------------------------------------------------------------------
\begin{VUpage}[124]{120}
\VUnotes{115.66mm}{%
(1) Certe nulu komprenos : homo dividenda.

(2) La kuntexto o la cirkonstanci impedas miskompreno. Cetere\nl
on sempre darfas tote precizigar per l'expresuro kompleta. (Videz\nl
l'apendico « Substantivigo dil adjektivo ».

(3) Videz la sufixo \textit{-al}.

(4) « Maniero » uzesas hike en sua maxim ampla senco. Ula\nl
adjektivi (\textit{quanta, tanta, multa, poka, kelka, omna, singla, plura,}\nl
la kardinala nombral adjektivi, e. c.) genitas, pro lia senco, adverbi\nl
qui implikas ideo di quanteso. Ma, mem en ta kazo, la supere donita\nl
regulo cd expliko restas vera, nam, relate ago, la quanteso o la\nl
nombro esas manieral cirkonstanco : \textit{li venis plure, single, omne}\nl
indikas reale « quamaniere » li venis. \VUgras{Il laboris tante} = \textit{en tanta}\nl
\textit{maniero}.

(5) \VUgras{Nokte} (de nokto) = \textit{dum la nokto, en nokto}. (Videz la\nl
manieral adverbi en l'unesma parto.)

(6) « Quaze delikto » tre diferas de « kom delikto ».%
}

\VUcontinue
substantivigita nemediate povas signifikar nur : \textit{to quo}\nl
\textit{esas...} (irge quon « to » reprezentas, \textit{ento} o \textit{kozo}) : \VUgras{belo,}\nl
\VUgras{malado} = \textit{homo bela, malada}; \VUgras{dividendo} = \textit{nombro divi}\cc
\textit{denda} (1); \VUgras{plano} = \textit{surfaco plana}; \VUgras{parazito} = \textit{parazita homo,}\nl
\textit{bestio o planto} (2).

Kontree, on ne darfas uzar adjektivo nemediate derivita\nl
vice genetivo (kazo dil komplemento posedala, t. e. vice la\nl
prepoziciono \VUgras{di}) : se \VUgras{gardena urbo} ya signifikas \textit{gardeno}\cc
\textit{urbo}, \VUgras{gardena pordo} ne povas signifikar \textit{pordo di gardeno}\nl
(ma on darfas dicar : \VUgras{gardenpordo} o \VUgras{pordo gardenala} (3).\parplein

3\textsuperscript{e} Adverbo nemediate formacita de adjektivo signifikas :\nl
\textit{en maniero} (tala) (4) : \VUgras{bele, rapide} = \textit{en maniero bela,}\nl
\textit{rapida}; \VUgras{noktale} (de l'adjektivo noktala) = \textit{en maniero nok}\cc
\textit{tala} (\VUgras{noktale vestizita} = \textit{per vesti di nokto}) (5).

Inverse, de adverbo darfas venar direte nur adjektivo\nl
signifikanta : \textit{qua esas talsorta, talmaniera} : \VUgras{frua} = \textit{qua}\nl
\textit{eventas frue} (frua veko, frua legumo); \VUgras{balda} = \textit{qua eventas}\nl
\textit{balde} (balda respondo); \VUgras{quaza delikto} = \textit{ago konsiderata}\nl
\textit{quaze delikto} (6).

Exemplo di vorto-familio : \VUgras{parolar, parolo, parola} (qua\nl
esas parolo) : \VUgras{parola promiso}; \VUgras{parole} (per parolo) : \VUgras{il pro}\cc
\VUgras{misis parole...} Ma tala familii esas rara. Mem importas ne\nl
opinionar, ke irga radiko povas \textit{nemediate} genitar la quar\nl
speci : verbo, substantivo, adjektivo, adverbo per nura chanjo\nl
di la dezinenco.
\end{VUpage}

% -------------------------------------------------------------------
%  feuillet 125 — folio 121.
%  PREMIER SOUS-TITRE DU VOLUME : sous « DERIVADO PER AFIXI. », la ligne
%  « (Mediata derivado.) » est en bas de casse italique, au corps du
%  texte courant (cap 21 px, comme les lignes de corps de la page, contre
%  29 px pour le titre). \VUtitre la composerait en demi-gras capitales :
%  d'ou \VUsoustitre, ajoute au preambule pour ce cas.
%  IRREGULARITE DE L'ORIGINAL, conservee : a la derniere ligne, la glose
%  francaise est composee en ROMAIN, alors que le livre met partout
%  ailleurs en italique ce qui est cite d'une autre langue.
% -------------------------------------------------------------------
\begin{VUpage}[125]{121}
\VUnotes{154.43mm}{%
(1) Tala akumulo de sufixi devas evitesar ordinare. Ol uzesas,\nl
quale hike, nur kande l'ideo tradukenda impozas olu.%
}

Precipue, \textit{verbo darfas derivesar nemediate de adjektivo}\nl
\textit{nur se la radiko di ta adjektivo esas verbala}. Verbo quale\nl
\VUgras{sanar} havus nula senco : \textit{sano} = ento sana, homo sana,\nl
ma ne la \textit{stando esar sana} o \textit{saneso}; on do devas dicar : \VUgras{sane}\cc
\VUgras{sar} = esar sana, e \VUgras{sanigar} = \textit{igar sana.}

Pro plu forta motivo, nula verbo darfas venar nemediate de\nl
adverbo primitiva o de partikulo nevarieba, nam evidente la\nl
radiko ne esas verbala. Verbi quale \VUgras{fruar, baldar, perar, cir}\cc
\VUgras{kumar, transar, proximar} havus nula senco. Inverse adjektivo\nl
nemediate formacita de verbo ne povas havar la senco di parti\cc
cipo aktiva o pasiva (altre la sufixi \VUgras{ant, at}, e. c., nule utilesus\nl
e ne havus senco). Exemple, \VUgras{instrukta} ne povas remplasar\nl
\VUgras{instruktanta}, o \VUgras{instruktiva} (nek mem \VUgras{instruktala}); \VUgras{obedia}\nl
ne povas, konseque ne darfas remplasar \VUgras{obedianta, obediema,}\nl
\VUgras{obediera}, e. c.

\VUsaut{5.48mm}
\VUtitre{10.27pt}{-17}{DERIVADO PER AFIXI.}
\VUsaut{4.60mm}

\VUsoustitre{10.2pt}{-22}{\textit{(Mediata derivado.)}}
\VUsaut{4.45mm}

Kontre ke en la derivado per dezinenci (o nemediata deri\cc
vado), suficas soldar dezinenco a radiko, quale ni vidis,\nl
por obtenar vorto kompleta, en la derivado per afixi (o me\cc
diata derivado), on devas unesme adjuntar e soldar a la\nl
radiko afixo od afixi, ante pozar la dezinenco an la fino\nl
dil toto. Ex. : \VUgras{arki-duk(o), para-pluv(o), pian-ist(o), dis}\cc
\VUgras{til-eri(o); bo-frat-in(o), retro-ag-iv-es(o).}

La sufixi insertesas inter la radiko e la dezinenco segun\nl
l'ordino en qua li aplikesas ica ad ita; to signifikas, ke\nl
singla de oli modifikas ta, qua preiras lu, ed aplikesas a\nl
lu quale a simpla radiko : \VUgras{lum-iz-iv-eso} = eso (\textit{qualeso di})\nl
\VUgras{iv} (\textit{to quo povas}) \VUgras{iz} (\textit{provizar per}) \VUgras{lum(o)}. Praktike on\nl
sequas l'ordino progresanta dil kompozo : \VUgras{lumizar, lumiziva,}\nl
\VUgras{lumiziveso} = \textit{kapableso lumizar}, France : pouvoir éclai\cc
rant (1).
\end{VUpage}

% -------------------------------------------------------------------
%  feuillet 126 — folio 122.
%  Le detecteur de lignes n'a PAS vu le folio courant de cette page :
%  la premiere bande du releve geometrique est le titre « AFIXI. ». Le
%  folio est pourtant bien imprime (verifie a l'agrandissement, base a
%  y = 147). Le champ folio est donc rempli.
%  LIGNE INVISIBLE RETROUVEE : « prefixo. », deuxieme et derniere ligne
%  de la note (2). Le depisteur l'annoncait « double net, aucun blanc » ;
%  lue a l'agrandissement, elle y est. La page porte 10 lignes de note.
%  IRREGULARITE DE L'ORIGINAL, conservee : « anti- » est en gras dans le
%  corps et en italique dans la note (3).
% -------------------------------------------------------------------
\begin{VUpage}[126]{122}
\VUnotes{127.25mm}{%
(1) Grande eroras Esperanto uzante sua afixi kom radiki : 1\textsuperscript{e} to\nl
duopligas neutile la vorto por la sama koncepto (pro quo \textit{ilo} apud\nl
\textit{instrumento}, od inverse?); 2\textsuperscript{e} to konstante ruptas por profani la\nl
kompreno di texto. \textit{Ema, eblo, inde, eta, aro} e. c. esas por ili tam\nl
misterioza kam \textit{edzo, fraulo} e. c. e la famoza korelativi : \textit{chiam, kie,}\nl
\textit{iom, ia, iel, chiu} e. c., qui ja per su ruptas la kompreno sat freque.

(2) En noto : Nam on havas e devas havar \textit{kontre} kom tala\nl
prefixo.

(3) En noto : Apartenos do a la ciencala e teknikala komitati\nl
determinar la kazi, en qui on uzas \textit{anti-}.%
}

\VUtitre{10.27pt}{-30}{AFIXI.}
\VUsaut{4.43mm}

La afixi esas partikuli, elementi tre kurta, quin on uzas\nl
por modifikar la koncepto expresata dal radiko. Exemple,\nl
en \VUgras{bo-frat-in-o} la partikuli \textit{bo}, \textit{in} esas afixi e modifikas la\nl
koncepto \VUgras{frat.}

Se li unionesas a la radiko avane (quale \textit{bo}) on nomizas\nl
li \textit{prefixi} e, se li unionesas a la radiko dope, on nomizas\nl
li \textit{sufixi.}

En Ido, la afixi esas nur to, quo li esas nature e grama\cc
tike : elementi (preske sempre monosilaba), partikuli mo\cc
difikanta, quale ni dicis. Konseque \textit{li generale ne darfas}\nl
\textit{uzesar kom radiki} (1).

Advere pri kelka vorti o radiki, quale l'adverbo \VUgras{ne-}, la\nl
radiko \VUgras{-es}, ula prepozicioni, on kustumas dicar, ke li uzesas\nl
kom afixi. Ma to ne signifikas, ke li reale e gramatike esas\nl
\textit{afixi.} Restas tote justa, ke ti, quin ni studios sube, ne darfas\nl
uzesar kom radiki, exter kazi sorge indikota.

\VUsaut{6.93mm}
\VUtitre{10.63pt}{-41}{PREFIXI.}
\VUsaut{4.49mm}

\VUgras{anti.} --- Ta prefixo esis l'objekto di la sequanta decido\nl
883 (\textit{Progreso}, VI, 113) :

« On repulsas \VUgras{anti-} kom prefixo generale aplikebla (2),\nl
ed adoptas ol kom prefixo rezervenda a la ciencala e tek\cc
nikala terminaro. » (3).
\end{VUpage}

% -------------------------------------------------------------------
%  feuillet 127 — folio 123.
%  COQUILLES DE L'ORIGINAL, conservees : « ciencal » et « matematikal »
%  pour ciencala / matematikala, alors que « teknikala » voisin garde
%  son -a ; et un guillemet ouvrant, a l'alinea « auto. », qui n'est
%  jamais referme.
% -------------------------------------------------------------------
\begin{VUpage}[127]{123}
\VUnotes{96.84mm}{%
(1) On questionas, kad ni devas admisar la prefixo \textit{anti-} kun la\nl
sama senco kam \textit{kontre-}. \textit{On devas distingar tri kazi} :

1\textsuperscript{a} En la teknikala kazi (ex. en la matematikal vortaro) on devas\nl
admisar la prefixo \textit{anti} kun specala senco. To nule koncernas la\nl
komuna linguo.

2\textsuperscript{a} En la linguo komuna, on devas admisar vorti integra, tote\nl
pronta, qui kontenas \textit{anti-}, kande li esas internaciona : pro ta motivo\nl
ni adoptis \textit{antipodo, antipatio, antisepta}, e. c. ed on devos probable\nl
adoptar, \textit{antipatriotismo, antimilitarismo} (same kam \textit{militarismo}),\nl
e. c. Ico havas nula detrimento, nam, o la radiko a qua aplikesas\nl
\textit{anti-} ne existas en nia linguo, o la kompozajo esas tante konocata,\nl
ke nulu povas dubar o nesavar lua senco.

3\textsuperscript{a} En la cetera kazi, on devas uzar kom prefixo la prepoziciono\nl
\textit{kontre}, qua sola devas expresar ta ideo (segun la principo di una\cc
senceso). Tale ni havas : \textit{kontre-pezo, kontre-veneno}, e. c. E singlu\nl
esas libera formacar simila kompozaji, kande li esas necesa e logi\cc
kala : \textit{kontre-konvulsa, kontre-reumatisma, kontre-religia, kontre}\cc
\textit{dinastia, kontre-konstituca}. To rezultas de la strukturo ipsa di la\nl
linguo; e se on uzus simile \textit{anti-}, to genitus nur duopla formi, e\nl
konseque konstanta hezito en l'uzado. (\textit{Progreso}, februaro 1911,\nl
pag. 689.)

(2) En noto : Sama remarki kam pri \textit{anti-}. L'ideo di \textit{auto-} povas\nl
tradukesar per \textit{su-} (ex. : \textit{su-defenso}) o per \textit{propra} (ex. : \textit{propr-amo}).\parplein
}

Ol signifikas \textit{kontre}, qua expresas la sama ideo en la lin\cc
guo komuna : \VUgras{kontreveneno, kontrefebra, kontreverma} (1).\parplein

\VUgras{arki.} --- Ta prefixo, pasable internaciona (Greka origine),\nl
signifikas : \textit{en eminenta grado}; ol do sempre indikas grado\nl
supera : \VUgras{arkianjelo, arkiduko, arkidiakono}, e. c. (\textit{Progreso},\nl
III, 419. --- Decido 649.)

\VUgras{auto.} --- Quale \VUgras{anti-}, ta prefixo repulsesis « kom prefixo\nl
generale aplikebla, ed adoptesis kom prefixo rezervenda a\nl
la ciencala e teknikala terminaro (2).

\VUgras{bi.} --- Ta ciencal o teknikala prefixo « signifikas ofte\nl
« dufoye », sed anke « duopla », « du kune » dicas S\textsuperscript{ro} \textsc{Cou}\cc
\textsc{turat} (\textit{Progreso}, III, 3); ol esas impozata da l'internacioneso;\nl
e pluse ol esas tre oportuna por supresar la dusencesi » : se\nl
on uzus, exemple, \textit{du} (vice \textit{bi-}) en \textit{biplani, bicikli}, on hezitus\nl
omnafoye, kande on parolus pri \textit{du plani, du cikli}. --- Ex. :\nl
\textit{biciklo, bipedo, bipeda, bilingua, bidenta, bikonkava, bikon}\cc
\textit{vexa, bimanua}, e. c.

\VUgras{bo-.} --- Ta prefixo, Franca per \textit{beau-} (pron. : \textit{bo}) esas\nl
heredita de Esperanto senkonteste. Ol indikas : \textit{parenteso per}\nl
\textit{mariajo} : \VUgras{bopatrulo} = \textit{patrulo di la spozino} (relate elua
\end{VUpage}

% -------------------------------------------------------------------
%  feuillet 128 — folio 124.
%  Le depisteur annoncait une derniere ligne manquante : c'etait de la
%  rousseur, verifiee a l'agrandissement.
%  COQUILLE DE L'ORIGINAL, conservee : « rimarijinta » sans le a, alors
%  que le meme mot est imprime « rimariajinta » deux lignes plus haut.
% -------------------------------------------------------------------
\begin{VUpage}[128]{124}
\VUnotes{134.15mm}{%
(1) \VUgras{Mala} e \VUgras{odiar} certe esas la preferinda kontreaji di \textit{bona} e\nl
\textit{amar;} e tale pri la vorti qui, sub formo simpla, esas la kontreaji\nl
di altri : \textit{lenta, mikra, povra} e. c., kontreaji di \textit{rapida, granda,}\nl
\textit{richa.} Ma, pro ke \textit{des} sempre facas ek vorto la kontreajo dil vorto\nl
uzata sen olu, on darfas dicar o skribar : \textit{deslenta, despovra, des}\cc
\textit{mala} e. c., se la memorado ne furnisas : \textit{rapida, richa, bona.}

(2) On imperis a me agar ulo, e ton me ne agas; lore me\nl
\textit{ne-obedias.} On interdiktis a me ulo, e tamen me agas lu; lore me\nl
\textit{desobedias.} Ed anke : on imperis a me agar ulo, e me agas lo kontrea,\nl
lore me \textit{desobedias.}%
}

\VUcontinue
spozo); \VUgras{bofratulo} = \textit{fratulo di la spozino} (relate elua spozo).\nl
Ta prefixo quik adoptesis.

La parenteso per \textit{rimariajo} indikesas, segun la kazi, per\nl
la prefixo mi- (V. ta prefixo), o per l'adjektivo \VUgras{stifa}\nl
(\textit{stief} D., \textit{step} E.).

Mea \textit{stifa} frato esas ta, qua divenis mea frato per rima\cc
riajo di mea patrulo kun lua matro, o di mea matro kun\nl
lua patrulo. Fakte ta \textit{stifa} frato (homulo o homino) havas\nl
altra patrulo ed altra matro kam me, do nula sango-paren\cc
teso kun me. Mea \textit{stifa} patrulo esas la spozo di mea matro\nl
rimariajinta su kun ilu, ma tote ne mea patrulo.

Mea \textit{stifa} matro esas la spozo di mea patrulo rimariajinta\nl
su kun elu, ma tote ne mea matro.

Mea \textit{stifa} onklulo esas la spozo di mea onklino rimari\cc
jinta su kun ilu, ma tote ne mea onklulo (\textit{Progreso}, IV, 105.\nl
--- \textit{Decido}, 745).

\VUgras{des-.} --- Ta prefixo internaciona indikas la kontreajo dil\nl
koncepto expresata sen olu. Exemple : \VUgras{desobedio, deshonoro,}\nl
\VUgras{desavantajo, deshonesta, desfacila, desbona} (1), \VUgras{desinfektar,}\nl
\VUgras{desorganizar, desamar} (1) esas la kontreaji di : \textit{obedio,}\nl
\textit{honoro, avantajo, honesta, facila, bona, infektar, organizar,}\nl
\textit{amar.}

On devas dicernar e distingar sorgoze la vorti \textit{kontrea},\nl
formacata per \VUgras{des-}, de la vorti \textit{kontredicanta} (simpla nego),\nl
formacata per \VUgras{ne-} : \textit{neafabla, neutila, neobedio.}

Por certe facar la dicerno (e pose la distingo per \VUgras{des-}\nl
o \VUgras{ne-}) suficas questionar su kad existas mezo inter la du\nl
extremaji; se ta mezo existas, l'extremaji esas kontrea, lore\nl
oportas uzar \VUgras{des-}; se ta mezo ne existas, lore oportas uzar\nl
\VUgras{ne-.} Ex. : \VUgras{afabla} e \VUgras{desafabla}, pro la mezo : \VUgras{neafabla}; \VUgras{obediar}\nl
e \VUgras{desobediar}, pro la mezo : \VUgras{neobediar} (2); \VUgras{richa} e \VUgras{desricha}
\end{VUpage}

% -------------------------------------------------------------------
%  feuillet 129 — folio 125.
%  Le depisteur annoncait une derniere ligne manquante : rousseur, comme
%  au feuillet 128. Verifie a l'agrandissement.
%  IRREGULARITE DE L'ORIGINAL, conservee : dans la note (2) l'italique
%  englobe le « e » de liaison et court jusqu'au point final, alors que
%  les notes (1) et (3) laissent ce meme « e » en romain.
% -------------------------------------------------------------------
\begin{VUpage}[129]{125}
\VUnotes{118.92mm}{%
(1) Kad esas tre evidenta, e precipue tre limitizebla, l'opozeso\nl
inter \textit{varma} e \textit{kolda} (des-varma), exemple?

(2) Kad esus tolerebla : \textit{yuni e desyuni, richi e desrichi, forti e}\nl
\textit{desforti, humili e deshumili, jenerozi e desjenerozi, boni e desboni,}\nl
\textit{felici e desfelici e. c. same finos per la morto.}

(3) Se on klamos a me : \textit{des-dextre} e me audos nur \textit{dextre}, kad\nl
me ne povos duktar mea biciklo juste an l'obstaklo e frakasar mea\nl
kapo? \textit{Sinistre e dextre} igas ne posibla ta desagreablajo.

(4) Kad, segun la cienco, \textit{frigoro} esas la kontreajo di \textit{kaloro?}\nl
E tamen Esperanto opozas \textit{malvarmeso} a \textit{varmeso}. Ido ne darfis\nl
imitar olu.

(5) Ni konocis ta fakto plura yari ante ke \textit{Progreso} tusheskis la\nl
punto. Ma motivo di « neoportuneso » deskonsilis la quika substituco\nl
di \textit{des} a \textit{mal}. Ol suficabus por impedar interkonsento kun l'Espisti :\nl
la stroko atingabus tro grave la \textit{netushebla} Fundamento.%
}

\VUcontinue
(se on ne uzas \textit{povra}), pro la mezo : nericha. Ta persono\nl
esas nericha; tamen ol ne esas povra (\textit{desricha}), nam ol\nl
posedas mikra domeno ek qua ol vivas.

Nultempe uzez : despovar, desdormar, desvolar, desexistar,\nl
desesar, desmovar, desvivo, desposedo, deshavanta, e. c., pro\nl
ke ta vorti ed altri analoga povas nur equivalar simpla\nl
nego : \textit{ne povar}, e. c. Multa kontreaji esas expresata en Ido\nl
(quale en nia lingui) da radiki specala. Exemple : jorno,\nl
nokto; lumo, tenebro; komencar, finar; audacar, timar;\nl
ganar, perdar; laborar, ociar; bona, mala; dextre, sinistre;\nl
supere, infre; alte, base; plu, min; ante, pos; avan, dop, e. c.\parplein

La motivi esas : 1\textsuperscript{e} ke la opozeso en ula kazi ne esas evi\cc
denta (1); 2\textsuperscript{e} ke serio kelke longa de kontreaji unike per\nl
\VUgras{des-} esus vere netolerebla (2); 3\textsuperscript{e} ke la uzo dil prefixo \VUgras{des-}\nl
povas ulkaze produktar grava detrimenti (3); 4\textsuperscript{e} ke lu donas\nl
vorti tote nejusta ulkaze e ne aceptebla da la cienci (4).\parplein

Nultempe uzez \VUgras{des} kom radiko, ma \VUgras{kontre} (kontree, kon\cc
trea, kontreajo).

\VUgras{Des-} adoptesis per la decido 251-252 pos longa e tote fun\cc
damentala diskutado en \textit{Progreso}, III, 19, 92, 150, 215, 216,\nl
282, 283, 284, 285, 287, 343, 344, 345, 483, 544, 618,\nl
dum qua on komparis ad olu por la sama rolo : \VUgras{mal-, non-,}\nl
\VUgras{an-.} --- \VUgras{Mal-} esis repulsata pro ke ol sustenesas nur da kelka\nl
Franca vorti, kontre ke \VUgras{des-} sustenesas da centi de vorti\nl
internaciona (5) e pro ke \VUgras{mal} esas la tre konocata radiko\nl
\textit{mal(a)}, kontreajo di \textit{bona}. Tala vorti quala esas en Esp.
\end{VUpage}

% -------------------------------------------------------------------
%  feuillet 130 — folio 126.
%  IRREGULARITES DE L'ORIGINAL, conservees : « (segun la \textit{kazo}) »
%  puis « (segun la kazo) », italique une fois et romain la suivante,
%  a une ligne d'intervalle ; et un point la ou la ponctuation attendue
%  est une virgule, dans la reference « (Progreso, VII, 65. e VI, 212.) ».
% -------------------------------------------------------------------
\begin{VUpage}[130]{126}
\VUnotes{121.58mm}{%
(1) La verbo « \textit{distributar} » rezervesas al idei teknikala : distri\cc
butar vaporo, gasi e. c.

(2) Ta prefixo quik adoptesis, ma pose ulu propozis ke on uzez lu\nl
por la kontreeso, alegante la Latina. Respondesis : « En la Latina\nl
quale en la cetera lingui, omna afixi havas plura senci. En la L. I.\nl
ni devas adoptar oli kun un sola e fixa senco. Do, se ni havas \textit{dis-}\nl
kun la senco di separo, ni ne povas uzar ol kun la senco di kontreeso\nl
(remarkez ke, en \textit{discordia} e \textit{dissentio} adminime, ol havas la senco di\nl
separo, eskarto). Cetere, \textit{des-} esas etimologie identa a \textit{dis-}; ni havas\nl
do la yuro atribuar ad ol l'altra senco di \textit{dis-}. » (\textit{Progreso}, III, 424.)\parplein

(3) Decido 1209, konsequo dil decidi 1089 e 1090. (\textit{Progreso}, VII,\nl
65. e VI, 212.) --- Ti qui prizus la duo : \textit{patro, matro} (kun exkluzo\nl
di \textit{patrino}), mustas agnoskar, ke se patro esus nur maskula, on ne\nl
povus nek darfus dicar \textit{gepatri}. E same pri \textit{frato, kuzo, onklo} e. c.\parplein
}

\VUcontinue
\textit{malbona, malfermi, malligi, malsupre, malsupreniri,} e. c., e. c.,\nl
esas etimologie ne justa e pluse groteska.

\VUgras{dis-.} --- Ta prefixo Greka e Latina origine, nun pasable\nl
internaciona, indikas : \textit{separo, disperso, dis-semo}, nam ol\nl
signifikas per su « hike ed ibe ». Ex. : \VUgras{disdonar, disdon(ad)o}\nl
= ago donar hike ed ibe, donar parto ad ica, parto ad ita\nl
(France \textit{distribuer} (1) \VUgras{dissendar} = sendar hoke ed ibe; \VUgras{dis}\cc
\VUgras{pozar} = pozar hike ed ibe, aranjar en diversa loki, manieri,\nl
sur diversa plasi (2).

\VUgras{ex-.} --- Ta prefixo, Latina origine, nun internaciona,\nl
signifikas : \textit{esinta..., qua havis antee la profesiono di..., la}\nl
\textit{ofico di...} Ex. : \VUgras{exoficiro, exprezidero, exnobelo, ex-vicerejo,}\nl
\VUgras{exadvokato.}

Quik adoptita, ica prefixo nultempe esis kontestata.

\VUgras{ge-.} --- Ta prefixo, Germana, uzesas por indikar « la en\cc
semblo di la du sexui (3) : \VUgras{la gepatri} = \textit{la patrulo e la}\nl
\textit{patrino} (matro); \VUgras{la gespozi} = \textit{la spozulo e la spozino}; \VUgras{mea}\nl
\VUgras{geonkli} = \textit{mea onklulo e mea onklino o mea onkluli e mea}\nl
\textit{onklini} (segun la \textit{kazo}); \VUgras{vua geavi} = \textit{vua avulo e vua avino}\nl
o \textit{vua avuli e vua avini} (segun la kazo). \VUgras{Patri, spozi, onkli,}\nl
\VUgras{avi} ne esus tam preciza, pro ke per su li ne indikas plu\nl
multe enti maskula kam enti femina.

Same \VUgras{frati} esas min preciza kam \VUgras{gefrati}, quale montras\nl
ica exemplo : \VUgras{kad vu havas frati, siorino?} --- \VUgras{Yes, me havas}\nl
\VUgras{gefrati : du fratuli ed un fratino.}

Ta prefixo esas komoda; l'exempli supera montras lo. Ma\nl
on devas evitar sorgoze trouzar o misuzar olu (Videz, en\parplein
\end{VUpage}

% -------------------------------------------------------------------
%  feuillet 131 — folio 127.
%  IRREGULARITES DE L'ORIGINAL, conservees : « bo- » est en gras a
%  l'alinea « mi-. » et en romain neuf lignes plus bas ; « stifa » est
%  en italique a une ligne et en romain a la suivante ; et un point la
%  ou la ligne parallele imprime une virgule, devant « e. c. ».
% -------------------------------------------------------------------
\begin{VUpage}[131]{127}
\VUnotes{157.57mm}{%
(1) La Franca prezentas lu kelke en : \textit{més-user}.%
}

\VUcontinue
l'unesma parto, la \textit{chapitro di la substantivi}, precipue la\nl
du alinei finala).

\VUgras{mi-.} --- Ta prefixo internaciona quik adoptesis. Ol signi\cc
fikas : \textit{duime}. Ex. : \VUgras{lasez la pordo miapertita; kun miklo}\cc
\VUgras{zita okuli; il parolis mivoce; laboro mifacita; mimortinta}\nl
\VUgras{pro} (o de) \VUgras{koldeso; pos mihoro.}

Per rezulto di rimariajo povas existar \VUgras{mifrati} o \VUgras{stifa frati}\nl
(Vid. \VUgras{bo-}). L'unesmi esas ti qui havas o la sama patrulo o\nl
la sama patrino. La duesmi esas la filii quin havas de sua\nl
unesma mariajo vua \textit{stifa} patrulo o patrino. Fakte nek li,\nl
nek la stifa patrulo esas parenta kun vu.

Do \textit{mi-patrulo} nek \textit{mipatrino} (o \textit{mimatro}) povas existar,\nl
nek anke \textit{mifilio}, en Ido, pro ke li ne existas en la naturo;\nl
e nun kande \textit{stifa} adoptesis, li ne plus darfas uzesar.

Remarkez, ke France \textit{beau-} (bo-) konstante produktas am\cc
bigueso : \textit{beau-père}, exemple, esas lore \textit{bopatrulo}, lore \textit{stifa}\nl
\textit{patrulo}, e nur posa expliko fixigas la senco.

\VUgras{mis-.} --- Ta prefixo esas Germana : \textit{miss} ed Angla : \textit{mis} (1).\nl
Ol quik adjuntesis dal Komitato ipsa di la \textit{Delegitaro}, dum\nl
lua kunsidi « por expresar : \textit{erore, nejuste} » : \VUgras{misaudar,}\nl
\VUgras{miskomprenar, misuzar, miskantar, mispronuncar, misduk}\cc
\VUgras{tar, misguidar, misfacar,} e. c.

On remarkez, ke esas difero inter \textit{misaudar} e \textit{audar male} o\nl
\textit{malaudar}. Ulu qua esas preske surda, o quan la sono preske ne\nl
atingas, ica \textit{male audas} o \textit{mal-audas}. Ma ulu qua audas \textit{boto},\nl
kande me dicas \textit{poto} (od inverse), ica \textit{misaudas}. Konseque,\nl
\textit{male} (o \textit{mal} uzata kom prefixo) ne esas sinonima kun \VUgras{mis-}.\parplein

Komprenende \VUgras{mis} ne darfas plear la rolo di radiko. Do\nl
ne dicez \textit{mise}, ma \textit{erore, nejuste} o \textit{male}, segun l'ideo.

\VUgras{ne-.} --- Ta vorto ne esas prefixo nature, nam ol esas negal\nl
\textit{adverbo}. Ma ol uzesas quale prefixo, kom elemento neganta\nl
relate la vorto a qua on unionas lu : \VUgras{neposibla, nevera,}\nl
\VUgras{nejusta, neyusta, nevidebla, nekohere, neprecizeso, nevide}\cc
\VUgras{blaji.} e. c.

Ne konfundez \VUgras{ne-} a \VUgras{des-} od a \VUgras{sen-} (prepoziciono). \VUgras{Ne-}\nl
indikas simpla nego, \VUgras{des-} kontreeso, \VUgras{sen-} manko : \VUgras{nericha,}\nl
\VUgras{desricha} = \textit{povra}; \VUgras{nebela, desbela} = \textit{leda}; \VUgras{senbarba, sen}\cc
\end{VUpage}

% -------------------------------------------------------------------
%  feuillet 132 — folio 128.
%  LIGNE INVISIBLE RETROUVEE : « IV, 68.) », deuxieme et derniere ligne
%  de la note (2). Le depisteur l'annoncait « double net, aucun blanc » ;
%  verifiee a l'agrandissement par l'assembleur. Le bloc de notes porte
%  donc 13 lignes la ou la geometrie n'en montrait que 12.
%  COQUILLE DE L'ORIGINAL, conservee : « automooilo » pour automobilo,
%  dans la note (4). Verifiee a l'agrandissement x 10 par l'assembleur :
%  la lettre litigieuse n'a AUCUNE encre au-dessus de la hauteur d'x,
%  quand la hampe du l voisin et le point du i s'impriment noirs et
%  pleins. C'est un o, sorte cassee ou coquille de composition ; on
%  transcrit ce qui est imprime.
%  Autre irregularite conservee : dans la note (1), « plena » reste en
%  romain quand pleno, pleneco, utilo, utileco et utila sont italiques.
% -------------------------------------------------------------------
\begin{VUpage}[132]{128}
\VUnotes{122.77mm}{%
(1) Pro quo ne \textit{senplena?} Esperantisto respondos : pro ke ni ne\nl
dicas \textit{pleno} (vice \textit{pleneco}). Ma pro quo vu dicas \textit{utilo} (vice \textit{utileco})?\nl
Ka plena ne esas adjektivo totsame kam \textit{utila?} Do, se on admisas\nl
\textit{senutila}, on devas admisar \textit{senplena.}

(2) \VUgras{Tralektar} signifikas nur : \textit{lektar rapide, flugante.} (\textit{Progreso},\nl
IV, 68.)

(3) De qua derivas \textit{parveninto}, dicas la decido 375. Ica vorto =\nl
\textit{parvenu} F., \textit{parvenyu} R., \textit{Parvenü, Emporkömmling} D., \textit{parvenu,}\nl
\textit{upstart} E.

(4) Decido 376. Vorto necesa e quan nek \textit{trakurar}, nek \textit{trairar}\nl
povas remplasar. Per extenso on dicas : \textit{X... parkuris en automooilo}\nl
\textit{100 kilometri en un horo. Y... facis en aeroplano parkuro de 35 kil.,}\nl
\textit{ye 100 metri di alteso.}%
}

\VUcontinue
\VUgras{kapa, senarma, senpova.} Konseque, \textit{senutila, senplena} (1),\nl
vice \VUgras{neutila, neplena} esus absurda, e \textit{nearma, nebarba, ne}\cc
\textit{kapa, nepova,} vice \VUgras{senarma, senbarba, senkapa, senpova} esus\nl
egale absurda.

Ne konfundez \VUgras{senmova} = \textit{qua esas sen movo}, a \VUgras{nemovebla}\nl
= \textit{qua ne esas movebla, quan on ne povas movar}; \VUgras{sendube}\nl
dicas, ke ne existas dubo pri..., \VUgras{nedubeble} : ke on ne povas\nl
dubar, ke lo dicata ne esas dubebla.

La vorti kompozita per \VUgras{sen-} opozesas kom simpla negi\nl
al derivaji havanta la sufixo \VUgras{-oz} (qua egalesas \VUgras{kun}) : \VUgras{sen}\cc
\VUgras{anma, sendolora} fakte esas la negi di \VUgras{anmoza, doloroza.}

\VUgras{sen-} = \textit{-los} Germana, \textit{-less} Angla.

\VUgras{par-.} --- Ta prefixo venas de la Latina e novlatina lingui :\nl
\textit{par-, per-.} Ol adoptesis per la decido 309 « \textit{por indikar la}\nl
\textit{perfekta fino o kompletigo} di ula ago » : \VUgras{parlernar} = \textit{lernar}\nl
\textit{komplete}; \VUgras{parlektar} = \textit{lektar komplete, til la fino} (dil\nl
libro) (2); \VUgras{parvenar} = \textit{venar til la fino, atingar la skopo} (3);\nl
\VUgras{parkurar} = \textit{kurar de la komenco til la fino (dil voyo)} (4).

\VUgras{para-.} --- Ta prefixo cience internaciona adoptesis per la\nl
decido 105 « \textit{por indikar shirmilo kontre-} » : \VUgras{parafalo, para}\cc
\VUgras{lumo, parafulmino, parafairo, paracintilo, parapluvo, para}\cc
\VUgras{suno, paravento, paragrelo, parapolvo, paramoskito,} e. c.

\VUgras{pre-.} --- Ta prefixo Latina e nun internaciona quik adop\cc
tesis. Kun formo plu kurta e plu internaciona ol remplasas\nl
\textit{ante} e havas la sama senco por la vorto a qua on unionas\nl
lu. \VUgras{Previdar} = \textit{ante-vidar} (o : vidar antee); \VUgras{predicar} = \textit{ante}\cc
\textit{dicar} (o : dicar antee); \VUgras{prematura} = \textit{matura ante} la tempo,\parplein
\end{VUpage}

% -------------------------------------------------------------------
%  feuillet 133 — folio 129. Premiere page d'un cahier : elle porte la
%  SIGNATURE « 5 » sous le bloc de notes, a gauche. Encre relevee de
%  y 1865 a 1890 px, x 103 a 119.
%  COQUILLES DE L'ORIGINAL, conservees : un guillemet ouvrant jamais
%  referme dans la citation de la decision 1629 ; et « 647 et 648 »,
%  avec le « et » francais au lieu du « e » ido.
% -------------------------------------------------------------------
\begin{VUpage}[133]{129}
\VUsignature{168.15mm}{5}
\VUnotes{154.01mm}{%
(1) Ne impedas \textit{avanirar}. Ex. : \VUgras{Unesme il preiris me, sequante}\nl
\VUgras{Petrus; ma balde il avaniris e Petrus ipsa trovesis dop ilu.}%
}

\VUcontinue
la sezono; \VUgras{prepozar, prehiere, prenomo, prehistorio, pre}\cc
\VUgras{dankar, preirar} (1)\VUgras{, prejudiko, preavo, preonklo.}

\VUgras{pseudo-.} --- Ta prefixo adoptesis unanime dal Akademio\nl
per la decido 1629 « kom prefixo internacione konocata\nl
e ja uzata en Ido, kun la senco : \textit{ne-vera, ne-reala}. Ex. :\nl
\VUgras{pseudo-profeto, pseudorepublikisto, pseudofilozofo, pseudo}\cc
\VUgras{amiko, pseudocienco, pseudomartiro, pseudolinguisto, pseu}\cc
\VUgras{donobelo, pseudoinventero,} e. c.

\VUgras{quadri-.} --- Ta ciencala prefixo signifikas : « \textit{qua havas}\nl
\textit{quar...} ». Exempli : \VUgras{quadripedo, quadrimanuo, quadripeda,}\nl
\VUgras{quadrimanua; quadriangulo, quadriangula.}

\VUgras{retro-.} --- Ta prefixo Latina, nun internaciona, esas quaze\nl
adverbo per sua senco; ol quik adoptesis por indikar \textit{la ago}\nl
\textit{inversa} di ta quan indikas sen olu la vorto fundamentala :\nl
\VUgras{retrocedar} = \textit{cedar ad ulu to quon ol cedabis a ni}; \VUgras{retroirar}\nl
= \textit{irar ad-dope, desavance}; \VUgras{retropulsar} = \textit{irigar desavance per}\nl
\textit{pulso}; \VUgras{retrovenar} = \textit{venar per iro inversa al loko de ube on}\nl
\textit{departabis}. Ex. : \VUgras{me livis la hemo ye duadek kloki e retro}\cc
\VUgras{venis en lu ye duadek e tri quarimi.} --- \VUgras{Retroagiva} = \textit{qua}\nl
\textit{povas agar, efikar pri kozi antea, pri lo pasinta} : \VUgras{la lego ne}\nl
\VUgras{havas efekto retroagiva.}

Pro sua senco intime adverbala, retro darfas uzesar izolite,\nl
kom adverbo : \VUgras{Retroe! ne plus avancez : esas danjero.}

\VUgras{ri-.} --- Ta prefixo, Angla sone, Italiana forme e sone,\nl
indikas nur \textit{ago quan on iteras, iteris od iteros}. Ol devas\nl
ne konfundesar a \textit{retro-}. \VUgras{Rihavar, ridicar, rifacar, riagar,}\nl
\VUgras{riskribar, rilektar} = \textit{havar, dicar, facar, agar, skribar, lektar}\nl
\textit{itere}. \VUgras{Nun me riesas kontenta, pro ke me rihavas mea per}\cc
\VUgras{dita pekunio.}

Ta prefixo unesme adoptita kun la formo \VUgras{re} chanjis pose\nl
ta formo a l'Italiana \VUgras{ri}, por evitar konfundi en multa radiki\nl
internaciona qui komencas per \textit{re}, ma en qui ta silabo ne\nl
havas senco di itero : \textit{rebel, recev, refuz, regard, regret,}\nl
\textit{reklam, rekolt}, e. c. (\textit{Decidi}, 647 et 648). Esperanto preferas\nl
erore alterar \textit{re} a \textit{ri} ed uzar \textit{re-} kun duopla senco.
\end{VUpage}

% -------------------------------------------------------------------
%  feuillet 134 — folio 130. La page OUVRE SUR UN TITRE.
%  COQUILLES DE L'ORIGINAL, conservees : un guillemet ouvrant jamais
%  referme dans le dernier alinea de la note ; et deux appels par
%  asterisque, devant « equi-. » et « ko-. », auxquels ne repond aucune
%  note en pied de page — le bas de page ne porte qu'une note numerotee.
%  Le bloc de notes compte trois alineas dont deux sans numero.
% -------------------------------------------------------------------
\begin{VUpage}[134]{130}
\VUnotes{110.47mm}{%
(1) Same kam pri \textit{anti}, on devas en la komuna linguo admisar\nl
vorti integra, tote pronta, qui kontenas \textit{auto}, kande li esas inter\cc
naciona : \textit{automobilo, autografo, autobiografio, autodidakta} e. c.

Kompreneble restas tote exter ico la uzo di \textit{su} en vorti quale\nl
\textit{su-levo, su-kusho, su-tenado} (on vere tenas su per la ago di la\nl
muskuli), \textit{su-ocido, su-amo, su-admiro} e. c.

« Same kam ni substantivigas sen skrupulo, e mem kun granda\nl
komodeso, la pasiva verbi (\textit{konstrukteso, konvinkeso}, e. c.), ni povas\nl
e devas substantivigar la reflektiva verbi, kun la reflektiva pronomo.\nl
Do on darfas uzar vorti quale \textit{su-levo, su-kusho}, por l'ago \textit{levar,}\nl
\textit{kushar su} (komparez \textit{su-ocido} a \textit{hom-ocido}). Remarko : Esperanto\nl
uzis \textit{kushi} en ne transitiva senco (la senco di \textit{jacar}), tale ke ol povis\nl
expresar \textit{kushar} nur per \textit{kushigi}, e \textit{kushar su} per \textit{kushighi!} Ed on\nl
ne havis la necesa analogeso inter \textit{levar} e \textit{kushar} (la du transitiva).\nl
\textit{Progreso}, V, 284.%
}

\VUtitre{10.27pt}{13}{PREFIXI TEKNIKALA.}
\VUsaut{4.94mm}

Ni vidis \textit{anti-}, \textit{auto-} (1), \textit{bi-}, \textit{quadri-}; restas studienda le\nl
sequanta pure teknikala.

*\VUgras{equi-.} --- Ol signifikas \textit{egala} (plu o min evidente) : \VUgras{equi}\cc
\VUgras{angula, equiaxa, equilaterala,} e. c., sen parolar pri \VUgras{equilibro,}\nl
\VUgras{equinoxo, equipolo, equivalo.} On ne povis uzar \textit{egal-} en ta\nl
omna vorti, qui esas maxim ofte internaciona.

*\VUgras{ko-.} --- Ol renkontresas en \VUgras{kosinuso, kolatitudo, kosekanto,}\nl
\VUgras{kofunciono, kotangento, kovarianto, konormala,} sen parolar\nl
pri \VUgras{kolineara, korelato} e \VUgras{korelativa, korespondo,} qui esas\nl
por ni primitiva vorti. En ta vorti, la prefixo \textit{ko-} ne havas\nl
la senco di kun, sed (ma) la senco : « komplemento di... »\nl
o « di komplemento ».

Ol esas do tote necesa. Ma ni rezervis ol a ta specala senco,\nl
ed uzis \textit{kun} o \textit{sam} en la kazi, en qui la senco postulis ta\nl
elementi : \textit{kun-junto, kun-kurar, kun-mezurebla; sam-axa,}\nl
\textit{sam-centra, sam-foka, sam-tempa, sam-plana...} (\textit{Progreso},\nl
III, 2.)

\VUgras{mono-.} --- Ol signifikas : \textit{un}. Ex. : \VUgras{monosilabo, monoko}\cc
\VUgras{tiledona, monokordo.} La linguo komuna uzas \textit{un-} : \VUgras{unlatera,}\nl
\VUgras{unnoma, unpersona,} ma ol admisas la vorti internaciona qui\nl
komencas per \textit{mono} : \VUgras{monogramo, monotona,} e. c.

Pro analogeso, e pro la sama motivi qui impozis \VUgras{mono-,}\parplein
\end{VUpage}

% -------------------------------------------------------------------
%  feuillet 135 — folio 131.
%  Les deux exces signales par le depisteur encadrent le titre centre :
%  91,4 px au-dessus, 52,6 px au-dessous. Aucune ligne ne manque, les
%  deux zones ont ete agrandies et sont vierges.
%  IRREGULARITES DE L'ORIGINAL, conservees : « mi- » cite en ROMAIN a
%  une ligne, quand « ne- » et tous les autres prefixes cites de la page
%  sont en italique ; et une parenthese ouvrante qui n'est refermee que
%  deux phrases plus loin.
% -------------------------------------------------------------------
\begin{VUpage}[135]{131}
\VUnotes{143.55mm}{%
(1) Plu precize : la teknikala subkomitato.

(2) N. B. : Kom generala \textit{prefixi} aplikebla ad irga radiko, nam\nl
lore on devus selektar e hezitar en singla kazo inter du prefixi. Ma\nl
on ne exkluzas la posibleso adoptar, kom integra radiki, vorti kom\cc
pozita per ta greka prefixi.%
}

\VUcontinue
\VUgras{bi-,} l'Akademio (1) adoptis : \VUgras{tri-, quadri-, quinqua-, sexa-,}\nl
\VUgras{septua-, okto-, nona-} (\textit{Progreso}, IV, 567).

L'Akademio repulsis (2) la serio de prefixi : \VUgras{di-, tetra-,}\nl
\VUgras{penta-, hexa-, hepta-, okto-, ennea-} (\textit{Progreso}, IV, 567).

Kompreneble on sequas la generala reguli di nia linguo\nl
omnube li esas suficanta a la justa expreso dil idei : mi-\nl
vice la tri prefixi : \textit{hemi-}, \textit{semi-} e \textit{demi-}, quo liberigas ni tre\nl
fortunoze de lia jenanta konkuro e de multa heziti (\textit{mi-axo,}\nl
\textit{mi-cirklo, mi-ombro, mi-plano, mi-sekanto, mi-sfero}. Same\nl
\textit{ne-} vice la Latina \textit{in-} o la D. E. \textit{un-} : \textit{nesimila, nevarianta}).\parplein

\VUsaut{7.74mm}
\VUtitre{9.92pt}{49}{REMARKO PRI LA AFIXI.}
\VUsaut{4.45mm}

Ula prefixi o sufixi ed ula prepozicioni neeviteble kolizionas\nl
kun kelka existanta radiki, exemple \textit{par-} kun \textit{pardonar}, \textit{par}\cc
\textit{fumo, partero, parturo}. On videz en nia vortolibri la radiki\nl
komencanta per \textit{de-, des-, ge-, mi-, mis-, per-, ri-, tra-, trans-.}\nl
Nula konfundo eventas praktike, pro ke on konocas ta\nl
radiki ed on komprenas li segun lia senco, e pro ke la sem\cc
blanta kompozajo havus nula senco, o senco evidente falsa\nl
en la kuntexto. Quon signifikus \textit{de-parto, des-tinar, ge-nitar,}\nl
\textit{mi-tologio, mis-tifikar, per-driko, ri-goro, trans-formo}? Mem\nl
kande dusenceso esus teorie posibla, ol ne eventas praktike,\nl
pro ke la kuntexto determinas klare la senco, e fakte on ne\nl
mem pensas ad altra senco, on ne hezitas mem dum un\nl
instanto. Mem pri \textit{pardonar}, qua homo komprenos lu kom\nl
« komplete o fine donar » en : \textit{pardonez me, il pardonis lu,}\nl
o mem : \textit{il pardonis omno a sua filio}? Adjuntez, ke pri\nl
\textit{donaco} o \textit{legaco}, on devus dicar : \textit{il donacis} o \textit{il legacis omno}\nl
\textit{a sua filio.}
\end{VUpage}

% -------------------------------------------------------------------
%  feuillet 136 — folio 132. DEUX titres centres.
%  LE VERDICT « double net, aucun blanc » EST ICI EN DEFAUT, pour la
%  premiere fois du volume. Le depisteur annoncait une ligne manquante
%  a la ligne 13 ; le profil d'encre sur toute la largeur de la page
%  donne ZERO pixel encre des rangees 768 a 819 (deux ou trois
%  poussieres exceptees). La cause est identifiee : le blanc qui
%  PRECEDE le titre « SUFIXI. » vaut presque exactement deux pas, ce que
%  le depisteur lit comme un double net. Voir le LISEZ-MOI.
%  IRREGULARITES DE L'ORIGINAL, conservees : « lernar komplete » compose
%  en romain ici et en italique au folio 128 pour le meme exemple ; le
%  mot allemand « aus-sprechen » laisse en romain quand le livre italise
%  partout ailleurs ce qui est cite d'une autre langue ; et, dans la
%  note, « transitiva » italique puis romain a six lignes d'intervalle.
% -------------------------------------------------------------------
\begin{VUpage}[136]{132}
\VUnotes{132.54mm}{%
(1) « Nia kunfrato ne prizas la regulo, qua igas \textit{transitiva} la\nl
verbi kompozita ek verbo \textit{netransitiva} e prepoziciono, quale \textit{en-irar,}\nl
\textit{ek-irar}; ed il preferas repetar la prepoziciono. To esas ya permisata,\nl
sed (ma) tote superflua, nam en : \textit{me eniras en la chambro}, l'ideo\nl
esas dufoye expresata. Tote logike, \textit{me eniras la chambro} = \textit{me iras}\nl
\textit{en la chambro}. L' « equivalo » esas « perfekta »; to ne esas « sub\cc
tilajo »; la verbo divenis transitiva, \textit{absorbante} la prepoziciono ipsa,\nl
qua « regas » la komplemento. (Komparez \textit{travidar, trairar}, e. c.)\nl
\textit{Progreso}, II, 304, noto.%
}

\VUtitre{10.27pt}{2}{PREPOZICIONI PREFIXA.}
\VUsaut{4.22mm}

Quale \VUgras{sen-}, vidita en antea chapitro, altra prepozicioni\nl
pleas la rolo di prefixi : \VUgras{adportar, cirkumpola, cirkum}\cc
\VUgras{skribo} (geom.), \VUgras{ekirar, enirar} (1), \VUgras{foreso, kontredicar, kun}\cc
\VUgras{voko, preterirar, subjacanta, vice-prezidero,} e. c.

Nultempe atribuez a prepoziciono, kompozante derivajo,\nl
senco altra kam ita, quan ol havas aparte. Exemple, \VUgras{eklernar}\nl
nule povas signifikar \VUgras{parlernar} = lernar komplete; \VUgras{ektrovar}\nl
o \VUgras{ekpensar} nule povas recevar la senco di \VUgras{inventar,} nek\nl
\VUgras{ekparolar} (Germana aus-sprechen) ta di \VUgras{pronuncar}; same\nl
\VUgras{ekportar} = \textit{portar extere} ne povas signifikar o \VUgras{suportar} o\nl
\VUgras{tolerar.}

\VUsaut{3.40mm}
\VUtitre{10.27pt}{-17}{SUFIXI.}
\VUsaut{4.24mm}

Quale ni dicis ye « Derivado per afixi », la sufixi inser\cc
tesas inter la radiko e la dezinenco, segun l'ordino en qua\nl
li aplikesas ica ad ita, se pluri uzesas.

\VUgras{-ab-.} --- Sufixo exkluzive verbala, quik propozita por in\cc
dikar l'anteeso, relate la tempi \textit{-is, -os, -us}, exemple. Ma la\nl
Komitato dil Delegitaro ne adoptis lu pro la motivi expozita\nl
en \textit{Progreso}, I, 567. Tamen ja lore dicesis : « se la sufixo \VUgras{ab}\nl
od irg altra procedo esas vere plu facila e komoda \textit{por la}\nl
\textit{maxim multi}, on devos preferar ica ad ita ». La prevido\nl
realigesis e, pro to, per la decido 407 « on admisis, kom\nl
provizora probado, la uzo di la sufixo \VUgras{-ab} por formacar la\nl
tempi antea : \VUgras{amabis} = \textit{esis aminta}, \VUgras{amabos} = \textit{esos aminta},\nl
\VUgras{amabus} = \textit{esus aminta} » (junio 1911). Noto dicas : « On\nl
devas remarkar, ke l'Akademio ne adoptis la formo\parplein
\end{VUpage}

% -------------------------------------------------------------------
%  feuillet 137 — folio 133.
%  TROIS LIGNES INVISIBLES RETROUVEES, dont une que l'assembleur avait
%  d'abord ecartee a tort :
%    - « tempi ». » (ligne 3) et « febla. » (ligne 29), toutes deux
%      annoncees « double net, aucun blanc » ;
%    - « arivis. » sous la derniere ligne du bloc de notes. L'assembleur
%      l'avait declaree jambage apres un agrandissement dont le cadrage
%      COMMENCAIT A x = 120 alors que la marge de la page est a x = 71 :
%      le mot etait coupe hors du cadre. Le releveur l'a mesure a
%      x 71..160 et l'a lu ; verifie ensuite a l'agrandissement x 8.
%  COQUILLE DE L'ORIGINAL, conservee : « dolcacha » pour dolcacho,
%  seul de sa serie a finir en -a.
% -------------------------------------------------------------------
\begin{VUpage}[137]{133}
\VUnotes{138.77mm}{%
(1) La formo \textit{-abas}, propozita (eventuale) por l'imperfekto, ne esis\nl
adoptata, pro ke l'imperfekto ipsa ne aceptesis, dum la diskuti di\nl
la Konstanta komisitaro (I, 569). On timis, ke ula populi trouzos,\nl
o misuzos ta tempo. Esis plu prudenta e simpla restriktar al formo\nl
\textit{-is} l'expreso dil pasinto (ne antea relate altra) : \textit{me vidis, il arivis,}\nl
e remplasigar ol per kompozita formi, se to esus necesa, ma nur\nl
takaze : \textit{il esas arivinta de du hori; me esis skribanta, kande il}\nl
\textit{arivis.}%
}

\VUcontinue
\textit{amabas} » (1). Fine, per la decido 784 (decembro 1912),\nl
« on adoptas unanime e definitive la sufixo \VUgras{ab} por l'antea\nl
tempi ».

La sufixo \VUgras{ab} inspiresis al propozinto da : \textit{amabam, ama}\cc
\textit{bas}, e. c. Latina, en la penso, ke on povos prizentar lu kom\nl
quaza radiko di D. \textit{haben}, E. \textit{have}, S. \textit{haber}, I. \textit{avere}, F. \textit{avoir}\nl
e docar : uzez \VUgras{ab} (\textit{-abis, -abos, -abus}) ube ta lingui uzas \VUgras{havar}\nl
en tempo kompozita.

Esas tre rekomendata ne uzar \textit{-abis} ube \textit{is} suficas, altra\cc
dice : ube ne esas necesa indikar, ke tempo pasinta preiris\nl
altra pasinto. Exemple, ne dicez : \VUgras{Me vidabis vua patrulo}\nl
\VUgras{hiere,} nam hike nula tempo pasinta preiris altra. Ma kon\cc
tree dicez : \VUgras{me finabis mea laboro kande vu arivis,} nam en\nl
ica exemplo la pasinto \textit{finabis} preiris la pasinto \textit{arivis}.

\VUgras{-ach-.} --- Ta sufixo, inspirita dal Italiana \textit{-accio}, quik\nl
substitucesis ad \VUgras{ac} propozita, kom \textit{pejorativo}. La Komitato\nl
di la Delegitaro judikis ol plu expresiva kam \textit{ac}. Per pejo\cc
rativa sufixo on expresas \textit{qualeso tre infra, nuanco di deses}\cc
\textit{timo, antipatio, repugneso} : \VUgras{kavalacho} = \textit{kavalo di infra,}\nl
\textit{di mala qualeso}; \VUgras{populacho} = \textit{infra populo, grosiera, male}\cc
\textit{dukita}; \VUgras{belacho} = \textit{di qua la beleso insipida, neexpresiva,}\nl
\textit{havas nula charmo}; \VUgras{dolcacha} = \textit{di dolceso sensapora} e, pri\nl
personi, \textit{di dolceso simulata}; \VUgras{bravacho} = \textit{falsa bravo qua}\nl
\textit{simulas braveso}; \VUgras{paperacho} = \textit{papero tote neutila, sen ul}\nl
\textit{valoro, ne konservinda}. \VUgras{Il ne kantas, ma kantachas. Quale}\nl
\VUgras{dormar sur ta litacho? Ico vino! ha no, certe, ma vinacho}\nl
\VUgras{nedrinkebla.}

Ta sufixo devas nultempe uzesar kom diminutivo.\nl
Exemple, \VUgras{brilacho} esas \textit{falsa brilo} e nule \VUgras{brileto} = \textit{brilo}\nl
\textit{febla.}

Same ol devas ne uzesar pri la organi o funcioni di la\nl
bestii, ecepte se on volas expresar repugneso : la bestii (quale\parplein
\end{VUpage}

% -------------------------------------------------------------------
%  feuillet 138 — folio 134.
%  IRREGULARITES DE L'ORIGINAL, conservees : « Fundamento » italique a
%  une ligne et romain six lignes plus bas ; la meme citation « indikas
%  duro en la ago » italique la premiere fois et romaine la seconde ;
%  et deux deux-points successifs dans une meme phrase.
% -------------------------------------------------------------------
\begin{VUpage}[138]{134}
\VUnotes{141.39mm}{%
(1) « L'iterado o frequeso di l'ago esas expresata en Ido per\nl
derivado, nome per la sufixo \textit{-ad} : \textit{pafar, pafadar, parolar, paro}\cc
\textit{ladar.} » (\textit{Progreso}, IV, 390, noto.)

(2) Quale \textit{pensado, volado, memorado, vidado, audado, gustado,}\nl
\textit{flarado, tushado} esas la fakultato \textit{pensar, volar} e. c. Ma por la 5 sensi\nl
la nomo plu justa esas : \textit{vidosenso, audosenso, gustosenso, flarosenso,}\nl
\textit{tushosenso.}%
}

\VUcontinue
ni) havas \textit{boko, pedi}, e. c., e (quale ni) \textit{manjas, drinkas}, e. c.\parplein

\VUgras{No, me ne povas asentar, ke la hipopotamo havas boko :}\nl
\VUgras{ol havas bokacho.}

\VUgras{-ad-.} --- Ta sufixo, quik adoptita, trovesas en la Latina\nl
lingui. \textit{Ol indikas} \VUgras{ago} \textit{durolonga o kontinue plurfoya}. Se per\nl
kanono o fusilo on pafas unfoye, la ago esas \textit{pafo}. Ma se on\nl
iteras plurfoye, se on duras pafar adminime kelke longe, la\nl
ago esas \textit{pafado} (1). Dicante tre laute : \textit{Ho!} me produktas\nl
\textit{klamo}. Ma, se homi tre forte parolas od interdisputas sur\nl
la merkato, sur placo, li produktas klamado.

\textit{Polko, valso}, e. c., esas dansi; ma dansado esas to, quon ni\nl
altre nomizas : dansarto : dansado esas praktikata da omna\nl
populi ed en omna epoki.

\VUgras{Parolo} esas frazo, frazeto pronuncata : il dicis a me :\nl
« \VUgras{vu esas idiota} »; il do insultis me per ta parolo. --- \VUgras{Parolado}\nl
esas la fakultato o la maniero parolar (2) ed anke longa\nl
uzo kontinua dil parolado : \VUgras{Parolado grantesis nur a la}\nl
\VUgras{homi. Ne ilua diskurso hierala fatigis lu, ma lua trihora}\nl
\VUgras{parolado kun l'advokato camatine.}

On evitez misuzar o trouzar ta sufixo. Ol esas uzenda nur\nl
kande ol esas nekareebla por la senco.

L'ex-esperantisti devas partikulare atencar por ne ripro\cc
duktar en Ido \textit{kronado} (vice \textit{kronizo}) ed altri analoga. Mem\nl
segun la \textit{Fundamento}, « ad » \textit{indikas duro en la} ago. Or certe\nl
\textit{krono} ne esas ago, ma objekto, quale \textit{glaso, stulo}, e. c., ed on\nl
ne darfas uzar \textit{kronado} plu multe kam \textit{glasado} o \textit{stulado}.\nl
Cetere Zamenhof ipsa konfirmis la senco di « indikas duro\nl
en la ago » per \textit{paf, danc}, e \textit{ir} qui certe esas verbal radiki.\nl
Do, mem judikante segun la defino e l'exempli donita da\nl
Zamenhof, \textit{kronado} violacas la Fundamento. E, por ni\nl
Idisti, ol violacus nia explicita principo, ke nul radiko povas\nl
donar verbo, se ol ne esas verbala. Or \textit{kronado} supozas \textit{kronar}\nl
e \textit{kron} esas ne verbal, ma \textit{nomal} radiko.
\end{VUpage}

% -------------------------------------------------------------------
%  feuillet 139 — folio 135.
%  DEUX LIGNES INVISIBLES RETROUVEES, toutes deux dans le bloc de notes
%  et toutes deux TRES COURTES, ce qui est la faiblesse connue du
%  detecteur : « ciono e. c. » (y 1573-1588) et « homi? » (y ~1815).
%  Le releveur les a trouvees par le PROFIL D'ENCRE, la seconde en
%  cadrant a 30 px a gauche de la marge ; l'assembleur les a verifiees
%  a l'agrandissement x 5.
%  Le verdict « double net » disait la premiere « sans blanc » : elle y
%  etait bien, mais si courte que la somme d'encre par rangee tombe a
%  30 px la ou une ligne pleine en donne 300.
%  COQUILLE DE L'ORIGINAL, conservee : dans la liste des verbes formes
%  directement sur l'instrument, « martelagar » est imprime la ou la
%  serie appelle « martelar » — et « martelagar » est justement la forme
%  en -ag- a laquelle le paragraphe l'oppose.
% -------------------------------------------------------------------
\begin{VUpage}[139]{135}
\VUnotes{111.93mm}{%
(1) \textit{Klov-izar} = garnisar per klovi (talono di shuo, exemple).

(2) Kad on konsilus (por evitar li) belaji quale : \textit{krucasiono,}\nl
\textit{klovasiono, agrafasiono, butonasiono, martelasiono?} Ma kad ica\nl
sufixo di ago ne devus anke soldesar ad omna radiko di ago, quale\nl
\textit{am, trov, kant, lekt} e. c.? Lore ni havus : \textit{amasiono, trovasiono,}\nl
\textit{kantasiono, lektasiono e mem parolasiono, diskursasiono}, se on ne\nl
preferus : \textit{krucaciono, klovaciono, agrafaciono, martelaciono, ama}\cc
\textit{ciono, kantaciono, diskursaciono, uzaciono, repozaciono, promena}\cc
\textit{ciono} e. c.

Ni rimemorigez, ke la direta verbigo di radiko nomala (quale\nl
\textit{kruc, klov, agraf, buton, martel}) esas ne nur kontrelogika, ma pluse\nl
kontrepraktika, quo esas milfoye plu grava. Ico demonstresis per\nl
\textit{plumar, dokumentar, bazardar, domar, bruetar}, klara por Franci\nl
(ecepte \textit{domar}) ma ne komprenata o \textit{diverse} komprenata da la ceteri.\nl
Ka L. I. devas esar rebusa divinendo por milioni e milioni de\nl
homi?%
}

On anke atencez, ke « ad » unionita ad \textit{-is} nule donas\nl
imperfekto. Nek \textit{is}, nek \textit{os}, nek \textit{us} chanjas ule sua valoro\nl
tempala per uniono kun « ad » : \textit{-adis, -ados, -adus} = \textit{is, os,}\nl
\textit{us}, quale \textit{adas} = \textit{as}, ma kun duro plu longa en la ago.

Kompreneble « ad » nultempe darfas uzesar kom vorto\nl
nedependanta. Do ne imitez \textit{ade} = longe, trovebla che l'Espe\cc
rantisti. Kad \textit{malade} = kurte, o kad \textit{eke} esas la kontreajo di\nl
\textit{ade}? Pro quo ne \textit{adeco} vice \textit{longeco}, e \textit{ekeco} vice \textit{mallongeco}?\parplein

\VUgras{-ag-.} --- Ta radiko dil verbo \textit{agar} uzesas kun vorti di\nl
instrumento por formacar verbi signifikanta : \textit{agar per}...\nl
Ex. : \VUgras{crucagar} = \textit{agar per kruco} (por mortigar ulu); \VUgras{klovagar}\nl
= \textit{agar per klovi} (por fixigar ulo, planko) (1); \VUgras{agrafagar} =\nl
\textit{agar per agrafi}; \VUgras{butonagar} = \textit{agar per butoni}; \VUgras{martelagar}\nl
fero, kupro, e. c. La substantivi : \textit{krucago, klovago, agrafago,}\nl
\textit{butonago, martelago} ne povas konfundesar a l'instrumenti :\nl
\textit{kruco, klovo, agrafo, butono, martelo}. Ma, se on formacus\nl
direte la verbo per l'instrumento : \textit{krucar, klovar, agrafar,}\nl
\textit{butonar, martelagar}, la nomo dil ago esus identa a ta dil\nl
instrumento : \textit{krucar, kruco}; \textit{klovar, klovo}, e. c., e to esus\nl
tam nelogikala kam nepraktikala, pro la konfundi e konfuzi\nl
qui rezultus de tala formaco (2).

\VUgras{-aj-.} --- Ta sufixo, quik adoptita, havas equivalanti en\nl
novlatina lingui. Ol indikas \textit{kozo facita ek ula materio} o\nl
\textit{posedanta ca o ta karaktero} : \VUgras{lanajo} = \textit{stofo ek lano}; \VUgras{silkajo}\parplein
\end{VUpage}

% -------------------------------------------------------------------
%  feuillet 140 — folio 136.
%  IRREGULARITES DE L'ORIGINAL, conservees : « meaopinione » agglutine ;
%  « grasajo » donne sans l'equivalence « = ... » que portent tous les
%  autres membres de sa serie ; et la meme formule composee « indikas la
%  objekto » en italique puis « indikas la subjekto » en romain sauf le
%  dernier mot.
% -------------------------------------------------------------------
\begin{VUpage}[140]{136}
\VUnotes{151.99mm}{%
(1) \textit{Progreso}, III, 482.

(2) Videz la sufixo \VUgras{-ur(o)}.

(3) Videz en l'unesma parto ye verbi mixita p. 52.%
}

\VUcontinue
= \textit{stofo ek silko}; \VUgras{kotonajo} = \textit{stofo ek kotono}; \VUgras{linajo} =\nl
\textit{texuro ek lino}; \VUgras{molajo} = \textit{mola parto} (ex. : mola parto di\nl
frukto); \VUgras{grasajo} (ex. : grasa parto di karno); \VUgras{ferajo} = \textit{ulo}\nl
\textit{ek fero} (ex. : il vendas olda ferajo, fripi); \VUgras{dolcajo} = \textit{ulo}\nl
\textit{dolca}; \VUgras{bitrajo} = \textit{ulo bitra}; \VUgras{bonajo} = \textit{ulo bona}; \VUgras{verajo} =\nl
\textit{ulo vera}; \VUgras{belajo} = \textit{ulo bela.}

Per extenso, e kun \textit{-al}, ta sufixo uzesas por expresar la\nl
ideo : \textit{ago, parolo, procedo} di... Ex. : \VUgras{amikalajo} = \textit{ago,}\nl
\textit{parolo, procedo di amiko}; \VUgras{pueralajo} = \textit{ago, parolo di puero}\nl
(\textit{puerala}).

Aplikata a radiko di verbo transitiva, \VUgras{aj} \textit{indikas la objekto}\nl
\textit{pasiva di la ago} (tote ne la rezultajo); lore ol equivalas\nl
\textit{-ataj} (1) e remplasas participo pasiva : \VUgras{drinkajo} =\nl
\textit{drink(at)ajo, kozo drinkata} (aquo, biro, vino, e. c.); \VUgras{manjajo}\nl
= \textit{manj(at)ajo, kozo manjata} (karno, pano, legumi, e. c.);\nl
\VUgras{vidajo} = \textit{vid(at)ajo, kozo vidata} (peizajo, panoramo, e. c.);\nl
\VUgras{sendajo} = \textit{send(at)ajo, kozo, kozi sendata}; \VUgras{piktajo} = (ne \textit{pik}\cc
\textit{turo} ma) \textit{pikt(at)ajo, to quo esas piktata} (personi, kozi, e. c.);\nl
\VUgras{skultajo} = (ne \textit{skulturo} ma) \textit{skult(at)ajo, to quo esas skul}\cc
\textit{tata} (viro, muliero); \VUgras{konstruktajo} = (ne \textit{konstrukturo} ma)\nl
\textit{konstrukt(at)ajo, to quo esas konstruktata} (la petri e ligni o\nl
fero, e. c., uzata por konstruktar); \VUgras{kopiajo} = (ne \textit{kopiuro} ma)\nl
\textit{kopi(at)ajo} (la texto, modelo, e. c.). \VUgras{Mea kopiuro esas ne tote}\nl
\VUgras{fidinda pro ke mea kopiajo esis apene videbla} (2).

Same, kun la verbi mixita, \VUgras{ajo} = \textit{a(ta)jo} e konseque havas\nl
la senco pasiva : \VUgras{chanjajo} = \textit{kozo chanjata}, ne : \textit{kozo chan}\cc
\textit{janta} (3).

Aplikata a radiko di verbo netransitiva \VUgras{-aj-} indikas la\nl
\textit{subjekto} (necese, pro la naturo dil verbo); ol equivalas\nl
\textit{-(ant)aj} e remplasas participo aktiva, pro ke ne existas pasivo\nl
en verbo netransitiva : \VUgras{rezultajo} = \textit{rezult(ant)ajo, to quo}\nl
\textit{rezultas, kozo rezultanta}; \VUgras{eventajo} = \textit{event(ant)ajo, to quo}\nl
\textit{eventas, kozo eventanta} (od eventinta). \VUgras{De longe me timis la}\nl
\VUgras{nuna eventajo, qua esas meaopinione nur la rezultajo di}\nl
\VUgras{lia intrigi e di vua inerteso.} (Dec. 541.)
\end{VUpage}

% -------------------------------------------------------------------
%  feuillet 141 — folio 137.
%  IRREGULARITES DE L'ORIGINAL, conservees : le titre francais du livre
%  de Couturat, en italique ailleurs dans l'ouvrage, est ici en ROMAIN,
%  capitale E toujours sans accent ; les citations en esperanto des
%  lignes 4 a 8 sont egalement en romain, quand le livre italise partout
%  ailleurs ce qui est cite d'une autre langue.
% -------------------------------------------------------------------
\begin{VUpage}[141]{137}
\VUnotes{127.08mm}{%
(1) Nam ol uzesas tale, mem en stilo tre autoritatoza, adminime\nl
en \textit{Fundamenta Krestomatio}.

(2) Kredeble pro ico l'Esperantisti ne volas lu.

(3) Videz en ta verko da S\textsuperscript{ro} \textsc{Couturat} la noto qua trovesas infre\nl
de la pagino 29.

(4) Espo havas : \textit{kriminala, ma krimino esas krimo!} Pro quo ne\nl
\textit{krimo, krima?} E pro quo \textit{universala} e ne \textit{universa}. Yes, pro quo ne\nl
\textit{universo, universa} quale \textit{brako, braka?} (En Ido : \textit{brakio, brakiala}).

Che ni \textit{brako} = F. chien braque, quan tre ortodoxa vortaro esp.\nl
tradukis per \textit{brakhundo} = hundo kun brakii!

(5) \VUgras{Muliero virgina} (o nur virga) = \textit{muliero} qua esas virgina.\nl
Ma \textit{velo} povas esar nur virginala, di virgino.%
}

Nultempe \VUgras{aj darfas} uzesar kom pejorativo (vice \textit{ach}),\nl
quale che l'Esperantisti (1). Nultempe anke ol darfas uzesar,\nl
quale che li, vice \VUgras{kozo} o \VUgras{afero}. Se \VUgras{ajo} sola = \textit{afero, kozo,}\nl
ici lasta esas neutila, e vice : « nia sankta afero » o : « mi\nl
havas multajn aferojn por fari », suficas dicar : « nia\nl
sankta ajho », « mi havas multajn ajhojn por fari ». Certe\nl
la profani komprenos quik e plu bone kam en Ido : « me\nl
havas multa kozi facenda ».

\VUgras{-al-.} --- Ja Latina (en la formo \textit{-alis}), la maxim interna\cc
ciona de la sufixi (en la formo \textit{al, el}) (2), la afixo \VUgras{-al-} pro\cc
pozesis por kompletigar la listo trovata en « Etude sur la\nl
Dérivation dans la Langue internationale » (3) ed ol quik\nl
adoptesis.

Lua signifiko esas : \textit{qua relatas, koncernas; qua apartenas}\nl
\textit{a..., dependas de..., konvenas por...} Ol juntesas a radiki no\cc
mala o verbala.

Exempli kun radiki nomala : \VUgras{universala, dorsala, nazala,}\nl
\VUgras{faciala, ventrala, brakiala} (4), \VUgras{nacionala, normala, grama}\cc
\VUgras{tikala, nazalo, labialo, guturalo, vertikalo, kolateralo}, e. c.\parplein

Kun verbala radiki : \VUgras{experimentala, finala, komencala,}\nl
\VUgras{administrala, baptala, mariajala, guvernala, ordinala, kon}\cc
\VUgras{fidencala}, e. c.

La sufixo \VUgras{-al-} darfas sustenar altra sufixi : \VUgras{universaleso,}\nl
\VUgras{universaligar; socialismo, socialista}, e. c.

Adjektivo kun \VUgras{-al-} ofte equivalas genitivo, to esas « \VUgras{di} »\nl
sequata da substantivo : \VUgras{hundo blindala} = \textit{hundo di blindo}\nl
(ne \textit{hundo blinda}); \VUgras{velo virginala} = \textit{velo di virgino} (5). De\nl
to la regulo praktikala : Por savar, kad adjektivo devas esar\parplein
\end{VUpage}

% -------------------------------------------------------------------
%  feuillet 142 — folio 138. Page a bloc de notes tres long : 29 lignes
%  de notes pour 14 de corps, filet a 76,88 mm.
%  IRREGULARITES DE L'ORIGINAL, conservees : « -al » en romain dans la
%  note (1) quand la meme note l'italise ailleurs ; « Geheimer Rat »
%  romain a une ligne et italique deux lignes plus haut ; « poemi da
%  Vergilius » romain quand les exemples paralleles sont en gras ; et
%  « oz » sans tiret initial quand la meme note ecrit « -oz » et « -al ».
% -------------------------------------------------------------------
\begin{VUpage}[142]{138}
\VUnotes{77.96mm}{%
(1) On devas tradukar « genitivo » per -al nur kande la sub\cc
stantivo en genitivo esas nedefinita (indikas objekto nedefinita)\nl
e nule kande ol esas definita (indikas objekto determinita, aparta,\nl
konocata) quon montras l'artiklo. Exemple : \textit{rejala kastelo} signi\cc
fikas nur : \textit{kastelo di rejo}, t. e. di ula, irga rejo, o \textit{digna de rejo}.\nl
Ma, se on volas parolar pri \textit{kastelo di} la \textit{rejo}, on ne darfas dicar\nl
rejala kastelo. On dicos exemple : « En ica lando existas multa\nl
rejala kasteli, ma nur tri kasteli di la rejo » t. e. : « la rejo ipsa\nl
havas nur tri kasteli ». (\textit{Progreso}, III, 413.)

(2) La sequanta noto di \textit{Progreso} (IV, 220) pruvas, inter mult\nl
altra kozi, la neceseso dil sufixo \textit{-al-}, kontre qua tante skribis adversi\nl
di Ido : « En ula Esp. broshureto me trovis traduko di D. \textit{Geheimer}\nl
\textit{Rat} per \textit{sekreta konsilisto}. Ca exemplo klare montras per sua absur\cc
deso la logikaleso di nia sufixo \textit{-ala}. \textit{Sekreta konsilisto} povas nur\nl
expresar, ke la konsilisteso di ta persono esas sekreta, t. e., ke on\nl
ne darfas savar, ke lu esas konsilisto. La konstanta publika uzado\nl
di ta titulo pruvas la kontreajo. La nomo Geheimer Rat, qua cetere\nl
nuntempe esas nur metaforo, volas expresar, ke konsili donenda da\nl
ta persono \textit{koncernas sekretaji}, do il esas \textit{sekretala} konsilisto. La\nl
uzo di \textit{sekreta} esas korekta en \textit{sekreta policisto} (D. \textit{Geheimpolizist}),\nl
nam tala policisto ipsa esas sekreta, t. e. ula personi ne darfas savar,\nl
ke il esas policisto. (A.-J. \textsc{Storfer}.)

« On povas questionar, kad on devas dicar \textit{regul-ala} o \textit{regul-oza}.\nl
Kelki tentesas rezonar tale : Nula de la du sufixi konvenas, nam\nl
\textit{-al} signifikas simple : « relativa ad », e \textit{-oz} signifikas « plena de ».\nl
On devus havar specala afixo por signifikar « konforma ad ». ---\nl
Ni respondas : Kontree, la du sufixi esas justa e legitima en ica\nl
kazo; --- \textit{oz}, pro ke ol signifikas « havanta » : reguloza = qua havas\nl
regulo; e \textit{-al}, pro ke ol expresas la maxim generala relato di kozo\parplein
}

\VUcontinue
formacata per la sufixo \VUgras{-al-} videz kad ol equivalas « \VUgras{di} »\nl
sequata da substantivo.

On devas ne trouzar ta formaco, quale agas ula lingui,\nl
exemple la Germana. Precipue kun la propra nomi on uzez\nl
prefere prepoziciono : \VUgras{la verki di Goethe} e mem plu bone :\nl
\VUgras{la verki da Goethe} (ne la Goeth-ala verki) (1).

Ma on uzas \VUgras{-ala-} por dicar : \textit{segun la maniero di, analoga}\nl
\textit{a la verki da, digna de} : \VUgras{poemi Vergiliala, Anakreonala,}\nl
\VUgras{dramo Shakespearala,} qui tote ne esas : poemi da Vergi\cc
lius, e. c. (2).

Kande l'adjektivo formacenda signifikas ne : \textit{qua relatas,}\nl
\textit{qua koncernas, qua apartenas a...} ma « qua esas... », lore\nl
nur la dezinenco \VUgras{-a} sola esas uzenda : \VUgras{letro responda} =\nl
letro \textit{qua esas} respondo; \VUgras{parolo konsolaca} = parolo \textit{qua esas}\parplein
\end{VUpage}

% -------------------------------------------------------------------
%  feuillet 143 — folio 139.
%  LE FILET DES NOTES ECHAPPAIT A L'OUTIL : il est si pale qu'il ne
%  devient continu qu'a un seuil de 210, la ou le detecteur seuille a
%  185. Trouve par le profil d'encre a y = 837 px, cale a la marge
%  (x = 69) et long de 210 px = 17,8 mm — les 212 px attendus. Ordonnee
%  posee : 76,12 mm. Le depistage n'ayant pas pu se faire, le releveur a
%  verifie pas a pas les lignes de base : corps 14 lignes au pas de
%  41,4 px, notes 30 lignes au pas de 32,7 px, aucun trou.
%  COQUILLES DE L'ORIGINAL, conservees : « mam » pour kam ou nam ;
%  « ia vorti » pour la vorti ; un guillemet fermant sans ouvrant ; et
%  un point-virgule la ou un deux-points serait attendu.
% -------------------------------------------------------------------
\begin{VUpage}[143]{139}
\VUnotes{76.75mm}{%
\VUcontinue ad altra; ni tradukas ol per « relativa ad », ma ni povus ofte tradu\cc
kar ol per « dependanta de ». On povas do tre bone uzar \textit{regulala} kun\nl
la senco dependanta de regulo », do « konforma a regulo ». Kun ana\cc
loga senco ni uzas; \textit{legala, normala, kustumala, konvencionala}, e. c.\nl
On povas tre klare uzar ia vorti : \textit{legaleso, legaligar}, e. c. impli\cc
kanta la kompreno di \textit{legala} kom « konforma a la lego ».

Cetere ti qui deziras (pro irga motivo) absoluta precizeso povas\nl
omnakaze uzar, sive \textit{lego-konforma, regul-konforma}, sive (quale ni\nl
ja indikis : VI, 28) : \textit{segun-lega, segun-regula;} same kam ni uzas :\nl
\textit{kontre-lega, kontre-regula} (\textit{Progreso}, VI, 595).

« On devas ne restriktar tro multe la senco di nia sufixi, e ne\nl
interpretar li strikte segun la defini o traduki, quin ni mustas\nl
donar pri li en la lingui nacionala, e qui necese esas kelke tro preciza,\nl
pro ke li uzas vorti preciza di nia lingui. Singla de ta sufixi repre\cc
zentas generala ideo, to e. relato, quan ta traduki sugestas ma ne\nl
limitizas. On devas konsiderar, ke singla nociono kontenas multa\nl
gradi e nuanci; inter du o plura nocioni existas kontinua gradaro,\nl
tale ke ula nuanci povas preske indiferente tradukesar per un od\nl
altra sufixo, inter qui ol jacas. On devas do admisar larja uzo di\nl
la sufixi, tale ke lia domeni esez kontigua e lasez nul kazo extere;\nl
e mem se lia domeni kelkafoye mixas su an la limiti, ico nule esas\nl
defekto di la linguo o kulpo di lua adepti, ma natural e necesa\nl
konsequanto di la dicita kontinueso (\textit{Progreso}, VI, 594).

L'exemplo \textit{regulala, reguloza}, ed altri simila posibla, esas inter la\nl
sufixi la analogajo di \textit{pro, de} en ula kazi, o \textit{pri, koncerne} en altri,\nl
inter la prepozicioni. Ex. : \textit{mortar pro o de hungro; pri co, o kon}\cc
\textit{cerne co} (e mem \textit{relate co}). La vicineso dil nuanci esas tante granda,\nl
ke on povas e darfas ulfoye selektar cis o trans la limito apene\nl
dicernebla. Lo importanta esas lore ke l'ideo restez klara e juste\nl
komprenebla.%
}

\VUcontinue
konsolaco; \VUgras{ento dea} = ento \textit{qua esas} deo; \VUgras{ruina kastelo} =\nl
kastelo \textit{qua esas} ruino.

Kande la vorto esas kompozita, ordinare on ne uzas \textit{-ala},\nl
ma nur \textit{a-} Ex. : \VUgras{omnadia, singladia, externaciona, super}\cc
\VUgras{naciona, internaciona, cadia, cayara, segunlega, kontrelega,}\nl
\VUgras{senaqua, senkapa, senviva,} e. c. Ma povas divenar necesa\nl
uzar \textit{-ala-} mem kun vorto kompozita, kande la klareso pos\cc
tulas lo. Exemple on dicos : \VUgras{la muskuli avanbrakiala} (o :\nl
dil avanbrakio) mam \textit{muskuli avanbrakia} signifikus : la\nl
muskuli \textit{qui esas} avanbrakio.

\VUgras{-an-.} --- Ta sufixo, \textit{-anus} en la Latina, \textit{an, ano, ain, ien} en\nl
altra lingui, quik adoptesis por indikar individuo apartenanta\nl
a klaso (urbo, lando, ensemblo) homo esanta « \textit{membro}\nl
\textit{di...} » : \VUgras{urbano, landano, senatano, akademiano, societano,}\parplein
\end{VUpage}

% -------------------------------------------------------------------
%  feuillet 144 — folio 140.
%  IRREGULARITE DE L'ORIGINAL, conservee : « Zamenhof » est compose en
%  romain ordinaire, quand les autres noms d'auteurs du livre sont en
%  petites capitales ; et « ossuaire », « ossuarium », « ossario »,
%  « ossuary » restent en romain quand « formulary » et « vocabulary »,
%  cites au meme titre deux lignes plus haut, sont en italique.
% -------------------------------------------------------------------
\begin{VUpage}[144]{140}
\VUnotes{91.86mm}{%
(1) Ne \textit{isto}, pro ke la Kristani, segun sua doktrino, dicas esar la\nl
membri di sama korpo, havanta kom kapo la Kristo. Li do esas\nl
Krist-ani.

(2) La origino di ta vorto, qua havis tanta fortuno inter ni esas\nl
pasable stranja. Me esis vizitanta S\textsuperscript{ro} \textsc{Gaston Moch} e me trovabis\nl
lu komencanta letro ad Esperantisto stranjera. Il ne volis uzar\nl
« kara samlinguano ». Lore spontane venis en mea spirito la vorto\nl
\textit{samideano} e me propozis lu al \textit{samideano}, qua aprobis ed uzis ta\nl
vorto. Me ipsa pose uzis lu skribante ad Esperantisti e rapidege ni\nl
imitesis. Se mea memorado esas justa to eventis en 1890 o 1891.

(3) Ne sen motivo me selektas ta du exempli. Questionita da me\nl
Zamenhof respondis : « \textit{-ar-} inspiresis a me da la Franca \textit{formulaire},\nl
\textit{vocabulaire} ». Ta du vorti esas en l'Italiana : \textit{formulario}, \textit{vocabu}\cc
\textit{lario}; en la Angla : \textit{formulary}, \textit{vocabulary}. La vorto « ossuaire »\nl
(L. ossuarium, I. ossario, E. ossuary) kelke sustenas \textit{-ar-}, nam ol\nl
indikas preske tam vere \textit{ostaro} kam \textit{osteyo}.

Kontre ta sufixo en Ido uli alegis posibla dusenceso. Yen to quon\nl
\textit{Progreso} respondis pri ca punto : « La dusenceso alegita (pro \textit{-ar}\nl
dil infinitivo) nultempe prizentis su fakte, ye nia konoco. Kad ol\nl
esas timenda? Ne semblas : 1\textsuperscript{e} pro ke la substantivi en \textit{-aro} esas\nl
preske nultempe uzata kom adjektivi (quo esus : \textit{homara}, \textit{navara},\nl
\textit{armara?} on havas okaziono uzar nur : \textit{homarala}, \textit{navarala}, e. c.);\nl
2\textsuperscript{e} pro ke mem en ta kazo, on ne esas \textit{obligata} elizionar, e tala\nl
eliziono certe esus evitenda; 3\textsuperscript{e} pro ke existas nula verbo,\nl
kun qua ta formi elizionita povus koincidar : ni povas havar nula%
}

\VUcontinue
\VUgras{Kristano} (1), \VUgras{samreligiano}, \VUgras{samideano} (2), \VUgras{samskopano},\nl
\VUgras{samluktano}, e. c.; \VUgras{Kanadano}, \VUgras{Italiano}, \VUgras{Portugalano}, \VUgras{Brazi}\cc
\VUgras{liano}, \VUgras{Boliviano}, e. c. (Videz ye la propra nomi en la Gra\cc
matiko.)

En la realeso (e segun \textit{Progreso}, III, 93) \VUgras{-an} definesas\nl
plu juste per : \textit{qua apartenas ad ula domeno} (lando, socio,\nl
e. c.). On do darfas ne restriktar olu a homi, ma uzar lu\nl
anke por kozi. Co explikas, ke ni parolas pri \textit{Amerikana}\nl
planti od animali, pri \textit{Chiniana} od \textit{Italiana} linguo, e. c. Nam\nl
la sufixo \VUgras{-an-} ne indikas esence personi, quankam pro lua\nl
senco, ol aplikesas precipue a personi.

Advere on darfus dicar \textit{Amerikala}, \textit{Chiniala}, e. c., ma lore\nl
ta vorti ne plus esus analoga a \textit{Franca}, \textit{Angla}, e. c., quin on\nl
uzas egale por personi e por kozi.

\VUgras{-ar-}. --- Ta sufixo, heredita de Esperanto, soldesas a nomal\nl
radiko por indikar \textit{kolektajo}, \textit{ensemblo} di la kozi o personi\nl
nomata dal radiko. Ex. : \VUgras{formularo} = kolektajo, ensemblo\nl
de formuli; \VUgras{vortaro} = kolektajo, ensemblo de vorti (3).

On devas gardar su kontre l'ambiguaji qui povas naskar
\end{VUpage}

% -------------------------------------------------------------------
%  feuillet 145 — folio 141.
%  DEUX LIGNES INVISIBLES RETROUVEES, toutes deux TRES COURTES : une
%  dans le corps (y ~742) et une dans les notes (y ~1524). Le depisteur
%  les annoncait « double net, aucun blanc » ; le releveur a montre par
%  le profil d'encre qu'elles portent 536 et 625 px d'encre, soit 8 %
%  d'une ligne pleine — assez pour passer pour du blanc a la somme, pas
%  assez pour l'etre. L'assembleur les a verifiees a l'agrandissement.
%  COQUILLES DE L'ORIGINAL, conservees : un blanc de deux cadratins
%  entre « kazi » et son point ; « L'Espisti » pour l'Esperantisti ;
%  et « -ul » italique quand « -ier », cinq mots plus loin, est romain.
% -------------------------------------------------------------------
\begin{VUpage}[145]{141}
\VUnotes{117.43mm}{%
\VUcontinue verbi simila a \textit{armar, homar, navar}. Generale la sufixo \textit{-ar} aplikesas\nl
nur a nomal radiki, or ni ne derivas verbi \textit{nemediate} de tala radiki :\nl
do la koliziono semblas ne posibla. --- To ne signifikas, ke ta sufixo\nl
\textit{-ar} ne esas kelke arbitriala; ma oportus trovar altra sufixo « plu\nl
internaciona » e min arbitriala; e to ne semblas facila. » (\textit{Progreso},\nl
II, 651.)

(1) Lore \textit{ul} ankore signifikis « karakterizata per », quale en\nl
Esperanto. Ma pose e plu juste, \textit{-ul} atribuesis a la enti \textit{maskula}, -ier\nl
recevinte la senco : \textit{karakterizata per} (ulo extera).

(2) Defino donita kun la decido (\textit{Progreso}, VI, 322).

(3) L'Espisti dicas sen skrupulo : \textit{adresato}, segun la Germana\nl
expreso \textit{adresat}. Ico esas duopla eroro : 1\textsuperscript{e} pro la nomal naturo di\nl
\textit{adreso}; ta radiko ne povas donar \textit{adresar} (ma \textit{adresizar}) ed on\nl
devas dicar : \textit{sendar, direktar, expediar} ulo ad ulu; 2\textsuperscript{e} pro ke,\nl
mem se \textit{adresar} existus, \textit{adresato} povus signifikar nur la kozo sen\cc
data, la \textit{sendajo}.%
}

\VUcontinue
de la neprecizeso di ta sufixo. Generale ol indikas \textit{la maxim}\nl
\textit{extensita} kolektajo : \VUgras{homaro} = \textit{la ensemblo di omna homi,}\nl
e ne irga grupo o societo de homi; \VUgras{vortaro} = \textit{la ensemblo}\nl
\textit{dil vorti di linguo}, e ne ica od ita grupo de vorti, vortolisto\nl
o mem lexiko; \VUgras{navaro} esas \textit{l'ensemblo di omna navi} (militala\nl
o komercala) quin posedas naciono, e ne eskadro, floto, e. c.;\nl
\VUgras{militistaro} = \textit{l'ensemblo di omna militisti di lando}, e ne\nl
armeo, regimento, kompanio, e. c. Ex. : \VUgras{L'unesma armeo}\nl
\VUgras{di nia militistaro nun komandesas dal ex-komandanto dil}\nl
\VUgras{duesma. --- Nia navaro havas tri floti. --- La vortaro Franca,}\nl
\VUgras{Germana od Angla trovesas komplete en preske nula vorto}\cc
\VUgras{libro o lexiko.} On do atencez en singla kazo uzar la justa\nl
vorto.

En altra kazi, la senco restas nepreciza; exemple la senco\nl
di \textit{arboraro} povas variar de arborgrupo til foresto. Esas do\nl
preferinda, segun quante on povas, uzar la specal vorti :\nl
\VUgras{foresto, bosko, bosketo.}

Kompreneble on ne uzas \VUgras{aro} aparte, kom radiko, ma :\nl
\textit{ensemblo, kolektajo, amaso, bando, trupo, serio}, e. c., segun\nl
la kazi .

\VUgras{-ari-.} --- Ta sufixo, Latina origine e novlatina, adoptesis\nl
dal Akademio, en junio 1913, per la decido 1143. Fakte la\nl
decido remplasigis la unionuro di \textit{du} sufixi (\textit{aj, ul}) (1), kom\cc
prenebla pos reflekto ma ne evidenta. Ol indikas « \textit{persono}\nl
\textit{a qua ula ago esas direktata od ula kozo esas destinata} (2) :\nl
\VUgras{depozario, legacario, sendario} (3), \VUgras{konfidencario, donacario,}\parplein
\end{VUpage}

% -------------------------------------------------------------------
%  feuillet 146 — folio 142.
%  IRREGULARITES DE L'ORIGINAL, conservees : « -atr » sans tiret final
%  quand le meme suffixe est ecrit « -atr- » trois fois sur la page ;
%  et, dans un meme paragraphe d'exemples, sponjatra / dolcatra /
%  dolcacha en gras quand rezinatra / silexatra / lanatra sont romains.
% -------------------------------------------------------------------
\begin{VUpage}[146]{142}
\VUnotes{115.55mm}{%
(1) En la vortolibro Franc-Idala, ye « confesseur », trovesas\nl
\textit{konfesigero, konfesiganto}, qui expresas l'ideo min simple. Yen\nl
l'expliko dil fakto. Kande me korektigis, dum la milito, ye l'explozo\nl
dil obusi, la duesma duimo di la lexiko Franc-Idala, l'unesma ja\nl
esis imprimita ed omna kayeri tote pronta. Do la vorto \textit{confesseur},\nl
tradukita en epoko kande ni ne havis \VUgras{-ario-}, trovesis en olu. Regre\cc
tinde on obliviis pozar \textit{konfesario} inter la vorti adjuntenda o korek\cc
tigenda. Ma nur olu esas vere bona ed ol tote korespondas al vorto\nl
Latina \textit{konfesarius}.

Ni remarkigez, ke \textit{confessor} Latina expresas altra ideo (pluse), ta\nl
di homo konfesanta (sua kredo Kristana) mem avan tormenti, e\nl
nule l'ideo dil sacerdoto recevanta la konfeso di peki. Ca detalo\nl
montras quale e quante on povas erorar prenante blinde Latina\nl
radiko (quale agas Latinisti en sua sistemi) sen atencar lua senci\nl
multopla, qui singla postulas \textit{logikale} e \textit{praktikale} sua traduko par\cc
tikulara en la helpolinguo.%
}

\VUcontinue
\VUgras{koncesionario, adjudikario, konfesario} (1), e. c. = \textit{qua re}\cc
\textit{cevas depozo, legaco, sendo, konfidenco, donaco, koncesiono,}\nl
\textit{adjudiko, konfeso}, la homo atingata da ica agi, e fakte dal\nl
\textit{depozajo, legacajo, sendajo}, e. c.

\VUgras{-atr-}. --- Ta sufixo, quik adoptita, trovesas en la Franca\nl
(\textit{-atre}) ed en l'Italiana (\textit{-astro} : \textit{verdastro, biancastro}). Ol jun\cc
tesas a nomala radiki por formacar adjektivi signifikanta :\nl
\textit{di la naturo di..., afina ad..., kelke simila a...} Ex. : \VUgras{sponjatra}\nl
fungo, ligno. Materio rezinoza, o plu juste rezinatra. Ta\nl
stono ne esas silexa ma nur silexatra. On vendis a vu ta\nl
texuro kom lana, on trompis vu : ol esas lanatra. Saporo\nl
\VUgras{dolcatra} (kelke, ma ne precize dolca); \VUgras{homo dolcacha} (homo\nl
ne dolca vere, ma qua simulas dolceso, por disimular sua\nl
nedolceso). Quale on vidas per la komparo di ta du exempli,\nl
\VUgras{-atr} tote ne havas la senco pejorativa; ica apartenas nur a\nl
la sufixo \VUgras{-ach}.

Kun radiko indikanta koloro, \VUgras{-atr-} expresas koloro afina\nl
ad olta quan enuncas la radiko : \VUgras{redatra, blankatra,}\nl
\VUgras{bluatra, verdatra}, di koloro afina a redeso, blueso, verdeso.\nl
Per la diminutivo \textit{et} on nur diminutas la koloro ed expresas\nl
altra nuanco dil ideo : on afirmas certeso pri la koloro\nl
(quon on ne agas per \VUgras{-atr-}) ma on afirmas anke olua\nl
extrema febleso.

En la vorti teknikala, ta sufixo darfas esar remplasata,\nl
segun la kazi, da la radiki \VUgras{-form}, \VUgras{-simil}, o dal sufixo tek\cc
\end{VUpage}

% -------------------------------------------------------------------
%  feuillet 147 — folio 143.
%  LIGNE INVISIBLE RETROUVEE que le depisteur n'avait PAS signalee :
%  « IV. 521). », cinquieme ligne de la note (1). Elle explique l'exces
%  de 28 px porte au compte de la ligne 36 : ce n'etait pas un blanc
%  d'alinea mais une ligne courte non detectee.
%  IRREGULARITE DE L'ORIGINAL, conservee : « (F. comestibles) » reste en
%  romain quand toutes les autres citations etrangeres de la page sont
%  en italique ; et « IV. 521 » porte un point la ou les autres renvois
%  au Progreso portent une virgule.
% -------------------------------------------------------------------
\begin{VUpage}[147]{143}
\VUnotes{139.87mm}{%
(1) « Exempli en qui \textit{-ea} sentesas kom granda alejo, esas : \textit{rozea},\nl
\textit{violea, lilacea, oranjea, orea, kremea} », dicis samideano qua adjuntis :\nl
« Hike la lingui, ofte mem D. ne adjuntas la elemento \textit{-kolor}, ex.\nl
Dana \textit{rosa, violet, lila, orange, gylden, kreme} (B. \textsc{Jonsson}, \textit{Progreso},\nl
IV. 521).

(2) Nule \textit{forsan}, qua expresas tote altra ideo, quale montras :\nl
\textit{to esas forsan posibla.}%
}

\VUcontinue
nikala \VUgras{-oid} : metaloido, elipsoido, konkoido, konoido,\nl
helicoido, e. c.

\VUgras{-e-}. --- Ta sufixo korespondas a \textit{at} (en \textit{rosata} L.), a \textit{at} (en\nl
\textit{rosato} I.), a \textit{ad} (en \textit{atigrado} S.), a \textit{é} (en \textit{rosé, tigré} F.). Ol\nl
adoptesis per la decido 758 kun ica signifiko : \textit{qua havas la}\nl
\textit{koloro od aspekto di...} Ex. : \VUgras{rozea} = \textit{rozkolora}; \VUgras{violea} =\nl
\textit{violkolora}; \VUgras{kastanea} = \textit{kastankolora}; \VUgras{oranjea} = \textit{oranjokolora},\nl
cindrea, orea, arjentea, kuprea, plombea; tigrea, zebrea felo;\nl
cindreo, sangeo. Per ta sufixo ni sparas centi de tedanta\nl
vorti kompozita per \textit{-kolor(a)} (1).

\VUgras{-ebl-}. --- Ta sufixo, internaciona en la formi \VUgras{abile, able,}\nl
\VUgras{ibile, ible} venas de la Latina \VUgras{abilis, ibilis}. Ol signifikas :\nl
\textit{qua povas esar ...ata} : \VUgras{facebla, kredebla, lektebla, videbla}\nl
= \textit{qua povas esar facata, kredata, lektata, vidata}. Do lu havas\nl
senco nature pasiva, quo explikas ke lua komplemento sem\cc
pre postulas \textit{da} : \VUgras{videbla da omni, kredebla da nulu.}

La linguo Angla inspiris (per la fonetismo) la formo\nl
-\textit{e}bl, meza inter -\textit{a}bl e -\textit{i}bl.

L'adjektivi quin ta sufixo formacas darfas esar substan\cc
tivigata nemediate : \VUgras{trompeblo} = \textit{homo trompebla}; \VUgras{manjebli}\nl
= \textit{nutrivi} (F. comestibles) \textit{kozi manjebla}; \VUgras{kombusteblo} =\nl
\textit{materio} (ligno, karbono) \textit{qui povas esar kombustata.}

Kun \VUgras{-eso} ta sufixo expresas la qualeso abstraktita : \textit{vide}\cc
\textit{bleso, kredebleso}, e kun \VUgras{ajo} lu donas substantivi konkreta :\nl
\VUgras{videblaji, kredeblajo.}

Ne konfundez ta sufixo, indikanta nur posibleso, a la\nl
sufixo \VUgras{-end-} indikanta \textit{obligo} o \textit{musto}, ed a \VUgras{-ind-} indikanta\nl
\textit{merito}. Ex. : \VUgras{omna kozi dicebla ne esas dicenda, nek dicinda.}\parplein

La sufixo \VUgras{-ebl-} ne darfas uzesar kom radiko o plu juste\nl
vorto izolita. Uzez do nur \VUgras{posibla} (qua povas eventar o\nl
facesar) e \VUgras{posible} (en maniero posibla) (2).
\end{VUpage}

% -------------------------------------------------------------------
%  feuillet 148 — folio 144.
%  COQUILLE DE L'ORIGINAL, conservee : « pladego » dans la liste des
%  exemples du suffixe -ed-, la ou la serie appelle « pladedo ».
%  IRREGULARITE conservee : « boko » est en gras quand « pedi », qui le
%  suit immediatement dans la meme enumeration, est en maigre.
% -------------------------------------------------------------------
\begin{VUpage}[148]{144}
\VUnotes{130.35mm}{%
(1) A la komuna raciono esas ne plu multe posibla konfundar a\nl
homo « \textit{posiblo} » kam dezerto. (Videz en l'apendico 4\textsuperscript{a}, la pagino 196.)\nl
Okazione di ta pagino ni memorigas, ke S\textsuperscript{ro} L. \textsc{Couturat} esis tre\nl
remarkinda specalisto pri filozofio, ke il iris aden Germania por\nl
facar longa e minucioza serchi en la manuskripti da \textsc{Leibniz}, ke il\nl
publikigis « L'Algèbre de la Logique » e « La Logique de Leibniz\nl
d'après des documents inédits » kun « Opuscules et fragments inédits\nl
de Leibniz », ke dum un yaro il supleis en \textit{Collège de France} la\nl
eminenta filozofiisto H. \textsc{Bergson}. Lua idei e judiki pri la kozi filo\cc
zofiala konseque havas valoro tote specala e certe plu granda kam\nl
olti di altra personi plu o min profana pri filozofio.%
}

\VUgras{Posiblo} = \textit{to quo esas posibla} : Kad vu forsan opinionus\nl
ke la homo povus transirar la limiti dil posiblo? (1).

\VUgras{-ed-.} --- Ta sufixo havas korespondanti en \textit{ée} F., \textit{at} I.,\nl
\textit{ad} S. Ol adoptesis per la decido 54 « \textit{por indikar la quanto}\nl
\textit{qua plenigas X} (la radiko) \textit{o korespondas ad X} » : \VUgras{bokedo,}\nl
\VUgras{glasedo, manuedo} = \textit{quanto de manjajo, de liquido, de ma}\cc
\textit{terio qua plenigas la boko, un glaso, un manuo} : \VUgras{bokedo de}\nl
\VUgras{pano, glasedo de vino, manuedo de sablo, tero, e. c.}

Altra exempli : \VUgras{kulieredo, pladego, bruetedo, charetedo,}\nl
\VUgras{chariotedo, kamionedo, nestedo, aguledo, pinchedo, glutedo,}\nl
\VUgras{brakiedo, kruchedo, potedo, e. c., e. c.}

Me drinkis nur un glutedo de vino. --- Il prenis ofte pin\cc
chedo di sniflotabako ek mea tabakuyo.

\VUgras{-eg-.} --- Ta sufixo, inter altri, venis che Ido de Espo. Ol\nl
uzesas por formacar \textit{augmentivi} qui indikas grado supera\nl
od extrema, \textit{chanjanta la nociono}. Ol juntesas a nomala ed a\nl
verbala radiki. Exempli : varmega = \textit{varma til la maxim alta}\nl
\textit{grado}. La aquo quan on pozas en kaserolo sur fairo por\nl
boliigar olu, esas unesme kolda, ma pose ol divenas var\cc
meta o tepida, varma, tre varma, e varmega kande ol bolias.\parplein

\textit{Kruela} esas min forta kam \textit{tre kruela}, ma ica esas min\nl
forta kam \textit{kruelega}, nam homo kruelega atingas la extrema\nl
\textit{grado}. Il esas ne nur tre bona, ma bonega e ciencozega. Il\nl
divenis ne nur tre richa ma richega. Il ploris e ploregis kun\nl
singluti. Ne tro bruisez ridante : ridego ne konvenas. Me\nl
tante sufris, ke me kriegis pro doloro.

On ne donez ad \VUgras{-eg-} la pejorativa senco di \VUgras{-ach-}, e ne\nl
aplikez lu a la bestii. La bestii havas \VUgras{boko,} pedi quale la\nl
homi, ma proporcione a sua grandeso e tipi. Same li manjas,
\end{VUpage}

% -------------------------------------------------------------------
%  feuillet 149 — folio 145. Le bloc de notes occupe les quatre
%  cinquiemes de la page : 35 lignes de notes pour 9 de corps, filet
%  a 59,94 mm.
%  IRREGULARITES DE L'ORIGINAL, conservees : « -eg- » cite en gras a une
%  ligne et en italique quatre lignes plus bas ; « -ug » et « -et » sans
%  tiret final quand les memes formes le portent ailleurs sur la page ;
%  et les gloses entre parentheses de la note (3) en romain quand les
%  formes qu'elles glosent sont en italique.
% -------------------------------------------------------------------
\begin{VUpage}[149]{145}
\VUnotes{62.13mm}{%
(1) Ni ne povis imitar D. e E., uzante \textit{plen-}, nam tre ofte tala\nl
vorti ne esas justa : \textit{plenboko} ne esas \textit{bokedo} (D. ipsa havas \textit{Bissen},\nl
peco demordita, apud \textit{Mund voll}); e quo esus : \textit{plenaghulo?} Pluse,\nl
li povas esar dusenca, nam \textit{plen-plado}, exemple, equivalas, pro l'eli\cc
ziono permisata, \textit{plena} plado. Ma donar \textit{plen-plado} o \textit{plena plado} de\nl
ulo ne esas lo sama : l'unesma expresuro indikas nur quanto, la\nl
duesma implikas la plado ipsa.

On ne povis prenar \textit{-ad} (veninta de \textit{-ata} L., che S. exemple, e\nl
vicina a \textit{-at} I.) pro ke \textit{-ad} havis altra rolo en Ido. On do kunfuzis \textit{ée}\nl
F. kun \textit{-ad}, quo donis \textit{-ed} e lasis \textit{-e} libera por altro.

(2) Ni sincere konfesez, ke \VUgras{-eg-} semblas tote arbitriala. On\nl
expresis pri olu reprocho plu grava : « Ube esas bruiso, o kande on\nl
telefonas, ta sufixo povas facile konfundesar al diminutivo \textit{-et}. »\nl
Se la futuro montros vere ta desavantajo, pos la periodo di stabi\cc
leso forsan on povus (por chanjar minim posible) chanjar \textit{-eg-} ad\nl
\textit{-ug-} qua certe ne povus konfundesar a \textit{-et}, segun la propozo quan me\nl
facis en \textit{Progreso}, III, 683. En la sama revuo, IV, 97, S\textsuperscript{ro} \textsc{Jonsson}\nl
judikis lu kom « rekomendinda » e remarkigis : « Kontraste, \textit{-ug}\nl
diferas perfekte de \textit{-et-} ed ol semblas aparte fortunoza pro la memo\cc
rigo por Angli pri E. \textit{huge} (imensa); \textit{animalugo} facile komprenesas\nl
quale E. \textit{a huge animal}; \textit{longuga} quale E. « \textit{hugely long} ».

Me lore agis nulo por ke l'Akademio adoptez mea propozo, e nun\nl
me nule solicitas futura adopto por olu. Sempre me esis e sempre me\nl
esos l'unesma obediero dil decidi, acepti o refuzi di nia Akademio,\nl
judikante kun la tante regretata \textsc{Couturat}, ke ne penvaloras elektar\nl
e laborigar akademio, se on ne aceptas lua decidi, o se mem on\nl
kritikas li. Singlu facez sua linguo e proklamez, ke nur olu esas\nl
digna aceptesar dal mondo. Pri Ido, quale olim pri Espo, me\nl
predikos la diciplino proprexemple, tante plu ke me ne havas, quale\nl
altri, la doturo di neeroriveso.

\VUgras{Ego} en Ido indikas « la \textit{me} filozofiala ».

(3) Se on atencos, en Esperanto, kompozaji quale : \textit{vertebrohava}\nl
(vice : vertebrohavanta), \textit{fruktodona} (vice fruktodonanta), on ne\nl
astonesos, ke Zamenhof konceptis \textit{-ema} kom modifikuro di \textit{-ama} =\nl
qua amas (\textit{-ama}, vice \textit{amanta}). Il departis do de \textit{lernama} = qua%
}

\VUcontinue
drinkas, e. c., quale la homi; konseque esas adminime neu\cc
tila uzar specala vorti por organi e funcioni quin la bestii\nl
havas komuna kun la homi (1).

\textit{Treege}, restajo Esperantala, fortunoze desaparis de che ni.\nl
Ni do ne uzez lu. Dicez : \VUgras{Me chagrenas extreme (en maniero}\nl
\VUgras{extrema) pro la tristiganta novajo quan vu savigas da me.}\nl
E nultempe uzez \textit{ege} o \textit{ega} (2).

\VUgras{-em-.} --- Ta sufixo soldesas a verbala radiki e signifikas :\nl
\textit{inklinita ad..., tendencanta ad...} : \VUgras{agema} (3), \VUgras{kredema} =
\end{VUpage}

% -------------------------------------------------------------------
%  feuillet 150 — folio 146.
%  LIGNE INVISIBLE RETROUVEE : « lema e. c. », derniere ligne de la note
%  en continuation. Le depisteur avait raison. Le releveur l'a etablie
%  par le profil d'encre — amas net de 16 rangees portant 30 px chacune,
%  soit 8 % d'une ligne pleine, entre deux blancs absolus — puis lue a
%  l'agrandissement.
%  IRREGULARITES DE L'ORIGINAL, conservees : « Fundamento » en romain a
%  une ligne quand il est en italique deux lignes avant et trois lignes
%  apres ; et « kredemo » en italique la ou les autres occurrences de la
%  meme vedette sont en gras.
% -------------------------------------------------------------------
\begin{VUpage}[150]{146}
\VUnotes{110.24mm}{%
\VUcontinue
amas lernar, e formacis \textit{lernema}; de \textit{babilama} il formacis \textit{babi}\cc
\textit{lema} e. c.

Cetere ni nule blamas Zamenhof venir ad \VUgras{-em-} per ta voyo e\nl
donir a sua linguo sufixo ne existanta en linguo naturala. La ques\cc
tiono ya ne esas, ka ta od ica sufixo esas naturala (uli ne mem\nl
povas esar naturala), ma kad li esas utila e klara. Se yes (e se\nl
altra plu bona ne existas), nule importas kad li havas o ne havas\nl
modelo en la « naturala » lingui. Se on volus esar plu « naturala »,\nl
on uzus (quale \textit{Neutral}) la sufixo \textit{-ator} por l'instrumento e l'aganto,\nl
la sufixo \textit{-orio} por la loko e (itere) por l'instrumento, ed on dronesus\nl
en la sama kaoso kam la « natural » lingui.

(1) On remarkez, ke ta vorti ne esas plu longa kam en la F. lia\nl
korespondanti : \textit{activité, crédulité, prodigalité}. \textit{Venjemeso} esas mem\nl
min longa kam la perifrazo : \textit{nature vindicative, caractère vindi}\cc
\textit{catif}, unika moyeno quan posedas la Franca por expresar l'ideo di\nl
\textit{venjemeso}. Esperanto, qua ne expresas \textit{es} (\textit{ec}) en ta vorti, posedas\nl
tre ofte kompozaji plu longa : \textit{venghemulino} (venjemino), \textit{malla}\cc
\textit{boremulo} (ociemo, ociero), \textit{mallaboremulino} (ociemino, ocierino).%
}

\VUcontinue
\textit{inklinita ad agar, a kredar}; \VUgras{venjemo} = \textit{homo venjema}; \VUgras{pro}\cc
\VUgras{digemo} = \textit{homo prodigema}; \VUgras{agemulo} = \textit{viro, homulo agema};\nl
\VUgras{agemino} = \textit{muliero, homino agema}; \VUgras{sentemino} = \textit{homino}\nl
\textit{sentema.}

La qualeso abstraktita qua korespondas a ta adjektivo\nl
devas esar expresata per la sufixo \VUgras{-es-} : \VUgras{agemeso, krede}\cc
\VUgras{meso, venjemeso, sentemeso, prodigemeso} (1). L'Esperan\cc
tisti, ne uzante sua \VUgras{-ec-} por ta vorti, duople violacas sua\nl
\textit{Fundamento} netushebla : 1\textsuperscript{e} pro ke, ye \VUgras{-ec-}, la santa libro\nl
dicas « ec » \textit{indikas (abstrakte)}; do, sen « ec », l'abstrakteso\nl
ne indikesas; 2\textsuperscript{e} pro ke, se \VUgras{kredema} = (quale trovesas en la\nl
Fundamento) \textit{kiu havas la inklino, la kutimo kredi}, \VUgras{kredemo}\nl
= (quale en Ido) indikas \textit{ne la qualeso} ipsa, \textit{ma la homo}, kiu\nl
havas la inklino, la kutimo kredi. Konkluze : per \textit{kredemo}\nl
ed analogi la \textit{Fundamento} esas violacata e la qualeso ne\nl
expresata.

Pro ke ol esas nur sufixo, \VUgras{-em-} ne darfas uzesar izolite.\nl
Do nultempe uzez \textit{emo} vice \textit{inklineso, tendenco.}

Ni rimarkigez la grosa eroro etikala quan facis Zamenhof,\nl
inkluzante en \VUgras{-em-} du senci tante diversa kam l'inklineso\nl
e la kustumo. Per la volado, homo povas frenagar, represar\nl
mem maxim forta inklineso naturala. Exemple on povas\nl
esar nature tre \VUgras{iracema, desobediema, ociema}, e. c., e tamen\nl
ne \VUgras{iracera, desobediera, ociera} (Videz la sufixo \VUgras{-er-}). Esas
\end{VUpage}

% -------------------------------------------------------------------
%  feuillet 151 — folio 147.
%  LIGNE INVISIBLE RETROUVEE : « defino. », large de 117 px seulement.
%  DOUTE CONSIGNE PAR LE RELEVEUR, non tranche : la graisse de « aj »
%  cite en vedette. L'agrandissement le donne gras, la largeur de fut
%  le donne romain, et la zone est delavee par la gouttiere. Compose en
%  gras, par analogie avec les autres suffixes cites de la page.
%  IRREGULARITE DE L'ORIGINAL, conservee : « iero » sans tiret initial
%  quand tous les autres membres de sa liste le portent.
% -------------------------------------------------------------------
\begin{VUpage}[151]{147}
\VUnotes{138.35mm}{%
(1) Ta sufixo esas tante necesa, ke mem l'Esperantisti pokope\nl
insinuis lu, e duras, en sua linguo (kun altra sufixi Idala), malgre\nl
la solena ed explicita refuzo da sua Akademio, qua deklaris li\nl
ne adoptenda nek aceptinda, pro ke ne mem utila, adminime en la\nl
komuna linguo. Li esas tante neutila, ke li sempre plu multe pene\cc
tras en ta linguo, mem kontre la \textit{Fundamento} netushebla. Icon me\nl
dicas nur por vidigar, per exemplo frapanta, quante ta sufixi \textit{neutila}\nl
esas en la realeso \textit{necesa} por la bona expresado dil pensi.%
}

\VUcontinue
do kolosal eroro konfundar inklineso e kustumo en un sola\nl
sufixo, quale Espo; nul moralisto od etikisto povas aceptar lu.\parplein

Granda difero anke inter \textit{ant} e \textit{em} : la homo nun \textit{almonanta}\nl
povas esar nule \textit{almonema.}

\VUgras{-end-.} --- Ta sufixo, existanta en la Latina e novlatina\nl
lingui, quik adoptesis. Ol soldesas a radiki verbala e signi\cc
fikas : \textit{qua devas o mustas esar ...ata} : \VUgras{facenda laboro} =\nl
\textit{laboro qua devas o mustas esar facata}; \VUgras{letro lektenda} o \VUgras{skri}\cc
\VUgras{benda} = \textit{letro qua devas o mustas esar skribata}; \VUgras{sumo}\nl
\VUgras{pagenda, problemo solvenda, korpo dissolvenda} forsan ne\nl
pagesos, ne solvesos o dissolvesos; li do tote ne equivalas :\nl
\textit{sumo pagota, problemo solvota} e \textit{korpo dissolvota.}

On pluse remarkez, ke \VUgras{-end-} tote ne egalesas \VUgras{-ind-} (videz\nl
ica) nam ol nule implikas merito : \VUgras{problemo solvenda}, e\nl
forsan \VUgras{solvota}, povas ne meritar solvo, do ne esar \VUgras{solvinda},\nl
nek mem \VUgras{solvebla.}

Remarkez la granda difero inter : \VUgras{me havas nulo skri}\cc
\VUgras{benda} (nulo devas o mustas skribesar da me) e : \VUgras{me havas}\nl
\VUgras{nulo por skribar} = \textit{me havas nula del kozi necesa por skri}\cc
\textit{bar} : \textit{papero, plumo,} e. c.

Ta sufixo darfas formacar substantivi : \VUgras{dividendo, multi}\cc
\VUgras{plikendo}, e. c. Ol darfas unionesar ad altra sufixo, nome\nl
\VUgras{aj} : \VUgras{kredendaji, facendaji} = \textit{kozi kredenda, facenda} (1).

\VUgras{-er-.} --- Ta sufixo (moderna formo : \textit{-ier, -er, -eur} F.;\nl
\textit{iero, -aio} I., \textit{ero} S., \textit{eiro} Port.) venas de \textit{arius} Latina. Ol quik\nl
adoptesis dal Komitato di la \textit{Delegitaro} por indikar la \textit{ama}\cc
\textit{toro di}... Ma pose la decidi 591 e 690 pluampligis lua senco.\nl
Konseque nun ol indikas : \textit{ta qua kustume o mem nur ofte}\nl
\textit{okupas su pri}... sen facar ek la kozo sua profesiono (Videz\nl
sufixo \VUgras{-ist}). Quale on vidas, la amatoro inkluzesas en ta\nl
defino.

La sufixo \VUgras{-er-} soldesas a radiki nomala ed a radiki ver\cc
\end{VUpage}

% -------------------------------------------------------------------
%  feuillet 152 — folio 148.
%  IRREGULARITE DE L'ORIGINAL, conservee : le meme suffixe est compose
%  en gras dans la note (1) et en italique dans la note (3).
% -------------------------------------------------------------------
\begin{VUpage}[152]{148}
\VUnotes{99.19mm}{%
(1) En noto dicesas : Ica decido chanjas nulo en l'antea decidi\nl
(quin on lektis supere) pri la sufixo \VUgras{-er-}. Ol relatas nur la propozo\nl
di sentempa participo, qua esis repulsata (decido : 428. --- IV, 322).\parplein

(2) \VUgras{er}, en Esperanto, formacas tre mikra nombro de vorti plu o\nl
min justa, ex. : \textit{fajrero} (cintilo), \textit{hajlero} (greluno). Lua uzo kom\nl
nedependanta vorto (\textit{ero} = elemento!) privacas ta linguo de l'inter\cc
naciona substantivo \VUgras{ero}, quan Esp. tradukas per \textit{tempo-kalkulo},\nl
\textit{jar(cent)aro}, en la maxim ortodoxa lexiko!!

(3) En \textit{Progreso} (V, 491) trovesas ico : « Nun kande ni havas la\nl
sufixo \textit{-er} por indikar ta qua facas \textit{kustume} ula ago, importas\nl
restriktar la sufixo \textit{-em} a la nociono di \textit{inklineso, tendenco}. (Me\nl
adjuntas hike : omna serioz etikisto e nome la teologio etikala tote\nl
ne povas aceptar la valoro ambigua di la \textit{-em} Esperantala.) Ulu\nl
povas esar \textit{babilema}, e ne \textit{babilera}, se lua profesiono od obligesi\nl
impedas lu \textit{babilar}. Same \textit{drinkemo} povas ne esar \textit{drinkero} (tale ke\nl
pri lu on darfus dicar : \textit{drinkemo ne drinkera}). \textit{Furtemo} povas ne\nl
esar \textit{furtero; inventemo, inventero}, e. c. Ni havas do en ta du sufixi\nl
moyeno expresar precize du diversa nuanci... Nulo impedas dicar :\nl
\textit{babilera, drinkera, furtera}... Nam ica vorti esas tote analoga a parti\cc
cipi; li equivalas : \textit{kustume babilanta, drinkanta, furtanta}. On\nl
remarkez bone la diversa expreso-moyeni quin nia linguo furnisas\nl
a ti, qui studias ol serioze e praktikas ol habile. »%
}

\VUcontinue
\VUgras{bala : biciklero (amatoro), gimnastikero, arkero, fumero,}\nl
\VUgras{kolektero, chasero, peskero. --- Il nultempe enoyas, pro ke}\nl
\VUgras{esante granda lektero, fotografero ed exkursero il sempre}\nl
\VUgras{okupesas. --- Dum sep yari il esis prezidero di nia societo}\nl
\VUgras{e nun il esas lua honor-prezidero.}

Ta sufixo aplikesas anke a nomi di bestii o di objekti qui\nl
esas karakterizata da special ago kustumala : \VUgras{reptero, rumi}\cc
\VUgras{nero, rodero, klimero; remorkero, krozero, flotacero} (Vid.\nl
sufixo \VUgras{-il}).

Per la decido 689 l'Akademio aprobis, ke la participal\nl
sufixo \VUgras{-ant-} expresez « la nuna ago » ed ol adoptis la sufixo\nl
\VUgras{-er-} por « la ago ne precize nuna » exemple : \VUgras{suskriptero} (1).\nl
Societo havas prezidero, ma la \textit{prezidanto} di ula kunsido\nl
povas esar sive la prezidero, sive la vice-prezidero, sive altra\nl
membro.

Komparez \VUgras{-er-} a \VUgras{-ist-} ed anke a \VUgras{-em-} por ne uzar ica od\nl
ita neoportune.

\VUgras{-ero} (2) darfas, kompreneble, recevar la formo adjekti\cc
vala. Ex. : \VUgras{Yes, me konfesas lo, nature me esas tre iracema}\nl
\VUgras{e venjema; ma me luktis kontre ta duopla inklineso, e me}\nl
\VUgras{ne divenis iracera e venjera} (3).
\end{VUpage}

% -------------------------------------------------------------------
%  feuillet 153 — folio 149.
%  LIGNE INVISIBLE RETROUVEE : « e. c. (1). », large de 161 px.
%  COQUILLE DE L'ORIGINAL, conservee : « la sufixo es formacas », ou le
%  suffixe est imprime « es » en romain nu la ou le contexte appelle la
%  forme « -es- ».
% -------------------------------------------------------------------
\begin{VUpage}[153]{149}
\VUnotes{146.73mm}{%
(1) Same kam \VUgras{videsar} = \textit{esar vidata} = \textit{vidatesar}, on darfas dicar,\nl
ke \VUgras{instrukteso} = \textit{instruktateso} o \textit{instruktiteso}, \VUgras{konverteso} = \textit{kon}\cc
\textit{vertateso} o \textit{konvertiteso}, e. c. Ma sen parolar pri la plu granda\nl
kurteso dil derivaji citita en la texto, li ne kontenas nek postulas la\nl
determino pri tempo, quan implikas necese \VUgras{at} e \VUgras{it}.%
}

\VUgras{-eri-.} --- Ta sufixo, internaciona, adoptesis per la decidi\nl
594 e 1091. Ol indikas : \textit{establisuro en qua on okupas su}\nl
\textit{irgamaniere} (ne necese per fabrikado) \textit{pri la kozo indikata}\nl
\textit{da la radiko.} Ol soldesas a radiki nomala : librerio,\nl
papererio, lakterio, od a radiki verbala : fabrikerio, impri\cc
merio, direkterio.

Ol indikas \textit{establisuro, firmo,} e. c. e distingesas de -ey,\nl
qua indikas \textit{irga loko} destinata a...; \VUgras{drinkerio} = \textit{establisuro}\nl
(taverno) \textit{en qua on drinkas}; \VUgras{drinkeyo} = precipue la \textit{loko ube}\nl
\textit{bestii drinkas.} Cetere, en \textit{...erio}, povas esar \textit{...eyo.} Exemple : \VUgras{On}\nl
\VUgras{dicis a me, ke vu rikonstruktigas vua imprimerio.} --- \VUgras{To ne}\nl
\VUgras{esas exakta, nam me rikonstruktigas nur l'imprimeyo ipsa;}\nl
\VUgras{on ne tushos la cetero dil imprimerio.}

\VUgras{-es-.} --- Ta radiko di la lingui Indo-Europana, qua donas\nl
la verbo \VUgras{es-ar}, pleas la rolo di sufixo en la sequanta kazi :\parplein

1\textsuperscript{e} Soldita al radiko di verbo transitiva, ol formacas\nl
pasiv-abreviuro por ta verbo : \VUgras{videsar, trovesar} = \textit{esar}\nl
\textit{vidata, esar trovata.}

Konseque ol uzesas por formacar substantivi qui expresas\nl
\textit{pasiva stando}, sen irga ideo tempala : \VUgras{instrukteso} = \textit{la fakto,}\nl
\textit{stando esir od esar instruktata}; \VUgras{izoleso} = \textit{la fakto, stando}\nl
\textit{esir od esar izolata}; \VUgras{konverteso} = \textit{la fakto, stando esir od}\nl
\textit{esar konvertata}; \VUgras{konvinkeso} = \textit{la fakto, stando esir od esar}\nl
\textit{konvinkata}; \VUgras{tenteso} = \textit{la fakto, stando esir od esar tentata};\nl
\VUgras{konstrukteso} = \textit{la fakto, stando esir od esar konstruktata,}\nl
e. c. (1).

2\textsuperscript{e} Soldita a nomala radiko, la sufixo es formacas subs\cc
tantivi, qui expresas la \textit{stando} od \textit{abstraktita qualeso} : \VUgras{ava}\cc
\VUgras{reso, beleso, blindeso, maladeso, utileso} = \textit{la stando esar}\nl
\textit{avara, bela, blinda, malada, utila.} Ol devas uzesar omnafoye\nl
kande on derivas de adjektivo l'abstraktita qualeso kores\cc
pondanta.

3\textsuperscript{e} Ta sufixo formacas anke verbi signifikanta : \textit{esar tala}\nl
\textit{qual indikas la radiko} : \VUgras{utilesar} = \textit{esar utila} a, por...; egale\cc
\end{VUpage}

% -------------------------------------------------------------------
%  feuillet 154 — folio 150.
%  COQUILLE DE L'ORIGINAL, conservee : dans la note (2), « di 1894. en
%  qua il lasis » — un point suivi d'une minuscule, la ou une virgule
%  est attendue.
% -------------------------------------------------------------------
\begin{VUpage}[154]{150}
\VUnotes{90.34mm}{%
(1) Nula konfundo esas posibla inter ica lasta uzo di la radiko \VUgras{es}\nl
ed olua uzo en la verbi pasiva, pro la dicerno inter la radiki verbala\nl
e la radiki neverbala (o nomala) quan omnu facas instinte. Nulu\nl
inklinesas komprenar : \VUgras{amar} = \textit{esar amo}, ma komprenas lu kom\nl
\textit{abreviuro} di \VUgras{am(at)esar.}

La Komitato dil Delegitaro quik adoptis la propozita radiko \VUgras{es}\nl
por \VUgras{esar} e la uzo supere expozita. On lektas lo yena, pag. 24 di\nl
\textit{Compte Rendu des Travaux du Comité} : « On varme aprobis la iden\cc
tigo dil sufixo \VUgras{ec} al radiko dil verbo \textit{esti} (\textit{esar}), e la logikal klareso\nl
qua rezultas de olu por la verbi dil tipo \VUgras{utilesar} e la korespondanta\nl
substantivo (\VUgras{utileso}). »

(2) \VUgras{-i\/j-} esis (quale en Espo) tote arbitriala e, pluse tro simila\nl
a \textit{ig}, fine ledega en l'anciena : \textit{rujijis, chanjijis, sajijis, ebriijis,}\nl
\textit{glaciijis} e. c.; ol desplezis pro la sono e mem pro la formo. No, vere\nl
tala vorti esis delico nek por la okulo, nek por la orelo, e lia kores\cc
pondanti en Esperanto esas digna kompani por \textit{kn, sc, kaj, -ajn,}\nl
\textit{-ojn, -ujn} e. c. Ido ne povis konservar to omna.

La \VUgras{ek-} Esperantala havas, kom inkoativo, nula susteno en nia\nl
lingui; ol esas tam arbitriala kam \textit{igh}. Zamenhof ipsa agnoskis lo,\nl
quale montras lua reformo (ne aceptita, e vere ne aceptinda) di 1894.\nl
en qua il lasis a \textit{ek-} nur la funciono indikar ago instantala (?) e\nl
remplasigis \textit{igh} per \textit{isk}. Reale \textit{ek-} esis la kontreajo di \textit{-ad}, en la\nl
mento di Zamenhof (\textit{ek-krii, kriadi}), e la du destinesis korespondar\nl
a to, quon on nomizas « aspekto » en la konjugado di la Rusa.\nl
Komence, sub Rusa plumi, preske omna personal modi di Espo%
}

\VUcontinue
\VUgras{sar} = \textit{esar egala a...}; \VUgras{similesar} = \textit{esar simila a...}; \VUgras{friponesar}\nl
= \textit{esar fripona} (o fripono); \VUgras{profetesar} = \textit{esar profeta} (o\nl
profeto); \VUgras{hipokritesar} = \textit{esar hipokrita} (o hipokrito); \VUgras{gastesar}\nl
= \textit{esar gasta} (o gasto); \VUgras{hostesar} = \textit{esar hosta} (o hosto); \VUgras{mala}\cc
\VUgras{desar} = \textit{esar malada} (o malado) (1).

On tote darfus uzar \VUgras{es(o)} mem izolita, nam ol esas radiko\nl
esence. Ma fakte \VUgras{stando} (de \textit{standar}) remplasas lu konstante.\nl
Tamen nule esus nekorekta dicar : \VUgras{Me trovis lu en bona eso.}\nl
Same on darfus dicar : \VUgras{Quale vu esas?} quo tote ne signi\cc
fikus : \VUgras{quale vu existas?} nam esar ne signifikas existar.

\VUgras{-esk-.} --- Greka e Latina, ta sufixo ritrovesas en novla\cc
tina lingui. On nomizas lu \textit{inkoativo}, to esas \textit{montranta}\nl
\textit{komenco}. Ol quik adoptesis del reformo-plano kun la senco :\nl
\textit{komencar}, kande ol juntesas a radiki verbala : \VUgras{dormeskar} =\nl
\textit{komencar dormar}; \VUgras{morteskar} = \textit{komencar mortar}; \VUgras{videskar}\nl
= \textit{komencar vidar}; \VUgras{sideskar} = \textit{komencar sidar}; \VUgras{iraceskar} =\nl
\textit{komencar iracar}; \VUgras{emoceskar} = \textit{komencar emocar}.

Pose l'Akademio per la decido 1221 remplasigis \textit{ij} (2)
\end{VUpage}

% -------------------------------------------------------------------
%  feuillet 155 — folio 151.
% -------------------------------------------------------------------
\begin{VUpage}[155]{151}
\VUnotes{137.61mm}{%
\VUcontinue
ornesis per \textit{ek-} o \textit{-ad}. Bela e multa traci di ta « aspekto » ankore\nl
trovesas en la duesma « santa libro », \textit{Fundamenta Krestomatio,}\nl
stilo-modelo por l'Esperantisti, segun explicita deklaro da Zamen\cc
hof. Ma sempre plu multe li oblivias la modelo por individuala\nl
preferaji. Evoluciono, evoluciono, tante influata da Ido!

(1) Diferanta de \textit{kampeyestro}, nam Ido havas \textit{kampar, kampo,}\nl
\textit{kampeyo}, por qui Espo donas nur ridinda vorti pro sua \textit{kampo, qua}\nl
\textit{egalesas agro.}

(2) Diferanta de \textit{pafeyestro.}%
}

\VUcontinue
(heredita de Esperanto) per \VUgras{-esk} unionebla ad adjektivi\nl
kun la senco : \textit{divenar}... Ex. : \VUgras{paleskar, febleskar, riches}\cc
\VUgras{kar, vireskar, muliereskar, adulteskar, amikeskar,} e. c. =\nl
\textit{divenar pala, febla, richa, vira} (o viro), \textit{muliera} (o muliero),\nl
\textit{amika} (o amiko). En la realeso, \VUgras{paleskar} = \textit{komencar esar}\nl
\textit{pala} e tale pri la ceteri. Do ne existas difero esencala inter\nl
l'uzo di \textit{esk} kun verbi e l'uzo di \textit{esk} kun adjektivi : ol in\cc
dikas : o \textit{komenco} di ago, o \textit{komenco} di qualeso.

Noto relatanta la decido 1221 dicis : « La sufixo \textit{-esk}\nl
uzesos do kun la verbal radiki (quale til nun) por signi\cc
fikar \textit{komencar} (\textit{dormeskar}) e kun la nomal radiki por signi\cc
fikar \textit{divenar} (\textit{paleskar}). L'apliko di ca regulo en la tota vor\cc
taro pruvis, ke ol genitas nula dusenceso. »

Por expresar l'ideo : \textit{divenar ...ata}, la du formi \VUgras{...ateskar}\nl
e \VUgras{...eskesar} tote darfas konsideresar kom regulala e kon\cc
seque legitima. Ex. : \VUgras{vidateskar, uzateskar, probateskar} o :\nl
\VUgras{videskesar, uzeskesar, probeskesar,} e. c. Ma on konsilas uzar\nl
prefere, kom plu simpla : \textit{divenar vidata, uzata, probata.}

Se la prima verbo esas transitiva, la verbo derivita per\nl
\VUgras{esk} esas anke transitiva, do povas havar objekto : \VUgras{lekteskar}\nl
\VUgras{jurnalo.}

Kompreneble \VUgras{esk} ne darfas uzesar kom radiko vice \textit{komenc.}\nl
Do nultempe dicez \textit{eskar, eskas}, quale ul Espisti dicas \textit{eki,}\nl
\textit{ekas}, vice \textit{komenci, komencas.}

\VUgras{-estr-.} --- Inspirita dal duesma parto dil vorti : \textit{Bürger}\cc
\textit{meister} D., \textit{burgmaster} E., \textit{borgomastro} I., \textit{burgomaestre} S.,\nl
\textit{burgo-mestre} Por., \textit{bourgmestre} F. Ol quik adoptesis kun la\nl
senco : \textit{chefo, mastro di}... Ol soldesas a radiki nomala :\nl
\VUgras{urbestro, kastelestro, chapelestro, hotelestro} (hotel-mastro),\nl
\VUgras{policestro, domestro, skolestro, navestro, statestro,} e. c.

Ol mem darfas okazione soldesar a verbal radiko :\nl
\VUgras{kampestro} (1), \VUgras{manovrestro, pafestro} (2), e. c.
\end{VUpage}

% -------------------------------------------------------------------
%  feuillet 156 — folio 152.
%  LIGNE INVISIBLE RETROUVEE : « -eri.) », derniere ligne d'un alinea.
%  COQUILLES DE L'ORIGINAL, conservees : « distinita » pour destinita,
%  et « religial » pour religiala, sur la meme ligne.
%  DOUTE CONSIGNE PAR LE RELEVEUR : la graisse de « direktisto », que
%  l'oeil donne romain et la mesure de fut donne gras, sur une ligne
%  moins encree que ses voisines. Compose en gras, par coherence de
%  liste.
% -------------------------------------------------------------------
\begin{VUpage}[156]{152}
\VUnotes{130.98mm}{%
(1) La \textit{chj} e \textit{nj} di Espo transformas la bapto-nomi til nerikono\cc
cebleso. Pro to li esas kondamninda e Zamenhof ipsa agnoskis to, nam\nl
en sua reformo di 1894, il remplasigis li, quale Ido, per la sufixo \textit{et},\nl
sen altero di la radiko : \textit{Paul(us), Pauletus; Petr(us), Petretus;}\nl
\textit{Luk(as), Luketas; Mari(a), Marieta; Sofi(a), Sofieta} e. c. En la\nl
tradukuri on darfas admisar : \textit{Ted, Dick, Sacha, Sonia,} e. c. quale la\nl
cetera propra nomi.

(2) Ol esas quaza modifikuro di \textit{uj} (uy), respondis Zamenhof\nl
questionita : \textit{uj} esas « recipiento por... », \textit{ej} esas « loko por... ».

(3) \textit{Manjo-chambro, dormo-chambro} en domi privata.

(4) Videz la sufixo \textit{-ier.}%
}

Kom vorto izolita ne uzez \VUgras{estro} nek \VUgras{estraro} quale Espe\cc
rantisti, ma \VUgras{chefo, chefaro, mastro, patrono, direktero,}\nl
\VUgras{direktisto,} e. c. segun la kazi.

\VUgras{-et-.} --- Ta sufixo trovesas en la Latina e novlatina lingui,\nl
kom \textit{diminutivo}. Ol quik adoptesis por indikar \textit{grado extrema}\nl
\textit{di mikreso o febleso qua chanjas la nociono}. Ex. : \VUgras{dometo,}\nl
\VUgras{rivereto; varmeta, beleta; ridetar, kantetar, saltetar, dor}\cc
\VUgras{metar,} e. c.

Ta sufixo formacas anke diminutivi di afeciono : \VUgras{filieto,}\nl
\VUgras{matreto, patruleto, fratuleto, fratineto,} e. c. (1).

\VUgras{-ey-.} --- Ta sufixo, veninta de Esperanto, quik adop\cc
tesis (2). Soldita a radiko nomala o verbala, lu formacas sub\cc
stantivi nomizanta la \textit{loko destinita} (ad objekto od ago). Ta loko\nl
esas edifico, chambro, e. c., loko en la maxim generala senco,\nl
tereno di kultivado por... Ex. : \VUgras{kavaleyo} = \textit{kavalostablo};\nl
\VUgras{boveyo} = \textit{bovostablo}; \VUgras{mutoneyo} = \textit{mutonostablo}; \VUgras{haneyo};\nl
\VUgras{manjeyo, dormeyo} (en kolegio, hospitalo) (3); \VUgras{koqueyo,}\nl
\VUgras{laveyo, tombeyo, pregeyo} (ne \textit{kirko}, nek \textit{kapelo} o \textit{templo},\nl
ma chambro distinita a religial exerci); \VUgras{pastureyo} (senco\nl
plu ampla kam \textit{prato}); \VUgras{viteyo, alneyo; querkeyo, ribiereyo,}\nl
\VUgras{roziereyo, kastaniereyo} (4).

La senco dil derivaji povas esar precizigata da altra sufixi\nl
o radiki : \VUgras{parfumeyo, parfumvendeyo; aer-kuracerio.} (Vid.\nl
-eri.)

Pro la signifiko kelke ne preciza di \textit{-ey} en ula derivaji,\nl
uzez vorti specala, kande la senco lo necesigas : \VUgras{gimnazio,}\nl
\VUgras{liceo, skolo, universitato} vice \textit{lerneyo}; \VUgras{baziliko, katedralo,}\nl
\VUgras{kirko, kapelo, templo,} vice \textit{pregeyo}.

Ta sufixo devas ne uzesar kom radiko izolita. Selektez,\nl
segun la kazo : \VUgras{edifico, konstrukturo, domo, loko, agro,} e. c.,\parplein
\end{VUpage}

% -------------------------------------------------------------------
%  feuillet 157 — folio 153.
%  COQUILLES DE L'ORIGINAL, conservees : une parenthese fermante sans
%  ouvrante (« kun la senco) : »), la parenthese ouverte a la ligne
%  precedente etant deja refermee ; et une rupture d'italique en plein
%  syntagme, « e quaze » restant en romain entre deux passages italiques.
% -------------------------------------------------------------------
\begin{VUpage}[157]{153}
\VUnotes{144.16mm}{%
(1) \textit{Episkopio} esas la diocezo, ed \textit{episkopeyo} esas la palaco, la\nl
domo dil episkopo. Sama difero inter \textit{prefektio, prefekteyo}; sen obli\cc
viar l'ofico : \textit{episkopeso, prefekteso}.

(2) Per la decido 590 la sufixo \VUgras{ier} remplasis la sufixo \VUgras{ul}, qua\nl
divenis libera por indikar la enti di maskulsexuo.

En komento dil decidi 589-593, \textit{Progreso} dicas pri \VUgras{ier} (IV, 565) :%
}

\VUcontinue
e ne \textit{eyo}, nekomprenebla da omna profano. Ni lasez l'Espe\cc
rantisti uzar obstine sua \textit{eyo, eta, uyo, emo, ilo, ebla, ero, ano,}\nl
\textit{ino,} e. c. Per tala vorti li atingas nur konfuzeso quan la L. I.\nl
devas evitar sorgoze.

\VUgras{-i-.} --- Ta sufixo pasable internaciona adoptesis per la\nl
decido 487 por indikar « \textit{lando, regiono, domeno dependanta}\nl
\textit{de...} » e konseque la resortiso dil autoritato di... Ex. : \VUgras{baro}\cc
\VUgras{nio, komtio, dukio, princio, rejio, episkopio} (1), \VUgras{parokio,} e. c.\parplein

\VUgras{-id-.} --- Greka, Latina e kelke internaciona, ta sufixo quik\nl
adoptesis ma nur por indikar : \textit{homa decendanto} (adminime\nl
en la komuna linguo) : \VUgras{Izraelido, la Semidi, la tribuo Rube}\cc
\VUgras{nida,} e. c. Por la homa genituro ne uzez \textit{ido} izolite, ma \textit{decen}\cc
\textit{danto, filio di} (\textit{filiulo, filiino, gefilii,} segun la kazo).

Por la bestia genituro uzez \VUgras{yun} (Vid. ta sufixo).

\VUgras{-ier-.} --- Ta sufixo quik adoptesis (ma en la formo \textit{yer})\nl
kun la senco) : \textit{etuyo kontenanta o sustenanta la kozo indi}\cc
\textit{kata da la radiko}; \VUgras{kandeliero} = \textit{utensilo sustenanta, portanta}\nl
\textit{un kandelo}; \VUgras{plumiero} = \textit{kurta stango portanta un plumo}\nl
(plumoskribilo); \VUgras{sigariero} = \textit{kurta tubo qua portas un sigaro}\nl
(sigarfumilo). Ne konfundez \VUgras{-ier-} a \VUgras{uy}. (Videz ica.)

Pro analogeso (lua senco esante « qua portas »), ta sufixo\nl
uzesas por indikar la \textit{arboro} o \textit{planto qua portas, produktas} :\nl
\VUgras{pomiero, figiero, kafeiero, teiero, palmiero, roziero, fragiero,}\nl
\VUgras{diantiero,} e. c.

Ma on darfas dicar ke, se \textit{pomiero, figiero, roziero,} e. c.\nl
portas, produktas \textit{pomi, figi, rozi,} li esas quaze karakterizata\nl
da ici. To komprenigas, ke l'Akademio per la decido 592\nl
adjuntis a l'unesma senco : « \textit{karakterizata per} » (ulo extera) :\nl
\VUgras{kurasiero} = \textit{soldato portanta kuraso, karakterizata da kuraso};\nl
\VUgras{lanciero} = \textit{soldato portanta lanco, armizita ye lanco} e quaze\nl
\textit{karakterizata da olu}; \VUgras{rentiero} = \textit{homo karakterizata da}\nl
\textit{renti}; \VUgras{miliardiero} = \textit{homo karakterizata adminime da un}\nl
\textit{miliardo} (2).
\end{VUpage}

% -------------------------------------------------------------------
%  feuillet 158 — folio 154.
%  LIGNE INVISIBLE RETROUVEE dans le bloc de notes : « per --- ». »,
%  qui ferme la note reportee du folio precedent.
%  COQUILLES DE L'ORIGINAL, conservees : un deux-points colle au mot
%  (« ol nun: ») la ou toute la page les espace ; et un deux-points
%  redondant devant le signe d'egalite.
% -------------------------------------------------------------------
\begin{VUpage}[158]{154}
\VUnotes{93.81mm}{%
\VUcontinue
« On povas definar ol nun: « qua portas ---, od esas karakterizata\nl
per --- ».

Ed en la sama revuo (VI, 595) dicesas : \VUgras{ier} esas rezervata por\nl
signifikar : « qua portas ulo extera a su ipsa », kontre ke \VUgras{-oz}\nl
signifikas : « qua havas en su nature ed esence ». Pro to on povas\nl
dicar : « \textit{kurasiero, rentiero}, ma ne \textit{gibiero} ».

Komence \VUgras{ier} skribesis \VUgras{yer,} ma l'Akademio modifikis ca lasta orto\cc
grafio per la decido 593, konsequante plu generala chanjo di \textit{y} a \textit{i.}\parplein

Ta sufixo inspiresis direte da la Franca linguo : \textit{figuier, rosier,}\nl
\textit{cuirassier, lancier, rentier} e. c. ed on aplikis lu anke a la radiki por\nl
qui Espo uzas la misterioza \textit{ing} : \textit{kandelingo} = kandeliero. Aplikata\nl
same ad omna radiki di planti, ol donas vorti certe plu komprenebla\nl
da multa homi kam lia Esperantala korespondanti. Komparez :\nl
\textit{figiero, pomiero, pruniero, roziero} a \textit{figujo, pomujo, prunujo, rozujo.}\parplein

(1) « Quankam la verbi derivita per \VUgras{if} esas esence netransitiva,\nl
pro ke li expresas generale kompleta ideo (li kontenas en su la\nl
nociono dil objekto : \VUgras{versifar} : = \textit{facar versi}), ni vidas nula detri\cc
mento en donar a li, en ecepta kazi, komplemento direta, exemple\nl
por dicar : \VUgras{urinifar sango.} To esas kozo tote analoga a \textit{vivar} (mize\cc
roza) \textit{vivo, dormar} (profunda) \textit{dormo}, e. c., ube l'ideo ipsa di la\nl
verbo (di l'ago) divenas komplemento direta. --- On povas anke\nl
dicar : \VUgras{urinifar kun sango,} nam reale la sango mixesas kun urino.\nl
La mediki havos la teknikala vorto \VUgras{hematurio.} » (\textit{Progreso}, V, 350.)\parplein

S\textsuperscript{ro} H. objecionas, ke \textit{sangifar, urinifar} devus ne havar la nuna%
}

\VUgras{-if-.} --- Ta sufixo, trovata en la Latina e novlatina lingui,\nl
quik adoptesis. Ol soldesas a radiko nomala (neverbala)\nl
por formacar netransitiva verbi qui signifikas : \textit{produktar,}\nl
\textit{genitar, sekrecar} : \VUgras{florifar, fruktifar} = \textit{produktar flori,}\nl
\textit{frukti}; \VUgras{sangifar} = \textit{produktar sango} : \VUgras{vunduro sangifanta tre}\nl
\VUgras{abunde}; \VUgras{ha! vu sangifas de la nazo}; \VUgras{versifar} = \VUgras{produktar}\nl
\VUgras{versi}; \VUgras{urinifar} = \textit{sekrecar, produktar urino}; \VUgras{sudorifar;}\nl
\VUgras{burjonifar; chapelifar, vestifar, shuifar; segifar, klovifar}\nl
(tre diferas de \textit{klovizar, klovagar}), \VUgras{martelifar} (altra kam \textit{mar}\cc
\textit{telagar, marteluzar, martelbatar}) (1).

Kompreneble de ta verbi darfas derivar substantivi signi\cc
fikanta la ago : \VUgras{florifo, fruktifo, sangifo, versifo, urinifo.}\parplein

A ta sufixo on darfas unionar altra sufixi : \VUgras{-ey}, \VUgras{-ist}, e\nl
lore \VUgras{ifeyo} signifikos \textit{la loko ube on facas}, fabrikas, e \VUgras{ifisto}\nl
esos ta qua fabrikas : \VUgras{la armifisti fabrikas la fusili en la}\nl
\VUgras{armifeyo dil arsenalo.}

On ne darfas uzar ta sufixo kom vorto izolita. Konseque\nl
dicez : \VUgras{facar, fabrikar, produktar,} e. c., ma nultempe \VUgras{ifar.}\parplein

\VUgras{-ig-.} --- Ta sufixo (modifikuro di \textit{ag}) quik adoptesis. Ol\parplein
\end{VUpage}

% -------------------------------------------------------------------
%  feuillet 159 — folio 155.
%  IRREGULARITE DE L'ORIGINAL, conservee : dans la note, « su levas »
%  est en italique et « su kushas » en romain, a trois mots d'ecart.
%  Le releveur signale aussi que la geometrie donne un faux alinea a
%  l'avant-derniere ligne : elle ouvre sur un signe « = » trop plat pour
%  le filtre de glyphes, et l'abscisse est mesuree sur le mot suivant.
% -------------------------------------------------------------------
\begin{VUpage}[159]{155}
\VUnotes{104.73mm}{%
\VUcontinue
senco, ma la senco : produktar sango, urino; exemple : \textit{la nutrivi}\nl
\textit{sangifas; la reni urinifas.}

Ni previdis ta objeciono kande ni kompozis la vortolibri. Yen nia\nl
respondo. La sufixo \textit{-if} (same kam la ceteri) esas destinata a la\nl
komuna uzado e devas uzesar segun la komuna raciono. Or en nia\nl
lingui, la verbi qui signifikas l'ekfluo od ekpulso di sango, urino,\nl
derivas de la vorti \textit{sango, urino}; ta derivo esas naturala e komoda,\nl
ed on darfas indikar ol per \textit{-if,} sequante la vulgara kompreno. Nur\nl
tre poka homi konocas la reala (interna) produktado di sango,\nl
urino; ma li povas expresar ol per ta vorti ipsa, o, se li preferos,\nl
per « ciencala » ed okulta vorti, quale \textit{hematopoyezo!} La L. I. ne\nl
esas facita nur por li, ma por la maxim multi, do devas sequar\nl
la vulgara koncepti. E ni devas ne havar plu multa skrupuli, dicante\nl
\textit{sangifar,} kam kande ni dicas, ke la suno \textit{su levas} o su \textit{kushas.}\nl
(\textit{Progreso}, IV, 282.)

(1) La radiko Latina \textit{ag} divenas \textit{ig} en kompozaji, e tale havas un\nl
del du senci di la sufixo \textit{ig}. \textit{Reinigen}, en la Germana, quale ni dicis\nl
supere, prizentas l'altra (\textit{net-igar, pur-igar}). Do \textit{ig, igar} ne esas\nl
arbitriala e, segun la justa remarko di S\textsuperscript{ro} O. \textsc{Jespersen}, nia lingui\nl
ofras a ni nulo plu bona praktike.%
}

\VUcontinue
trovas susteno en la Latina (\textit{ab-igo}) ed en la Germana (\textit{rein}\cc
\textit{ig-en}) (1).

Ol soldesas a nomala ed a verbala radiki.

Kun radiko nomala ol signifikas : \textit{donar la qualeso}\nl
\textit{expresata da la radiko, igar tala quala ol dicas, transformar}\nl
\textit{a...} Ex. : \VUgras{fortigar} = \textit{igar forta}; \VUgras{richigar} = \textit{igar richa}; \VUgras{plu}\cc
\VUgras{grandigar} = \textit{igar plu granda}; \VUgras{virigar} = \textit{igar vira} (o viro),\nl
\textit{transformar} (adultulo) \textit{a viro}; \VUgras{mulierigar} = \textit{igar muliera}\nl
(o muliero), \textit{transformar} (adultino) \textit{a muliero.}

Kompreneble de ta verbi derivas substantivi : \VUgras{fortigo,}\nl
\VUgras{richigo, plugrandigo, virigo, mulierigo,} e. c.

Kun radiko verbala \VUgras{-ig-} signifikas : \textit{esar kauzo, ke eventez}\nl
\textit{lo dicata da la radiko} : \VUgras{dormigar} = \textit{esar kauzo di dormo,}\nl
\textit{produktar dormo} : \VUgras{opiumo dormigas}; \VUgras{vidigar ulo da ulu} =\nl
\textit{igar, ke ulo videsez da ulu.}

Se la verbal radiko esas netransitiva (\textit{venar}) la sufixo \VUgras{-ig-}\nl
havas la senco : \textit{igar ...anta,} e konseque la verbo derivita\nl
per \VUgras{igar} havas kom objekto (komplemento direta) to quo\nl
esus la subjekto dil verbo primitiva : \VUgras{venigar} = \textit{igar}\nl
\textit{venanta} : \VUgras{me venigis la mediko} = \textit{me igis la mediko venar}\nl
= \textit{me igis ke la mediko venis}; \VUgras{iracigar} = \textit{igar iracanta} : \VUgras{vua}\nl
\VUgras{sencesa kontredici iracigas lu.}
\end{VUpage}

% -------------------------------------------------------------------
%  feuillet 160 — folio 156.
%  DEUX COQUILLES DE L'ORIGINAL, conservees et signalees par le releveur
%  comme non tranchables entre coquille de composition et sorte abimee :
%  « scndar » pour sendar (un e prive de sa barre) et « Icannes » pour
%  Ioannes (un o ouvert a droite). Deux accidents du meme type sur la
%  meme page rendent l'hypothese du defaut de fonte plausible ; on
%  transcrit ce qui est imprime et l'on signale.
% -------------------------------------------------------------------
\begin{VUpage}[160]{156}
\VUnotes{141.40mm}{%
(1) Se on volus aplikar severega logiko, \textit{nula} afixo povus sold\cc
esar samtempe a verbal ed a nomal radiki, nam lu ne povas havar\nl
exakte la sama senco en la du kazi. To pruvesas per la fakto, ke\nl
on devas definar \textit{ig} kun verbal radiki per \textit{ig... anta} o \textit{ig...ata}. On\nl
povas transirar de un kazo a l'altra nur per \textit{konvencional} abrevio,\nl
qua supresas la sufixo \textit{ant} od \textit{at}, e ta abrevio, qua esas tre komoda,\nl
povas genitar nul \textit{reala} miskompreno (segun \textit{Progreso}, II, 402).%
}

Kontraste, \VUgras{igar} = \textit{igar ...ata}, se la verbal radiko esas tran\cc
sitiva (\textit{vidar, scndar}). Lore, la verbo derivita per \VUgras{igar} havas\nl
kom objekto (komplemento direta) la objekto dil verbo\nl
primitiva : \VUgras{me sendigos a vu ta libro da Petrus} = \textit{me igos}\nl
\textit{ta libro sendata} (sendesar) \textit{a vu da Petrus}; \VUgras{me vidigis mea do}\cc
\VUgras{mo da Icannes} = \textit{me igis vidata} (videsar) \textit{mea domo da I.}; \VUgras{me}\nl
\VUgras{instruktigas mea filii da la skolestro} = \textit{me igas instruktata}\nl
(instruktesar) \textit{mea filii da la skolestro}; \VUgras{me imprimigas nova}\nl
\VUgras{libro da mea imprimisto} = \textit{me igas imprimata} (imprime\cc
sar), e. c.

Segun ta regulo, \VUgras{me manjigas mea kavalo} povas signifikar\nl
nur : \textit{me igas mea kavalo manjesar, me igas ke lu esez}\nl
\textit{manjata.}

Se me volus dicar : \textit{me donas nutrivo a mea kavalo}, me\nl
dicus : \VUgras{me igas mea kavalo manjar.}

Se la verbal radiko esas mixita (transitiva e ne transi\cc
tiva), \VUgras{igar} kompreneble signifikas \textit{igar ...anta} o \textit{igar ...ata},\nl
pro la senco ipsa dil mixita verbi (1). Ma to ne havas detri\cc
mento, quale montras l'exempli sequanta : \VUgras{me komencigas}\nl
\VUgras{la laboro} = \textit{me igas la laboro komencanta} o \textit{komencata}; \VUgras{me}\nl
\VUgras{durigos la laboro da altra laboristi} = \textit{me igos la laboro}\nl
\textit{durata da altra laboristi}; \VUgras{cesigez la infanti ludar} = \textit{igez la}\nl
\textit{infanti cesanta ludar.}

Cetere, se to esas necesa o mem nur utila, on sempre dar\cc
fas uzar explicite : \textit{igar ...anta} o \textit{igar ...ata.}

On tote darfas preferar : \VUgras{igar ulu vidar ulo} a : \textit{vidigar}\nl
\textit{ulo da ulu}; \VUgras{me igos vu konstatar la kozo} a : \textit{me konstatigos}\nl
\textit{da vu la kozo}; e tale pri multa kazi analoga.

Quale on vidis en la supera exempli, la sufixo \VUgras{ig} uzesas\nl
kom radiko dil verbo \VUgras{igar} : \VUgras{me igas vitro neruptebla} = \textit{me}\nl
\textit{traktas vitro} (ja fabrikita) \textit{tale, ke lu divenas neruptebla}, tre\nl
diferas de : \VUgras{me facas neruptebla vitro} (o \VUgras{vitro neruptebla})\parplein
\end{VUpage}

% -------------------------------------------------------------------
%  feuillet 161 — folio 157.
%  DEUX LIGNES INVISIBLES RETROUVEES, toutes deux tres courtes : une
%  dans le corps (~130 px) et une dans les notes (~75 px).
%  COQUILLE DE L'ORIGINAL, conservee : « Dicernez do bone igar do
%  facar. » — le second « do » la ou le sens appelle « de ». Verifie a
%  l'agrandissement x 10.
%  A NOTER : « utensilo » est compose avec sa finale en ITALIQUE (le
%  suffixe cite dans le mot meme), et non en gras comme la vedette qui
%  suit deux lignes plus bas.
% -------------------------------------------------------------------
\begin{VUpage}[161]{157}
\VUnotes{135.13mm}{%
(1) La nomi di instrumenti esas punto extreme desfacile fixi\cc
gebla. Ne sempre on povas departar de la verbo o de la instrumento\nl
ipsa. L'Akademio ja studiis atencoze la temo e kredeble ristudios\nl
olu.

Pri la origino di \textit{-il-} yen ulo interesanta :

« La sufixo \textit{-il-} perfekte konkordas kun la sufixi \textit{al-}, \textit{il-} uzata da\nl
omna Slava lingui, e. c. : Slovena \textit{gladilo} = glatigilo, \textit{komilo} =\nl
direktilo, \textit{belilo} = blankigilo, \textit{rezilo} = tranchilo, \textit{ovalo} = plugilo,\nl
\textit{pisalo} = skribilo (J. \textsc{Kovacic}, \textit{Progreso}, V, 622).%
}

\VUcontinue
qua signifikas : \textit{me fabrikas}, e. c. (Dicernez do bone \VUgras{igar}\nl
\VUgras{do facar.})

On anke vidis, ke la verbo \VUgras{igar} darfas havar propoziciono\nl
kom objekto : \VUgras{me igas mea kavalo manjar aveno.} Quale on\nl
vidas, importas (por la justa expreso e kompreno) sequar\nl
la « normal ordino » en ta propozicioni infinitivala : \VUgras{Il igis}\nl
\VUgras{Petrus ocidar Alexander} = \textit{il igis Alexander ocidesar da}\nl
\textit{Petrus.}

Ne konfundez \VUgras{igar} a \VUgras{koaktar} od a \VUgras{lasar}. \VUgras{Il igis sua spo}\cc
\VUgras{zino manjar la kordio di elua amoranto} (el ne savis lo);\nl
\VUgras{il koaktis sua spozino,} e. c. (el ne volis, ma il impozis ad\nl
\VUgras{elu koakte sua volo}). \VUgras{Lasez me enirar} = \textit{ne impedez me}\nl
\textit{enirar}; \VUgras{igez ni enirar} = \textit{ayez por ke ni enirez.}

La uzo di \VUgras{lasar} en la senco di \VUgras{igar} esas Germanajo.

\VUgras{-ik-.} --- Internaciona, ta sufixo adoptesis per la decido\nl
566 kun ica senco : \textit{malada per} : \VUgras{kolerika, koleriko; ftiziika,}\nl
\VUgras{ftiziiko; leprika, lepriko; skrofulika, skrofuliko; hidrop}\cc
\VUgras{sika, hidropsiko; pestika, pestiko; alkoholika, alkoholiko;}\nl
\VUgras{kordiika, kordiiko,} e. c.

La sufixo \VUgras{oz} (videz ica) ne povas expresar juste ta ideo.\nl
Exemple, \textit{kordioza}, \textit{kordiozo} nule dicus, ke la persono esas\nl
malada ye la kordio. Fakte omna homo esas \textit{kordioza}, quale\nl
ol esas \textit{racionoza}; ma, fortunoze, omna homo ne esas \textit{kordiika},\nl
ne sufras de la kordio. La \textit{gibozo} ne sufras (ordinare) de sua\nl
gibo; ol ne esas \textit{gibiko}.

\VUgras{-il-.} --- Finalo dil vorto utens\textit{ilo}, ta sufixo quik adoptesis.\nl
Soldita a verbal radiko, ol donas substantivi signifikanta :\nl
\textit{instrumento por}... Ex. : \VUgras{pektilo, brosilo, entravilo, razilo,}\nl
\VUgras{fotografilo, telegrafilo, telefonilo,} e. c. (1).

Nultempe uzez \VUgras{ilo} kom vorto izolita, ma selektez inter :\nl
\VUgras{aparato, instrumento, utensilo, mashino,} e. c. segun la kazi.\parplein
\end{VUpage}

% -------------------------------------------------------------------
%  feuillet 162 — folio 158.
%  LE FILET ETAIT PERDU POUR UN PIXEL. Le releveur l'a trouve a la main
%  a y = 1407 px, puis a diagnostique la cause dans outils/filet.py : le
%  test « rien d'autre sur la ligne du filet » sommait jusqu'au bord de
%  l'image, ou traine la tranche du feuillet voisin. Corrige (marge de
%  25 px retiree a droite), l'outil rend exactement la meme ordonnee que
%  la mesure a la main, et recupere aussi le feuillet 164.
%  IRREGULARITE DE L'ORIGINAL, conservee : « (homulo) » en romain et
%  « (homino) » en italique, a trois mots d'ecart ; et les memes suffixes
%  cites en italique a une ligne, en gras a la suivante.
% -------------------------------------------------------------------
\begin{VUpage}[162]{158}
\VUnotes{129.07mm}{%
(1) On dicas anke : \VUgras{tiposkribilo} (tiposkribar, tiposkribisto).

(2) Quale on savas, omna nomi di animali (konseque anke di\nl
homa enti) esas \textit{epicena} en Ido, t. e. komuna a la du sexui (Dec.\nl
1089). La maskuleso indikesas per \VUgras{ul} (homulo), la femineso per \VUgras{in}\nl
(\textit{homino}). Pro ke nature e per su, la sexuo ne esas ulo plu grama\cc
tikala kam altra qualeso, ni agus tote nelogike, se ni indikus ol per\nl
dezinenci (\textit{a, o, u, e, i} exemple). La gramatikal dezinenci devas\nl
uzesar nur por kozi gramatikala (modi, tempi e. c.), pro ke vere\nl
on ne vidas per quo esar maskula o femina esas plu gramatikala\nl
kam esar yuna, olda, o richa, povra. La sexuo devas esar expresata\nl
per \textit{maskula, femina} o per lia abreviuri sufixa : \VUgras{-ul, -in.}%
}

\VUcontinue
Pro la neprecizeso di ta sufixo, on uzas radiki specala, vice\nl
olu, kande on volas indikar objekto preciza. Exemple, \VUgras{kul}\cc
\VUgras{telo, hakilo, serpo, siklo, sabro,} e. c. esas \VUgras{tranchili;} kon\cc
seque \VUgras{tranchilo} ne nomizus ica plu reale kam ita. Same :\nl
\VUgras{fusilo, karabino, revolvero, kanono} esas \VUgras{pafili;} se do vu\nl
saciesus da \VUgras{pafilo} por nomizar un de oli, on tote ne savus\nl
exakte pri qua vu parolas.

Se traktesas pri mashino on devas uzar prefere la vorto\nl
« mashino » ipsa : \VUgras{skribmashino} (1), \VUgras{sutmashino, steb}\cc
\VUgras{mashino, drashmashino,} e. c.

\VUgras{-il-} indikas \textit{instrumento} (o implemento) : on brosas \textit{per}\nl
brosilo, on pektas \textit{per} pektilo; ma \VUgras{-er} indikas \textit{aganto} : remor\cc
kero remorkas per \textit{su ipsa}.

\VUgras{-in-.} --- Sufixo Latina (\textit{eroina, heroina}), Franca (\textit{héroïne,}\nl
\textit{Albertine, Pauline}), Italiana (\textit{eroina}), Hispana (\textit{heroina}), Ger\cc
mana (\textit{Heldin, Lehrerin}), Rusa (\textit{grafinja}). Ol quik adoptesis\nl
por indikar la ento femina : \VUgras{rejino, heroino, profesorino,}\nl
\VUgras{instruktistino, docistino, baronino,} e. c. (2).

Pro ke la sexuo devas indikesar per \textit{-ul} (m.) od \textit{-in} (f.)\nl
nur kande to esas necesa, on ne adjuntas \VUgras{-in} a radiki expre\cc
santa per su la femineso : \VUgras{Amazono, Parco, megero, subreto,}\nl
\VUgras{matrono, primadono.} Pro quo on indikus per sufixo la\nl
femineso ja nature indikata da la radiko? Ma \VUgras{feino} esas\nl
uzebla, pro ke on darfus tote logike dicar \VUgras{feulo} (feo mas\cc
kula). On adoptis \VUgras{matro} apud \VUgras{patrino} (Dec., 662) pro ca\nl
motivi : etimologie la radiko \VUgras{matr} esas tam internaciona\nl
kam la radiko \VUgras{patr;} historie la nociono \VUgras{matr} esas adminime\nl
tam anciena kam ta di \VUgras{patr;} fiziologie ed etike la relato\parplein
\end{VUpage}

% -------------------------------------------------------------------
%  feuillet 163 — folio 159.
%  FILET ENCORE PLUS PALE QU'AU FEUILLET 143 : il ne devient continu
%  qu'a un seuil de 220, quand l'outil seuille a 185 et que le feuillet
%  143 demandait 210. Trouve a y = 839 px, long de 216 px, cale a la
%  marge. Ordonnee posee : 79,93 mm.
%  COQUILLES DE L'ORIGINAL, conservees : « iaudigna » pour laudigna ;
%  un point la ou la phrase appelle une virgule ; « la spozo » en romain
%  quand les autres occurrences du mot sont en italique ; et un
%  point-virgule devant deux des quatre numeros d'une enumeration dont
%  les deux autres portent un deux-points.
% -------------------------------------------------------------------
\begin{VUpage}[163]{159}
\VUnotes{79.93mm}{%
(1) Eminenta sociologi, arkeologi, yurohistoriisti e. c. demonstris,\nl
ke ante \textit{patrul-yuro} existis \textit{matro-yuro}, ke mem la nociono « matro »\nl
esas plu anciena kam la nociono « patrulo ». En la primitiva kul\cc
turo on ne konocas « patro », qua esas produkturo di plu tarda\nl
epoki kun developita yurala koncepti, kontre ke la relato inter la\nl
parturanta muliero e l'infanto naskanta produktesas dal naturo.\nl
Konseque omna ciencisti, qui okupas su per la mencionita materio,\nl
judikus, kom tre neoportuna, se li, uzante Ido, devus dicar e\nl
skribar \textit{patrino}, juste kande li pruvas la existo di \textit{matro} ante \textit{patro}\nl
(segun \textit{Progreso}, IV, 143).

« Pro ke dicesas pri la sufixo \VUgras{-in}, quale anke pri la sufixo \VUgras{ul}.\nl
ke on devas uzar ol nur \textit{en kazo di bezono}, uli asertas, ke on devas\nl
dicar : 1\textsuperscript{e} « Mea \textit{filiino} esas la spozo di mea \textit{amiko} », ma ke on ne\nl
darfas dicar; 2\textsuperscript{e} « Mea \textit{filiino} esas la \textit{spozino} di mea \textit{amiko} » e mem\nl
min multe : 3\textsuperscript{e} « Mea \textit{filiino} esas la \textit{spozino} di mea \textit{amikulo} », pro ke\nl
\textit{filiino} implikas ya, ke \textit{spozo} esas femina e ke konseque \textit{amiko} esas\nl
maskula. On asertas mem, ke on povus dicar; 4\textsuperscript{e} « Mea \textit{filio} esas\nl
la \textit{spozo} di mea \textit{amikulo}, pro ke \textit{amikulo} determinas inverse la\nl
sexuo di \textit{spozo} e di \textit{filio}.

Omna ta rezoni esas tre bona... kom espritaji e logikala ludeti.\nl
Ma praktike on ne darfas esar tante severa... e subtila; on ne povas\nl
kalkular tale singla frazo quale per algebro. La formo 1 esas bona\nl
e suficanta; ma la formo 2 esas anke bona; la formo 3 esas nur\nl
kelke pleonasma; e la formo 4, quankam « logikala » esas praktike\nl
absurda, nam on ne devas vartar la lasta silabo di frazo por indikar\nl
(nedirete!) ke \textit{filio} esas \textit{filiino} e ne \textit{filiulo}. »

(2) Vice \textit{ign} (de \textit{digna}) : iaudigna = laudinda, qua propozesis.\nl
\VUgras{-ind} versimile esas modifikuro di \textit{and}(us), \textit{end}(us) latina kun\nl
chanjo di senco.%
}

\VUcontinue
inter \VUgras{matro} e \VUgras{filio} tote ne esas identa al relato existanta inter\nl
\VUgras{patrulo} (patro maskula) e \VUgras{filio} (1).

\VUgras{-ind-}. --- Ta sufixo esis riestablisata dal Komitato ipsa\nl
dil \textit{Delegitaro} (2). Ol juntesas a verbal radiki kun la senco :\nl
« \textit{qua meritas esar ...ata} » : \VUgras{aminda, estiminda, kredinda},\nl
\VUgras{respektinda} = \textit{qua meritas esar amata, estimata, kredata, res}\cc
\textit{pektata.}

Evidente ta sufixo povas ofte esar remplasigata per \VUgras{end},\nl
nam « on devas » estimar la homi estiminda, kredar la\nl
homi kredinda, e. c. Ma lua senco esas plu restriktita e plu\nl
preciza.

Kom vorto izolita ne uzez \textit{inda}, \textit{indeso}, ma \VUgras{digna}, \VUgras{digneso}.\parplein

Remarkez bone, ke \VUgras{em}, \VUgras{er}, \VUgras{iv} (videz ica) havas senco\nl
aktiva, ma \VUgras{ebl}, \VUgras{end}, \VUgras{ind} senco pasiva. Do la tri lasta ne\nl
darfas aplikesar a verbi netransitiva. Ma, se la verbo esas\parplein
\end{VUpage}

% -------------------------------------------------------------------
%  feuillet 164 — folio 160.
%  L'ordonnee du filet, 124,38 mm, vient de la correction apportee a
%  outils/filet.py apres le feuillet 162.
%  DEUX LIGNES INVISIBLES RETROUVEES par le releveur, qu'aucun depistage
%  n'avait signalees puisque la page n'avait pas pu etre depistee : une
%  de 85 px dans le bloc de notes, et une derniere qui tombait hors de
%  la plage exploree par l'outil.
%  IRREGULARITE DE L'ORIGINAL, conservee : deux listes d'exemples de
%  meme fonction, l'une en romain et l'autre en gras, dans le meme
%  article.
%  DOUTE CONSIGNE : dans la note (1), le a final de « Kristana » porte
%  un accent grave que l'ido ne connait pas et qui ne reparait nulle
%  part ; compose sans l'accent, faute de pouvoir distinguer la coquille
%  du defaut d'impression.
% -------------------------------------------------------------------
\begin{VUpage}[164]{160}
\VUnotes{124.38mm}{%
(1) De \textit{Kristana} (Videz la noto 1 p. 140).

(2) Ido : \textit{socio, sociala, socialismo}; Espo : \textit{societo, societa, socia}\cc
\textit{lismo.}

(3) Ido : \textit{turo} (ne \textit{turmo}), \textit{turismo}; Esp. : \textit{turo} (la \textit{turmo} di\nl
Ido), \textit{turismo} kom vorto specala.

Til nun \textit{ismo} ne uzesas en Espo kom vorto aparta!

(4) On bone remarkez, ke la defino dil sufixo \textit{-ist} dicas « profe\cc
sione ». Or la profesiono implikas multa okupi qui ne esas mestiero\nl
propre dicata. Mestiero esas profesiono, ma tre ofte profesiono ne\nl
esas mestiero : \textit{ciencisto}.

(5) Do \textit{Idisto} ne \textit{Idano}, tante plu ke \textit{-an} signifikas esence :\nl
\textit{membro} di : \textit{societano, senatano}, ed instinte omna profani nomizas\nl
ni \textit{Idisti}.%
}

\VUcontinue
mixita (trans. e netrans. : \textit{variar}, \textit{chanjar}) ol recevas ebl\nl
o iv segun l'ideo expresenda : \textit{variiva} o \textit{variebla}; \textit{chanjiva} o\nl
\textit{chanjebla} : ta formo esas variiva o variebla; mea sentimenti\nl
esas nun tote nechanjiva o chanjebla.

\VUgras{-ism-.} --- Ta sufixo internaciona quik adoptesis \textit{por in}\cc
\textit{dikar doktrino, partiso, sistemo} : Kristanismo (1), katoli\cc
kismo, Luterismo, Kalvinismo, Protestantismo, Mahome\cc
dismo, pozitivismo, imperialismo, kapitalismo, empirikismo,\nl
Epikurismo, stoikismo, socialismo (2), turismo (3), skepti\cc
kismo, asketismo, e. c.

\VUgras{-ist-.} --- Ta sufixo internaciona quik adoptesis \textit{por indikar}\nl
« \textit{la homo qua profesione okupas su pri...} » : artisto, pianisto,\nl
flutisto, kantisto, muzikisto, komercisto, telegrafisto, poli\cc
cisto, misionisto, ciencisto, e. c. (4).

Per extenso, qua trovesas en nia lingui, ta sufixo indikas\nl
\textit{adepto} o \textit{adherinto di partiso, skolo} (doktrino) (5); \VUgras{socialisto,}\nl
\VUgras{materialisto, idealisto, e. c.;} ecepte se la radiko dil vorto\nl
en \VUgras{-ism} ja posedas ta senco : \VUgras{Kristano, skeptiko, mistiko,}\nl
\VUgras{katoliko, protestanto, e. c.}

Por distingar la vendisto de la produktisto, on uzas (se\nl
to esas necesa od utila) -vendisto por l'unesma e -ifisto por\nl
la duesma (vid. \textit{-if}) : shu-ifisto (lu fabrikas), shu-vendisto\nl
(lu vendas); od on precizigas altre : flor(vend)isto, flor\cc
(kultiv)isto.

Kompreneble on uzas \VUgras{-er} (studiita) ne \VUgras{-ist,} se la prakti\cc
kanto ne esas vere profesionano : fotografero, biciklero,\nl
dansero, peskero, e. c. (ne profesione); fotografisto, bici\cc
klisto (kurista), dansisto, peskisto (profesione).
\end{VUpage}

% -------------------------------------------------------------------
%  feuillet 165 — folio 161. Premiere page d'un cahier : elle porte la
%  SIGNATURE « 6 » sous le bloc de texte, a gauche. Encre relevee de
%  y 1895 a 1920 px, x 103 a 118.
%  COQUILLE DE L'ORIGINAL, conservee : une virgule entre sujet et verbe
%  dans la note (5).
%  IRREGULARITE conservee : dans la glose francaise, « le fait d'etre
%  mortel, » est en romain quand le reste de la parenthese est italique.
% -------------------------------------------------------------------
\begin{VUpage}[165]{161}
\VUsignature{169.50mm}{6}
\VUnotes{76.96mm}{%
(1) Diferanta de \textit{instruktala} (libro) e mem de \textit{instruktanta}\nl
(libro). Espo dicas \textit{instrua}, por la tri idei, e co per plumi modela.\nl
Advere ol uzas anke la du nejusta vorti : \textit{instruiga, lerniga}.

(2) Tre diferanta de \textit{sentebla} (qua povas sentesar) e de \textit{sentema}\nl
(qua facile, forte sentas) : \textit{inter la sentivi, ici esas multe plu}\nl
\textit{sentema pri la mondo sentebla.}

(3) \textit{Fero esas bona konduktivo di elektro.}

(4) \textit{Esar responsiva pri..., responsar pri...} dicesas en Esperanto :\nl
\textit{respondi pri...}, deplorinda imituro di la Franca, quan on akuzis ni\nl
sequar, e quan reale tro sequis Zamenhof, ankore nun imitata da\nl
sua fidela o nefidela dicipuli en ta eroro.

La sufixo \textit{-iv} esas konforma a l'uzado di nia lingui : \textit{lénitif,}\nl
\textit{sédatif, carminatif}, e. c. Nula sufixo esas plu « naturala » kam lu,\nl
quankam la Esp-isti ne volas agnoskar lo ed opozas ad ol omnaspeca\nl
ne sincera shikani.

(5) Pri ta sufixo, on lektas en \textit{Progreso} (VI, 597) : Ni definas\nl
ol generale per : « \textit{qua povas...} », « \textit{kapabla...} ». Ma la praktiko\nl
duktas ni ad extensar e samtempe precizigar lua senco. Exemple, ula\nl
samideano objecionis, ke irga substanco qua en irga cirkonstanci\nl
povas \textit{purgar}, darfas qualifikesar \textit{purgiva}, do ke ca epiteto ne povas\nl
restriktesar a la substanci nomizita kustume \textit{purgivi}. No, co esus\nl
misuzo di la vorto, t. e. di la sufixo \textit{-iv}. Se ula substanco purgas\nl
acidente e hazarde ulu, ol esas simple \textit{purganta} (o \textit{purginta}) fakte,\nl
en ta kazo; ol ne esas \textit{purgiva}. Do la sufixo \textit{-iv} signifikas exakte :\nl
« \textit{qua havas nature}, esence, la proprajo di... » e konseque : « qua\nl
havas la destino di...; nam la uzi a qui ni destinas singla kozo rezul\cc
tas evidente de lua natural ed esencala propraji. Remarkez cetere, ke\nl
la sufixo \textit{-iv}, esas tre utila por indikar la materii qui esas apta\nl
produktar ula efekto, e qui nule esas utensili (a qui rezervesas \VUgras{-il.})\nl
Ex. : multa korpi esas \textit{konduktiva} (dil elektro) sen esar propre\parplein
}

\VUgras{-iv-.} --- Ta sufixo (internaciona en la formi \VUgras{if, iv}) quik\nl
adoptesis. Ol soldesas a radiki verbala e signifikas : \textit{kapa}\cc
\textit{bla..., povanta...} Ex. : \VUgras{instruktiva} (libro) (1); \VUgras{nutriva, nu}\cc
\VUgras{trivo}; \VUgras{konsolaciva} (paroli); \VUgras{sentiva} (2), \VUgras{konduktivo} (3);\nl
\VUgras{rezistiva} (qualeso); \VUgras{responsiva} (qua povas responsar) (4);\nl
\VUgras{mortiva} (omna \textit{mortivi}), e. c.

La qualeso abstraktita devas esar expresata per la soldo\nl
di \VUgras{-es} al sufixo \VUgras{-iv}, do : \VUgras{iveso.} Ex. : \VUgras{konduktiveso} (F. \textit{con}\cc
\textit{ductibilité}), \VUgras{rezistiveso} (F. \textit{capacité de résistance, résistivité}),\nl
\VUgras{responsiveso} (F. \textit{responsabilité}); \VUgras{mortiveso} (F. \textit{mortalité}, le\nl
fait d'être mortel, \textit{mortiva}); \VUgras{ne-mortiveso} (F. \textit{immortalité}).\parplein

Nultempe uzez \textit{iva, ivo} o \textit{iveso} kom izolita vorti. Dicez :\nl
\textit{povanta, kapabla...; povo, kapableso...} Ex. : \VUgras{la povo, la kapa}\cc
\VUgras{bleso lumizar} esas : \VUgras{lumiziveso} (Vid. \VUgras{-iz}) (5).
\end{VUpage}

% -------------------------------------------------------------------
%  feuillet 166 — folio 162.
%  DOUTE CONSIGNE PAR LE RELEVEUR, en bloc : la graisse des exemples en
%  -izar et -izo. La structure « mot ido = glose italique » appelle le
%  gras partout ailleurs dans le corpus, mais la comparaison directe
%  avec un gras certifie du feuillet 162 les donne romains, et la mesure
%  de fut ne les separe pas du romain voisin. Composes en ROMAIN, comme
%  le fac-simile les donne ; c'est le seul point du folio qui puisse
%  basculer en bloc.
% -------------------------------------------------------------------
\begin{VUpage}[166]{162}
\VUnotes{75.45mm}{%
\VUcontinue
\textit{konduktili}; \textit{explozivo} (t. e. materio exploziva) esas tote diversa de\nl
\textit{explozilo} (qua povas esar amorcilo o simpla perkutilo). Generale,\nl
la sufixo \textit{-iv} esas tre utila por omna materii qui esas definita per\nl
lia efekto ed uzo (medicinala, farmaciala, e. c.) e por distingar li,\nl
sive de la utensili, instrumenti, aparati (sufixo \textit{-il}), sive de la\nl
vivanta (homa) aganti (sufixo \textit{-er}). Existas en la vapor-mashini\nl
aparati nomata \textit{purgili}. E \textit{purgero} esus... \textit{Monsieur Purgon!}

(1) On questionas ni kad on darfas uzar la sufixo \VUgras{-iz} kun verbala\nl
radiki... Ni respondas, ke l'ideo generala di \VUgras{-iz} esas aplikebla a verbal\nl
radiki, mem se lua \textit{vortala} defini, necese tro strikta, semblas restrik\cc
tar ol a la nomala radiki. Ni ja havis l'okaziono montrar, ke \textit{respon}\cc
\textit{dizar} povas esar necesa, o adminime utila (VI, 296). Plu frue ni\nl
anke montris, ke on bezonas \textit{notizar} apud \textit{notar} : se me \textit{notas} mea\nl
reflekti en libro quan me lektas, me \textit{notizas} la libro, me garnisas ol\nl
per noti. Forsan on objecionos, ke me misuzas hike \textit{noto} vice \textit{noturo},\nl
e ke on devus dicar rigoroze : \textit{noturizar}. Ma co esas nur subtila\nl
shikano : nam on garnisas reale la libro per l'ago ipsa notar, same\nl
kam on \textit{respondizas} letro per l'ago ipsa respondar (nam la letro per\nl
qua on respondas nule juntesas materiale a la letro « respondizita »).\nl
Simile \textit{spricizar} signifikas ago derivata de l'ago spricar, e ne de la\nl
kozo « spricuro » : ol signifikas quaze efekto e rezulto dil ago\nl
\textit{spricar}, aplikata al objekto. Ni havas ja (e nulu astonesis o shokesis\nl
da lo) \textit{lumar} e \textit{lumizar}. Ni havas anke, derivata de \textit{fum-ar}, \textit{fum}\cc
\textit{agar} e \textit{fum-izar} (*). Do ni povus definar maxim generale la sufixo\nl
\textit{-iz} tale : « aplikar objekto od \textit{ago} ad (altra) objekto ». Kompre\cc
neble, ica defino esas kelke tro abstraktita por la novici, ed on\nl
sempre bezonos la defini plu konkreta, ma plu strikta, quin on uzis\nl
til nun (\textit{Progr.}, VI, 596).

(*) Simile on povas dicar \textit{bavizar} (D. \textit{begeifern}), \textit{vomizar} (D.\nl
\textit{bespeien}), \textit{desegnizar} (libro), forsan mem \textit{skribizar} (papero).%
}

\VUgras{-iz-.} --- Ta sufixo Greka, Latina, internaciona (ed ankore\nl
nun tre produktanta) quik adoptesis. Ol soldesas a radiki\nl
kun la senco : \textit{provizar, garnisar, indutar, impregnar per...}\nl
Ex. : armizar, lumizar, limitizar, regulizar = \textit{provizar per}\nl
\textit{armi, lumo, limito, regulo}; kolorizar, stukizar = \textit{indutar}\nl
\textit{per koloro, stuko}; salizar, sulfizar = \textit{provizar od impre}\cc
\textit{gnar per salo, sulfo} : karni salizita o sulfizita koruptesas\nl
min rapide; vu ne salizis nek piprizis ta supo.

La substantivi armizo, alkoholizo, elektrizo havas aktiva\nl
senco. Se on volas la senco pasiva, on adjuntas \VUgras{es} a \VUgras{iz}, do\nl
izeso : armizeso, alkoholizeso, elektrizeso = \textit{la fakto esar}\nl
\textit{od esir armizata, alkoholizata, elektrizata} (1).

Kom radiko uzez : provizar, garnisar, impregnar, e. c.,\nl
segun la senci expresenda, ma nultempe \textit{izar} od \textit{izeso}.
\end{VUpage}

% -------------------------------------------------------------------
%  feuillet 167 — folio 163.
%  IRREGULARITES DE L'ORIGINAL, conservees : « kalkoza » en italique
%  maigre au milieu d'une enumeration dont tous les autres membres sont
%  en gras ; « kontinue » en romain au milieu d'un exemple gras ; un
%  point la ou la phrase appelle une virgule ; et un guillemet ouvrant
%  jamais referme.
%  Le releveur signale que outils/graisse.py est inexploitable sur cette
%  page : l'ombre de reliure couvre le tiers gauche et l'outil declare
%  gras des mots manifestement romains.
% -------------------------------------------------------------------
\begin{VUpage}[167]{163}
\VUnotes{110.90mm}{%
(1) Ni definas la sufixo \textit{-oz} per : « plena de, richa de », o plu\nl
larje : « \textit{kontenanta, havanta} ». Nu, ca lasta expresuro, la maxim\nl
generala, esas kelkafoye la sola justa. La \textit{nuanco} specala di la\nl
sufixo dependas de la radiko, a qua ol aplikesas, ed esas suficante\nl
determinata per olu. Exemple, esas evidenta, ke parieto \textit{poroza}\nl
havas multa, multega \textit{pori}; ma mem se ol havus nur una, ol esus\nl
ankore poroza. (Ico respondas a la sofisma objecioni di ula kriti\cc
kantachi, qui alegas, ke \textit{kronizar} esas garnisar per (multa) kroni;\nl
or la komuna raciono indikas, ke kande on kronizas rejo, on garnisas\nl
lu per un sola krono!) Ma kande on parolas pri homo, \textit{kurajoza}, ico\nl
signifikas. ke lu esas plena de kurajo, o simple \textit{havas kurajo}. [Ni\nl
adjuntas ke nur la duesma parto dil expliko esas rigoroze justa;\nl
nam, se \textit{kurajoza} signifikus strikte « plena de kurajo », Ido ne\nl
povus dicar, totsame kam nia lingui : \textit{poke kurajoza, ne tre}\nl
\textit{kurajoza}, o kontree : \textit{extreme} kurajoza] simile on darfas dicar\nl
\textit{giboza} equivalas simple : \textit{gibo-havanta}. (La Espisti devus shamar\nl
facar ula objecioni, li qui pro manko dil sufixo \textit{-oz} esas obligata\nl
dicar, en simila kazi : \textit{gibo-hava!}) \textit{Progr.}, VI, 595.%
}

No konfundez \textit{saligar}, \textit{saligo} (kemio) a \textit{salizar}, \textit{salizo}.\nl
L'unesma esas F. \textit{salifier} e la duesma \textit{saler}. Karno esas \textit{sali}\cc
\textit{gebla} ma ne \textit{saligebla} (F. salifiable).

\VUgras{-oz-.} --- Latina ed internaciona (en la \textit{os, ous, eux} (\textit{euse}),\nl
e. c.) ta sufixo quik adoptesis. On soldas lu a nomala radiki\nl
kun la senco : \textit{qua havas}... Ex. : \VUgras{kurajoza} = \textit{qua havas}\nl
\textit{kurajo}; \textit{kalkoza} (tereno, aquo) = \textit{qua havas kalko}; \VUgras{poroza} =\nl
\textit{qua havas pori}; \VUgras{glorioza}, \VUgras{danjeroza} (konduto) = \textit{qua havas}\nl
\textit{glorio, danjero}; \VUgras{vigoroza} = \textit{qua havas vigoro} (1).

Quale omna adjektivi, le formacata per la sufixo \VUgras{-oz} dar\cc
fas esar substantivigata direte : \VUgras{ambiciozo}, \VUgras{kurajozo}, e. c.

La difero inter \VUgras{-anta} e \VUgras{-oza} esas ke l'unesma expresas \textit{ago}\nl
kun ideo tempala, kontre ke la duesma expresas \textit{qualeso} sen\nl
ul ideo tempala. Exemple, \VUgras{la viro amoroza ne sempre esas}\nl
\VUgras{amoranta;} same \VUgras{persono favoroza} (ad ulu od ulo) \VUgras{ne esas}\nl
kontinue \VUgras{favoranta.}

La difero inter \VUgras{-oz} e \VUgras{-iz} esas, ke l'unesma expresas ulo\nl
naturala, e la duesma expresas ulo artificala. Exemple,\nl
tereno \VUgras{sabloza} nature kontenas sablo; ma korto \VUgras{sablizita}\nl
havas sablo nur artifice, per laboro qua kovris lu ye sablo.\parplein

\VUgras{-ul-.} --- Adoptita per la decido 1090, pos tote fundamen\cc
tala studiado dil questiono. Yen la texto di la decido : « On\nl
adoptas la sufixo -ul \textit{por indikar maskula sexuo} (decido,\nl
714 : V, 65). Specala noto dicas : « Komprenende, ico ne
\end{VUpage}

% -------------------------------------------------------------------
%  feuillet 168 — folio 164.
%  COQUILLE DE L'ORIGINAL, conservee : un guillemet ouvrant jamais
%  referme. Le bloc de notes porte un seul appel et trois alineas, dont
%  deux sans numero.
% -------------------------------------------------------------------
\begin{VUpage}[168]{164}
\VUnotes{96.90mm}{%
(1) La sama argumento (alegita olim da S\textsuperscript{ro} Jespersen) qua\nl
valoras por \textit{lu}, kom sengenra pronomo, valoras anke por la sengenra\nl
nomi. Or esus evidente absurda serchar od inventar aparta radiki\nl
por expresar l'ideo komuna a \textit{spozulo} e \textit{spozino}, a \textit{patrulo} e \textit{patrino},\nl
a \textit{fratulo} e \textit{fratino}, e. c. Ta ideo povas expresesar logikale nur per\nl
la radiko ipsa.

Ti qui questionas, kad on devas nun dicar \textit{soldatulo}, \textit{parokulo}, ne\nl
mem atencis ica fakto evidenta, ke \textit{soldato}, \textit{paroko} esas esence mas\cc
kula, per sua propra senco, e konseque nule bezonas la sufixo \textit{-ul}; ol\nl
esus mem nekorekta, \textit{kontre-regula!} Simile ti qui objecionas, ke\nl
patrulo = \textit{père mâle}, siorulo = \textit{monsieur mâle}, facas absurda\nl
rezono (o plu juste mala joko). Nam li oblivias precize la nuna\nl
regulo, segun qua la radiki \textit{patr}, \textit{sior} esas neutra, sengenra; e li\nl
persistas atribuar a li l'anciena senco maskula. Li intermixas\nl
ridinde la du sistemi; e li ne remarkas, ke lia propra rezono kon\cc
damnas (plu juste) l'anciena sistemo, nam ol implikas ica absurdaji :\nl
patrino = \textit{père femelle}, siorino = \textit{monsieur femelle}. (\textit{Prog.}, VI,\nl
588-589.)

Pro ke \textit{sior} esas sengenra, on darfas uzar olu ad o pri \textit{homulo}.\nl
Ma, se en la sama domo lojas \textit{homulo} e \textit{homino} havanta sama nomo,\nl
on agos plu prudente uzar \textit{siorulo} pri la homulo. Cetere mem\nl
sempre, on tote darfas uzar \textit{siorulo} pri homulo, quale on uzas\nl
\textit{siorino} pri homino. (Videz l'apendico « \textit{Genro e maskulismo} ».)%
}

\VUcontinue
exkluzas la formaco di kompozaji quale \textit{porko-maskulo}; ma\nl
\textit{pork-ulo} esas equivalanta e plu kurta. La sufixo -ul ya esas\nl
abreviuro di \textit{(mask)ulo}, quale -in esas abreviuro di \textit{(fem)ino}. »\parplein

La decido pri ul esabis preirata da ica (1089) : « On\nl
decidas definitive, ke la substantivi generale signifikez ne\cc
determinita sexuo (decido, 714 : V, 65). Exempli : \VUgras{avulo,}\nl
\VUgras{patrulo, onklulo, fratulo, kuzulo, siorulo; bovulo, katulo,}\nl
\VUgras{hundulo; finkulo, sturnulo; luciulo, karpulo, tenkulo, per}\cc
\VUgras{kulo; leonulo, panterulo, tigrulo, hienulo,} e. c., e. c. (1).

Komprenende \textit{ulo} ne darfas uzesar sola, quale \textit{ino} uzesas\nl
sola (vice \textit{femino}) dal Esperantisti.

\VUgras{-um-.} --- Ta sufixo, quan Ido konservis de Esperanto,\nl
havas nedeterminita signifiko. Ol esas inter la sufixi simi\cc
lajo di ye inter la prepozicioni. Ol formacas (en mikra\nl
nombro) derivaji, di qui la senco havas relato nedefinita\nl
kun l'ideo expresata dal prima radiko. La ligilo qua restas\nl
inter li helpas kelke la memorado, ma, kande on ne savas\nl
la senco di la derivajo, nur la lexiko povas furnisar lu\nl
certe. Ex. : \VUgras{kolumo} = \textit{parto di kamizo (vesto) qua cirkondas}\nl
\textit{la kolo}; \VUgras{bordumo} = \textit{to quo garnisas bordo} : \VUgras{la tablo havas}\parplein
\end{VUpage}

% -------------------------------------------------------------------
%  feuillet 169 — folio 165.
%  FILET PERDU POUR UNE TOUTE AUTRE RAISON QUE LE SEUIL. Aux feuillets
%  143 et 163 il etait trop pale ; ici il passe deja le test de longueur
%  au seuil ordinaire (186 px a 185), et c'est le TEST DU RESTE qui le
%  rejette : deux saletes a x 1301-1335, tranche du feuillet voisin,
%  donnent reste = 23 pour un seuil de 6. La marge BORD = 25 px posee
%  apres le feuillet 162 ne suffit pas ici. Filet trouve a la main a
%  y = 1791 px, long de 219 px, continu au seuil 220. Ordonnee : 157,40 mm.
%  COQUILLES DE L'ORIGINAL, conservees : « solvor » pour solvar ; et le
%  meme mot compose romain puis italique a la meme ligne.
%  Le releveur etablit que la liste des derives en -uro est ROMAINE et
%  non grasse, contrairement a l'usage du chapitre : meme fut que les
%  mots romains voisins, quand les vedettes de la page sont visiblement
%  plus epaisses. C'est le meme doute que le feuillet 166 avait laisse
%  ouvert, tranche cette fois.
% -------------------------------------------------------------------
\begin{VUpage}[169]{165}
\VUnotes{157.40mm}{%
(1) Ol avantajoze remplasas la \textit{-er-} di Espo : \textit{sablero, hajlero.}%
}

\VUcontinue
\VUgras{kupra bordumo}; \VUgras{foliumar} = \textit{turnar la folii} (di libro) \textit{hal}\cc
\textit{tante tempope por lektar}; \VUgras{formikumar} = \textit{agitar su, movar}\nl
\textit{grandanombre quale formiki}; \VUgras{krucumar} = \textit{dispozar, aranjar}\nl
\textit{krucoforme} : \VUgras{krucumar sua brakii; la voyi interkrucumas}\nl
\VUgras{ye ta punto.}

Per la decido 1556 (\textit{Mondo}, marto 1922, pag. 68) « On\nl
decidas, ke derivaji per -um darfas formacesar nur dal\nl
Akademio » (\textit{Mondo}, X, 54). Dicesas en noto : « La Aka\cc
demio ne indikis til nun \textit{-um}, quale l'altra sufixi \textit{definita},\nl
kom uzebla segunvole. Precize pro ke ol esas nek definita\nl
nek definebla e pro ke lua derivaji konseque povas inter\cc
diferar, ol postulas aparta kompetenteso. --- Komprenende\nl
la supra decido ne interdiktas la propozo di vorti formacita\nl
per \textit{-um}; ma la definitiva decido apartenas al Akademio.

\VUgras{-un-.} --- Ta radiko \textit{darfas uzesar}, dicas la « Grammaire\nl
Complète », p. 71, \textit{por indikar un elemento} : \VUgras{sabluno, gre}\cc
\VUgras{luno} (1). « Kande la elemento havas nula unajo naturala,\nl
on uzas prefere \VUgras{peco} : \VUgras{monetpeco, sukropeco} ». Kompreneble\nl
nulo interdiktas (tote kontree) : \VUgras{aven-grano, hordeo-grano,}\nl
\VUgras{frumento-grano} e tale pri la cetera cereali (\textit{maizo-grano,}\nl
\textit{sekalgrano, rizograno}, e. c.). Vice \VUgras{nivuno} esas preferebla\nl
\VUgras{nivo-floko.}

\VUgras{-ur-.} --- Ta sufixo Latina e nun pasable internaciona quik\nl
adoptesis. Ol soldesas al radiko \textit{por distingar} (kande to esas\nl
necesa) \textit{la produkturo dil ago de la ago ipsa, expresata dal}\nl
\textit{verbo}. Ex. : konstrukturo, pikturo, skulturo, desegnuro, kal\cc
quuro, imprimuro, fotografuro, suturo, solvuro, imituro,\nl
aperturo, fenduro, bavuro, inventuro, skriburo, e. c., la kozo\nl
\textit{obtenita} per konstruktir, piktir, skultir, e. c.

Ni dicis : \textit{kozo obtenita} per la ago, nam produkturo existas\nl
nur \textit{pos obtenir} olu. Pri Ido qua esas produktita (ne pro\cc
duktata o produktota) me darfas dicar, ke ol esas \textit{solvuro}. Ma\nl
pri problemo, di qua me serchas la solvo (quale solvor lu) me\nl
ne uzos la vorto solvuro. Se me uzus \textit{solvuro}, la senco esus\nl
altra; to signifikus, ke la problemo esas nun solvita (ne\nl
solvenda) e ke me serchas la « solvuro » en mea paperi,\nl
exemple.
\end{VUpage}

% -------------------------------------------------------------------
%  feuillet 170 — folio 166.
%  IRREGULARITE DE L'ORIGINAL, conservee : « sendata » compose romain a
%  une ligne et italique deux lignes plus bas.
%  Le releveur signale qu'une rayure diagonale en bord de gouttiere
%  imite un tiret de coupure en fin de ligne ; verification faite, les
%  deux lignes concernees finissent sur un mot entier.
% -------------------------------------------------------------------
\begin{VUpage}[170]{166}
\VUnotes{149.93mm}{%
(1) To, quo preexistis, relate l'edifico konstruktita, esas la\nl
petri, ligno, fero e. c.; fakte li esas la \textit{konstruktajo,} quale la kozi\nl
sendata esas la \textit{sendajo.} Nur pos ke li esos konstruktita en formo di\nl
domo, palaco, e. c., li esas \VUgras{konstrukturo.}%
}

Remarkez ke \VUgras{-uro} povas venar o rezultar mem de verbo\nl
netransitiva. Nam de la fakto, ke la ago di tal verbo ne\nl
atingas direte ula persono o kozo, nule konsequas, ke ta\nl
verbo ne produktas ulo, ke ol ne donas produkturo. Exemple,\nl
la verbo netransitiva « abortar » certe ne atingas ulu od\nl
ulo; la ago « abortar » ne pasas a persono od a kozo, quale\nl
« frapar ». Ma tamen, kad ol ne donas produkturo? Kad ol\nl
ne donas ulo nova, til lore ne existinta : \textit{aborturo?} Same,\nl
on ne dicas « emersar » ulu, o « emersar » ulo. Ma, kad\nl
de la ago « emersar » ne rezultas \textit{emersuro?} Ex. : Multa\nl
insuli esas emersuri ek l'oceano; li produktesis o subite da\nl
tertremo, o lente da madrepori.

Koncerne la difero inter \VUgras{-ajo} e \VUgras{-uro}, la questiono konsistas\nl
en \textit{preexisto} o \textit{ne-preexisto} di la kozo konsiderata. Tam longe\nl
kam \textit{domo}, \textit{procedo}, exemple, ne esas konstruktita, li ne\nl
existas. Li do ne preexistas relate lia konstrukteso od in\cc
venteso, quale la kozi sendata, qui preexistas relate lia sen\cc
deso (1). La \textit{domo}, la \textit{procedo} esas do \VUgras{konstrukturo, inven}\cc
\VUgras{turo;} ma la kozi \textit{sendata,} qui preexistis, esas \VUgras{sendajo.}

Objekto trovita, deskovrita existas ante la ago trovar,\nl
deskovrar olu, do : \VUgras{trovajo, deskovrajo : montrez vua des}\cc
\VUgras{kovrajo.}

Vino ed aquo mixita en un botelo ne esas \textit{mixajo,} ma\nl
\VUgras{mixuro.} Advere la du liquidi, vino ed aquo, preexistis, ma\nl
separite, aparte, e la triesma liquido obtenita ne preexistis :\nl
ol naskis de la mixo, ol esas lua produkturo, do \VUgras{mixuro.}

En mult okazioni la rezulto di la ago ne distingesas reale\nl
del ago ipsa : \VUgras{dekreto, parolo, diskurso, pretendo, abandono,}\nl
\VUgras{aserto, afirmo,} e. c. Lore la sufixo \VUgras{-ur} evidente ne esas\nl
uzenda. Ta sufixo existas nur por la kazi en qui la produk\cc
turo esas evidenta e distingenda, quale en \VUgras{pikturo, skulturo,}\nl
\VUgras{imprimuro, imituro, fotografuro, kopiuro,} e. c. qui certe\nl
esas altro kam \textit{lia modelo} od objekto : \textit{piktajo, skultajo} (per\cc
sono, kozo, e. c.), \textit{imprimajo} (to quo imprimesas), \textit{imitajo,}\nl
\textit{kopiajo} (to quo imitesas, kopiesas), \textit{fotografajo} (to quo foto\cc
\end{VUpage}

% -------------------------------------------------------------------
%  feuillet 171 — folio 167.
%  LIGNE INVISIBLE RETROUVEE : « tradukir. », derniere ligne de la
%  note (1), au fer et tres courte.
%  COQUILLES DE L'ORIGINAL, conservees : un « i » prive de son point ;
%  et une rupture d'italique en plein syntagme dans la note (3).
% -------------------------------------------------------------------
\begin{VUpage}[171]{167}
\VUnotes{122.01mm}{%
(1) \VUgras{Skriburo} esas la produkturo dil ago skribar; \VUgras{skribajo} (la\nl
texto skribata) dicas « Grammaire Complète », p. 55, quale \textit{citajo}\nl
esas texto citata. \textit{Me recevis de ilu tre longa skriburo, qua esas vere}\nl
\textit{interesanta studiuro pri la temo; ma la citaji e la skribajo dil}\nl
\textit{pagino 5 meritas reprochi pro ne kompleta exakteso.} --- \VUgras{Tradukajo}\nl
esas la texto tradukata, \VUgras{tradukuro} la produkturo rezultanta de\nl
tradukir.

(2) Diferas de \textit{kafekrucho, tekrucho} qui kontenas l'infuzuro di\nl
kafeo, di teo.

(3) Jokacheri tre mokis ica lasta. Li objecionis ke \textit{butonagar}\nl
« agar per butono o butoni » ne signifikas necese la Franca verbo\nl
\textit{boutonner.} Ni agnoskas lo tante plu volunte, ke la Franca vorto\nl
\textit{havas du} senci : \textit{haveskar butoni} (parolante pri vejetanto) e \textit{klozar}\nl
\textit{per butoni} (parolante pri vesto). Ica vorto suaparte montras la eroro%
}

\VUcontinue
grafesas) (1) e qui esas anke altro kam la ago : \textit{pikt(ad)o,}\nl
\textit{skult(ad)o, imprim(ad)o, imit(ad)o,} e. c.

Esas notenda, ke la kompozaji (longa e pezoza) \VUgras{-antaj,}\nl
\VUgras{-intaj;} \VUgras{-ataj, -itaj} ne povus remplasar \VUgras{-uro,} pro ke li ne\nl
expresus l'ideo. Exemple, \textit{fendantajo, fendintajo} o \textit{fendatajo,}\nl
\textit{fenditajo} ne equivalas \VUgras{fenduro} = \textit{marko, traso, lineo} pro\cc
duktita da la fendo, rezultanta de olu.

\VUgras{-uy-.} --- Anke konservita de Esperanto, ta sufixo darfas\nl
konsideresar kom abreviuro di \textit{etuyo.} Ol soldesas a radiki\nl
neverbala por formacar substantivi signifikanta : \textit{recipiento,}\nl
\textit{etuyo, buxo,} e. c. : \VUgras{inkuyo} = \textit{mikra vazo, botelo por inko;}\nl
\VUgras{sukruyo} = \textit{vazo por sukro;} \VUgras{supuyo} = \textit{vazo en qua on pozas}\nl
\textit{la supo;} \VUgras{kafeuyo} = \textit{buxo por la kafeo-grani o la kafeo mue}\cc
\textit{lita;} \VUgras{teuyo} = \textit{buxo por la tefolii} (2); \VUgras{tabakuyo} = \textit{mikra buxo}\nl
\textit{por tabakpulvero,} Fr. : \textit{tabatière.}

Dicernez \VUgras{uy} de \VUgras{ier} (Vid. ica) : \VUgras{sigaruyo} = \textit{etuyo por}\nl
\textit{cigari;} \VUgras{sigariero} = \textit{kurta tubo por fumar sigaro.}

Kom radiko ne uzez \VUgras{uyo,} quale l'Esperantisti, ma selektez\nl
inter \textit{etuyo, gaino, buxo, vazo, kesto, recipiento,} e. c.

\VUgras{-yun-.} --- La radiko \VUgras{yun} uzesas kom sufixo por la « \textit{yuni} »\nl
l'infanti dil bestii : \VUgras{bovyuno} = \textit{yuno di la bova speco;} \VUgras{bov}\cc
\VUgras{yunulo} = \textit{yuna bovulo;} \VUgras{bovyunino} = \textit{yuna bovino;} \VUgras{hanyuno,}\nl
\VUgras{porkoyuno, katyuno, leonyuno,} e. c.

Por la homa genituro, quale on ja savas, nur filio, in\cc
fanto, puero esas uzenda, e kun indiko pri la sexuo : filiulo,\nl
filiino; infantulo, infantino; puerulo, puerino.

Ultre \VUgras{yun} la radiki \VUgras{es,} \VUgras{ig} e \VUgras{ag} (3) uzesas kom sufixo.
\end{VUpage}

% -------------------------------------------------------------------
%  feuillet 172 — folio 168.
%  LE DETECTEUR N'A PAS VU LE FOLIO COURANT : la premiere bande de la
%  geometrie est le titre centre « KOMPOZADO. ». Le releveur l'a etabli
%  en comparant la page a ses voisines — ligne 0 a y = 253 et large de
%  252 px, quand les feuillets voisins ont leur folio a y 125-153 et
%  large de 50 a 158 px. Verifie par outils/apparat.py, qui voit bien
%  les deux lignes distinctes (folio a y = 143, titre a y = 253).
%  LIGNE INVISIBLE RETROUVEE : « freno. », derniere ligne du deuxieme
%  alinea de la note, longue de 90 px et tout en italique.
% -------------------------------------------------------------------
\begin{VUpage}[172]{168}
\VUnotes{114.98mm}{%
\VUcontinue
dil autori qui verbigas direte radiki nomala. La F. \textit{boutonner} povus\nl
mem signifikar pluse : \textit{garnisar per butoni}, quale \textit{fleurir} F., havanta\nl
la senco \textit{florizar} apud la senco \textit{floreskar}. Ma, se \textit{garnisar per butoni}\nl
donas en Ido : \textit{butonizar, se haveskar butoni} esas \textit{butoneskar}, quon\nl
do povas signifikar la verbo \textit{butonagar}, se ne \textit{klozar} (vesto) per\nl
\textit{butoni?} Simile la posibla verbi \textit{frenizar, krucizar, martelizar, pedal}\cc
\textit{izar = garnisar per...} lasas al \textit{frenagar, krucagar, martelagar, peda}\cc
\textit{lagar} e. c. nur un senco posibla : \textit{agar per...} E tale la mokeri povas\nl
vidar la neceseso logikala e praktikala dil sufixi \textit{ag} e \textit{iz}.

Fakte en la derivado per la radiki \VUgras{es, ig, yun, ag} (uzata sufixe)\nl
ni esas sur la limito inter la derivado e la kompozado, quale pruvas\nl
la deskompozo : \textit{maladeso = eso di malado} (\textit{o malada}); \textit{bonigar =}\nl
\textit{igar bona} (\textit{o bono}); \textit{bovyuno = yuno di bovo}; \textit{frenagar = agar per}\nl
\textit{freno.}

(1) En la tri unesma klasi la kompozajo komencas per substan\cc
tivo; en la tri lasta per vorto altraspeca.%
}

\VUtitre{10.27pt}{-106}{KOMPOZADO.}
\VUsaut{5.03mm}

En Ido, quale en altra lingui, la kompozado esas l'uniono\nl
di elementi diversa en un sola vorto, qua dicesas kompozajo\nl
e prizentas signifiko komplexa.

Semblas, ke on povas reduktar la Ido-kompozaji al sis\nl
klasi sequanta :

1\textsuperscript{a} Substantivo kun substantivo : \VUgras{fer-voyo, vapor-navo,}\nl
\VUgras{dormo-chambro, skribo-tablo, kap-kuseno,} e. c.

2\textsuperscript{a} Substantivo kun verbo : \VUgras{sabro-frapar, hers-agar, parto}\cc
\VUgras{prenar, tipo-skribar,} e. c.

3\textsuperscript{a} Substantivo kun adjektivo (\textit{o radiko adjektivigita}) :\nl
\VUgras{ciel-blua, graco-plena, lenso-forma, simio-kapa,} e. c. (1).

4\textsuperscript{a} Prepoziciono kun verbo : \VUgras{en-irar, ek-irar, preter-pasar,}\nl
\VUgras{ad-portar, ek-pulsar, apud-pozar,} e. c.

5\textsuperscript{a} Prepoziciono, o vorto nombrala, kun adjektivo o\nl
substantivo : \VUgras{sub-mara, super-natura, inter-naciona, du}\cc
\VUgras{yara, dek-e-du-yara, dua-dek-yara, un-dia, cent-yara, cent}\cc
\VUgras{yaro, inter-tempo, sub-grupo, sub-genero.}

6\textsuperscript{a} Irga adjektivo (qual., demons., nedefin., posedala)\nl
kun radiko igita adjektivo od adverbo : \VUgras{granda-nombra,}\nl
\VUgras{blu-okula, gros-kapa, ta-speca, tala-speca; sam-tempa, sam}\cc
\VUgras{tempe; ca-dia, ca-die; ta-kaza, ta-kaze; plur-foya, plur-foye;}\nl
\VUgras{omna-monata, omna-monate; lua-parte, lia-latere; mea}\cc
\VUgras{nome,} e. c.
\end{VUpage}

% -------------------------------------------------------------------
%  feuillet 173 — folio 169.
%  FILET TROP PALE, panne (a) : il ne devient continu qu'au seuil 210
%  quand l'outil seuille a 185 ; au seuil ordinaire le segment partant
%  de la marge vaut 0 a 1 px et le test de LONGUEUR echoue avant meme
%  celui du reste. Le releveur distingue donc nettement cette panne de
%  celle du feuillet 169. Filet a y = 1303 px, long de 210 px.
%  DEUX LIGNES INVISIBLES RETROUVEES, absentes de la geometrie :
%  « horlojo. » (corps) et « di vokalo. » (derniere ligne des notes).
%  COQUILLE DE L'ORIGINAL, conservee : « l'enfonio » pour l'eufonio, le
%  mot etant correctement compose deux lignes plus bas.
% -------------------------------------------------------------------
\begin{VUpage}[173]{169}
\VUnotes{115.49mm}{%
(1) Certe plu fluanta, plu dolca kam : \textit{postkarto, postmarko,}\nl
\textit{mondlinguo}. Ni ne obliviez, ke ula sucedi di konsonanti, sen la\nl
« bufro » di vokalo, esas facile pronuncebla da nordani, ma preske\nl
ne pronuncebla da sudani.

(2) La kompozaji esas abreviuri. To explikas e justifikas, ke on\nl
darfas uzar kom unesma elemento nur simpla radiko, e ne vorto\nl
kompleta, se nur la klareso o l'enfonio ne sufras de lo.

Nature Ido povas havar kompozado preske tam flexebla e richa\nl
kam la Greka. Ma importas, ke ta kompozado satisfacez justa\nl
expreso dil idei e posedez klara komprenebleso, kun pronunco facila\nl
ed agreabla eufonio. On do evitez tro longa kompozaji ed akumulo\nl
de konsonanti, o ruda vicinaji. Ex. : \textit{blankabrakia, larjafronta, kur}\cc
\textit{vabeka, blankadenta, laktosuganta, arjentoklova} esas nedubeble\nl
preferinda kam : \textit{blankbrakia, larjfronta, kurvbeka, blankdenta,}\nl
\textit{laktsuganta}, en qui intershokas tro multa konsonanti sen la bufro\nl
di vokalo.%
}

Se l'eufonio postulas o nur konsilas lo, on lasas en\nl
l'unesma elemento la vokalo dezinenca, quan lu havas kom\nl
izolita vorto : \VUgras{granda-nombra, longa-hara, omna-monate,}\nl
\VUgras{lenso-forma, dormo-chambro, skribo-tablo, posto-marko,}\nl
\VUgras{posto-karto, mondo-linguo} (1).

Ma, se l'eufonio permisas lo, on omisas la vokalo dezi\cc
nenca en l'unesma elemento : \VUgras{ciel-blua, samtempe, plurfoye,}\nl
\VUgras{aernavo, fervoyo,} e. c. (2).

Ordinare on darfas omisar la streketo, quan ni pozis inter\nl
la parti dil vorto kompozita. Ma on uzas lu quik kande on\nl
timas, ke la dicerno di la elementi ne facesos sat facile o\nl
sat rapide : \VUgras{mar-aquo, filtr-etamino, kredit-institucuro, posh}\cc
\VUgras{horlojo.}

Cetere on sempre darfas ne uzar la kompozo e dicar :\nl
\VUgras{aquo marala} o \VUgras{di maro;} \VUgras{institucuro kreditala} o \VUgras{di kredito;}\nl
\VUgras{etamino por filtrar; postal marko, postal karto,} o \VUgras{marko}\nl
\VUgras{postala, karto postala, mondal linguo} o \VUgras{linguo mondala,} e. c.\parplein

\VUsaut{6.90mm}
\VUtitre{10.27pt}{-45}{REGULO DI ANALIZO O DESKOMPOZO.}
\VUsaut{4.27mm}

Se, pos lekto atencoza, vu ne komprenas bone la kom\cc
pozajo, procedez tale :

\textit{Deskompozez lu}, se ol apartenas al 1, 2, 3 e 4esma klasi,\nl
\textit{irante de lua lasta elemento al unesma} : \VUgras{fer-voyo} = \textit{voyo} (ek)\parplein
\end{VUpage}

% -------------------------------------------------------------------
%  feuillet 174 — folio 170.
%  IRREGULARITES DE L'ORIGINAL, conservees : la glose de « simio-kapa »
%  est en romain quand celle de « lenso-forma », ligne au-dessus et de
%  meme fonction, est en italique ; le meme compose est en gras au corps
%  et en italique dans la note ; deux mots francais cites restent en
%  romain ; et deux guillemets ouvrants ne sont jamais refermes.
% -------------------------------------------------------------------
\begin{VUpage}[174]{170}
\VUnotes{117.24mm}{%
(1) Komparez \textit{lenso-forma} a F. pisciforme, lentiforme, e. c.

(2) \textit{Cent yari} = \textit{yarcento}, quale \textit{dek yari} = yardeko.

« \textit{Centyaro} ne povas esar sinonima di \textit{yarcento}; ol esas la sub\cc
stantivo di \textit{centyara} = qua havas centyari e konseque tradukas\nl
F. \textit{centenaire} (sive homo, sive altra ento, ex. querko). Se on volas\nl
indikar ne ambigue qua speco di ento, on devas simple nomar ol :\nl
\textit{centyara homo, centyara querko}. Evidente (remarko generala e tre\nl
importanta) on ne povas postular, ke kompozita vorto expresez\nl
omno; e se kompozajo ne suficas por expresar la penso, \textit{on ne uzez ol}.\nl
Nultempe on esas obligata uzar kompozita vorto! (\textit{Progr.}, VII, 496.)\nl
Ma sempre la helpolinguo \textit{devas} expresar klare ed exakte la penso.\nl
Se ol ne povus lo, ol esus refuzenda e forjetenda, nam ol adportus\nl
al mondo, ne solvuro, ma vera mistifiko.

(3) « On devas dicar \textit{omnadia, ne omnadiala}; nam quo esas \textit{omna}\cc
\textit{dio}? Nula dio! (\textit{Progr.}, VII, 497.)%
}

\VUcontinue
\textit{fero}; \VUgras{vapor-navo} = \textit{navo} (movata per) \textit{vaporo}; \VUgras{dormo-cham}\cc
\VUgras{bro} = \textit{chambro} (por) \textit{dormo}; \VUgras{skribo-tablo} = \textit{tablo} (por)\nl
\textit{skribo}; \VUgras{kap-kuseno} = \textit{kuseno} (por la) \textit{kapo}; \VUgras{sabro-frapar} =\nl
\textit{frapar} (per) \textit{sabro}; \VUgras{hers-agar} = \textit{agar} (per) \textit{herso}; \VUgras{parto}\cc
\VUgras{prenar} = \textit{prenar parto} (en); \VUgras{tipo-skribar} = \textit{skribar} (per)\nl
\textit{tipo}; --- \VUgras{ciel-blua} = \textit{blua} (quale) la \textit{cielo}; \VUgras{graco-plena} =\nl
\textit{plena} (de la) \textit{graco}; \VUgras{lenso-forma} = \textit{qua havas la formo di}\nl
\textit{lenso}, \VUgras{simio-kapa} = qua havas kapo di simio (1); --- \VUgras{en}\cc
\VUgras{irar} = \textit{irar en}, \VUgras{ek-irar} = \textit{irar ek}, \VUgras{preter-pasar} = \textit{pasar preter},\nl
\VUgras{ad-portar} = \textit{portar ad}, \VUgras{ek-pulsar} = \textit{pulsar ek}.

La motivo dil supera regulo e di tala iro, por la deskom\cc
pozo di ta 4 klasi de kompozaji esas ke, en li, la elemento\nl
determinata esas sempre la lasta, ed olua senco determinesas\nl
dal unesma elemento. En \VUgras{dormo-chambro}, \textit{chambro} deter\cc
minesas da \textit{dormo}; en \VUgras{simio-kapa}, \textit{kapa} determinesas da\nl
\textit{simio}.

Ma por la 5 e 6esma klasi de kompozaji, sequez l'ordino di\nl
la elementi quan vu trovas lektante li. Cetere maxim ofte ta\nl
ordino korespondas ad olta quan li havas en multi de nia\nl
lingui, e pro to on komprenas li ordinare per simpla lekto.\nl
\VUgras{Ex.} : \VUgras{sub-mara} = \textit{qua esas sub maro}, \VUgras{super-natura} = \textit{qua}\nl
\textit{esas super la naturo}, \VUgras{inter-naciona} = \textit{qua esas inter nacioni};\nl
\VUgras{un-dia} = \textit{qua evas} (od evis) \textit{un dio}, \VUgras{cent-yara} = \textit{qua evas}\nl
(od evis) \textit{cent yari} (2); \VUgras{ca-dia} = \textit{di cadio}, \VUgras{omnadia} (3) =\nl
\textit{di omna dio}, \VUgras{taspeca} = \textit{di ta speco}, \VUgras{talaspeca} = \textit{di tala}\nl
\end{VUpage}

% -------------------------------------------------------------------
%  feuillet 175 — folio 171.
%  FILET TROP PALE, panne (a) : au seuil ordinaire aucun pixel sombre
%  dans la fenetre de depart, le test de LONGUEUR echoue avant celui du
%  reste. Continu au seuil 210. Filet a y = 792 px, long de 209 px.
%  COQUILLES DE L'ORIGINAL, conservees : « direte » pour direkte ; un
%  double point final ; « groskapa » sans trait d'union au corps quand
%  la note ecrit « gros-kapa » ; et un point parasite dans « (5. e
%  6esma klasi) ».
%  Le releveur note que les composes cites sont ici en ROMAIN alors que
%  les memes etaient en gras au folio 170.
% -------------------------------------------------------------------
\begin{VUpage}[175]{171}
\VUnotes{74.17mm}{%
(1) « Existas nula nekonsequo inter \textit{samtempa} e \textit{samtempano} :\nl
du eventi esas \textit{samtempa}, ne \textit{samtempana}, pro ke li eventas \textit{sam}\cc
\textit{tempe} (en la sama tempo). --- On devas analizar direte la kompozita\nl
vorti per la simpla komuna raciono, e ne serchar reguli subtila e\nl
neutila, qui nur komplikus ed obskurigus kozi simpla e klara per\nl
su. » (\textit{Progreso}, VII, 497.)

« On ne povas dicar, ke la derivado en Ido dependas de la \textit{speco}\nl
(gramatikala) di la radiko, nam kontree, ni emendis ica defekto di\nl
Esp. : la derivado dependas nur de la \textit{senco} di la radiko. Omna\nl
subtila reguli propozita por ico semblas tote neutila. Esas anke\nl
nejusta asertar, ke por komprenar vorto derivita de « radiko kompo\cc
zita » on devas serchar la senco di ca « radiko » : nam ofte ca\nl
« radiko » havas nula senco. Exemple \textit{samtempano} havas tote klara\nl
senco, dum ke \textit{samtemp} havas nula senco : nam quo esas \textit{samtempo?}\nl
Nula speco di \textit{tempo!} (\textit{Progreso}, VII, 497.)

(2) \textit{Blu-okula, long-hara, gros-kapa} e. c. memorigas tre justa\nl
remarkigo en \textit{Progreso}, VII, p. 496 : « En tala kompozaji on ne\nl
povas parolar pri elemento precipua o determinata : la du elementi\nl
esas ne subordinita, ma koordinita, ed on povus preske permutar\nl
li kelkafoye, ex. \textit{okul-blua}. Lia senco esas tote klara sen irga sufixo.\nl
Vane on objecionas, ke on povas substantivigar li : \textit{bluokula}, e. c.\nl
esas esence adjektiva. Duesme, mem se on substantivigas li, eventos\nl
nula bisenceso, se on sorgas egardar ica regulo di komuna raciono :\nl
\textit{Nultempe formacez kompozajo por expresar ideo qua povas esar}\nl
\textit{expresata tam kurte per la simpla elementi.} D\textsuperscript{ro} Talmey bone\nl
montris, ke on ne darfas dicar \textit{redvango}, vice \textit{reda vango}, \textit{grandurbo}\nl
\textit{vice granda urbo} (kulpo quan facas ofte nia Germana samideani). Do\nl
on povas sen ula detrimento dicar \textit{Red-pelo} por F. \textit{Peau-Rouge}, nam\nl
\textit{peau rouge} dicesus \textit{reda pelo.} »

Pose S\textsuperscript{ro} \textsc{Couturat} dicas pri \textit{triangulo, trimasto, centyaro, mil}\cc
\textit{pedo} : nu, anke ici ne esas bisenca..%
}

\VUcontinue
\textit{speco}, \VUgras{samtempa} = \textit{di sama tempo} (1), \VUgras{granda-nombra} =\nl
\textit{di granda nombro, qua prizentas granda nombro}, \VUgras{blu-okula}\nl
= \textit{qua havas blua okuli}, \VUgras{groskapa} = \textit{qua havas grosa kapo} (2).\parplein

Praktike memorez, ke on deskompozas komencante per la\nl
fino, se la kompozajo unionas \textit{substantivo} ad altro (1, 2,\nl
3esma klasi) o \textit{prepoziciono} a verbo (4esma klaso). Altrakaze\nl
on sequas l'ordino di la elementi (5. e 6esma klasi).

La relato inter la elementi dil kompozajo esas tacata :\nl
\textit{quale} en ciel-blua, \textit{por} en dormo-chambro, \textit{per} en sabro-fra\cc
par; ma, quale on vidis supere, la spirito trovas lu facile.\parplein

En ula kazi, la relato esas la liganta konjunciono \textit{e}. Ex. :\nl
\VUgras{sekretario-kasisto} = sekretario \textit{e} kasisto; \VUgras{kapitano-instruk}\cc
\VUgras{tero} = kapitano \textit{e} instruktero.
\end{VUpage}

% -------------------------------------------------------------------
%  feuillet 176 — folio 172.
%  LIGNE INVISIBLE RETROUVEE : « arma. », en italique, au fer, large de
%  155 px, qui achevait le mot coupe de la ligne precedente.
%  Le second signalement, « deux lectures » a la ligne 2, est tranche :
%  c'est un BLANC, celui qui suit le titre centre — la zone est vierge
%  a l'agrandissement.
%  Le releveur signale une neuvieme confirmation de la meme regle : une
%  liste qui paraissait grasse sur la bande est romaine a fort
%  grossissement.
% -------------------------------------------------------------------
\begin{VUpage}[176]{172}
\VUnotes{154.60mm}{%
(1) Videz la 4\textsuperscript{a} klaso dil kompozaji p. 168.%
}

\VUtitre{10.27pt}{-56}{KOMPOZADO PER PREPOZICIONI.}
\VUsaut{5.79mm}

La kompozado per prepozicioni formacas la transito inter\nl
la derivado e la kompozado dil radiki; ma ol similesas plu\nl
multe la derivado (pro l'analogeso inter la afixi e la par\cc
tikuli). Pro to lu ne esas submisata al generala regulo di\nl
kompozado qua postulas, ke la vorto determinanta esez avan\nl
la vorto determinata. Konseque, kande prepoziciono esas\nl
unionita a verbal radiko, l'analizo o deskompozo generale\nl
facesas per inversigar l'ordino di la elementi, to esas komen\cc
cante de la radiko verbala (1). Ex. : \VUgras{en-irar}, \textit{irar en}; \VUgras{ek-irar},\nl
\textit{irar ek}; \VUgras{ad-portar}, \textit{portar ad}; \VUgras{kontre-dicar}, \textit{dicar kontre};\nl
\VUgras{inter-mixar}, \textit{mixar inter} (li); \VUgras{inter-veno}, \textit{veno inter}; \VUgras{ek-pul}\cc
\VUgras{sar}, \textit{pulsar ek}; \VUgras{apud-pozar} = \textit{pozar apud.}

Ma, kande la prepoziciono esas unionita a nomal radiko,\nl
on devas generale komencar la deskompozo de la prepozi\cc
ciono ipsa; altravorte : on sequas l'ordino di la elementi.\nl
Ex. : \VUgras{sen-hara} = \textit{sen hari}; \VUgras{sur-tera} = \textit{qua esas sur tero};\nl
\VUgras{sub-mara} = \textit{qua esas sub maro}; \VUgras{super-natura} = \textit{qua esas}\nl
\textit{super la naturo}; \VUgras{inter-naciona} = \textit{qua esas inter nacioni};\nl
\VUgras{inter-popula}, \VUgras{inter-homa} = \textit{qua esas inter la populi, la homi};\nl
\VUgras{inter-yuna} = \textit{qua esas inter yuni}, e. c., e. c.

Konseque on darfas dicar : intertempo, intertempe; inter\cc
akto; subtaso, subvesto, subjupo, e. c.

Ta duopla regulo posibligas interpretar korekte la verbi\nl
kompozita per \VUgras{-igar} e prepoziciono; on devas komencar de\nl
la verbo \VUgras{igar} : \VUgras{enter-igar} devas analizesar \VUgras{igar entera}, e ne\nl
\VUgras{en-terigar}, qua havus nula senco; \VUgras{sen-arm-igar} = \textit{igar sen}\cc
\textit{arma.}

On atencez, ke la kompozado per prepozicioni ne chanjas\nl
la direta komplemento dil verbo (kande ica esas transitiva) :\nl
\VUgras{en limpida aquo, on travidas la stoneti dil fundo; on ne}\nl
\VUgras{travidas la aquo ipsa.} Konseque l'adjektivo \VUgras{travidebla} povas\nl
aplikesar nur a la stoneti, e ne a la aquo ipsa; ol ne povas\nl
signifikar \VUgras{diafana}, ma nur (segun sua elementi) \textit{videbla tra.}\nl
Simile on \VUgras{supervarsas aquo sur flori}; quo esas \VUgras{supervarsita}?\parplein
\end{VUpage}

% -------------------------------------------------------------------
%  feuillet 177 — folio 173.
%  FILET TROP PALE, panne (a) : au seuil ordinaire la fenetre de depart
%  ne contient aucun pixel sombre, le test de LONGUEUR rend n = 0 et le
%  filet est ecarte avant le test du reste — que le releveur a d'ailleurs
%  verifie : il valait 0, il n'y a pas de salete a droite. Filet a
%  y = 1177 px, long de 207 px, continu au seuil 210.
%  DEUX LIGNES INVISIBLES RETROUVEES, toutes deux dans les notes et
%  toutes deux tres courtes : « p. 132. » (85 px) et « tiko. » (53 px).
%  COQUILLES DE L'ORIGINAL, conservees : « p.497 » sans espace apres le
%  point abreviatif, quand la meme page ecrit « p. 132. » et « p. 168 » ;
%  et une liste composee en romain quand la liste strictement parallele
%  de la ligne precedente est en gras.
% -------------------------------------------------------------------
\begin{VUpage}[177]{173}
\VUnotes{105.24mm}{%
(1) Videz l'expliko en la noto (1), ye « Prepozicioni prefixa »,\nl
p. 132.

(2) Okazione di to ni dicez, ke la helpolinguo devas ne imitar la\nl
verbi Latina kun du akuzativi : \textit{doceo pueros grammaticam}. On\nl
dicas : \textit{docar la gramatiko a pueri} o \textit{instruktar pueri pri la grama}\cc
\textit{tiko}.

(3) Okazione di ta vorti ni citez \textit{Progreso}, VII, p.497 : « On ne\nl
darfas dicar, ke \textit{unkorniero} esas ento karakterizata per la qualeso\nl
\textit{unkorna} : ico apartenas a la logiko Esperantala, e duktus ad ico :\nl
ento bela = \textit{beliero!} Se \textit{unkorna} esas suficanta, \textit{unkorno} anke esas\nl
tala. E \textit{tripedo, trimasto} anke esas bona e suficanta. Ma on adoptis\nl
\textit{unikorno}, pro ke ol esas specala animalo, ne irga animalo havanta\nl
un korno (ex. rinocero); \textit{uniforma}, pro ke ca vorto nule esas sino\cc
nima di \textit{un-forma}; \textit{centifolio} nule signifikas : qua havas cent folii,\nl
ma speco di planto; e simile on ne dicas \textit{milpedo}, ma \textit{miriapodo}.\nl
\textit{Bipunto} esas ensemblo di \textit{du} punti superpozita, do plu preciza kam\nl
\textit{du-punto} o \textit{duopla punto}. Reale omna \textit{reala} bisencesi esas evitita.

Cetere, quale ni ja dicis, la prefixi \textit{mono-, bi-, quadri-}... esas tre\nl
utila, nam li signifikas precize : \textit{havanta un, du, quar}... On devas\nl
dicar ne \textit{du-pedo, quar-pedo}, ma \textit{bi-pedo, quadri-pedo}. »%
}

\VUcontinue
La aquo, e ne la flori. Do on ne darfas dicar (quale Espo) :\nl
\VUgras{vango supervarsita da lakrimi}, ma juste : \VUgras{balnita, inundita}\nl
\VUgras{da lakrimi}.

Kontraste, nul detrimento en ke prepoziciono-prefixa igas\nl
transitiva netransitiva verbo : \VUgras{enirar, ekirar la chambro} (1).\nl
Konseque ula verbi netransitiva (quale \textit{irar}) darfas en la\nl
kompozado divenar transitiva e posedar pasivo. Se on darfas\nl
dicar : \VUgras{trairar la rivero}, la rivero ipsa darfas dicesar \VUgras{trairata}.\nl
Ma de to on ne darfas konkluzar, ke en la kompozado, verbo\nl
transitiva darfas analoge recevar duesma komplemento\nl
direta, nam la posedo di du komplementi direta esas kontrea\nl
a la logiko ed al klareso (2).

\VUsaut{6.05mm}
\VUtitre{10.27pt}{-17}{LA KOMPOZAJI E LA SUFIXI.}
\VUsaut{4.28mm}

En la kompozaji dil klasi 5 e 6-ma (vid. p. 168) on \textit{darfas}\nl
omisar, se nur la senco restas klara, afixi quin on expresus\nl
en derivaji. Mem en ula kazi on \textit{devas} ne expresar li. Exemple,\nl
on dicas \VUgras{tadia, cadia, omnadia, singladia} malgre la derivajo\nl
\textit{diala}. Same : internaciona, dukapa, trifolia (3), longhara,\nl
malgre la derivaji \textit{nacionala, kapala, folioza, haroza}. Pro quo
\end{VUpage}

% -------------------------------------------------------------------
%  feuillet 178 — folio 174.
%  COQUILLES DE L'ORIGINAL, conservees : deux guillemets ouvrants
%  jamais refermes ; et les initiales de l'auteur en capitales romaines
%  ordinaires quand les autres noms d'auteurs sont en petites capitales.
%  Le releveur signale deux fois le meme piege, tranche a
%  l'agrandissement : des mots qui paraissent gras sur la bande sont
%  romains, et le gras s'interrompt a l'interieur des parentheses.
% -------------------------------------------------------------------
\begin{VUpage}[178]{174}
\VUnotes{138.58mm}{%
(1) « On ne vidas quon povus signifikar la substantivi neutra\nl
\textit{longharo, senfingro}; e la substantivo \textit{internaciono} povus signifikar :\nl
naciono inter (la ceteri), do esas multe plu klara dicar \textit{interna}\cc
\textit{cionajo} = ulo internaciona. Do existas nula praktikala bisenceso, se\nl
on ne fabrikas artificala exempli. Se \textit{senfingri} povas dicernesar de\nl
\textit{sen fingri, senfingro} (homo) povas tam bone dicernesar de \textit{sen}\nl
\textit{fingro}. E, quale ni dicis (p. 217-8), on darfas uzar por plua klareso :\nl
\textit{senfingra homo} (\textit{Progreso}, VII, 495, noto).%
}

\VUcontinue
on adjuntus a ta derivaji elementi neutila? Nulo ya mankas\nl
en \textit{internaciona, dukapa, trifolia, longhara} (1).

Ma kontraste, on \textit{devas} ne omisar la sufixo, se lua manko\nl
povus efektigar miskompreno o se vera ambigueso rezultus\nl
praktike de ta omiso. Exemple, se me dicus : \VUgras{esas agreabla}\nl
\VUgras{vivar kun bonhumoro e desagreabla vivar kun malhumoro,}\nl
preske omni komprenus : \textit{esas agreabla vivar kun bona}\nl
\textit{humoro e desagreabla vivar kun mala humoro}. Nur pos re\cc
flekto on divinus, ke me intencis dicar : \VUgras{esas agreabla vivar}\nl
\VUgras{kun bonhumorozo} (homo bonhumora) \VUgras{e desagreabla vivar}\nl
\VUgras{kun malhumorozo} (homo malhumora). Do uzez \textit{bonhumo}\cc
\textit{rozo} e \textit{malhumorozo} por preventar omna miskompreno.

Pri \textit{sejorno} qua duris dum non dii me dicos : mea \textit{nondia}\nl
sejorno, pro ke vere ta sejorno esas \textit{nondia}. Ma pri soldati\nl
qui restis kaptita dum non dii me ne dicos : la \textit{nondia} kaptiti,\nl
ma : la \textit{nondiala} kaptiti, nam certe li ne esas \textit{nondia} quale\nl
mea sejorno.

Same pri la muskuli dil \textit{avanbrakio}, me ne dicos : la mus\cc
kuli \textit{avanbrakia}, nam li ne esas « avanbrakio », li nur kon\cc
cernas olu. Me do dicos : la muskuli \textit{avanbrakiala}.

En \textit{Progreso}, VII, p. 217, lektesas ica tre utila remarko\nl
pri la temo : « Se l'adjektivi kompozita divenas dusenca\nl
per substantivigo, la kauzo ne esas sempre la manko di ula\nl
elemento. Exemple, l'adjektivo \textit{sen-hara} esas kompleta : nula\nl
elemento mankas : \textit{homo sen-hara} esas \textit{homo sen haro}. Ni\nl
povas do substantivigar ol, e dicar : \textit{sen-haro} (tre prefere :\nl
\textit{kalvo} de \textit{kalva} L. B.); ne nur \textit{senharozo} o \textit{senhariera} ne esas\nl
necesa, ma li esus nelogikal e nejusta. E tamen \textit{sen-haro, sen}\cc
\textit{hari} povas esar dusenca, komprenesar kom \textit{sen haro, sen hari}.\nl
Same \textit{sen-laboro, sen-labori} ed omna simila kompozaji (tre\nl
multa!)
\end{VUpage}

% -------------------------------------------------------------------
%  feuillet 179 — folio 175.
%  DEUX LIGNES INVISIBLES RETROUVEES, toutes deux dans les notes :
%  « e. c. (L. B.) » et « kalvini. ».
%  LA SECONDE VALIDE UNE LECON DE L'OUTIL : le signalement « encre sous
%  la derniere ligne » s'etait revele trois fois de la rousseur ; ici
%  c'est un mot entier. Le releveur a cadre a 47 px a gauche de la marge
%  et l'a lu sans peine — c'est exactement le cadrage que l'assembleur
%  avait manque au feuillet 137, ou il avait pris un mot pour un jambage.
%  PARTICULARITE DE L'ORIGINAL, conservee : le sigle « DFIS » en
%  capitales romaines pleines, quand les autres titres cites sont en
%  italique.
% -------------------------------------------------------------------
\begin{VUpage}[179]{175}
\VUnotes{115.57mm}{%
(1) On rilektez pri ca temo la tre saja artiklo : \textit{Kompozo per}\nl
\textit{prepozicioni}, da S\textsuperscript{ro} \textsc{Jespersen} (III, 402). Komparez l'artiklo da\nl
S\textsuperscript{ro} \textsc{Pivet} pri la sama temo (I, 403).

(2) Me adjuntas : Or ne semblas esar la skopo di L. I. fabrikar\nl
vorti « stranja e poke rikonocebla ». Ol destinesas ad altro kam ad\nl
esar ludilo por kelki, o materio por subtilaji ne atingebla facile\nl
dal komuna raciono. L. B.

(3) Do kompozaji tala kam : \textit{ir-ekar, ir-transar, ir-traar, ir-enar,}\nl
\textit{ir-preterar; dic-kontrear, port-adar} (ad-portar), \textit{jet-forar} (for-jetar)\nl
e. c. (L. B.)

(4) Quankam, segun III-201, \textit{kalveso} ne esas kompleta senhareso\nl
(tre vere), tamen en la vivo ordinara \textit{kalva, kalvo, kalveso, kalvulo,}\nl
\textit{kalvino} solvas sat bone la questiono, e juste pro ico me propozigis\nl
\textit{kalv} DFIS, qua adoptesis. Ex. : \textit{la kalvuli esas plu multa kam la}\nl
\textit{kalvini.}%
}

Se on volus evitar omna posibleso di dusenceso en tala\nl
kazi, oportus, sive uzar specala sufixo equivalanta a sen,\nl
quale D. \textit{-los}, E. \textit{-less}, sive pozar la prepoziciono \textit{sen} dop la\nl
substantivo, segun la propozo di S\textsuperscript{ro} \textsc{de Janko} (II, 152) (1) :\nl
\textit{har-sena}, genitus \textit{har-seno, har-seni}, qui ne esus ambigua;\nl
ma li esus stranja e poke rikonocebla (2). Pluse on devus\nl
aplikar (logikale) la sama regulo ad omna prepozicioni (3),\nl
e dicar : \textit{mar-suba, ter-ena, nacion-intera!} Ma se on ne volas\nl
irar til ta extrema konsequanti (or nulu aprobis li o mem\nl
semblas pronta aceptar li), on devas aceptar la nuna sistemo,\nl
qua esas suficante klara e logikala, e qua genitas nula dusen\cc
ceso \textit{en la praktiko.} »

En noto S\textsuperscript{ro} \textsc{Couturat} adjuntis :

« Existas moyeno tre simpla por evitar la dusenceso di \textit{sen}\cc
\textit{haro, senhari}, se ol semblos ulfoye timinda : nome, uzar la\nl
vorto \textit{homo}, e dicar : \textit{senhara homo} (o \textit{senharhomo}), \textit{senhara}\nl
\textit{homi} (o \textit{senharhomi}). Pluse, se nula sufixo esas necesa\nl
por indikar \textit{homo} (sengenra), ni havas \textit{-ul} e la sufixo \textit{-in} por\nl
indikar la homuli e homini rispektive. On povus do dicar tre\nl
bone, exemple : « La \textit{senharuli} esas plu multa kam la \textit{sen}\cc
\textit{harini} » (ma esas plu bona uzar \textit{kalvo}, do \textit{kalvuli, kalvini}) (4).\nl
On memorez bone, ke la derivaji e kompozaji esas generale\nl
\textit{abreviuri} : on nule \textit{devas} uzar li kande li esas neklara o\nl
dusenca; en ta kazi on \textit{devas} uzar la expresuri kompleta.\nl
Exemple, se \textit{petrolnavo} divenas praktike dusenca, on devas
\end{VUpage}

% -------------------------------------------------------------------
%  feuillet 180 — folio 176. FIN DE LA DEUXIEME PARTIE.
%  TROISIEME CAUSE DISTINCTE POUR « FILET NON TROUVE ». Ce n'est ni le
%  seuil (143, 163, 173, 175, 177) ni la salete a droite (169) : ici le
%  filet est a y = 377 px, au-DESSUS de la fenetre de recherche de
%  filet.py, qui part de 0,25 H = 520. La page est presque vide et son
%  bloc de notes remonte tres haut. Le filet passe les deux tests ; il
%  n'est jamais examine. Consigne au LISEZ-MOI.
%  Le releveur corrige aussi la lecture de la page : la ligne 4 n'est
%  pas une ligne de corps renfoncee mais LA NOTE elle-meme.
%  Le grand blanc qui precede « FINO. » est VIDE — pas d'ornement,
%  contrairement a la fin de la premiere partie. Verifie au profil
%  d'encre a quatre seuils et a l'agrandissement.
%  PARTICULARITE DE L'ORIGINAL, conservee : pas de point apres « ideo »
%  devant les initiales de l'auteur (blanc net de 18 px, lu au pixel).
% -------------------------------------------------------------------
\begin{VUpage}[180]{176}
\VUnotes{37.85mm}{%
(1) \textit{Progreso}, VII, p. 218, noto.%
}
\VUcontinue
dicar, sive \textit{navo portanta petrolo}, sive \textit{navo movata per petrolo}.\nl
Simile, se \textit{senhari} esas dusenca, on devas uzar \textit{senhara}\nl
\textit{homi} » (1). (O \textit{kalvi}, se ta vorto suficas por l'ideo L. B.)

\VUcentreA{89.24mm}{10.63pt}{-23}{FINO.}
\end{VUpage}

% -------------------------------------------------------------------
%  feuillets 181 et 182 — folios 177 et 178 : BLANCS.
%  Verifies a l'agrandissement par l'assembleur : les quelques bandes
%  que le detecteur y releve sont des artefacts du pli et du bord de
%  page, aucune n'est du texte.
% -------------------------------------------------------------------
\begin{VUpage}[181]{}
\end{VUpage}

\begin{VUpage}[182]{}
\end{VUpage}

% -------------------------------------------------------------------
%  feuillet 183 — folio 179. OUVERTURE DES APPENDICES.
%  La page ne porte PAS de folio courant : la premiere bande est le
%  titre. Elle porte DEUX lignes centrees — un titre en capitales
%  demi-grasses et, au-dessous, un sous-titre en bas de casse italique
%  entre parentheses, d'un corps plus petit que le texte courant
%  (hauteur de capitale 17 px contre 21 pour le corps).
%  AUCUNE NOTE : le releveur a verifie le profil d'encre aux seuils 185,
%  210 et 220 sous la derniere ligne — pas de filet. Les pages
%  d'appendice n'en portent pas.
%  LIGNE INVISIBLE RETROUVEE : « ad li. » », fin du premier alinea.
%  Le second signalement, « deux lectures » a la ligne 2, est tranche :
%  c'est le blanc qui suit le bloc de titre, la zone est vierge.
% -------------------------------------------------------------------
\begin{VUpage}[183]{}
\VUtitre{10.27pt}{-27}{L'ACENTIZO EN IDO.}
\VUsaut{1.50mm}
\VUsoustitre{6.02pt}{227}{\textit{(Apendico 1ma.)}}
\VUsaut{4.54mm}

Quale dicesas tre juste en \textit{Progreso}, I, 383, « La questiono\nl
di l'acento esas un ek le maxim desfacila, quan on devas\nl
solvar en L. I. artificala, se on volas konciliar maxim bone\nl
posible la regulozeso e la naturaleso. Unlatere, on devas\nl
havar regulo generala e simpla pri acentizo, e ne, quale en\nl
omna lingui, fidar a l'uzado kom sola regulo o guidilo.\nl
Altralatere, on devas respektar maxim ofte posible la « na\cc
tural » acentizo di la vivanta lingui, por ne ofensar la fone\cc
tikal kustumi di lia adepti, ed igar la L. I. leda e desaminda\nl
ad li. »

En I, 479, on remarkigis, ke « l'ortografio di l'interna\cc
ciona vorti restas generale la sama en plura lingui, dum\nl
ke l'acentizo multafoye varias : ex. I. \textit{ortografîa}, R. \textit{ortho}\cc
\textit{grâphiya}, E. \textit{orthôgraphy}; Hungara : \textit{ôrtografia}.

D. \textit{Agonîe}; I. \textit{agonîa}; R. \textit{agôniya}; E. \textit{âgony}.

D. \textit{Geologîe}; I. \textit{geologîa}; R. \textit{geolôgiya}; E. \textit{geôlogy}; Hun\cc
gara : \textit{gêologia}.

Pos altra exempli, on propozis fixigar, ke : l'acento ne\nl
povos falar sur la i, mem ne en la vorti : \textit{agonio}, \textit{alegorio},\nl
\textit{amfibio}, e. c. Ed on adjuntis : « Por enduktar to, on devus\nl
establisar aparta regulo, nome, ke on ne povas acentizar la i\nl
en la songrupi \VUgras{io}, \VUgras{ia} (anke \VUgras{ie}). »

En II, 8, pri la regulo dil acentizo en \textit{neutral} (acento sur\nl
la vokalo avan la lasta konsonanto) on igis atencar, ke ta\nl
regulo, se ni adoptus lu, devus aplikesar anke ad altra\nl
finali kam \VUgras{ia}, \VUgras{io}, nome ad omna radiki qui finas per vokalo :\nl
\textit{aleo}, \textit{ideo}, \textit{lineo}, \textit{livreo}, \textit{tineo}, \textit{manuo}, \textit{sexuo}, \textit{rituo}, \textit{precipue};\nl
ed ol transportus, en ta omna vorti, l'acento sur la \textit{prelasta}\nl
silabo di la radiko.

Pluse la regulo esas kelke ambigua : « la lasta kon\cc
\end{VUpage}

% -------------------------------------------------------------------
%  feuillet 184 — folio 180. Page d'appendice, SANS NOTES : le releveur
%  a verifie le profil d'encre aux seuils 185, 210 et 220 sous la
%  derniere ligne — rigoureusement nul.
%  PARTICULARITES DE L'ORIGINAL, conservees : un mot d'une serie
%  accentuee laisse sans accent ; l'alternance d'un accent circonflexe
%  et d'un accent aigu sur le meme mot a deux lignes d'ecart ; et une
%  forme cumulant l'accent orthographique francais et l'accent tonique.
% -------------------------------------------------------------------
\begin{VUpage}[184]{180}
\VUcontinue
sonanto » ... kad \textit{di la radiko}, o \textit{di la vorto}? Se : « di la\nl
vorto », to chanjus la loko di l'acento en l'akuzativo (\textit{homôn}\nl
vice \textit{hómo}) e precipue en la verbal formi : \textit{amâs, amîs, amôs}...\nl
Nur l'infinitivi esus lore konforma a la general regulo. Se :\nl
« di la radiko », on devas definar la radiko, plu juste la \textit{radi}\cc
\textit{kalo} di singla vorto : ex. en \textit{amanta, amesos, am} esus la\nl
radiko, e \textit{amant, ames}, la radikalo; en \textit{lumizo, lumiziva, lumi}\cc
\textit{ziveso, lum} esus la radiko, e \textit{lumiz, lumiziv, lumizives} la radi\cc
kali. Kad to ne esus komplikajo e desfacilajo por la lernanti?\parplein

En II, 583, S\textsuperscript{ro} \textsc{Otto Jespersen}, per fundamental e deta\cc
loza artiklo klarigis admirinde la questiono. Il dicis :

La exempli alegata da S\textsuperscript{ro} \textsc{de Janko} tre bone demonstras,\nl
ke esus komplete neposibla trovar acentizo naturala por\nl
omna populi, pro ke omna imaginebla sistemi esas repre\cc
zentata en nia lingui, ed ultree, omna speci de nesistemo o\nl
de to quo aparas adminime tala, exemple l'Angla. Mem la\nl
maxim internaciona vorti ofte prizentas variiva acento, quale\nl
E. \textit{châracter}, D. \textit{Charâkter}, F. \textit{caractêre}; isl., hung., e finl.\nl
\textit{fîlosofi(a)}, E. \textit{philôsophy}, R. \textit{filosôfiya}, F. D. \textit{philosophîe}; e\nl
mem en geografial nomi ni trovas simila diferi, ex. : isl., hung.\nl
e finl. \textit{Amerika}, D. E. \textit{Amêrika}, F. \textit{Amérîque}...

En linguo havanta strukturo simila a nia, on ne povas\nl
strikte e sole aplikar la unesala principo, acentizante sempre\nl
la lasta silabo, pro ke ol esas, en preske omna vorti, sen\cc
valora finalo (vokalo di nura gramatikala valoro), e tama\cc
niere on ruptus komplete la valorprincipo.

Pezigante la prelasta silabo, on kontentigas grandaparte\nl
la unesala tendenco, ed en granda nombro de vorti anke\nl
la valorprincipo, pro ke \textit{nul} esas la maxim grava parto di\nl
\textit{nule}; tale anke \textit{kâra, vôrto, mûlti, sêmblas, atrâkto}, e. c. En\nl
altra vorti, quankam la prelasta ne povas dicesar la maxim\nl
valoranta od importanta por la signifiko, la sama acentizo\nl
esas tote ne evitebla pro la karaktero esence romanika quan\nl
la linguo posedas; nulu propozus altra pronunco kam \textit{acênto},\nl
\textit{detrimênto, exêmplo, hiacînto, exâkta, krokodîlo, pelikâno} e\nl
sennombra altri.

Esas do naturala, ke la maxim multi favoras acento sur\nl
la prelasta; la questiono esas nur, kad on devas establisar\nl
ecepti. »

Citinte 34 vorti havanta per ica regulo acento « sur silabo\parplein
\end{VUpage}

% -------------------------------------------------------------------
%  feuillet 185 — folio 181. Page d'appendice, sans notes.
%  LIGNE INVISIBLE RETROUVEE : « tionas : », tres haut dans la page,
%  fin de l'alinea commence au folio precedent. Le releveur precise
%  qu'il ne s'agit PAS de l'exception connue du verdict « double net »
%  (un titre centre) : il n'y a pas de titre sur ce folio.
%  IRREGULARITE DE L'ORIGINAL, conservee : une forme grecque en « ph »
%  a cote de trois formes en « f », le passage citant la graphie latine.
% -------------------------------------------------------------------
\begin{VUpage}[185]{181}
\VUcontinue
qua aparas nenaturala a multa homi », S\textsuperscript{ro} \textsc{Jespersen} \textit{ques}\cc
\textit{tionas} :

« Kad on devus modifikar la regulo ed acentizar la triesma\nl
de la fino? Me respondas : no. Unesme pro ke to komplikus\nl
la linguo. Duesme, pro ke omni ti qui savas la Franca kono\cc
cas ja acentizi quale consûl, gondôle, publîc, Amerîque,\nl
lyrîque, splendîde; e. c., sen esar shokata per oli. Triesme,\nl
pro ke en multa kazi vorti derivata de ca vorti aquiras natu\cc
rale duesmagrada acento sur la silabo « qua devas portar\nl
l'acento » : pôpulôza, sîmulânta, sîmilêso, pûblikîgo, Amêri\cc
kâno, spîritîsmo; e. c. »

Pri la acento il indikas plura reguli inter qui on povus\nl
selektar. Ma il judikas la sequanta kom la maxim bona :

La prelasta vokalo, ecepte ke en \textit{plursilaba} radiki, \textit{i} e \textit{u}\nl
nemediate avan altra vokalo ne povas portar acento...

La vorto « plursilaba » esas necesa, pro ke sen ol on\nl
havus nula regulo pri \textit{pîa, nîa}, e. c...

Ta regulo « konkordas en maxim multa kazi kun la valor\cc
principo. Ol evitas acento sur \textit{patrôn, niân, amâs}.

Ol esas facile memorenda e aplikenda sen presupozar ula\nl
ciencala konoci pri « mivokali » o pri latina od altra\nl
acentoreguli. Esas vera ke en klaso 3 (\textit{irônio, akadêmio, sim}\cc
\textit{fônio, filosôfio, melôdio, teôrio}, e. c.) ni havas acento qua\nl
semblas nenaturala, ma on devas memorar, ke on havas\nl
kelka nenaturalaji, irge qua regulo adoptesas, ed anke, ke\nl
per ca regulo la ritmo konservas acento (miforta) sur la\nl
sama silabo kande on adjuntas finalo :

melôdio : melôdiôza.

sôcio, sôciâla, sôciêto.

irônio, irôniâla.

gênio, gêniâla.

sêxuo, sêxuâla.

Multa ek ta derivaji konkordas kun la acentizo en natu\cc
rala lingui, kp. E. \textit{melôdious}, D. \textit{melôdisch}; E. \textit{sôcial}, D.\nl
\textit{sôcial}; E. \textit{irônic}, D. \textit{irônisch}; E. \textit{ceremônious, sêxual, mânual},\nl
e. c. Notez anke ke \textit{filozôfo, filolôgo} e \textit{filolôgio} ligesas per\nl
acento, e ke \textit{philosôphia, ceremônia, litûrgia}, e. c., esas la\nl
latina acentizo, quan la Franci e segun lia exemplo anke\nl
kelka altra nacioni chanjis...
\end{VUpage}

% -------------------------------------------------------------------
%  feuillet 186 — folio 182. Page d'appendice, sans notes.
%  PARTICULARITE DE L'ORIGINAL, conservee et signalee par le releveur
%  comme voulue : sur la meme page, les finales nominales citees sont
%  en GRAS et les finales verbales en ITALIQUE. Tranche au grossissement
%  x 8, c'est le seul point qui demandait ce grossissement.
%  Un guillemet ouvrant n'est referme nulle part sur ce folio : la
%  citation court au-dela.
% -------------------------------------------------------------------
\begin{VUpage}[186]{182}
Pos tante bona demonstro l'Akademio povis nur adoptar\nl
la regulo quan S\textsuperscript{ro} \textsc{Jespersen} indikis klare kom preferinda.\nl
Ol agis lo unanime minus un voco per la decido 57 :\nl
L'acento esas sur la lasta silabo di l'infinitivi e sur la pre\cc
lasta silabo di la cetera vorti. Ma en plursilaba radiki, i e u\nl
nemediate avan vokalo ne povas ricevar (recevar) l'acento.\nl
Konseque omna radiki finanta nun per y pos konsonanto\nl
chanjos ol ad i. (\textit{Progreso}, III, 322.)

Pose aparis en III, 392, la sequanta konsideri :

« Unu de nia amiki astonesis pri la nuva (nova) regulo\nl
di acentizo, qua semblas ad il komplikita; il preferus regulo\nl
tote uniforma, exemple : l'acento sempre sur la prelasta\nl
silabo, ed il qualifikas \textit{ecepti} omna deviaco de ta simplega\nl
regulo. Altru remarkas nur, ke la nova regulo donas a kelka\nl
vorti acentizo ne « naturala ».

A to ni respondos : 1\textsuperscript{e} ke la regulo di acentizo esis longe\nl
diskutata hike, ante la decido di l'Akademio, qua rezultis\nl
de ta diskuto ipsa; 2\textsuperscript{e} ke nula \textit{regulo} di acentizo povas\nl
sequar la naturala acentizo, pro ke ca acentizo esas \textit{senregula}\nl
(t. e. senuniforma regulo); 3\textsuperscript{e} ke nula regulo di acentizo\nl
povas kontentigar omna populi, pro ke l'acento varias, por\nl
la sama vorti, segun la lingui. On devas agnoskar, ke la\nl
nova regulo konstitucas granda simpligo en l'ortografio, per\nl
ke ol unigas la finali \VUgras{-io}, \VUgras{-yo}, inter qui on ofte hezitis e\nl
dubis, segun l'atesto di multi. Ol esas do reala simpligo ed\nl
uniformigo.

Ti qui dezirus simpligar ed uniformigar « komplete »\nl
la regulo di acentizo esas Franci, e quale Franci nule kom\cc
prenas l'importo di l'acento por la cetera populi. Semblas\nl
a li indiferenta acentizar ca o ta silabo, pro ke li reale acen\cc
tizas nula, o (nekoncie) acentizas \textit{sempre la lasta}. Ico esas\nl
grava kulpo e danjeroza eroro. L'acentizo maxime importas,\nl
ne nur a la beleso e harmonio di nia linguo, ma a lua kom\cc
prenebleso por omna (o maxim multa) populi. Or de l'inter\cc
naciona vidpunto, esas nule indiferenta, ke la finali \textit{-ar}, \textit{-ir},\nl
\textit{-or} dil infinitivo esez acentizata o senacenta. Simile, esas\nl
tote \textit{neposibla}, fonetike, acentizar \textit{aqûo} e \textit{lingûo}, \textit{filîo} o\nl
\textit{familîo}, \textit{sekretarîo} e \textit{notarîo} (quankam Esperanto facas la\nl
lasta kulpo) ankore (mem) min \textit{qûa}, \textit{qûo}. Ne esas arbitriala\nl
regulo, ma fonetikal fakto universala, ke la vokali \textit{i} e \textit{u}\nl
\end{VUpage}

% -------------------------------------------------------------------
%  feuillet 187 — folio 183. OUVERTURE DE L'APPENDICE 2.
%  Trois elements d'apparat, dans cet ordre : un ORNEMENT (filet nu de
%  93 px = 7,9 mm, centre), puis le titre en capitales demi-grasses,
%  puis un sous-titre italique entre parentheses, d'un corps plus petit
%  que le texte courant.
%  FILET DES NOTES : panne (a), trop pale. Le releveur l'a etablie en
%  rejouant le detecteur — un grain isole a x = 95 fixe le depart, puis
%  neuf pixels clairs que la fermeture ne soude pas, et la longueur
%  tombe a 1 ou 2 px. Il ecarte (b) par mesure (reste = 4 et 6, sous le
%  seuil) et (c) par position (le filet est bien dans la fenetre).
%  PARTICULARITE DE L'ORIGINAL, conservee : les lettres citees changent
%  d'enrichissement sur la meme page — s en gras trois fois puis en
%  italique, n en italique puis en gras.
% -------------------------------------------------------------------
\begin{VUpage}[187]{183}
\VUnotes{127.59mm}{%
(1) Pro to, en la linguistiko indo-europana, li esas la sola vokali,\nl
qui genitas diftongi. Per to li distingesas esence de la ceteri.

(2) Videz l'Apendico « Konjugo-sistemo di Ido ».

(3) Vere kad esus tolerebla e facile pronuncebla : \textit{quasn vu}\nl
\textit{prenos?} vice : \textit{quin vu prenos?}

O : \textit{la homos quasn me renkontris?} O fine : \textit{mokosn, persekutosn,}\nl
\textit{mem tormentosn dum yaros il subisis...} vice : \textit{mokin, persekutin,}\nl
\textit{mem tormentin dum yari il subisis...?}

Or la \VUgras{n} inversigala esas necesa, nome por \textit{qua}, \textit{qui}, \textit{quo}, quale on\nl
vidis en la gramatiko. Senco klara, nedubebla povas obtenesar nur\nl
per olu, en ula kazi. Do ol devas esar konservata, ed ol ne esas\nl
admisebla pos \textit{s} en finalo,%
}

\VUcontinue
avan altra vokali divenas generale mi-konsonanti, ne nur en\nl
la finali, ma interne di la vorti (ex. : \textit{koliaro}, \textit{kuliero}, \textit{buli}\cc
\textit{ono}, \textit{liuto}) (1). La nova regulo simple konstatas ica fakto,\nl
nur kelke generaligante ol, pro ke omna regulo devas esar\nl
generala.

On ne devas parolar pri « ecepti », pro ke ni adoptas\nl
regulo diversa por diversa speci de vorti o de vokali, qui\nl
esas reale diferanta. Existas nula neceseso, mem nula mo\cc
tivo, por ke l'acentizo esez « simpla ed uniforma », t. e. falez\nl
sempre sur la sama silabo \textit{en omna vorti}; tala acentizo me\cc
ritus prefere la nomo di « arbitriala regulo ». Se Zamenhof\nl
adoptis tala regulo por Esperanto, to ne venas de irga pro\cc
funda intenco o misterioza sajeso; il simple imitis l'acentizo\nl
di sua linguo, la Polona. Se ulu havas la yuro reprochar,\nl
ke nia regulo ne esas « simpla », to certe ne esas l'Esperan\cc
tisti, qui, fanfaronante pri la simpleso di lia gramatiko en\nl
« 16 reguli », kontas kom un sola regulo... la tota konjugado! »\parplein

\VUsaut{4.61mm}
\VUfilet{7.9mm}
\VUsaut{3.20mm}
\VUtitre{10.27pt}{26}{LA PLURALO EN IDO.}
\VUsaut{0.87mm}
\VUsoustitre{5.67pt}{234}{\textit{(Apendico 2ma.)}}
\VUsaut{3.60mm}

Ta punto esas dominacata da du kozi : 1\textsuperscript{e} la verbal rolo\nl
di \VUgras{s} (2), unesma motivo decidiganta por ne adoptar ta kon\cc
sonanto kom plural-marko; 2\textsuperscript{e} ne posibleso (adjuntante\nl
\textit{n} a \VUgras{s}) obtenar finalo facile pronuncebla ed eufonioza,\nl
duesma motivo decidiganta por eskartar \VUgras{s} kom plural\cc
marko (3).
\end{VUpage}

% -------------------------------------------------------------------
%  feuillet 188 — folio 184.
%  LIGNE INVISIBLE RETROUVEE : « olu. », au fer, fin du deuxieme alinea
%  de la note (2). Le verdict « double net » garde son score.
%  COQUILLES DE L'ORIGINAL, conservees : « niaopinione » agglutine et
%  « lon » pour l'on.
% -------------------------------------------------------------------
\begin{VUpage}[188]{184}
\VUnotes{123.11mm}{%
(1) Quan havis \textit{Neutral} e \textit{Reform-Neutral}.

(2) On reprochis a Ido (por justifikar Esperanto) formacar sua\nl
pluralo per \textit{substituco}. \VUgras{O}, on dicas, esas la marko distingiva di la\nl
substantivo. Konseque on devas konservar ol en la pluralo.

Unesme \VUgras{o} ne esas nur la marko dil substantivo : ol esas anke\nl
la marko dil substantivo \textit{singulara}. On do havas nul motivo por\nl
adjuntar la marko dil pluralo ad olta dil singularo; kontree, on\nl
devas remplasigar ica per ita; nam l'ideo dil pluralo ne \textit{adjuntesas}\nl
al ideo dil singularo : ol kontredicas lu absolute. La marko dil plu\cc
ralo konseque devas remplasar olta dil singularo e ne adjuntesar ad\nl
olu.

Altraparte, se \VUgras{o} esus la marko generala di la substantivo, on\nl
devus adjuntar \textit{un} specal marko por la singularo, ed \textit{un altra} marko%
}

Un del argumenti uzata kontre Ido esas opozar a lua plu\cc
ralo per -i, la pluralo per -s dil Angla, di la Hispana e di\nl
la Franca. Tre rapide on deduktas ica konkluzo : Ido vio\cc
lacas l'internacioneso, formacante sua pluralo per -i.

Pri ta reprocho un del maxim eminenta partisani dil\nl
naturaleso en L. I. ed anke dil pluralo per -i (1), S\textsuperscript{ro} \textsc{Rosen}\cc
\textsc{berger}, mortinta, suflas a ni la respondo : « On devas ne\nl
regardar unike l'internacioneso (en la gramatiko), ma anke\nl
la regulozeso, simpleso e komodeso. Or la finalo -i esas plu\nl
komoda kam -s (-es), qua postulas komplikita regulo ».

Ni adjuntez, ke nur ta finalo plurala posibligas l'akuza\cc
tivo (\textit{-n} inversigala) e co esas decidiganta niaopinione. Se\nl
la homo poke instruktita sucesas nur tre desfacile aplikar\nl
l'akuzativo (lon montras l'Esperantisti ed ante li la plebeyi\nl
Latina), ta homo komprenas frazo qua aplikas ol saje. En\nl
la vivo omnadia, ta kazo esas kompliko neutila, e jenanta,\nl
quale omna kompliko. Ibe l'inversigo di la subjekto povas\nl
e devas proxime reduktesar al kazo ube \textit{qua, qui, quo} divenas\nl
komplemento direta e preiras la subjekto : \textit{quan} me vidas?\nl
evidente diferanta de \textit{qua} vidas me? o : \textit{qua} me vidas?

Ma, exter la domeno omnadia, ne esas certa ed on ne\nl
pruvis, ke l'akuzativo esas neutila por unionar plu rapide\nl
e plu evidente la du parti di rezono filozofiala, teologiala,\nl
o mem matematikala, e. c. Yen pro quo Ido ne forjetis tote\nl
l'akuzativo. Or, ni lo dicis e repetas, l'akuzativo esas posibla\nl
nur kun la pluralo per -i, ecepte se on inventus ula sub\cc
tilajo qua fakte pruvus nia aserto (2).
\end{VUpage}

% -------------------------------------------------------------------
%  feuillet 189 — folio 185.
%  FILET DES NOTES : panne (a), trop pale, etablie par le meme procede
%  qu'au feuillet 187 — l'extremite gauche du trait tombe sous le seuil,
%  la fermeture de 7 px ne soude pas le trou de 12 px, et la longueur
%  retombe a 1. Filet a y = 1038 px, ordonnee 93,98 mm.
%  COQUILLES DE L'ORIGINAL, conservees : « specala » pour speciala ;
%  « rispektive » pour respektive ; et « e. c. e. c. » repete.
%  PARTICULARITE conservee : dans une meme enumeration de mots francais,
%  deux d'entre eux sont en italique quand tous les autres sont en gras.
% -------------------------------------------------------------------
\begin{VUpage}[189]{185}
\VUnotes{93.98mm}{%
\VUcontinue
specala por la pluralo. Tal esas la konsequo di ta postulo super\cc
logika. Ma lore la kritiko atingas tam Esperanto kam Ido.

Altraparte, se \VUgras{o} esas la marko di la substantivo generale, qua do\nl
esas la marko dil verbo generale? Qua esas la komuna elemento en\nl
ta omna finali \textit{as, is, os, us, u, i} di Esperanto? Nula, ne mem un\nl
litero! On certe respondos ke suficas savar, ke ta omna finali esas\nl
finali \textit{verbala} e karakterizas la verbo. Tre bone! ma same suficas\nl
savar, ke \VUgras{o} e \VUgras{i} esas finali substantivala e karakterizas rispektive la\nl
substantivi singulara e la substantivi plurala. La rezono esas\nl
rigoroze identa en la du kazi. Se l'objeciono havus la maxim mikra\nl
justeso, on devus konservar la marko dil verbo (\VUgras{i} dil infinitivo\nl
en Esperanto, exemple) en omna modi e tempi, e dicar : \textit{amias,}\nl
\textit{amiis, amios..., amianta...} Fakte, por defensar la pluralo di Esper\cc
anto, on kondamnas lua tota konjugo! On do koaktesas admisar, ke\nl
se la konjugo di Esperanto ne pekas kontre logiko, la pluralo\nl
per \VUgras{i} ne pekas, anke lu, kontre logiko. O l'objeciono ne valoras, od ol\nl
atingas la kara linguo quale « la pavo dil urso »!

Finante, ni donez la decidiganta motivo di la substituco dil Ido\cc
pluralo ad olta di Espo; ta motivo esas, ke l'Ido-pluralo esas tam\nl
facile pronuncebla da omni (ye la nominativo o ye l'akuzativo) kam\nl
ne esas tala por tota populi la pluralo di Espo, precipue ye l'akuza\cc
tivo. Ex. : \textit{Me amas chiujn tiujn, kiujn vi amas, kaj mi malamas}\nl
\textit{chiujn tiujn, kiujn vi malamas.} En Ido : \textit{me amas ti omna quin vu}\nl
\textit{amas, e me odias ti omna quin vu odias.}%
}

Pos dicir to, ni examenez la susten-punti dil pluralo\nl
per \VUgras{-s} (\VUgras{-es}).

Fakte \textit{du}, ne \textit{tri} existas. On alegas \textit{tri} por plufortigar\nl
l'argumento, o pro distrakteso.

Nam la Franca devas esar sustracionata pro la sequanta\nl
motivi : 1\textsuperscript{e} \VUgras{s} en ta linguo finas pasable granda nombro de\nl
vorti \textit{singulara} : \VUgras{corps, puits, temps, fils, remords, mors,}\nl
\VUgras{buis, pus, pouls, pas, tas, os, dos, cas, jus, bris, accès, succès,}\nl
\textit{décès, procès}, e. c.; 2\textsuperscript{e} \VUgras{x} esas kun olu tre ofte la plural\cc
marko : \VUgras{émaux, agneaux, beaux, baux, rideaux, caveaux,}\nl
\VUgras{cerceaux, marteaux, tombeaux,} e. c.; \VUgras{jeux, cheveux, vœux,}\nl
\VUgras{bi\/joux, cailloux, choux, genoux, joujoux, hiboux, poux,}\nl
\VUgras{creux,} e. c.; 3\textsuperscript{e} on skribas \VUgras{s}, ma on ne pronuncas lu : \VUgras{clous,}\nl
\VUgras{sous, livres, tables, récits, pages, mers, mères, pères, sœurs,}\nl
e. c. e. c. Do on ne darfas kontar la Franca, od on devas\nl
pozar kom principo : quale en la Franca, on skribos \VUgras{s} por\nl
la pluralo, ma on ne pronuncos lu. Or ta principo esas\nl
posibla en la Franca pro la determinanti \textit{les, mes, tes, ses}, e. c.,\nl
ma la L. I. ne indikas la pluralo en sua determinanti : \textit{la}\nl
\end{VUpage}

% -------------------------------------------------------------------
%  feuillet 190 — folio 186.
%  ASTERISME, ornement neuf dans le volume : trois asterisques en
%  triangle entre deux alineas, la ou la premiere partie se contentait
%  d'un filet. Verifie a l'agrandissement par l'assembleur ; d'ou la
%  macro \VUasterismo, ajoutee au preambule.
%  LIGNE INVISIBLE RETROUVEE : « à 'tush. », troisieme ligne de la
%  note (1), large de 85 px seulement.
%  Le releveur signale que la geometrie est fausse a une ligne : une
%  salete isolee gonfle la boite et donne un « fin » trompeur.
%  PARTICULARITE DE L'ORIGINAL, conservee : la meme desinence citee est
%  tantot en gras, tantot en romain, cinq fois sur la page.
% -------------------------------------------------------------------
\begin{VUpage}[190]{186}
\VUnotes{129.54mm}{%
(1) On ne audacis til nun alegar la Portugalana. Ma ni ne deses\cc
perez. Pro quo ne? Quale en ta linguo, ni skribus \textit{actos} e pronuncus\nl
\textit{à 'tush.}

Ma on ne obliviis citar kom susteno, por \VUgras{s} plurala, linguo klasika,\nl
la Latina, en qua ta vokalo finas omnainstante nominativo o genitivo\nl
\textit{singulara}, tante vere ol esas marko di \textit{pluralo!}

(2) Ol ne esas nekonocato en la Greka moderna.

(3) Zamenhof adoptis lu en sua reformo di 1894, e lore il dicis,\nl
ke ta pluralo esas plu harmonianta kun la fizionomio di Esper\cc
anto. Il chanjis a \textit{a} sua infinitivo \VUgras{-i}, e donis \VUgras{-e} kom karakterizivo a\nl
l'adjektivo ed a l'adverbo.%
}
\VUasterismo{48.68mm}

\VUcontinue
(homi), \textit{mea} (frati). Do on ne darfas kontar la Franca kom\nl
susteno (por la helpo-linguo) dil pluralo per \VUgras{s}, ca \VUgras{s} ne\nl
sonante en la vorti plurala di ta linguo.

Restas vere por la pluralo per \VUgras{-s} nur \textit{du} susteni : la Angla\nl
e la Hispana (1). Nam esas quaza joko alegar altra lingui,\nl
en qui \VUgras{s} renkontresas preske tam ofte en la singularo kam\nl
en la pluralo.

\VUsaut{6.99mm}
Nun ni videz la susteno dil pluralo per -i.

La Slava lingui (susteno extreme importanta) e la Ruma\cc
niana kun l'Italiana (2).

En ica lasta linguo, -i havas ne nur la maskuli, ma mem\nl
femini : \textit{madri} (matri). E, fakto remarkinda, \VUgras{s} qua des\cc
aparis en la paso dil Gotika al Germana, transformesis\nl
a i (kun \textit{e} di \textit{es}) pasante de la Latina al Italiana : \textit{uomini}\nl
= \textit{homines}; \textit{madri} = \textit{matres}; \textit{cani} = \textit{canes}.

Fine, la pluralo per -i esas konocata da omna Latinisti\nl
(kad on negus lo?), tale ke \textit{boni, mali} di Ido esas exakte\nl
\textit{boni, mali} di la Latina (kad en ta linguo on dicas : \textit{bones,}\nl
\textit{males?}); ol esas konocata da multega altra homi per Italiana\nl
plurali quin on darfas dicar internaciona por omna Euro\cc
pano od Amerikano kelke lerninta la gramatiko di sua lin\cc
guo : \textit{carbonari, ciceroni, condottieri, dilettanti, lazzaroni,}\nl
\textit{quintetti}.

Ka do ni exajeras dicante : la pluralo per -i esas prefe\cc
rinda por la linguo helpanta, e Ido saje agis adoptante lu (3).\parplein

L'Italiana qua facis \textit{uomini} de \textit{homines}, \textit{duca} de \textit{dux}\nl
(ducis), \textit{mano} de \textit{manus}, e. c., inspiras ta reflekto : se ula\nl
\end{VUpage}

% -------------------------------------------------------------------
%  feuillet 191 — folio 187.
%  FILET DES NOTES : panne (a), trop pale. Le releveur l'a etablie en
%  rejouant le detecteur seuil par seuil — a 185 et 190 la longueur vaut
%  1 px, a 195 elle vaut 208 et le reste 0. Il ecarte (c) par position
%  (le filet est a 0,57 H, la fenetre part de 0,25 H) et (b) par mesure
%  (reste = 0). Filet a y = 1184 px, ordonnee 105,16 mm.
%  COQUILLE DE L'ORIGINAL, conservee : « diciplino » pour disciplino.
%  Le releveur signale que outils/graisse.py est inutilisable sur le
%  corps de cette page : il declare GRAS les vingt-deux lignes sans
%  exception, verdict uniforme donc suspect. Tranche a l'oeil : aucun
%  gras dans le corps.
% -------------------------------------------------------------------
\begin{VUpage}[191]{187}
\VUnotes{106.34mm}{%
(1) On reprochis a \textit{-as}, \textit{-is}, \textit{-os}, \textit{-us} dil konjugo produktar tro ofta\nl
sisi, obliviante ke li esas nekompareble plu multa en la Hispana,\nl
per olua pluralo, quan on asertas esar la maxim bona por L. I.\nl
Cetere ni komparez : \textit{omnasorta legumos, terpomos, karotos, napos,}\nl
\textit{kaulos, pizos, fazeolos, artichokos e diversa fruktos : pomos, piros,}\nl
\textit{vitberos, nucos} e. c., a : \textit{omnasorta legumi : terpomi, karoti, napi,}\nl
\textit{kauli, pizi, fazeoli, artichoki e diversa frukti : pomi, piri, vitberi,}\nl
\textit{nuci,} e. c. (*). Forsan mea orelo esas trompata da l'Italiana, ma\nl
me preferas kom dezinenci le i sucedanta a le \textit{os} tante sika e tante\nl
desagreable sisanta di la Hispana pluralo.

Ne esas sen utileso remarkigar la regulozeso dil pluralo per \VUgras{-i} en\nl
Ido. Olua personal pronomi ipsa revelas sua pluralo per \VUgras{i} : \textit{ni, vi,}\parplein

(*) En Espo : \textit{chiuspecaj legomoj : terpomoj, karotoj, napoj,}\nl
\textit{brasikoj, pizoj, artishokoj, kaj diversaj fruktoj : pomoj, piroj,}\nl
\textit{vinberoj, juglandoj} k. t. p. Akuzative : \textit{chiuspecajn legomojn : ter}\cc
\textit{pomojn, karotojn, napojn, brasikojn, pizojn, artishokojn, kaj diver}\cc
\textit{sajn fruktojn; pomojn, pirojn, vinberojn, juglandojn} e. c. Uf!

Ka, koram tal exempli e centi de altri posibla, la pluralo per \VUgras{-i}\nl
ne semblas la maxim bona?%
}

\VUcontinue
sistemo di helpolinguo agas plu o min tale, la diletanti di\nl
naturaleso e di latineso plendas pro barbarismi, klozante\nl
la okuli pri olti quin li facas, ed energioze li rimemorigas\nl
la prototipo Latina. Ni do ne trublesez da lia kritiki e ni\nl
ne shamez pri tala chanji en Ido, qua ne vizas neutil eru\cc
diteso, ma praktikala simpleso. Kun ta exkuzo, ni forsan\nl
darfas uzar tanta libereso kam l'\textit{evoluciono naturala} en\nl
vivanta linguo, e tante plu, ke ica tro ofte uzas blinde ta\nl
libereso, kontre ke L. I. bone konstitucita havas kom freno\nl
la simpleso e facileso necesa, e kom guidilo la koncio dil\nl
skopo vizita.

Ni lasez nia kontrediceri skandalesar en Ido pri kozi quin\nl
li admiras en vivanta lingui. Li unesme laboras por su, en\nl
la realeso, por havar \textit{sua} linguo internaciona propre indi\cc
viduala e partikulara. Nur ica povas satisfacar lia senti\cc
mento e gusto linguala, quin li ne povas sakrifikar a l'in\cc
teresto di la homaro. On asertas, ke ta marcho diverganta\nl
donos fine la konvergo e helpolinguo unika. Ni kontraste\nl
laboras por la maso homala, en la diciplino e obediado lin\cc
guala; kompreneble la procedi ne povas esar identa e sama.\parplein

La tempo e la homaro facos la selekto, e me esas tre certa\nl
ke fine li decidos e rezolvos por ni (1).
\end{VUpage}

% -------------------------------------------------------------------
%  feuillet 192 — folio 188. OUVERTURE DE L'APPENDICE 3.
%  Titre en capitales demi-grasses et sous-titre italique entre
%  parentheses, sans ornement au-dessus — contrairement a l'appendice 2,
%  qui s'ouvrait sur un filet. Verifie : l'espace entre le folio et le
%  titre est entierement blanc.
%  Le signalement « deux lectures » a la ligne 3 est tranche : c'est le
%  blanc qui suit le bloc de titre, la zone est vierge.
%  PARTICULARITE DE L'ORIGINAL, conservee : le meme nom d'auteur est
%  introduit par « Sioro » en toutes lettres a une ligne et par
%  l'abreviation en exposant huit lignes plus bas.
% -------------------------------------------------------------------
\begin{VUpage}[192]{188}
\VUnotes{114.15mm}{%
\VUcontinue
\textit{li, ili, eli, oli}; lua posedali anke : \textit{mei, tui, sui} e. c., lua relativo\nl
questionanta : \textit{qui}, pluralo di \textit{qua}; lua demonstrativi : \textit{(i)ci}, \textit{(i)ti},\nl
pluralo di \textit{(i)ca}, \textit{(i)ta}; lua nedefiniti : \textit{uli, nuli, kelki, pluri,}\nl
\textit{multi}, e. c.; lua substantivi : \textit{homi, viri, mulieri, infanti}; lua\nl
adjektivi : \textit{yen grosa e mikra pomi; prenez la grosi e lasez la mikri,}\nl
tam klara e reguloza kam \textit{le grosa, le mikra}.

En Espo la \textit{j} plurala ne trovesas \textit{en ni, vi, ili}, tamen plurala; e \textit{li,}\nl
\textit{shi, ghi} qui rimifas kun \textit{ni, vi, ili} ne esas plurala, ma singulara.

(1) Esas utila remarkigar, ke la nov-Latina lingui havas nur \textit{du}\nl
\textit{genri} : la maskulalo e la feminalo, kontre ke lia matro, la Latina,\nl
pluse havis la neutro. E la Angla, qua nun preske ne plus havas\nl
genro, venas de lingui posedanta la \textit{tri genri}.

(2) \textit{Progreso}, IV, 572.

(3) Fakto simila ne indijas exemplo mem en la Franca : \textit{de bons}\nl
\textit{pères, de bonnes mères}. La genro aparas nur en \textit{bons, bonnes}.%
}

\VUtitre{10.63pt}{-32}{GENRO E MASKULISMO.}
\VUsaut{1.10mm}
\VUsoustitre{5.31pt}{283}{\textit{(Apendico 3-ma.)}}
\VUsaut{4.27mm}

Maxim ofte la genro indikesas, en la lingui posedanta olu,\nl
da gramatikal formo partikulara, ma ulfoye, mem en ta\nl
lingui, ol esas expresata dal vortolibro, altradice dal nomo\nl
ipsa dil anmoza ento (ex. : F. \textit{père, mère; oncle, tante;}\nl
\textit{bœuf, vache, cheval, jument}). Sioro A. \textsc{Meillet} lo remar\cc
kigis en sua kurso pri Gramatiko generala, che la \textit{Sorbonne}.\nl
Il anke memorigis, ke multa lingui ignoras la genro.\nl
Exemple, la Angla posedas lu nur en la pronomi \textit{he, she, it}.\nl
La genro ne existas en la grupo (Finno-ugriana), en la Turka.\nl
Ula lingui perdis lu, qui lu posedis, exemple la lingui\nl
Iraniana, l'Armeniana, qua ne mem havas feminal sufixo\nl
por formacar \textit{rej-ino}, me supozas (1). « Do, konkluzas\nl
S\textsuperscript{ro} \textsc{Meillet}, la genro \textit{gramatikala} nule esas universala e\nl
necesa kategorio... ». « L'Europani tendencas atribuar a la\nl
genro tro granda importo, pro ke en lia lingui l'artiklo\nl
konstante akompanas la substantivo e memorigas lua genro\nl
(ecepte en E.). » (2).

En Indo-Europana, la genro ne indikesas en la substan\cc
tivo, ma en l'adjektivo : la substantivo posedas nur genro\nl
latenta (\textit{pater, mater} havas simila formo), quan revelas\nl
l'adjektivo akompananta (\textit{bonus, bona}) (3). Do nur l'akordo\nl
indikas la genro.
\end{VUpage}

% -------------------------------------------------------------------
%  feuillet 193 — folio 189.
%  COQUILLES DE L'ORIGINAL, conservees : « special » sans -a devant un
%  substantif, deux fois ; « a ni » puis « ad ni » dans deux tournures
%  identiques ; et un guillemet ouvrant jamais referme.
%  Le releveur signale deux fausses lectures evitees a l'agrandissement,
%  toutes deux sur la meme ligature serree.
% -------------------------------------------------------------------
\begin{VUpage}[193]{189}
\VUnotes{156.21mm}{%
(1) Ni bone remarkez, ke en ica frazo : « \textit{La hundo esas la amiko}\nl
\textit{di la homo} », nek \textit{hundo}, nek \textit{amiko}, nek \textit{homo} implikas la sexuo,\parplein
}

« La klasifiko dil genro maxim konocata da ni, ma ne\nl
la maxim importanta, esas la \textit{sexuo}. Ol uzesas preske nur\nl
en la min multa lingui, nome le Indo-Europana e le Semi\cc
dala. Ol esas do multe min importanta kam semblas a ni\nl
kom Europani. Mem en ta lingui ol aplikesas nur a parto\nl
di la nomi, a la nomi di l'enti sexuoza; e mem por multa\nl
animali on ne distingas maskulo e femino, adminime per\nl
chanjo di formo. Exemple en la Franca, por indikar la\nl
sexuo en : \textit{elefanto, hieno, strucho, buvrelo, skarabeo,}\nl
\textit{heliko}, e. c., e. c., on koaktesas adjuntar sive « maskula »,\nl
sive « femina ». Tale ke, se on kontus, on trovus, ke nur\nl
ecepte la genro esas markizata, mem en la Franca.

Semblis a ni esar utila memorigar, kun la susteno di lin\cc
guistikal autoritato :

1\textsuperscript{e} ke la genro \textit{gramatikala} nule esas universala e necesa\nl
kategorio;

2\textsuperscript{e} ke la klasifiko dil sexuo esas multe min importanta\nl
kam semblas ad ni.

Pos agir ico, esos plu facil ad ni expozar la procedo di\nl
Ido por la genro, o plu juste por la sexuo e konocigar olua\nl
motivi.

Unesme la questiono dil maskulalo o feminalo povas kon\cc
cernar en Ido nur \textit{la nomi di enti anmoza pri qui on volas}\nl
\textit{indikar la sexuo}. Or, maxim ofte, esas nek necesa, nek mem\nl
utila indikar olu. De ico konsequas : la dezinenco \VUgras{-o} dil\nl
substantivo suficas tam bone por la enti anmoza kam por\nl
la kozi. Nula special dezinenco esas necesa por l'uni o por\nl
l'altri. On ya ne povas vidar animalo en vorto indikanta\nl
kozo. Konseque, pro quo ni havus du dezinenci, kande un\nl
suficas?

Altra konsequo : substantivo di ento anmoza qua havas\nl
nul special marko di sexuo ne esas plu multe maskula kam\nl
femina, ol indikas nur la \textit{speco}, quale F. \textit{enfant, conjoint}\nl
e mem \textit{homme}, qui ne implikas plu multe un sexuo kam\nl
l'altra, quale \textit{bous} Greka, \textit{bos} Latina e mem \textit{agnus} en l'an\cc
tiqua Latina, ube on dicis \textit{agnus femina} por F. \textit{agnelle}, ante\nl
dicar \textit{agna} (1).
\end{VUpage}

% -------------------------------------------------------------------
%  feuillet 194 — folio 190.
%  FILET DES NOTES : panne (c), trop haut — le cas le plus net du
%  volume. Le filet est a y = 450 px quand la fenetre de recherche part
%  de 0,25 H = 516. Le releveur a etabli par mesure que les deux autres
%  tests passent (longueur 213 px, reste 2) et a verifie en rejouant le
%  detecteur avec la seule fenetre elargie : il le trouve aussitot.
%  Ordonnee 44,79 mm.
%  LIGNE INVISIBLE RETROUVEE : « (frato). », huit signes, absente de la
%  geometrie.
%  COQUILLE DE L'ORIGINAL, conservee : « evas » pour esas, verifiee a
%  l'agrandissement x 6.
% -------------------------------------------------------------------
\begin{VUpage}[194]{190}
\VUnotes{44.79mm}{%
\VUcontinue
mem en la Franca. Ni remarkez anke, en la Franca, ta stranja fakto :\nl
\textit{enfant} ne implikas sexuo, e \textit{père} esas maskula, adminime en la\nl
singularo! Nam, en la plurala, ol kontenas mem la \textit{matri}, quale\nl
\textit{homme} kontenas la homini : \textit{nos pères nous ont appris} (nia patri\nl
docis ad ni); \textit{l'homme est mortel} (la homo esas mortiva). --- La\nl
Madyara, plu logikoza, donas nula genro a \textit{père} (patro) nek a \textit{frère}\nl
(frato).

L'Angla lasas sen genro \textit{kuzo} ed \textit{amiko}. Pro to lu dicas : \textit{male}\nl
\textit{cousin, female cousin; male friend, female friend.}

La Rusa havas multa sengenra nomi, inter altri : \textit{sirota} (orfano).\nl
Ulkaze ol havas tri nomi diversa : 1 por la sengenra substantivo (di\nl
speco) : \textit{sobaka} (hundo), 1 por la maskulo : \textit{kobel} (hundulo), 1 por\nl
la femino : \textit{suka} (hundino); \textit{chelovek} (homo), \textit{mujchina} (homulo),\nl
\textit{jenschina} (homino).

On atencez ico : se \textit{patro}, \textit{frato} e. c. esus maskula (homuli),\nl
\textit{ge-patri} e. c. esus absurdajo, nam on ne povas unionar la du sexui,\nl
kande existas nur una.

(1) \textit{Omna} maskula nomi en nia lingui esas dusenca : oli signi\cc
fikas, lore maskula enti, lore (e maxim ofte) la du sexui sen\nl
distingo. La exempli esas sennombra; ni citez nur du. Kande ni\nl
parolas pri nia \textit{parenti} e nia \textit{amiki}, kad ni komprenas per to nur\nl
\textit{maskula} parenti ed amiki, o kad ni inkluzas anke nia \textit{parentini} ed\nl
\textit{amikini}? Certe la duesma senco esas preske sempre la justa. Do ca\nl
plurali esas dusenca en nia lingui : ni ne povas distingar klare,\nl
gramatikale nia parentuli ed amikuli inter nia parenti ed amiki\nl
generale.

Ma, se \textit{parento}, \textit{amiko} esas preske sempre sengenra en la pluralo,\nl
li devas evidente esar sengenra anke en la singularo; nam altre li\nl
\textit{chanjus sua senco} pasante del singularo al pluralo! Esas do absolute\nl
necesa, se ni volas parolar pri parento, amiko \textit{maskula}, ke ni indikez\nl
lo irgamaniere. Pro to, vice uzar la tota vorto \textit{maskula}, ni uzas\nl
la sufixo \textit{-ul} qua esas lua abreviuro, same kam \textit{-in} esas l'abreviuro\nl
di \textit{femino}. Do ni parolas pri \textit{parentulo}, \textit{amikulo}, same kam ni parolas\nl
pri \textit{parentino}, \textit{amikino}, e pro la sama motivo : a la \textit{sengenra} nociono\nl
di \textit{parent}, \textit{amik}, ni devas adjuntar, lore la nociono di maskulo (\textit{-ul}),\nl
lore la nociono di femino (\textit{-in}), ma \VUgras{nur kande ni bezonas li}, t. e.\nl
kande ni havas li en la mento, e deziras expresar li : \textit{on nultempe}\nl
\textit{obliviez ica grava restrikto!} Nam se, quale ni dicis, 95 foyi ek 100\nl
\textit{amiko} (en nia lingui) signifikas egale \textit{amikulo} od \textit{amikino}, 95 foyi\nl
ek cent ni devas uzar simple \textit{amiko}, e ne \textit{amikulo} nek \textit{amikino}. La\nl
suprega lego di la L. I. esas : « Expresez vua tota penso, ma nur\nl
vua penso. » (\textit{Progreso}, VI, 588.)%
}

Ne sen fundamental studiado di la questiono Ido adoptis\nl
sua nuna sistemo por la enti anmoza. Nam ol heredabis\nl
de Esperanto la eroro pri \textit{maskulismo}, qua konfundas im\cc
plicite la nomo di speco kun ta dil maskulo (1) e formacas\nl
la feminalo ek ta lasta. Ica procedo stranja « evas de ta\parplein
\end{VUpage}

% -------------------------------------------------------------------
%  feuillet 195 — folio 191.
%  LIGNE INVISIBLE RETROUVEE : « p. 37) : », qui clot un alinea de la
%  note (3), large de 180 px.
%  PARTICULARITES DE L'ORIGINAL, conservees : un blanc franc entre la
%  parenthese et le point final a la derniere ligne du corps ; et, dans
%  la citation en esperanto, les digrammes et l'accent grave de
%  l'orthographe ancienne, conserves tels quels.
% -------------------------------------------------------------------
\begin{VUpage}[195]{191}
\VUnotes{90.85mm}{%
(1) \textit{Progreso}, VI, 592.

(2) On remarkez bone, ke ca sistemo esas tote analoga e kon\cc
forma a la sistemo di nia pronomi personala : same kam ni havas\nl
sengenra \textit{lu} apud ed exter \textit{il}, \textit{el}, \textit{ol}, ni devas havar \textit{homo} exter\nl
\textit{homulo} e \textit{homino}, e. c. : nam ta formi devas necese korespondar e\nl
harmoniar inter su. L'Angla ipsa ne havas la sengenra pronomo e\nl
mustas dicar : \textit{he or she} (il od el), quo esas inkombranta e tedanta.\parplein

(3) S\textsuperscript{ro} \textsc{Couturat} finas per ica noto :

Lasta vorto. Se la fanatiki di la « Zamenhofa lingvo », qui kus\cc
tumas uzar kontre ni omnaspeca argumenti, bona o ne, sincera o\nl
ne, esus tentata kritikar la sistemo di Ido, esus facila klozar lia\nl
boko per la sequanta deklaro di la Maestro ipsa (\textit{Pri reformoj en}\nl
\textit{Esperanto}, p. 13-14; texto aparinta en \textit{Esperantisto}, marto 1894,\nl
p. 37) :

« Kelkaj amikoj proponis, ke ni enkonduku apartan sufikson\nl
por substantivoj speciale viraj » (Remarkez la ridinda uzo di \textit{vira}\nl
\textit{virina} vice \textit{maskula, femina} : ica uzo duras ankore nun : \textit{bovo-viro!})\nl
« Pripensinte tiun chi proponon, mi trovis ke \textit{ghi estas ne sole tre}\nl
\textit{logika sed ankaù tre oportuna.} \VUgras{Fratiro} ekzemple signifus tiam\nl
speciale fraton, kaj \VUgras{fratino} speciale fratinon, dum \VUgras{frato} signifus\nl
simple infanon de tiuj samaj gepatroj (au frato, au fratino). »

Do ni povas gratular ni (kun legitima fiereso!) ke ni havas\nl
l'aprobo mem di la « kara Maestro », quan la « fideluloj » blinde\nl
veneracas. Ti qui objecionas, ke ni riprenas nun la sufixo \textit{-ul}, quan\nl
ni tante kritikis e mokis en Esperanto, ne atencas, ke ni nule\parplein
}

\VUcontinue
prehistoriala tempi, en qui la homulo dominacis per la vio\cc
lento e teroro lua kompatinda kompanino, o plu juste lua\nl
sklavino », remarkigis S\textsuperscript{ro} \textsc{Couturat} (1). Ni liberigas ni \textit{en}\nl
\textit{Ido}, de ta tiranatra prejudiko e de ta obsedanta konsidero\nl
di la genri, qua imprimis a nia omna lingui neefacebla\nl
karaktero (2). Nam ni repetas lo ed insistas pri lo : la nova\nl
sistemo ne tendencas impozar a ni la konstanta e penoza\nl
distingo di la genri e konseque la selekto inter \textit{-ul} e \textit{-in},\nl
ma, inverse, dispensar ni de ta neutila e jenanta konsidero\nl
o selekto, en omna kazi en qui ol ne esas necesa ed esencala.\nl
La linguistiko montras, ke l'evoluciono di la lingui ten\cc
dencas sempre simpligar la gramatikala kategorii : or inter\nl
li la minim necesa ed esencala esas certe ta di la genro.\nl
Nu, Ido anticipas nun la rezultajo finala di ca evoluciono,\nl
ed atingis per un stroko la sola solvo logikala e definitiva.\nl
Ico konstitucas granda pazo ad avane en l'evoluciono di la\nl
lingui, e decidigiva avantajo por nia linguo kompare a la\nl
ceteri, natural od artificala (3) .
\end{VUpage}

% -------------------------------------------------------------------
%  feuillet 196 — folio 192.
%  PARTICULARITES DE L'ORIGINAL, conservees : le meme mot cite est en
%  italique puis en romain dans la meme phrase ; et la particule
%  interrogative alterne entre deux formes sur la meme page.
% -------------------------------------------------------------------
\begin{VUpage}[196]{192}
\VUnotes{136.06mm}{%
\VUcontinue
kritikis lua formo, ma lua \textit{absurda uzo} : \textit{blindulo} ne esas persono\nl
« karakterizata da \textit{blind} », ma « qua esas blind », do \textit{blindo}. La\nl
sufixo \textit{-ul} signifikis nulo en Esp., ma por ni ol signifikas \textit{maskulo} :\nl
do kande ni dicas \textit{blindulo}, \textit{richulo}, ni donas a ca vorti senco quan\nl
li ne havis en Esp. (opozante li a \textit{blindino}, \textit{richino}). Por komprenar\nl
la difero, suficas komparar la feminala nomi : Esp. \textit{blindulino}, \textit{vir}\cc
\textit{gulino} ad Ido : \textit{blindino}, \textit{virgino}. (\textit{Progreso}, VI, 593.)

(1) Doktoro Zamenhof uzis ipsa plurfoye \textit{homino}, vice \textit{virino}, to\nl
quo supozas en \textit{homo} la duopla valoro segun qua ni rezonas.%
}

Ni bone remarkez, ke la konservo dil maskulismo en Ido\nl
esabus kontrea al \textit{principo dil unasenceso}. Nam, se \textit{homo}\nl
opozesas a \textit{homino}, to venas evidente de ke ol havas du senci :\nl
1\textsuperscript{e} \textit{homa ento}. 2\textsuperscript{e} \textit{homa ento maskula}. Ta du senci, de ube\nl
ol recevas li? Certe ne de la dezinenco \textit{-o}, simpla karakteri\cc
zivo di la substantivo e quan ni ritrovas en hom-\textit{in}-o (homo\nl
femina) (1). Pro ke nul altra marko uzesis por \textit{homo} (kon\cc
traste a hom-\textit{in}-o), la kauzo esas evidente, ke on donas a \textit{homo}\nl
la du senci supere indikita por la radiko \textit{hom}; on do kon\cc
seque atribuas reale du valori ad ica : ol reprezentas sam\cc
tempe la homo \textit{speco} e la homo \textit{maskulo}, ek qua on facas\nl
\textit{homino}. Do tal procedo esas nekontesteble la maskulismo en\nl
sua kompleta nelogikaleso. Esperanto havas lu, e Ido havus lu\nl
ankore (sen la sufixo \textit{-ul} dil maskuli) por omna nomi di\nl
enti anmoza.

Altraparte ka la raciono komuna ne dicas ke, se la radiko\nl
ne povas sen adjunto, expresar la ento femina, ol ne povas\nl
plu multe, sen adjunto, expresar la ento maskula.

Pluse, \textit{homo} signifikante homo maskula, kad on ne dar\cc
fas dicar, ke on obtenas de olu « androgino, hermafrodito »,\nl
kande on derivas \textit{homino} de olu? Nam lore on unionas lua\nl
feminal valoro a la maskulal valoro di \textit{homo}.

Ma pro quo la du sufixi \textit{-ul} e \textit{-in} por markizar la sexuo,\nl
prefere kam dezinenci? Pro du motivi :

1\textsuperscript{e}. Genro e sexuo, quale ni vidis, nule havas l'importo\nl
quan ni inklinesas atribuar ad oli. Do esas plu racionala\nl
expresar per sufixo la sexual ideo qua adjuntesas a la radiko\nl
di speco e rezervar al dezinenci ofico plu importanta. Nam,\nl
se la formi gramatikala havas kom esencal skopo « indikar\nl
la rolo dil vorti en la frazo » (ka la sexuo esas rolo?), lia\nl
uzo maxim utila certe esas karakterizar la diversa speci di\parplein
\end{VUpage}

% -------------------------------------------------------------------
%  feuillet 197 — folio 193. Premiere page d'un cahier : elle porte la
%  SIGNATURE « 7 » sous le bloc de notes, a gauche. Encre relevee de
%  y 1882 a 1908 px, x 95 a 111, soit un renfoncement de +38 px — le
%  meme que les alineas de note.
%  Le releveur signale une lecture qui ne se voit qu'a l'agrandissement :
%  un mot compose MOITIE ROMAIN MOITIE ITALIQUE, que la taille reelle
%  donne pour tout italique.
%  Il signale aussi que les « ind » de la geometrie sont trompeurs sur
%  les lignes ouvertes par un cadratin : le detecteur accroche la
%  capitale et non le tiret.
% -------------------------------------------------------------------
\begin{VUpage}[197]{193}
\VUsignature{170.01mm}{7}
\VUnotes{136.31mm}{%
(1) Reale la Latina havis por sua substantivi nur dezinenci\nl
kazala. Ol nultempe havis sexual marki por ta vorti.

(2) Ni memorigez, ke la Provencala havas multega \textit{feminala} nomi\nl
ed adjektivi finanta per \textit{-o}. E kad l'Italiana ne dicas : \textit{la mia mano}\nl
(mea manuo)?

Ne esas sen utileso remarkigar, ke povas esar nur \textit{neutra} la\nl
substantivo qua havas nek la karakterizivo dil maskuli \textit{-ul}, nek la\nl
karakterizivo dil femini \textit{-in}, se per naturo, quale \textit{matro}, o mestiere\nl
quale \textit{maristo}, lua sexuo ne indikesas.%
}

\VUcontinue
vorti (substantivo, adjektivo, verbo, e. c.), nam to esas\nl
l'indiko maxim utila por la « konstrukto » di la frazo e\nl
por lua kompreno quika e certa.

2\textsuperscript{e} Qua dezinencin on adoptus?

--- Nu, \textit{-a}, exemple por la feminal nomi.

--- Ne pro la Latina, me pensas; nam me trovas en lu\nl
sennombra neutri plurala ed anke multa maskuli : \textit{conviva},\nl
\textit{geometra}, \textit{collega}, \textit{agricola}, \textit{homicida}, \textit{poeta}, e. c. (1). Ido\nl
multe plu preferas la dezinenco \textit{-a} por l'adjektivo, pro ke\nl
ta vokalo finas la plumulto dil feminal \textit{adjektivi} en la\nl
Latina, l'Italiana, e. c., e mult adjektivi en la Sueda. Ol\nl
valoras tre certe plu multe kam \textit{-i} (marko di pluralo en\nl
tanta lingui) e quan malgre co ula sistemi adoptis kom\nl
karakterizivo dil adjektivo.

--- Kad on ne povus adoptar \textit{o} por la maskulalo?

--- Co esus certe posibla. Ma, pro la dicita motivi, esas\nl
preferinda atribuar specal sufixo al nomi maskulala, quale\nl
on agas (per \textit{-in}) por la nomi feminala, ed atribuar \textit{-o}, kom\nl
dezinenco pure gramatikala, ad \textit{omna} substantivo (2).

--- Ma \textit{-e} por la substantivi indikanta kozo, \textit{scientie},\nl
exemple?

--- Se me ne eroras, scien\textit{tia} impozesas da la Latina por\nl
la fideli di ca linguo. Omnakaze, substantivo esas substan\cc
tivo, sive ol indikas ento anmoza, sive ol indikas kozo, e lua\nl
senco ipsa impedas omna konfundo. Sive me uzos \textit{scientia},\nl
\textit{scientie} o mem \textit{cienco}, nulu konfundos la vorto ad ento\nl
anmoza. Lore pro quo du etiketi, du karakterizivi, kande\nl
un sola tote suficas? Pluse, kad la vokalo \textit{-e} ne konvenas\nl
perfekte por l'adverbo? \textit{Forte}, \textit{dolce}, tante konocata kom\nl
muzikal vorti, donas quik la valoro di \textit{-e} e semblas adminime\nl
tam bona kam \textit{fortiter}, \textit{fortim}, \textit{fortemen}, \textit{forti}, e. c.
\end{VUpage}

% -------------------------------------------------------------------
%  feuillet 198 — folio 194.
%  LIGNE INVISIBLE RETROUVEE, et c'est LA PREMIERE DU CORPS — le cas le
%  plus consequent du volume, puisqu'une premiere ligne absente decale
%  toute la page. Elle est courte (162 px), ouverte par un cadratin, et
%  forme a elle seule un alinea. Aucun depistage ne l'avait signalee :
%  le releveur l'a trouvee parce que son compte de lignes ne tombait pas
%  juste. Verifiee a l'agrandissement x 6 par l'assembleur.
%  PARTICULARITE DE L'ORIGINAL, conservee : une virgule la ou la phrase
%  strictement parallele, quatre lignes plus haut, porte un point-virgule.
% -------------------------------------------------------------------
\begin{VUpage}[198]{194}
\VUnotes{147.15mm}{%
(1) Ni pasante remarkigez, ke la Franca ed altra lingui ne povas\nl
simple e sen ambigueso expresar ica ideo : \textit{Mea filii ne plus esas}\nl
\textit{infanti.}%
}

--- E \textit{-u}?

--- Quaresma vokalo (ek 5) nur por la substantivo! Kad\nl
ico ne esas troa? Vice facar ek olu la marko di irga genro\nl
(segun ula autoro) o la signo dil akuzativo (segun altra),\nl
ka ne esas plu bona uzar \textit{-u} por markizar la pronomi indi\cc
viduala : \textit{ulu}, \textit{nulu}, \textit{omnu}, e. c.?

Ni do lasez la dezinenci a lia vera rolo, precizigita supere.\nl
La sexuo esas gramatikala, en ula lingui, nur misuze; ni\nl
omna sentas lo ed ofte pruvas per la vorti « maskula » e\nl
« femina », quin ni tante freque juntas al nomo di speco,\nl
vice o apud dezinenci partikulara. Se ni kelke reflektos, ni\nl
asentos, ke esar \textit{maskula} od esar \textit{femina} ne plu multe klasi\cc
fikas nomo di animalo en kategorio gramatikala kam esar\nl
\textit{yuna} od esar \textit{olda}. To esas nociono partikulara qua devas esar\nl
expresata per vorto partikulara, sive kompleta (\textit{maskulo},\nl
\textit{femino}) sive abreviita en sufixo (\textit{-ulo}, \textit{-ino}). Certe min racio\cc
nala e min praktikala esas, por expresar ol dezinence, uzar\nl
4 vokali ek 5, quale propozesis, precipue se li esas multe plu\nl
necesa od utila por vera kategorii gramatikala.

Rezume, Ido solvas la questiono dil genro ed emendas la\nl
maskulismo di nia lingui en la maniero sequanta :

1\textsuperscript{e} Radiko di speco, konvenanta egale bone por la mas\cc
kuli kam por la femini : \textit{elefanto} (maskula o femina), \textit{homo}\nl
(ento homa, sive maskula, sive femina), \textit{filio} (la genitito sive\nl
maskula, sive femina) (1).

2\textsuperscript{e} Sufixo \textit{-ul} preiranta la dezinenco ed indikanta la\nl
sexuo maskulala, se oportas enuncar olu : \textit{elefant-ulo} (ele\cc
fanto maskula); \textit{hom-ulo} (homa ento maskula); \textit{fili-ulo} (geni\cc
tito maskula).

3\textsuperscript{e} Sufixo \textit{-in} preiranta la dezinenco ed indikanta la sexuo\nl
feminala, se oportas enuncar olu : \textit{elefant-ino} (elefanto\nl
femina); \textit{hom-ino} (homa ento femina), \textit{fili-ino} (genitito\nl
femina).

Kompreneble, se altra vorto, exemple pronomo personala,\parplein
\end{VUpage}

% -------------------------------------------------------------------
%  feuillet 199 — folio 195. OUVERTURE DE L'APPENDICE 4.
%  Ornement (filet de 96 px = 8,13 mm, centre), titre en capitales
%  demi-grasses, sous-titre italique entre parentheses.
%  DEUX LIGNES INVISIBLES RETROUVEES : « tikala. » (fin d'alinea, avant
%  l'ornement) et « homino. » (fin de la note 3).
%  Le releveur a tranche les QUATRE signalements de la page : la ligne 5
%  cachait a la fois une ligne manquante ET l'ornement ; la ligne 7 est
%  le blanc d'apres-titre ; la ligne 18 est le filet des notes.
%  PARTICULARITES DE L'ORIGINAL, conservees : un proverbe francais et
%  une citation latine composes en ROMAIN, quand le volume italise
%  partout ailleurs ce qui est cite d'une autre langue ; et une
%  metathese dans un mot repete correctement sur la meme ligne.
% -------------------------------------------------------------------
\begin{VUpage}[199]{195}
\VUnotes{107.74mm}{%
(1) E ne \textit{infantulo}, pro ke la sexuo ja indikesas en \textit{filiulo}.

(2) E ne \textit{sekretariino}, pro ke la sexuo ja indikesas dal pronomo \textit{el}.

(3) Til nun la \textit{ministri}, la \textit{soldati} ed altra profesionani apartenas\nl
al forta sexuo; do nul konfundo esas timebla e \textit{ministro}, \textit{soldato}\nl
suficas. Ma supozez, ke divenus necesa indikar la sexuo, nu, ni\nl
havus : \textit{ministrulo}, \textit{soldatulo} e \textit{ministrino}, \textit{soldatino}.

La Franca proverbo : « La paresse est la mère de tous les vices »\nl
implikas nula sexuo e \textit{mère} hike = genitanto, patro (sive maskula,\nl
sive femina). Konseque on tradukas : \textit{L'ociereso esas la genitanto}\nl
(o : \textit{la patro}) \textit{di omna vicii}. Genitantulo, genintantino, o patrulo,\nl
patrino esus hike groteska, nam \textit{ociereso} esas nek homulo, nek\nl
homino.

(4) E vere, F. \VUgras{grand} ne equivalas \textit{grandeur}, \VUgras{grandeso}, nek \VUgras{pauvre}\nl
esas identa a \textit{pauvreté}, \VUgras{povreso}. Li do ne indikas la qualeso ipsa, ma\nl
indikas, nomizas la ento karakterizata da olu : \textit{la povra} o \textit{povro}, \textit{la}\nl
\textit{granda} o \textit{grando}. Nam, dicis Leibniz : « Discrimen adjectivi et\nl
substantivi in lingua rationali non est magni momenti. » \textit{Granda} e\nl
\textit{grando}, \textit{povra} e \textit{povro} fakte nomizas ento konkreta e diferas nur\nl
gramatikale.%
}
\VUornement[8.13mm]{45.17mm}

\VUcontinue
ja indikas la sexuo, la radiko di speco suficas : \textit{mea filiulo}\nl
\textit{ne plus esas infanto} (1); \textit{el esas sekretario} (2) \textit{di ministro} (3).\parplein

Semblas a ni, ke ta sistemo sustenesas da la raciono, e\nl
ke ol havas la merito esar tam logikala e simpla kam prak\cc
tikala.

\VUsaut{11.17mm}
\VUtitre{10.63pt}{-21}{SUBSTANTIVIGO DIL ADJEKTIVO.}
\VUsaut{2.29mm}
\VUsoustitre{6.02pt}{185}{\textit{(Apendico 4-ma.)}}
\VUsaut{4.30mm}

Unesme importas precizigar bone la justa signifiko dil\nl
vorto « adjektivo » (qualifikanta). Yen to, quon skribis\nl
pri ta punto la ciencoza filozofiisto L. \textsc{Couturat}, nia regretata\nl
pioniro :

« Adjektivo (plu precize : adjektival radiko) tote ne\nl
indikas la abstraktita qualeso, ma la konkreta ento qua\nl
« posedas » ta qualeso, e qua tam multe diferas de ta qualeso,\nl
kam la habitanto diferas de sua lando o la mestieristo\nl
de sua mestiero (4). La radiko \VUgras{blind} expresas ulo \VUgras{blinda},\nl
e ne la blindeso; konseque \VUgras{blindo} = \VUgras{ento} (praktike homo)\nl
\VUgras{blinda}, e tote ne la \VUgras{blindeso}. Same \VUgras{belo}, \VUgras{bono} ne povas\parplein
\end{VUpage}

% -------------------------------------------------------------------
%  feuillet 200 — folio 196.
%  LIGNE INVISIBLE RETROUVEE : « vorti : », large de 110 px, fin
%  d'alinea. Ce n'etait pas l'exception connue du verdict « double
%  net » : il n'y a pas de titre a cet endroit.
%  Le releveur signale trois fausses lectures de graisse evitees a
%  l'agrandissement, toutes dans le meme sens : la bande faisait
%  paraitre gras des mots romains.
% -------------------------------------------------------------------
\begin{VUpage}[200]{196}
\VUnotes{120.31mm}{%
(1) \textit{Progreso}, I, 389.

(2) Do on substitucus un konvenciono ad altra, nam, per su \VUgras{-o}\nl
indikas nur substantivo, quale \VUgras{-a} indikas nur adjektivo. E kad ni ne\nl
devus serchar anke marko specala por la substantivi (nekontebla) di\nl
ago, di stando, quale \VUgras{amo, sufro, emoco, laboro, angoro, anxio,} e. c.\nl
Pro quo nur la \VUgras{kozi,} o la \VUgras{homi,} e ne la \VUgras{agi} e la \VUgras{standi} havus la pri\cc
vilejo di marko propra e tote specala? Se on alegos la logiko, me\nl
sen hezito respondos : Por ni la logiko devas esar \VUgras{moyeno, ne skopo.}\nl
Icon me ja dicis plurfoye. Kun superlogiko, forsan ni obtenus linguo\nl
plu logikoza. Ma me tote ne esas konvinkita, ke ol ne esus tro\nl
komplikita e neutile min simpla por la maxim multi. Or ni ne\nl
laboras por la logikozi unesme, e mem min multe por la super\cc
logikozi, ma por omna posibla uzonti di la helpolinguo. (L. B.)

(3) \textit{Progreso}, I, 565.%
}

\VUcontinue
indikar la \VUgras{beleso,} la \VUgras{boneso} (quale se la sufixo \VUgras{es,} o \VUgras{ec,}\nl
havus nula senco e nula valoro), nek felico la feliceso, nek\nl
fiero la fiereso, e. c., e. c., ma nur la enti, qui posedas ta\nl
qualesi. » (1).

Reale la substituco di o a a substantivigas la adjektivo,\nl
ma ol ne homigas lu. Icon agnoskis S\textsuperscript{ro} \textsc{Couturat} per ica\nl
vorti :

« Certe, \VUgras{substantivigar} adjektivo ne esas necese \VUgras{personigar}\nl
ol. \VUgras{Belo} nur signifikas « bela ento » (persono o kozo indi\cc
ferente), nam la finalo \VUgras{-o} di la substantivo aplikesas tote\nl
egale a personi ed a kozi. Ma hike aparas motivo, \VUgras{ne plus}\nl
\VUgras{di logiko,} ma di simpleso : se on bezonas distingar la kazo\nl
di persono e la kazo di kozo, on havas ja la sufixo \VUgras{-aj} por\nl
la kozi : do ne esas necesa uzar altra sufixo : to esas nur\nl
praktikal konvenciono di « komuna raciono ». Remarkez\nl
bone, ke la substantivigo di adjektivo ne implikas plu l'ideo\nl
di kozo kam l'ideo di persono; \VUgras{belo} vice \VUgras{belajo} esus tam\nl
konvencionala kam \VUgras{belo} vice \VUgras{bela viro} (2). Ma l'ideo ipsa\nl
dil adjektivo maxim ofte implikas, sive ideo di persono,\nl
sive ideo di kozo : \VUgras{acido, dezerto, vakuo,} e. c., povas (natu\cc
rale) indikar nur kozi; \VUgras{avaro, blindo, malado,} e. c. povas\nl
indikar nur personi (o vivanti). Do en omna kazi en qui\nl
povas aparar praktike nula dusenceso, on darfas substan\cc
tivigar \textit{direte} l'adjektivo : to esas simpligo e sparado di neu\cc
tila sufixo. Nur en kazo di reala dubo on uzos \VUgras{-aj} por dis\cc
tingar la kozi (3).
\end{VUpage}

% -------------------------------------------------------------------
%  feuillet 201 — folio 197.
%  Le releveur a corrige DEUX erreurs de son propre bloc avant de le
%  rendre — une petite capitale posee a tort et un groupe de gras
%  inutilement coupe — plutot que de les laisser passer.
%  COQUILLE DE L'ORIGINAL, conservee : deux guillemets ouvrants de
%  suite, dont aucun n'est ferme.
% -------------------------------------------------------------------
\begin{VUpage}[201]{197}
\VUnotes{146.56mm}{%
(1) Videz ye la « Pronomo « lo ». --- \VUgras{Lo} esas \VUgras{pronomo} per sua\nl
senco intima : \textit{to quo esas...}

(2) Hike on povas remarkar, ke persono anke esas \VUgras{nutrivo} che\nl
la kanibali; ke persono povas esar \VUgras{konduktivo} (ex. en elektrala\nl
cirkuito), e. c.%
}

Ni remarkigas, ke « \VUgras{le beau, le vrai} », e. c., \VUgras{en su}, segun\nl
la koncepto di \textsc{Platon} ed \textsc{Aristoteles} (qua darfas esar\nl
expresata ankore per \VUgras{la belajo, la verajo}) nun esas plu ofte\nl
expresata dal Idisti per la pronomo \VUgras{lo} e l'adjektivi : \VUgras{lo bela,}\nl
\VUgras{lo vera,} e. c. (1).

Altraloke S\textsuperscript{ro} \textsc{Couturat} dicabis :

« L'ideo di \VUgras{persono} o di \VUgras{kozo} esas nur implicita, tacite\nl
o spontane adjuntita; kompreneble ol impedas la strikta\nl
renversebleso, t. e. devas desaparar kande on rivenas de la\nl
substantivo a la adjektivo. Ma esus same arbitriala, same\nl
kontrea a la strikta renversebleso, pozar : \VUgras{belo} = \VUgras{bela kozo}\nl
(\textit{belajo}) kam pozar : \VUgras{belo} = \VUgras{bela homo o persono}. La du\nl
esas egale kontre la logiko, o plu juste, exter la logiko.\nl
Segun la strikta logiko, on devas tacar nulo, ed expresar\nl
omna elementi di l'ideo : l'ideo di persono, per la vorto\nl
\VUgras{homo}; l'ideo di \VUgras{kozo}, per la vorto \VUgras{kozo} (o per la sufixo \VUgras{-ajo}).\nl
Apene on bezonus alterar \VUgras{hom} en \VUgras{-om} (sufixo propozita da\nl
S\textsuperscript{ro} \textsc{de Janko}), nam, danke la posibleso di eliziono, \VUgras{bela}\nl
\VUgras{homo} = \VUgras{bel' homo}. Rezume, segun strikta logiko, on devus\nl
sempre uzar sufixo, sive por la personi, sive por \VUgras{kozi}...

« Pro quo ni ne adoptis tala sufixo? Simple pro ke prak\cc
tike ol ne esas necesa. Nam, pro ke on havas sufixo specala\nl
por indikar la kozi, nome \VUgras{-aj}, on ne bezonas sufixo specala\nl
por indikar la \VUgras{personi} : l'absenteso di sufixo suficas por to\nl
(ne obliviez, ke la questiono ne plus esas di strikta logiko,\nl
sed (ma) di praktikal komodeso...

« Mem la sufixo \VUgras{-aj}, por la kozi, devas uzesar nur en\nl
kazo di bezono, t. e. nur kande l'adjektivo substantivigita\nl
povus esar dusenca, e signifikar persono. Nam existas tre\nl
multa vorti, precipue en ciencal e teknikala linguo, qui\nl
esas adjektivi substantivigita ed indikas kozi : \VUgras{astringivo,}\nl
\VUgras{nutrivo, konduktivo} (2), e. c., esas esence ed evidente nur\nl
kozi; esas do tote neutila adjuntar a li la sufixo \VUgras{-aj}. Same\parplein
\end{VUpage}

% -------------------------------------------------------------------
%  feuillet 202 — folio 198.
%  PARTICULARITE DE L'ORIGINAL, conservee : sur la meme page, la meme
%  desinence citee est tantot en gras, tantot en romain — le releveur
%  l'a tranchee au grossissement 10 en comparant les delies.
% -------------------------------------------------------------------
\begin{VUpage}[202]{198}
\VUnotes{135.38mm}{%
(1) Kad tala vorti esus per \VUgras{-o} (substitucita a \VUgras{-a}) substantivi \textit{Aris}\cc
\textit{totelala}, e povus kom tala uzesar da la cienco? (L. B.). --- S\textsuperscript{ro} \textsc{Cou}\cc
\textsc{turat} tre remarkinda specalisto pri filozofio, skribis tre vere :\nl
« Kande Platon (ante Aristoteles) volis expresar \VUgras{l'ideo}, \VUgras{l'esenco}\nl
di la bela kozi (belaji), il uzis l'expresuro : \VUgras{la bela en su}. Ma to\nl
esas nur altra nomo di la beleso (kun la cirkonstanco filozo\cc
fiala, e ne linguala, ke il konceptis la \VUgras{beleso} kom existanta « en su »,\nl
exter la belaji, do exter la \VUgras{beleso} di ca o ta belaji)...

(2) \textit{Progreso}, I, 555-56.%
}

\VUcontinue
en matematiki on parolas pri \VUgras{kompozanto, rezultanto, deter}\cc
\VUgras{minanto, varianto}... (1).

« Nule on bezonas adjuntar -aj a tala vorti, tante plu ke\nl
li ne indikas propre kozi plu kam personi... (On darfas tre\nl
bone dicar \VUgras{konsequanto}, nam ta vorto ne povas indikar\nl
homo, ma nur propoziciono o fakto; e. c.) » (2).

Nun kande ni expozis la questiono tal quala nia Aka\cc
demio mantenis ol, pos diskuti e studiado, ni konocigos\nl
propozo, quan facis, la unesma, ed unesme en formo di\nl
questiono, S\textsuperscript{ro} \textsc{de Janko}, qua mortis dum la milito pos esir\nl
vice-sekretario dil Akademio :

« Pro ke la pronomi : \VUgras{unu, ulu, altru, nulu, omnu, irgu,}\nl
omni signifikas homi, kad on povus (darfus) analoge,\nl
aplikar l'afixo -u a l'adjektivi por donar ad li senco simila?\nl
exemple \VUgras{la blondu} --- la blonda homo, \VUgras{la modestu} --- la\nl
modesta homo; e. c.? » (\textit{Progreso}, I, 307.)

S\textsuperscript{ro} \textsc{Couturat} respondis : « Ca sugesto esas interesanta e\nl
meritas studio. Nun on uzas en ca senco simple la substan\cc
tivo : \VUgras{blondo, modesto}, e. c. E co havas nula malavantajo\nl
(desavantajo), nam, kande on parolas pri kozi, on uzas la\nl
sufixo \VUgras{-aj}, specala por ca kazi. »

Ni dicez quik, ke la propozo di S\textsuperscript{ro} \textsc{de Janko} rezultis de\nl
fundamental eroro pri la valoro dil dezinenco -u. Il ipsa\nl
agnoskis lo pose.

La dezinenco -u selektesis da « Ido » (me) en la reformo\cc
plano, \textit{tote ne} por homigar \VUgras{un, ul, altr, nul, omn, irg}, ma\nl
unike por pronomigar li, quo esas tre diferanta. Mem me\nl
propozabis \VUgras{icu, istu}, \textit{tote ne} por homigar \VUgras{ic, ist} (diveninta\nl
\VUgras{it}), ma por pronomigar li, per la substituco di -u al dezi\cc
nenco adjektivala (a) di \VUgras{ica, ista}. La -u esis ed esas ankore\nl
nur pronomal dezinenco, quale -a esis ed esas ankore\nl
\end{VUpage}

\begin{VUpage}[203]{199}
\VUnotes{140.72mm}{%
(1) Forsan ta konstato inspiris a lu la penso ofrar pose \VUgras{-u} kom\nl
specala karakterizivo dil prepozicioni.

(2) \textit{Progreso}, III, 278.

(3) E la agi? (Videz pag. 196, noto 2) Se on respondos, ke la\nl
senco ipsa dil radiko indikas, ka la vorto expresas ago substantive,\nl
me replikos, ke la senco dil radiko indikas, ka la vorto expresas homo\nl
o kozo sammaniere.%
}

\VUcontinue
adjektival dezinenco. -u formacas individuala, ma ne homala\nl
\textit{pronomi}. Kad ankore nun ni ne uzas omnainstante l'unu,\nl
l'altru, kom \textit{pronomi, pri enti nule homa, mem nule bestia?}\parplein

Cetere S\textsuperscript{ro} \textsc{Couturat} dicis : « Ni memorigas insiste, ke la\nl
finalo- u nule indikas la personi, ma karakterizas la \VUgras{pro}\cc
\VUgras{nomi} (\textit{Progreso}, II, 602).

Sualatere S\textsuperscript{ro} \textsc{de Janko} skribis : « Til nun semblis a me,\nl
ke la finalo -u esis institucata por indikar personeso; sed\nl
(ma) nun aparas, ke ol indikas simple pronomo (1).

Semblas do, ke se on komprenabus bone la valoro pro\cc
nomala di -u, on ne pensabus ad olu por homigar adjektivi,\nl
o plu juste adjektival radiki.

Tamen, pos dicir : « Sed (ma) en ta rolo (di pronomo)\nl
ol semblas a me tote ne utila, nam omnu vidas, kad la pro\cc
nomo esas o ne sequata da nomo » (ma nultempe la\nl
pronomo sequesas da nomo, pro ke ol \textit{remplasas} nomo),\nl
S\textsuperscript{ro} \textsc{de Janko} mantenis e mem generaligis sua propozo pri la\nl
personigo (homigo) dil adjektivo dicante :

« Uzez la finalo -u por personigar omna adjektivi indi\cc
ferente : \textquotesingle{} ulu venis. Quu venis? Icu esas genio. Anke la bonu\nl
pekas. L'amantu esas blinda \textquotesingle{} » (2).

Pri l'unesma propozo (videz p. 198) S\textsuperscript{ro} \textsc{Couturat} res\cc
pondis :

« La finalo -u indikas personi (exakte : \textit{individui}) nur\nl
en la pronomi, ne en la substantivi. Teorie, on devus adoptar\nl
ed uzar du sufixi, un por la personi, altra por kozi (3);\nl
ma praktike, un sola suficas por distingar la du kazi, e mem\nl
ol esas necesa nur en kazi di posibla dusenceso.

« Lasta remarko : en la pronomi la plurala formo (per \VUgras{-i})\nl
sempre korespondas a la singulara per \VUgras{-u}, nam la singularo\nl
per \VUgras{-o} indikanta nedeterminita kozo, ne povas havar pluralo\nl
(se on vizas plura kozi, ta kozi esas pro to ipsa determi\cc
\end{VUpage}

\begin{VUpage}[204]{200}
\VUnotes{150.03mm}{%
(1) \textit{Progreso}, I, 556-557.

(2) Me repetas, ke \VUgras{-u} indikas tote ne direte \VUgras{homi}, ma \VUgras{individui},\nl
quo esas tre diferanta.

(3) \textit{Progreso}, I, 554.%
}

\VUcontinue
nita, e konseque on devas dicar : \VUgras{ta kozi, multa kozi, altra}\nl
\VUgras{kozi,} e. c., e ne simple, \VUgras{to, multo, altro}...).

Ma, se on adoptus la propozo di S\textsuperscript{ro} \textsc{de Janko}, la formo\nl
per \VUgras{-o} indikanta determinita kozi, povus e devus havar plu\cc
ralo, qua esus necese la formo per \VUgras{-i}, segun la generala\nl
regulo : \VUgras{bono, boni} (L. \VUgras{bonum, bona}); ed on devus inventar\nl
altra pluralo por la formo per \VUgras{-u}. Ma to ipsa montras, ke\nl
ta formo ne esas admisebla, e ke on devus uzar specala\nl
sufixo konservenda en pluralo; ex. : \VUgras{bonomo, bonomi}. La\nl
rezultajo esus nur renversar la nuna regulo, t. e. uzar ula\nl
sufixo por personi, e nula sufixo por la kozi, quo esus « ar\cc
bitriala » o konvencionala. On dicus \VUgras{bonomo, bonomi}, vice\nl
nia \VUgras{bono, boni}, e \VUgras{bono, boni} vice nia \VUgras{bonaji} (1).

« Konseque, la distingo di la personi (per \VUgras{-u}) (2) esas\nl
utila nur en la pronomi, e ne en la substantivi nek en la\nl
adjektivi; nam la senco ipsa di substantivo indikas persono\nl
o kozo, ed adjektivo sempre referas a substantivo expresita\nl
o tacita. Cetere ta distingo esas nule logikal, sed (ma) pure\nl
praktikala : la logiko ne distingas personi o kozi, omni esas\nl
logike « objekti » od « enti ». Pro to, la finalo \VUgras{-o} di la\nl
substantivi ne signifikas plu kozo kam persono, sed (ma)\nl
simple « \VUgras{objekto} od « ento » (3).

Pri la duesma propozo di S\textsuperscript{ro} \textsc{de Janko}. (Vid. p. 199.) S\textsuperscript{ro} \textsc{Cou}\cc
\textsc{turat} respondis « ... l'uzado di \VUgras{-u} quale (kom) finalo dil\nl
adjektivi substantivigita (personigita) konstrastus krude kun\nl
la finalo \VUgras{-o} di la propra (primitiva) substantivi, quan il\nl
intencas konservar. Videz la frazo, quan il citis quale (kom)\nl
exemplo (\textit{Progreso}, II, 499), ube on trovas intermixita :\nl
\VUgras{malsaju, intrigantu, koketinu, kun kurtano, kompano}. L'ana\cc
logeso instigus nerezisteble la parolanto dicar, automate :\nl
\VUgras{kurtanu, kompanu.}

« La solvo maxim simpla semblas esar, tote kontree, la\nl
totala supreso di la finalo \VUgras{-u} en determiniva pronomi. Quale\nl
S\textsuperscript{ro} \textsc{de Janko} ipsa remarkas, on ne uzas ol en la demonstra\cc
tiva ed questionala-relativa pronomi e til nun aparis nula\nl
\end{VUpage}

\begin{VUpage}[205]{201}
\VUnotes{154.69mm}{%
(1) \textit{Progreso}, III, 279-280.

(2) Videz \textit{Progreso}, III, p. 280.%
}

\VUcontinue
detrimento. ` Qua venis? ' signifikas evidente : ` Qua per\cc
sono ', e se on volus dicar ` qua kozo ', on uzus, sive ca\nl
ipsa vorti, sive ` quo '. Same pri \VUgras{ica, ita}. Se do on volas res\cc
taurar perfekta analogeso, suficas supresar la formi en \VUgras{-u}\nl
o vicigar (remplasigar) oli en omna kazi per la formi en \VUgras{-a}.\nl
La exempli sube citata montras, ke nula dusenceso esas\nl
timenda pro ke omna pronomo vicas (remplasas) ula nomo,\nl
qua esas antee expresata o bone konocata, on savas sempre\nl
a qua speco de objekti ol aplikesas... « Logikale la substi\cc
tuco di la finalo \VUgras{-o} a la finalo \VUgras{-a} signifikas certe nur la\nl
substantivigo, do expresas nur l'ideo generala di objekto,\nl
sen indikar ka ta objekto esas persono o kozo (Vid. \textit{Pro}\cc
\textit{greso}, I, 555). \VUgras{La belo, la bono} indikas generale « to quo\nl
esas bela, bona ». La senco neutra o nedeterminita di ta\nl
expresuri korespondas perfekte a ta di la pronomi en \VUgras{-o} :\nl
\VUgras{ico, ito, omno, nulo, irgo,} e. c. Segun remarko da S\textsuperscript{ro} \textsc{de}\nl
\textsc{Janko} ipsa, ta senco generala kontenas egale personi e kozi :\nl
kande on dicas : ` \VUgras{omno} ', on komprenas anke : ` omna per\cc
soni '; kande on dicas : ` Nulo venis ', on komprenas anke :\nl
` nula persono '. Ma praktike, la senco interna di la adjek\cc
tivo substantivigita suficas por indikar, kad parolesas pri\nl
homo o pri kozo : \VUgras{la yusto} esas evidente persono e \VUgras{la dezerto}\nl
esas kozo. Same pri la participi substantivigita : \VUgras{la paro}\cc
\VUgras{lanto, la arestito} esas evidente homi; \VUgras{la konsequanto, la}\nl
\VUgras{dicito} esas kozi (plu exakte, pensi). Esus tote neutila, do\nl
pedantala, postular formala distingo di ta kazi per specala\nl
sufixi. Tamen ni havas la sufixo \VUgras{-ajo}, qua indikas sempre\nl
kozo; en la kazi dubebla (e nur en ta kazi) on devas uzar\nl
ol por indikar la kozi, e konseque lua absenteso indikas la\nl
personi. E se ico ne semblas suficanta (quo eventos tre rare),\nl
on esas sempre libera adjuntar a l'adjektivo la substantivi\nl
\VUgras{ento, homo, viro,} e. c. Li esos omnakaze plu klara e komoda\nl
kam specala sufixo, quan on inventus por ica funciono, e\nl
quan on riskus trouzar o misuzar, nam on uzus ol tre rare,\nl
se on uzus ol nur segun reala bezono. » (1).

Ca konsideri apogita sur tre multa exempli (2) granda\cc
\end{VUpage}

\begin{VUpage}[206]{202}
\VUnotes{155.03mm}{%
(1) \textit{Progreso}, III, 281.

(2) \VUgras{Vera salvo} en la texto (ideo \VUgras{salvar}),%
}

\VUcontinue
parte aprobesis da S\textsuperscript{ro} \textsc{de Janko} en repliko (1) ube on\nl
lektas :

« ... Do por esar klara e ne fidar a la kuntexto, semblas\nl
oportuna indikar specale la nedeterminita neutro. Altraparte\nl
me agnoskas la justeso dil objecioni facita en la Respondo\nl
kontre mea tilnuna propozo, ja plurfoye repetita, ke -u signi\cc
fikez persono ed -o la neutro. » Il adjuntis :

« Ula ecelanta konocanto di nia linguo recente remarkigis\nl
l'eventualajo permutar la -o ad -u, t. e. uzar la finalo -u\nl
(qua ne havas pluralo) por la nedeterminita neutro (qua\nl
ne bezonas pluralo). To semblas a me tre sagaca ideo e vera\nl
salvo (2) de ta desfacilaji. »

Konseque il propozis fine :

« Por substantivigo (signifikanta, segun la senco, per\cc
sono o kozo) uzez la finalo -o. Exempli ek la Respondo :\nl
\VUgras{la blondo, la yusto} (pers.) \VUgras{la dezerto} (kozo); me adjuntas :\nl
\VUgras{Irgo venis...} « Por la nedeterminita neutro uzez la finalo -u :\nl
\VUgras{La belu esas l'imajo di la bonu. Me prizas irga bonu} =\nl
\VUgras{irgu bona...} » E pos reflekti e propozo pri la pluralo dil\nl
adjektivo, il finis per : « Tale per malgranda (mikra)\nl
retusho ni advenos (arivos) a la maxim simpla e belsona\nl
formi propozita da S\textsuperscript{ro} \textsc{Couturat}, sen sakrifikar en kazi di\nl
bezono la precizeso di nia linguo. »

Yen do quala esis la \textit{definitiva} judiko e koncepto di\nl
S\textsuperscript{ro} \textsc{de Janko} pri ta punto gravega, ed on ne povas negar,\nl
ke ol esis injenioza, simpla, e chanjabus minime posible.

Ponderinte omna opinioni e propozi, l'Akademio decidis :\nl
« On repulsas la propozo substantivigar la qualifikiva adjek\cc
tivi sammaniere kam la determinivai (determinivi) (ex. :\nl
\VUgras{bonu} indikus persono, \VUgras{bono} indikus kozo o neutro), de\cc
cido 482, unanime per 8 voci. »

S\textsuperscript{ro} \textsc{de Janko}, quale omni, aceptis ta decido, qua fakte nur\nl
konfirmis la tillora uzado, duranta la substantivigo dil\nl
adjektivo tal quala establisabis lu la « Grammaire Com\cc
plète » aprobita dal Komisitaro Konstanta di la Delegitaro.\parplein

Ni adjuntez ica konsideri.
\end{VUpage}

\begin{VUpage}[207]{203}
\VUnotes{123.02mm}{%
(1) Pro lua graveso ed importo, ni ne povus tro lumizar pri ta\nl
punto dil substantivigo dil adjektivi, o dil adjektivigo dil substantivi.\nl
Quante plu on vidos olu plenlume, tante plu on prizos la decido\nl
motivizita solide dil Akademio, pos longa studiado. On do ne\nl
astonesez, ke ni donas ankore la sequanta konsideri. Ni extraktas li\nl
ek specal artiklo da S\textsuperscript{ro} \textsc{Couturat} « Pri l'adjektivo substantivigita »\nl
en \textit{Progreso}, IV, 84.

« On dicas a ni : Ni havas specala bonaji, e ni havas lia generala\nl
qualeso : \VUgras{boneso.} Ma ni havas anke la generala ideo di omno, quo\nl
esas bona; to ne esas bonajo, ma l'ensemblo o ideo di omna bonaji.\nl
Or ideo ne esas identa a specala kozo; singla bonajo esas un bonajo;\nl
ma on ne povas dicar \VUgras{la bonajo} generale. »

Me respondas, ke on povus aplikar la sama argumento ad omna\parplein
}

Vice un sola kategorio de substantivi, la procedo \VUgras{bono,}\nl
\VUgras{bonu} establisabus du kategorii, ica per \VUgras{-o}, ita per \VUgras{-u}.\nl
Omna substantivi per \VUgras{-o} nun havas sua pluralo per \VUgras{-i} :\nl
tablo, tabli; doloro, dolori. Ma, ecepte, segun altra pro\cc
pozo, omna substantivi per \VUgras{-o} venanta de adjektivo ne havus\nl
pluralo, o formacus ol per altra procedo : \VUgras{le, uli, kelki.}\nl
Ma se \VUgras{le, uli, kelki} ne esos uzebla? Ex. : purgivi tro\nl
frequa esas nociva, quo tre diferas de uli purgivo tro fre\cc
qua, e. c. \VUgras{Arabi, Franci, Italiani} = homi \VUgras{Franca, Araba,}\nl
\VUgras{Italiana.} Ma : \VUgras{Arabo, Franco, Italiano} egalesus : \VUgras{irgo Araba,}\nl
\VUgras{Franca, Italiana} (kavalo, kultelo, armo exemple); nur\nl
\VUgras{Arabu, Francu, Italianu} esus homo. Simile : \VUgras{parolanto, babi}\cc
\VUgras{lemo, amanto} ne esus \VUgras{homo,} ma \VUgras{irgo;} nur \VUgras{parolantu, babi}\cc
\VUgras{lemu, amantu} esus \VUgras{homo.} Nun, por direte substantivigar\nl
adjektivo, o adjektivigar substantivo, ni remplasigas \VUgras{-a}\nl
per \VUgras{-o}, \VUgras{-o} per \VUgras{-a} tote simple; de \VUgras{katolika} ni facas \VUgras{katoliko,}\nl
o de \VUgras{katoliko} ni facas \VUgras{katolika.} Ma kun la sistemo \VUgras{-o} = \VUgras{irgo,}\nl
\VUgras{-u} = \VUgras{homo,} oportus esar certa antee, kad la radiko \VUgras{katolik}\nl
esas adjektiva, o substantiva. Nam, se ol esas adjektiva, ni\nl
havos : \VUgras{katoliko} = \textit{irgo katolika} (exter homo), e \VUgras{katoliku} =\nl
nur \textit{homo katolika}. Ma, se ol esas substantiva, \VUgras{katoliko} esos\nl
nur \textit{homo katolika} (ne : irgo katolika). Ta distingi, qui ne\nl
mem havas la apogo di nia lingui, semblas vere tro super\cc
logika por helpolinguo destinata a spiriti di ordinara for\cc
teso e ne a superlogikozi.

Ni darfas joyar ke nia Akademio en sua decidi sucias nur\nl
l'unesmi (1).
\end{VUpage}

% -------------------------------------------------------------------
%  FEUILLET 208 — LA SEULE PAGE DU VOLUME ENTIEREMENT AU CORPS DES
%  NOTES. La note (1) ouverte au folio 203 deborde et prend le folio
%  204 tout entier : quarante-trois lignes. Trois mesures le montrent
%  et non le supposent — hauteur d'x 15-16 px sur les quarante-six
%  lignes (19-20 px au corps ordinaire) ; pas des lignes de base
%  33,76 px = 8,13 pt de bout en bout ; et le raccord de sens, la
%  derniere ligne du folio 203 s'achevant sur « ad omna » que la
%  premiere du 204 continue par « ideo abstraktita ».
%  C'est ce pas unique qui rendait « pas -1.00 » a la geometrie : le
%  detecteur mesurait y0, qui saute de dix pixels selon les hampes.
%
%  IL N'Y A PAS DE FILET. Cherche au pixel : aucune encre entre
%  y=1688 et y=1735, quand celui du feuillet 207 se leve net a 213 px.
%  Le compositeur n'avait rien dont il eut du distinguer sa note (*) :
%  la page etait deja tout entiere au corps des notes. D'ou
%  \VUnotes[0pt].  Ordonnee : le filet manquant serait a y=1685
%  (1736 moins les 51 px qui separent filet et premiere ligne de note
%  au feuillet 207), soit 24,30 + (1685-248) x 25,4/300 = 145,97 mm.
%
%  UNE LIGNE MANQUAIT A LA GEOMETRIE : « bela. », 65 px de large,
%  sous le seuil du detecteur — c'est la panne connue des lignes tres
%  courtes. Le compte reel est de 46 lignes, non 45.
% -------------------------------------------------------------------
\begin{VUpage}[208]{204}
\VUnotes[0pt]{145.97mm}{%
(*) Se on postulas, ke \VUgras{la bonajo} esez expresata altre kam \VUgras{la}\nl
\VUgras{bonaji,} lore on devos postular por \VUgras{la homo} altra nomo kam por\nl
\VUgras{la homi,} do \textit{duopligar} omna substantivi ed adjektivi.%
}
\VUpageNote{8.13pt}

\VUcontinue
ideo abstraktita, ex. : a la nociono \VUgras{homo.} « Ni havas aparta homi,\nl
ed anke la qualeso : \VUgras{homeso.} Ma ni havas anke la generala ideo di\nl
omna homi, to ne esas un homo, ma l'ensemblo o koncepto di omna\nl
homi. Singla homo esas ya un homo; ma ni ne povas parolar pri\nl
\VUgras{la homo generale.} » Tamen ni sempre parolas tale : ex. ni dicas :\nl
« la hundo esas l'amiko di la homo », sen vizar un aparta hundo nek\nl
un aparta homo, ma komprenante \VUgras{la hundo} e \VUgras{la homo} generale.

On objecionas : « Ideo ne esas identa ad aparta kozo ». Tre vera;\nl
tamen, to ne impedas, ke en omna nia lingui omna abstraktita idei\nl
indikesas per la sama vorto kam la singla individui, qui apartenas\nl
a li. To esas fakto linguistikala tre generala, e bone konocata da la\nl
logikisti, nam ta ambigueso di omna generala nomi esas fonto di\nl
multa erori e sofismi : ma ni ne povas supresar ol, nam nia tasko\nl
ne esas emendar la homal spirito (*). Do ni darfas e mem devas\nl
aplikar la sama regulo o kustumo a \VUgras{la bonajo, la belajo,} e. c.\nl
\VUgras{La belajo} esas l'ideo generala di omna belaji, exakte same kam\nl
\VUgras{la homo} esas l'ideo generala di omna homi. Dicar : « me amas la\nl
belajo » esas dicar : « me amas la (omna) belaji », same kam dicar :\nl
« Me odias la (omna) homi ». On ne trovos irga logikala difero inter\nl
ta du kazi. Remarkez, ke \VUgras{belajo,} en sua maxim generala senco,\nl
kontenas la bela enti kun la bela kozi, do anke la \VUgras{beli} (beluli) e la\nl
\VUgras{belini;} e pro to on darfas uzar \VUgras{belajo} mem pri \VUgras{belino.} Kande on\nl
parolas pri la \VUgras{belaji} di la naturo, on ne ekskluzas de li la bela enti,\nl
homi od animali! Do \VUgras{la belajo} esas la maxim generala nociono di\nl
« omno bela ».

« On dicos forsan, ke belajo esas kelkafoye ambigua o ne sat\nl
preciza : certe, same kam \VUgras{pneumatiko} citata supere. Ma anke \VUgras{belino}\nl
esas ambigua, nam ol povus, segun rigoro, aplikesar a bela kavalino\nl
same kam a bela virino (muliero). Tamen on savas sempre, a qua ta\nl
vorto aplikesas.

On facis objeciono preske kontrea : \VUgras{bonajo} ne povus tradukar ex.\nl
D. \VUgras{das Gute an einer Sache,} F. \VUgras{le bon de cette affaire} (est que...).\nl
Pro quo ne? On uzas la sufixo \VUgras{-ajo,} ne nur por kozo kompleta, ma\nl
anke por irga parto di kozo : me dicos exemple : donez la molajo\nl
(mola parto) di ca melono ». Pro quo lore on ne darfus dicar :\nl
« \VUgras{La bonajo di ca afero?} »

La artiklo (quan ni tre konsilas rilektar) finas per tre pruvanta\nl
frazi Germana tradukita ad Ido. Yen la unesma : \VUgras{Die Ethik lehrt}\nl
\VUgras{das Gute, die Wissenschaft sucht das Wahre, die Kunst pflegt das}\nl
\VUgras{Schöne.} La etiko docas la bonajo, la cienco serchas la verajo, l'arto\nl
kultivas la belajo. Nun la frazo povas e darfas anke tradukesar :\nl
\VUgras{La etiko docas lo bona, la cienco serchas lo vera, l'arto kultivas lo}\nl
\VUgras{bela.}
\end{VUpage}

\begin{VUpage}[209]{205}
\VUnotes{148.50mm}{%
(1) A. \textsc{Meillet}, \textit{Introduction à l'étude comparative des langues}\nl
\textit{indo-européennes}, p. 293-599.

Cf. \textsc{Brugmann}, \textit{Abrégé de grammaire comparée}, § 495.%
}
\VUtitre{10.63pt}{1}{CA, TA e QUA.}
\VUsaut{1.29mm}
\VUsoustitre{5.85pt}{211}{\textit{(Apendico 5ma.)}}
\VUsaut{4.43mm}

Esas eroro kredar, ke nia demonstrativi esas arbitrie\nl
selektita; sed (ma), pro ke en ta fako mankas internaciona\nl
formi, on devis serchar l'internacioneso en la primitiva\nl
(Indo-europana) formi. Or la Indo-europana demonstrativi\nl
havas la sequanta radiki : 1\textsuperscript{e} \VUgras{sa, ta}; 2\textsuperscript{e} k, qua divenis \textit{h} o \textit{c},\nl
ed indikas proxima objekto : « la preciza senco di \textit{k} esas\nl
bone definita per la fakto, ke ol esas la demonstrativo qua,\nl
juntita kun la vorto \textit{dio}, donas la senco di \textit{cadie} »(1). 3\textsuperscript{e} por\nl
la fora objekti, on havas tri radiki : \VUgras{w}, \VUgras{n}, \VUgras{l}. La radiko \VUgras{w}\nl
esas orientala; la radiko \VUgras{n} trovesas precipue en la Slava\nl
lingui ed en D. \textit{jene}; la radiko \VUgras{l} trovesas en L. \textit{ille}. Nu, ta\nl
tri lasta radiki ne povus uzesar, nam la sola qua esas sat\nl
internaciona, \VUgras{l}, uzesas ja por l'artiklo e la pronomi (per\cc
sonala). La k ne povas uzesar sub ta formo, qua karakterizas\nl
la relativ-questionala pronomi, nek sub la formo h, des\cc
facile pronuncebla e dicernebla; restas do la tri formi :\nl
\textit{sa, ta, ca}. Ma, se on adoptas \VUgras{s} por la posedala pronomo di\nl
3\textsuperscript{a} persono (qua esas pasable internaciona), restas nur la\nl
du radiki \VUgras{c} e \VUgras{t}, de qui \VUgras{c} necese koncernas la proximeso\nl
(ex. : L. \textit{hic} = \textit{hi-ce}, Romanala lingui : \textit{cil}, cist, nia \textit{cis}\nl
(L.), e. c. (\textit{Progreso}, II, p. 25.)

\VUblanc{6.28mm}
\VUcentre{10.2pt}{0}{\VUgras{Qua.}}
\VUblanc{2.01mm}

Ni ja dicis, por justigar (justifikar) la quaze-digramo \VUgras{qu},\nl
ke ol explikesas per la Indo-europana fonetiko, ube ol re\cc
prezentas la antiqua \textit{labio-velala} artikulado. On respondis a\nl
ni, ke ni ne devas sorgar pri la antiqua o prehistorial fone\cc
tiko, ma nur pri la moderna. To semblas tre justa, ma nur\nl
per ta « prehistorial » serchado on povas trovar l'origino\nl
komuna di ula importanta vorti di nia moderna lingui, e\nl
\end{VUpage}

\begin{VUpage}[210]{206}
\VUnotes{150.88mm}{%
% LE BLANC ENTRE FILET ET PREMIERE NOTE EST PLUS LARGE ICI qu ailleurs
% dans le volume : 48 px au fac-simile, quand le feuillet 209 en donne
% 27 et que \VUsautFilet en pose 19. Le filet, lui, est bien a sa place
% (y=1744, 213 px, verifie a l oeil et par outils/filet.py). On ajoute
% donc le surplus DANS le bloc, plutot que de descendre le filet de
% 20 px pour caler le texte : le filet est visible, on ne le deplace pas.
\VUblanc{1.69mm}
(1) \textit{Corespondense international} n\textsuperscript{o} 9 (\textit{Progreso}, n\textsuperscript{o} 11, p. 681).

(2) A. \textsc{Meillet}, \textit{Introduction}, p. 295; K. \textsc{Brugmann}, \textit{Abrégé de}\nl
\textit{grammaire comparée}, § 497.%
}
\VUornement[7.87mm]{60.62mm}

\VUcontinue
konseque lia formo \textit{internaciona}. To esas aparte vera pri la\nl
pronomi questionala. La Indo-europana radiko esas (duopla\nl
formo) \textit{quo, quei}, qua divenis en E. \textit{what, which} (qui pro\cc
nuncesas reale \textit{hwat, hwich}), ed en D. \textit{wer, was} (la palatal\nl
o velal elemento, ankore reprezentata en E. da \textit{h}, tote des\cc
aparis). Quale on remarkis (1), E. \textit{what} = L. \textit{quod}. Tale la\nl
maxim internaciona formo di la radiko esas vere \textit{qu}, e tote\nl
ne \textit{ki} (Esperanto). Kompreneble, la u qua sequas q esas kon\cc
sonanto, quan la fonetikisti reprezentas ofte per \textit{w} (2).

\VUsaut{13.30mm}
\VUtitre{10.63pt}{-20}{KONJUGO-SISTEMO DI IDO.}
\VUsaut{2.94mm}
\VUsoustitre{5.67pt}{240}{\textit{(Apendico 6ma.)}}
\VUsaut{4.95mm}

Pri \VUgras{-r} por l'infinitivo, \VUgras{-nt} por la participo aktiva, e \VUgras{-t} por\nl
la participo pasiva, nula kontesto esas posibla, ye la nomo\nl
dil linguistiko. Altraparte, ne pri oli kontestas nia adversi,\nl
ma pri \VUgras{s} di nia modi personala, pro ke li generale prenis\nl
ta letro por la pluralo.

Ni devas do montrar pro quo Ido havas \VUgras{s} kom karak\cc
terizivo di sua modi personala :

\VUblanc{0.92mm}
1\textsuperscript{e} Ta konsonanto esas la sol verbal litero quan ofras por\nl
l'organo internaciona kelka lingui naturala.

\textit{a}) \VUgras{S}, en la Greka e Latina lingui, renkontresas (artiku\cc
lata) ye omna duesma personi dil singularo, se on eceptas\nl
\textit{-isti} di la perfekto; la Latina lu prizentas (artikulata) anke\nl
ye l'unesma e la duesma personi dil pluralo.

\textit{b}) Ta konsonanto renkontresas (artikulata) ye la duesma\nl
persono, ne nur dil singularo, ma anke dil pluralo en omna\nl
verbo Hispana; l'unesma persono dil pluralo anke prizentas\nl
lu (artikulata).

\textit{c}) Ni anke trovas lu (artikulata \textit{z}, ligante al vokalo\nl
\end{VUpage}

% -------------------------------------------------------------------
%  FEUILLET 211 — L'EXCES DE L'INDICE 30 N'EST PAS LE FILET, C'EST UN
%  ASTERISME. Le filet est a y=1626 (211 px), entre les indices 31 et
%  32 ; le bloc des notes occupe donc les indices 32 a 38, non 30 a 38.
%  Les trois asterisques en triangle sont a y 1434-1471, verifies a x8.
%  Ordonnee calculee comme celle du feuillet 190, dont \VUasterismo
%  reproduit la valeur exactement : 24,30 + (1435-243) x 25,4/300.
%  La note (2) s'acheve sur « p. 266, », parenthese non fermee : le
%  folio 208 ne porte AUCUNE note (verifie au pixel sur toute la
%  hauteur, et tous ses pas valent 41 px), la note est donc coupee
%  dans l'original. Conserve tel quel.
% -------------------------------------------------------------------
\begin{VUpage}[211]{207}
\VUnotes{141.39mm}{%
(1) En qui la liganta vokali \textit{a, i, o, u} reprezentas rispektive la\nl
prezento, la pasinto, la futuro, la kondicionalo.

(2) En special artiklo « Verbala rolo dil \VUgras{s} en la helpo-linguo »\nl
ni pluse citis la susteno quan la Sanskrita linguo, la Zenda, l'Arme\cc
niana, la Gotika e l'ancien alta Germana donas a \VUgras{s} kom verbal\nl
finalo, en la modi personala : \textit{-asi, -as, ats, -ais, -mas} e. c. (Videz\nl
\textit{Bulletin français-Ido de la Langue auxiliaire} n\textsuperscript{o} 69-70, p. 266,%
}
\VUasterismo{125.14mm}

\VUcontinue
sequanta), ye la duesma persono dil singularo, e ye l'unesma\nl
dil pluralo, en la verbo Franca.

\textit{d}) Ol esas, on darfas dicar, l'unika dezinenco dil tota\nl
Angla konjugo, en la modo personala. La homi qui su\nl
apogas sur l'Angla por postular \VUgras{s} kom plural-marko di la\nl
L. I. devus vere ne obliviar ca lasta fakto.

Ni do vidas la Hispana ed Angla lingui ipsa, qui uzas \VUgras{s}\nl
kom plural-marko, unionar su al lingui supere citita por\nl
apogar \VUgras{s} kom verbal finalo. Koram to, on bezonas lin\cc
guala sentimento tre specala por konsiderar kom pure arbi\cc
triala la dezinenci \textit{-as, -is, -os, -us} di Esperanto e Ido (1).\nl
Ma, se \VUgras{s} atribuesas al verbo, on ne povas atribuar ol pluse\nl
al pluralo : ta kumuli darfas toleresar nur en nia lingui.\nl
En la L. I. li absolute violacus la principo di l'unasenceso.\parplein

2\textsuperscript{e} La sisanta \VUgras{s}, tre facile pronuncebla ed audebla, kon\cc
venas tre bone kom tala por la verbo, qua pleas rolo tante\nl
importanta en la frazo, e mem rolo plu grava kam la pluralo,\nl
se on ponderas bone la kozi.

3\textsuperscript{e} Altraparte, pro ke la siso di la konsonanto \VUgras{s} facile\nl
tedas, kande ol multafoye repetesas sucede, ol esabus quaze\nl
netolerebla, se on selektabus lu kom plural-marko. Nam\nl
multa substantivi intersequas plu freque kam verbi (2).\nl
To esis un plusa motivo por lasar a \VUgras{s} lua rolo verbala.\nl
L'Italiana montras eloquente, ke multa \textit{i} intersequanta\nl
fluas sen tedar : \textit{i miei buoni amici}, kontre ke le tro multa \VUgras{s}\nl
finala di la Hispana certe detrimentas l'eufonio di ca bela\nl
linguo. Or me sempre opinionis, ke la belsoneso importas\nl
en la L. I., ed esas atraktilo. Ido montras, ke me ne eroris\nl
pri ca punto

\VUsaut{7.04mm}
On dicis, ke la finali \VUgras{-as}, \VUgras{-is}, \VUgras{-os}, \VUgras{-us} meritas la reprocho\nl
esar adminime kelke arbitriala. Oportus remplasigar li per
\end{VUpage}

% -------------------------------------------------------------------
%  FEUILLET 212 — PAGE SANS AUCUNE NOTE. Cherche au pixel sur toute
%  la hauteur, hors de la fenetre de outils/filet.py et avec un seuil
%  de longueur abaisse a 60 px : aucun candidat. Les seuls segments
%  longs et propres sont le bord noir du scan, 150 px au-dessus du
%  folio. Tous les pas de ligne de base valent 41 a 45 px, jamais les
%  ~33 px d'un bloc de notes.
% -------------------------------------------------------------------
\begin{VUpage}[212]{208}
\VUcontinue
altra finali \textit{pure romanala}, e komprenebla da un o du popu\cc
li adminime. Ma on obliviis, ke por la cetera populi li ne\nl
esus min arbitriala, e forsan li esus min facile rikonocebla\nl
e memorebla. Altraparte la finali propozita koincidus kun\nl
altra gramatikal finali, o plu juste dezinenci (\VUgras{-o}, \VUgras{-a}, \VUgras{-e}, \VUgras{-i}),\nl
e konseque esus ne konciliebla kun la sistemo di la grama\cc
tikal finali, indikanta quaze per etiketi la naturo dil vorti\nl
e lia kategorio o rolo gramatikala. Do la chanjo di ta verbal\nl
dezinenci necesigus la \textit{total} reformo di la finali karakteri\cc
ziva di Ido, qui esas tante utila kom amortisanti eufoniala,\nl
e kom indikili por la lernanti e novici, qui esas ed esos\nl
ankore dum longa yari la maxim multi. Ica konsideri potente\nl
influis la Komitato dil Delegitaro pri la konservo dil finali\nl
supere citita, ed oli ne esis min konvinkanta pose che l'Aka\cc
demio Idista.

Al verbal dezinenci di Ido on reprochis esar nek \textit{naturala}\nl
nek \textit{ciencala}. Ni darfus saciesar da ca respondo : multe plu\nl
importas, ke li esez koheranta, simpla, facile merkebla ed\nl
aplikebla. Ma ni examenez la reprocho.

Pro ke ni ne savas (e kredeble anke vu, lektero) quo esas\nl
dezinenci ciencala por la konjugo, ni okupos ni nur pri \textit{le}\nl
\textit{naturala}.

Do, la konjugo-dezinenci di Ido ne esas naturala.

Se la aserto esas justa, ni certe ne trovos uli tala en la\nl
Latina ed en la Hispana, exemple, pro ke ta lingui esas\nl
nekontesteble naturala.

On bone remarkez, ke ni examenas la verbal dezinenci\nl
nur pri la \textit{formo} naturala, nule pri lia korespondo kun olti\nl
di la Latina o di la Hispana relate modo, persono o tempo.

\VUblanc{1.27mm}
Se \textit{-as, -is, -os, -us} (ni vidos \textit{-ez} plu fore) ne esas konjugo\cc
finali naturala, tote certe ni ne trovos li, nek en la verbo\nl
Latina, nek en la verbo Hispana.

Tamen, kad la Latina ne posedas, en l'unesma konjugo :\nl
\textit{am-as}, stranje simila a (tu) \textit{am-as} di Ido? E kad la Hispana\nl
ne havas anke (tu) \textit{am-as}? Kad en altra tempi li ne posedas :\nl
\textit{amabas} (L. S.), (tu) \textit{habras, amaras, habias} (S.), \textit{legas, audias,}\nl
\textit{amaveras} (L.)?

Do \VUgras{-as}, agnoskez lo, esas vere final sono verbal \textit{naturala}.\parplein

E pri \VUgras{-is}, ka ne existas en la Latina, ye la duesma per\cc
\end{VUpage}

% -------------------------------------------------------------------
%  FEUILLET 213 — CINQUIEME CAUSE DE « FILET INTROUVABLE ». Le filet
%  est parfait : y=1372, 214 px, cale sur la marge, franc a l'oeil.
%  Dix pixels de crasse a x = 1073-1082, AU MILIEU de la colonne,
%  portaient le « reste » a 10 pour un seuil de 6. La marge BORD, posee
%  pour la tranche du feuillet voisin, ne protege que du bord droit.
%  Corrige dans outils/filet.py : on juge desormais sur la taille du
%  plus gros amas, non sur le seul total.
%  Ordonnee : 24,30 + (1372-227) x 25,4/300 = 121,24 mm.
%  Le bloc de notes est UNE note (1) en trois alineas : le corps
%  n'appelle qu'un seul « (1) ».
% -------------------------------------------------------------------
\begin{VUpage}[213]{209}
\VUnotes{120.56mm}{%
(1) Kad \VUgras{audis} L. ne esas identa a \textit{audis} di Ido? Ma on respondas :\nl
« \textit{audis} L. esas prezento, e vu facas barbarismo, uzante lu por\nl
indikar pasinto ». Ne koaktez me, kara kritikeri, respondar a vi per\nl
cito di barbarismi multe plu grosa, quin vi facas en konjugo asertita\nl
kom Latina, kom natural e ciencala. Ma kad me ne precizigis, ke\nl
me parolas nur pri la \textit{formo}. Nu, kad \textit{audis, dormis, venis, sentis}\nl
ne esas, en Ido, \textit{forme} simila ad \textit{audis, dormis, venis, sentis} di la\nl
Latina?

Relate \VUgras{i,} karakterizivo dil pasinto en Ido, kad ol ne havas bela\nl
defenso en omna perfekti Latina, che qui ta vokalo uzesas kinfoyi\nl
ek sis? (\textit{Amavi, amavisti} e. c.)

E kad \VUgras{a} por la prezento ne bone defensesas dal infinitivo prezenta\nl
dil unesma konjugo Latina, Italiana, Hispana (\textit{-are, -ar}), dal indik\cc
ativo prezenta (\textit{amas, amat} e. c., \textit{ama, parla, aman}), fine dal parti\cc
cipo prezenta (\textit{amant...., amando, aimant} F.)?%
}

\VUcontinue
sono singulara; \textit{audis} (1), \textit{capis}, \textit{amabis}, e. c., ed en la His\cc
pana, ye la duesma plurala persono, kad on ne trovas : \textit{amais},\nl
\textit{amabais, amesteis, habais, habiais, amarais, temiais, temis}\cc
\textit{teis, partiais, partis}.

Do \VUgras{-is} esas vere final sono verbal \textit{naturala}.

Nun ni videz la dezinenco \VUgras{-us.}

La questiono esas ed esas nur : kad \VUgras{-us} trovesas en la\nl
Latina kom final sono di konjugo, identa (sone) a la verbal\nl
dezinenci \VUgras{-us} di Ido?

Negar lo esus negar l'evidenteso ipsa. Nam, de la verbo\nl
\textit{sum} (esar) til la lasta di sua verbi, la Latina prezentas ta\nl
sono en omna unesma personi di la pluralo.

Do \VUgras{-us} esas vere final sono verbal \textit{naturala}, pro ke on\nl
trovas lu dek foyi en omna verbo Latina.

Ma \VUgras{-os,} kad on darfas dicar lu final sono verbal naturala?\parplein

Yes, certe; nam la finalo \textit{mus} dil verbo Latina divenis\nl
\textit{mos} en la Hispana, ye la sama persono, en singla tempo.\nl
De to rezultas, ke la verbo Hispana, ecepte eroro, audigas\nl
dek e non foyi la final sono \VUgras{os} en sua konjugo. Konseque\nl
la dezinenco \VUgras{os} di Ido esas vere final sono verbal naturala.\parplein

On forsan alegos objecione, ke ta finalo, en la Hispana,\nl
pluse markas substantivo plurala.

Ma kad ol agas lo en Ido? E juste pro ke \textit{os} en la Hispana\nl
renkontresas 19 foyi kom final sono en la verbi, esas eroro\nl
selektar \VUgras{s} kom plural signo di substantivi, citante la His\cc
pana kom susteno. Nam fakte la verbi, en ica linguo, pre\cc
zentas \textit{os} kom final sono plu ofte kam la substantivi.
\end{VUpage}

% -------------------------------------------------------------------
%  FEUILLET 214 — LA GEOMETRIE DIT « exces : aucun » ET C'EST FAUX :
%  146 px separent l'indice 34 de l'indice 35, et le filet est a
%  y=1698 entre les deux. Le detecteur d'exces a manque le saut ; ce
%  sont les y0 qui font foi.
% -------------------------------------------------------------------
\begin{VUpage}[214]{210}
\VUnotes{148.00mm}{%
% Meme blanc large qu au folio 206 : 52 px du filet a la premiere
% note, quand le feuillet 209 en donne 27. Le filet est bien a sa
% place (y=1698, 212 px) ; le surplus va DANS le bloc.
\VUblanc{2.12mm}
(1) L'imperativo (o plu bone \textit{volitivo}) \textit{tale formacata} \textit{(ven-ez)}\nl
semblas a ni nek min ciencala (nam on sempre alegas la cienco) nek\nl
min naturala, nek min klara kam le \textit{venidum} o \textit{ples venir} di altra\nl
sistemi,%
}

Ma on questionas : pro quo la vokalo o kom karakteri\cc
zivo dil futuro?

Pro ke, se ni atribuas karakteriziva vokalo al prezento\nl
(\textit{a}), \VUgras{altra} (\textit{i}) al pasinto, ne esas kontrelogika o nekoheranta\nl
atribuar anke specal vokalo (\textit{o}) a la futuro. Facar ek ta\nl
vokalo la simbolo dil futuro en la tota konjugo ne esas plu\nl
arbitriala, ed alteras mine la Latina kam selektar (quale ula\nl
Latinisti) la partikulo \textit{pos} kom signo di ta tempo, o trans\cc
portar sur la finalo la acento dil infinitivo latina por facar\nl
ek lu kondicionalo : \textit{amarè, audirè}, e. c., kande on povabus\nl
evitar omna barbarismo, uzante \textit{amarem, audirem}, e. c.\nl
Stranja maniero, vere, restar ciencoza e naturala!

E la dezinenco \VUgras{-ez} dil imperativo? Unesme ni questionos,\nl
kad ol esas naturala od artificala? Me opinionas, ke la\nl
Franca linguo sate pruvas lua « naturaleso » per omna\nl
duesma personi plurala di sua verbo, ed aparte di sua im\cc
perativo. Ica fakto, pro ke nulo plu bona existis, igis me\nl
propozar ol vice \VUgras{-u} di Esperanto (markizanta anke pro\cc
nomi : \textit{kiu, chiu}, e. c.) e qua, se me ne eroras, posedas kom\nl
apogo nur imperativo Hebrea, nekonocata da preske omni.\nl
Remarkez ke \VUgras{-ez} reale apartenas al sistemo \textit{-as, -is, -os, -us}\nl
(per la substituco di la febla \textit{z} a la forta \textit{s}). Se me ne pro\cc
pozis \textit{-es}, la kauzo esas, ke \textit{venez, rakontez, dormez, preparez,}\nl
\textit{pensez, pardonez, donez, manjez}, e. c., e. c., esas quik kom\cc
prenata da irgu qua savas kelkete la Franca, e certe plu\nl
facile e bone kam \textit{venes, rakontes, dormes}, e. c. (1).

Irge quon opinionas kontrediceri, fakto esas nekonteste\cc
bla : la Franca linguo pruvas, ke \VUgras{-ez}, por l'imperativo, esas\nl
vere verbal finalo naturala. Pri la sono, le \textit{venez a ni, rakontez}\nl
\textit{ankore}, e. c., esas la preske kompleta riprodukturo di F. \textit{venez}\nl
\textit{à nous, racontez encore.}

Justifikinte \textit{-as, -is, -os, -us, -ez}, restas justifikenda \VUgras{-ir},\nl
\VUgras{-or} dil infinitivi pasinta e futura.

Fakte nia kontrediceri \textit{chikanas} nur pri l'infinitivo futura,\parplein
\end{VUpage}

% -------------------------------------------------------------------
%  FEUILLET 215 — LIGNE INVISIBLE RETROUVEE : « nitivo. », 106 px de
%  large, entre les indices 4 et 5 de la geometrie. L'exces de 41,6 px
%  qu'annoncait le detecteur n'etait pas un blanc mais un interligne
%  entier : la ligne y etait, sous le seuil d'encre. Verifiee a x7.
%  Le corps compte donc 32 lignes, non 31.
% -------------------------------------------------------------------
\begin{VUpage}[215]{211}
\VUnotes{139.28mm}{%
(1) Che li ta infinitivo esas kompozita per la verbo \textit{havar} o \textit{esar}\nl
e pasinta participo; che ni ol esas simpla do plu kurta : \textit{-ir},\nl
quale la prezenta \textit{-ar}. Yen la tota difero. Kad on volus donar solid\nl
expliko pro quo l'infinitivo prezenta devas esar simpla, ma la\nl
pasinta, kompozita.

(2) Pri l'infinitivo futura, ni remarkigos ke ol trovas kelka\nl
susteno en la Latina : \textit{fore} = \textit{esor}, di qua l'\textit{e} finala audesas apene\nl
pro la tonika acento.%
}

\VUcontinue
nam generale li ipsa uzas infinitivo pasinta (1) kun infi\cc
nitivo prezenta. Li anke ne kontestas pri -r kom karakterizivo\nl
dil modo infinitiva. Do reale kontestesas nur la vokalo sim\cc
bolizanta la tempo (\textit{-i, -o}) e la bezono ipsa di futura infi\cc
nitivo.

Ma, se on aceptas ke \VUgras{a} simbolizas la prezento : \textit{amas, amar,}\nl
pro quo \VUgras{i} ne simbolizus la pasinto : \textit{amis, amir} e \VUgras{o} la futuro :\nl
\textit{amos, amor}? Kad ica tri karakterizivi tempala (qui en omna\nl
modi pleas la sama rolo) esas kompliko o simpligo? Kad to\nl
ne esas simpla, praktikal e homogena? Kontraste, kad esas\nl
tala konjugo-sistemo sintezala en un parto, ed analizala en\nl
l'altra parto, quale olti di nia kontrediceri? Ka tal konjugo\nl
esas tre koheranta, tre homogena e tre ciencala, por parolar\nl
quale li?

Altraparte \textit{-ar, -ir} en Ido semblas a me tam natural sono\nl
verbala kam en la infinitivi Latina, Italiana, Hispana,\nl
Franca en \textit{-ire, -ir}, od en l'infinitivi Germana en \textit{-ieren},\nl
ankore tante produktiva (2).

Ma on objecionas, ke li esas prezenta en la lingui supere\nl
mencionita. Me agnoskas lo; ma li restas adminime infini\cc
tivi, en Ido; e ke \textit{i} simbolizas en li la pasinto, quale en la\nl
tota konjugo, certe ne povas jenar la adepto, nek super\cc
charjar lua memorado. Kad esus por lu plu facila memorar,\nl
ke \textit{du} simbolizas la pasinto, \textit{pos} la futuro, \textit{fe} l'imperfekto,\nl
\textit{vell} la kondicionalo, \textit{dum} (pos-pozita) l'imperativo, quale\nl
en altra sistemi plu o min laborita? Kad vere Ido ne darfas,\nl
kun adminime yuro tam granda e cienco egala, donar ica\nl
principo : en \textit{omna} dezinenco verbala, \VUgras{a} simbolizas la pre\cc
zento, \VUgras{i} la pasinto, \VUgras{o} la futuro? Kad ol esas per to min\nl
simpla o min koheranta kam altri?

Ma, on objecionas ankore, pro quo \textit{sis} participi, kande\nl
\textit{2} suficas a multa lingui? On ne dicas, kompreneble, ke pro\parplein
\end{VUpage}

% -------------------------------------------------------------------
%  FEUILLET 216 — TROISIEME ASTERISME DU VOLUME, apres les feuillets
%  190 et 211. Le detecteur ne l'a pas signale : l'intervalle qui le
%  porte vaut exactement quatre pas (166,9 px pour 4 x 41,40), si bien
%  qu'aucun « exces » n'en sortait. Mesure : rang haut y 1232-1247,
%  rang bas y 1254-1269, ecart 22 px, asterisques bas a 61 px l'un de
%  l'autre, centre x 750 quand le milieu de la justification tombe a
%  749,5. Au pixel les memes valeurs que les deux autres.
%  Ordonnee : 24,30 + (1232-248) x 25,4/300 = 107,61 mm.
%  PAGE SANS NOTE, la deuxieme du volume apres le feuillet 212 : les
%  trois segments qu'un balayage grossier signalait (y 639, 764, 1629)
%  sont des lignes de texte soudees par la fermeture morphologique,
%  longues de 106 a 127 px quand un filet en mesure 213 ; celle de
%  y=1629 a ete lue a l'agrandissement et porte du texte.
% -------------------------------------------------------------------
\begin{VUpage}[216]{212}
\VUasterismo{107.61mm}

\VUcontinue
la jus mencionita principo pri \textit{a}, \textit{i}, \textit{o} en dezinenco verbala,\nl
la lerno e memoro pri la participi reduktesas a \textit{2} kozi :\nl
nt participo aktiva, t participo pasiva. Yen fakte la \textit{sis} par\cc
ticipi! On ne dicas, ke la participi futura (\textit{-ont, -ot}), multe\nl
min frequa, lernesas pos la ceteri, okazione. On ne dicas\nl
de quanta dubi, miskompreni o perifrazi liberigas ta parti\cc
cipi. On ne dicas ke li genitas multega substantivi tre pre\cc
ciza, quin nia lingui kareas nur koakte e kun multa des\cc
avantaji. Ton omna on tacas por blamar mokante kun sem\cc
blanta justeso e yusteso. Ma bone ridos, qua laste ridos.\nl
Yes; e kande Ido vere sucesabos, on laudos me ne propozir\nl
lu \textit{tro} simpla, e ne sakrifikir la just expreso dil pensi a\nl
simpleso exajerita.

La nuna Greka kareas omna infinitivo. Ka ni devus imitar\nl
olu? No, tre certe. La Angla e la Germana havas nur \textit{1} per\cc
sonal pronomo : \textit{they, sie} por la triesma persono plurala, e\nl
ni havas \textit{tri}. Ka ni esas pro to blaminda? Certe no. La Franca\nl
posedas nur \textit{son, sa, ses} kom posedal adjektivo por la triesma\nl
persono. Ka ni devus imitar lu, quale agis granda kritikero?\nl
Ha certe no, nam ico esas ne simpleso, ma deplorinda mizero.\nl
En « Les vrais principes de la Langue Auxiliaire » me pozis\nl
la limito netranspasebla dil simpleso por ta linguo : \textit{la}\nl
\textit{klara e justa expreso dil pensi.}

\VUsaut{10.63mm}
\VUcentre{10.2pt}{29}{\VUgras{Pro quo ne} \textit{havar} \VUgras{en la tempi kompozita?}}
\VUsaut{3.02mm}

Ca punto di nia konjugo meritas exameno aparta.

Unesme ni konstatez ico : Judikata dal komuna raciono,\nl
la verbo \textit{havar} kom helpanto konjugala en la kompozita tempi\nl
esas \textit{absurdajo}. Nam verbo, qua esence indikas la posedo,\nl
quale \textit{havar}, ne darfas, mem ne povas logike divenar simpla\nl
marko di tempo, e fakte analoga a dezinenco, por indikar\nl
l'anteeso : j'\textit{avais} aimé (lore), j'\textit{aurai} aimé (lore), o \VUgras{simple}\nl
la pasinto : j'\textit{ai} aimé (me amis). Esas neposibla kompre\cc
nigar ed admisigar ta gramatikal idiotismo dal populi qui\nl
ne havas lu; or li esas nekontebla. La L. I. qua destinesas\nl
ad omni, devas konseque evitar ta eroro (nule naturala reale,\nl
pro ke ol esas absurda); pro to Ido evitis lu.
\end{VUpage}

% -------------------------------------------------------------------
%  FEUILLET 217 — TROISIEME PAGE SANS NOTE, et une ligne invisible de
%  plus : « jugo! », 90 px de large, entre les indices 30 et 31. Elle
%  ne porte que 8 % de l'encre d'une ligne pleine — la signature
%  connue. Le corps compte 41 lignes, non 40.
%  Aucun segment horizontal de 90 px ou plus sur toute la page : ici
%  « filet non trouve » est la bonne reponse, non une panne.
% -------------------------------------------------------------------
\begin{VUpage}[217]{213}
Artiklo aparinta, dum aprilo di 1910, en \textit{Progreso}, expozas\nl
tre bone la questiono. Ol debesas a M. J. \textsc{Kovacic}, instruk\cc
tisto en \textit{St. Veit i. Jauntal} (Karintia). Yen olu :

« On propozis (n\textsuperscript{o} 19, p. 398) formacar la kompozita\nl
tempi di l'aktivo per la verbo \textit{havar} e la pasinta participo,\nl
pro ke la lingui D. E. F. I. S. uzas kune la helpanta verbo\nl
\textit{havar}. Sed (ma), kad la mencionita lingui formacas la\nl
kompozita tempi di l'aktivo tute (tote) e sen ecepto per \textit{havar?}\nl
No, nam oli havas multa ecepti, ex. : D. por \textit{mortar}, \textit{vekar},\nl
\textit{venar}, \textit{irar}, \textit{saltar}, di qui la tempi kompozita formacesas per\nl
\textit{esar}. On devus do admisar en Ido ecepti qui esas, quale on\nl
savas, granda malfacileso (desfacileso) por omna lernanto.\parplein

Pluse, quale on devus tradukar D. \textit{Ich habe geschlafen,}\nl
\textit{Wir haben uns gefreut, Du bist tapfer gewesen?} Frazi quale :\nl
\textit{Me havas dormita, Ni havas joyita, Tu havas brava esita,}\nl
semblas certe nelogikala ad omnu. Sed (ma) on ne darfus\nl
uzar \textit{esar} vice \textit{havar} en la alegita frazi, pro ke \textit{esar} esus defi\cc
nota (destinota) nur por formacar la tempi di la pasivo.

Ankore, quale on tradukus D. \textit{der gefallene Soldat, die}\nl
\textit{verflossene Nacht, der verstorbene Vater?} Dicar : \textit{la falita}\nl
\textit{soldato, la pasita nokto, la mortita patro,} esus tote nelogikala!\parplein

Ni videz quale altra Indogermana lingui, quin ni devas\nl
konsiderar pri ca questiono, formacas la kompozita tempi\nl
di l'aktivo. La Slava lingui formacas oli per la helpanta\nl
verbo \textit{esar} e la aktiva participo, exakte same kam en Ido.\nl
Konseque Ido en la formaco di la komp. tempi ne agas\nl
arbitriale, sed (ma) ol apogas su sur granda parto di l'Indo\cc
germana linguaro, nome sur la Slava lingui. Pluse la nuna\nl
participi esas tre aplikebla e praktikala quale atributi, kozo,\nl
quan on devas bone konsiderar ante enduktar altra kon\cc
jugo!

Fine, kad on inventos plu bona, plu internaciona e min\nl
artificala konjugo, kam esas la nuna? Me dubas. La kon\cc
jugo propozita en n\textsuperscript{o} 22 semblas nulkaze esar facila, pre\cc
cipue en la pasivo. Ni videz nur quanta sufixi esas por ol\nl
necesa : \textit{-as, -avas, -iras, -it, -avit, -irit, -olas, avolas, -irolas,}\nl
do triople multa kam en nuna Ido! Mem se ta sufixi igus\nl
omna nuna analizala verbformi superflua, la konjugo esus\nl
malgre to min facila kam la nuna. \textit{Amesavolas} o \textit{amesirolas}\nl
ne esas plu simpla kam : \textit{esos amita} o \textit{esos amota}. E quala\nl
\end{VUpage}

% -------------------------------------------------------------------
%  FEUILLET 218 — le bloc de notes prend plus de la moitie de la page :
%  vingt-deux lignes en sept alineas. Le filet, lu a y=1142 par le
%  releveur ET par outils/filet.py, long de 212 px : les deux lectures
%  concordent, ce qui n'etait pas le cas au folio 195. Blanc filet ->
%  premiere note : 26 px, le regime ordinaire.
% -------------------------------------------------------------------
\begin{VUpage}[218]{214}
\VUnotes{100.67mm}{%
(1) L'akademio repulsis ta ed altra chanji. Videz ye la fino di la\nl
chapitro pri « \textit{verbo} », pag. 55.

(2) Kad, en la Germana, \textit{fiel} e \textit{fiele,} \textit{fällt} e \textit{fielt,} \textit{werde} e \textit{würde,}\nl
\textit{waren} e \textit{wären,} \textit{sprecht} e \textit{spricht} e. c., diferas inter su, por la orelo,\nl
plu forte kam \textit{amas, amis, amos, amus, amez} o \textit{amar, amir, amor?}\parplein

Kad, en l'Angla, \textit{drink, drank, drunk} --- \textit{ring, rang, rung} --- \textit{sing,}\nl
\textit{sang, sung} diferas per altro kam per vokalo, o kad la difero audesas\nl
plu bone kam en \textit{esanta, esinta, esonta?}

Kad, en la Franca, e por l'orelo, \textit{fais} e \textit{fis,} \textit{lus} e \textit{lis,} \textit{marcha,}\nl
\textit{marchait, marché, serai} e \textit{serais} e. c. diferas plu dicerneble kam\nl
\textit{esus, esez} o \textit{amita, amata?}

Kad, en la Latina, olim internaciona ed uzita, parolita dal maxim\nl
diversa nacioni, on ne dicernis e nun ankore ne dicernas aude \textit{leges}\nl
de \textit{legis, audiunt} de \textit{audiant, audies} de \textit{audias, sunt} de \textit{sint, erant}\nl
de \textit{erunt, fueras} de \textit{fueris, ames} de \textit{amas, amabas} de \textit{amabis} e. c.?\nl
Kad on plendis o plendas pro ke li diferas inter su nur per un sola\nl
vokalo quale \textit{amos, amus} o \textit{aminte, amante} en Ido?

Se en linguo patriala, plu o min neglijata dal nacionani, pri la\nl
pronuncado, on dicernas e komprenas vorti o verbal formi diferanta\nl
nur per un vokalo, pro quo on ne atingus sama rezultajo en Ido,\nl
linguo stranjera a l'uzanti e kom tala plu atence pronuncata?\parplein

Cetere, ka nia renkontri e kongresi ne sat konvinkive respondis?\parplein
}

\VUcontinue
longa rubandatra vorti! \textit{Amesavolas} ed \textit{amesirolas} ankore\nl
konvenas, ma ne plus la vorti : \textit{diminutesavolas, kalum}\cc
\textit{niesavolas} o mem \textit{familiarigesirolas}. On sentas justa angoro\nl
en pronuncar tala longaji! (1).

Kelka kritikanti anke propozas chanjar la verbal finali :\nl
\textit{-as, -is, -os, -us, -ez,} dicante, ke oli ne esas sufice distingebla\nl
(dicernebla). Me opinionas, ke ta aserto ne esas justa. La\nl
finali alegita diferas inter su per la bone audebla vokali\nl
multe plu bone kam en multa nacionala lingui. Exemple :\nl
en D. \textit{er lernt -lernte, wir erwachen -erwachten} diferas per\nl
un litero facile misaudebla (2).

Me finas do kun la propozo : « Nia estimata Akademio\nl
konservez la nuna, facila e tre praktikala konjugo! »

L'artiklo esis sequata da la yena remarko : « Ne nur D.\nl
ed F. havas verbi qui formacas la kompozita tempi per \textit{esar}\nl
vice \textit{havar,} ma inter la korespondanta verbi di la du lingui,\nl
sat (pasable) multi uzas \textit{havar} en un linguo ed \textit{esar} en l'altra,\nl
exemple : \textit{esar, kurar, saltar,} e. c. Do ne nur l'adepti di singla\nl
linguo hezitus ofte inter \textit{esar} e \textit{havar,} sed (ma) la Franci ne\nl
uzus la sama helpanta verbo en la sama kazi kam la Ger\cc
mani. Adjuntez ke, en la reflektiva verbi, la F. vicigas (rem\cc
\end{VUpage}

% -------------------------------------------------------------------
%  FEUILLET 219 — OUVERTURE DE L'APENDICO 7, « VORTORDINO ». Page sans
%  note : les deux segments qu'un balayage grossier signalait (y 932 et
%  989, 151 et 122 px) sont des lignes de texte soudees, lues a
%  l'agrandissement. C'est la troisieme page du volume sans note.
%  LE FILET COURT DE SECTION ECHAPPE AU DETECTEUR comme il echappait au
%  feuillet 187 : 91 px de large, 5 px d'epais, a y 485-489, x 583-674,
%  centre a 628,5 pour un milieu de justification a 631,5.
% -------------------------------------------------------------------
\begin{VUpage}[219]{215}
\VUcontinue
plasigas) absurde \textit{havar} per \textit{esar}, dum ke la D. uzas (logike)\nl
ankore \textit{havar} : \textit{ich habe mich gewaschen} = \textit{je me} suis \textit{lavé} :\nl
altra frequa kauzo di erori, di heziti, di miskompreni. Tale\nl
en ta sola punto questionesas pri la logikala karaktero di\nl
nia linguo, do pri sa (lua) vera internacioneso.

\VUsaut{4.61mm}
\VUfilet{7.7mm}
\VUsaut{3.20mm}
\VUtitre{10.27pt}{-73}{VORTORDINO.}
\VUsaut{2.39mm}
\VUsoustitre{6.02pt}{228}{\textit{(Apendico 7-ma.)}}
\VUsaut{4.02mm}

Quale ni ja vidis en la sintaxo, la \textit{normala vortordino}\nl
esas : 1\textsuperscript{e} subjekto, 2\textsuperscript{e} verbo, 3\textsuperscript{e} komplemento direta, singla\nl
ek oli esante akompanata da sua omna komplementi. La\nl
komplementi nedireta povas esar lokizata irgaplase, ma li\nl
prefere sequas la verbo. Ex. : \VUgras{La filiulo di nia olda vicini}\nl
\VUgras{retrovenis de Paris, ube il kompris multa ludili por sua in}\cc
\VUgras{fantino.}

Omna rupto dil ordino normala nomizesas \textit{inversigo}, sive\nl
pro ke la komplemento direta preiras la verbo : \VUgras{elun ne ilun}\nl
\VUgras{me vokis,} sive pro ke la subjekto sequas lu : \VUgras{venez, venez,}\nl
\VUgras{klameskis l'instruktisto.} L'inversigi povas esar postulata dal\nl
bezono insistar sur vorto, pozante lu avane, quale \textit{elun} en\nl
l'unesma exemplo, o dal deziro (en la traduki) sequar\nl
l'ordino dil texto originala.

\textit{La questionanta vorti} sempre komencas la propoziciono e\nl
ne diplasas la subjekto : \VUgras{Kad vu venos; dicez kad vu venos.}\nl
\VUgras{--- Quo eventis? Quon vu vidas? --- Quantin me kontis?}\nl
\VUgras{me ne plus memoras.}

\textit{L'atributo} devas quik sequar la verbo esar : \VUgras{Esez mizerikor}\cc
\VUgras{dioza por omni,} e ne : \textit{esez por omni mizerikordioza.} --- Il \textit{ne}\nl
toleras to, quo esas erora segun sua konvinkeso, e ne : \textit{il ne}\nl
\textit{toleras to, quo esas segun sua konvinkeso erora,} en qua \textit{erora}\nl
semblas relatar \textit{konvinkeso}. E tale kun omna verbo havanta\nl
atributo : \VUgras{Il mortis tre povra de pekunio, ma richa de vertui,}\nl
e ne : \textit{il mortis de pekunio tre povra, ma de vertui richa.}

\textit{L'adjektivo} avan o dop lua substantivo (Videz en la gra\cc
matiko, ye adjektivo). Ma ne sistematre avan la substantivo,\nl
quale en la Germana e l'Angla,\nl
\end{VUpage}

% -------------------------------------------------------------------
%  FEUILLET 220 — quatrieme page du volume sans note. Les quatre
%  segments du balayage grossier (y 528, 861, 1540, 1664) ont ete lus a
%  l'agrandissement : ce sont des lignes de texte soudees.
%  COQUILLE DE L'ORIGINAL, conservee : « è » pour « e » a la ligne 22,
%  accent grave de plomb et non salissure (noir plein, pente
%  descendante a droite, verifie a x14).
%  PARTICULARITE conservee : a la ligne 9 le contre-exemple est compose
%  en GRAS, alors que les douze autres de la page sont en italique.
% -------------------------------------------------------------------
\begin{VUpage}[220]{216}
\textit{L'adverbo} preiras o sequas nemediate la vorto quan ol\nl
modifikas. \textit{Ne, tre, nur} sempre avan olu. Ofte l'adverbo bone\nl
lokizesas inter la subjekto e la verbo : \VUgras{Il sempre plendas}\nl
\VUgras{pri omno.}

\textit{La participo,} en la kompozita tempi dil verbo (aktiva o\nl
pasiva), sempre devas sequar la helpanto \textit{esar,} e darfas\nl
separesar de olu nur da adverbo relanta la verbo : \VUgras{Ni esas}\nl
\VUgras{arivinta de du hori,} e ne : \textit{ni esas de du hori arivinta.} ---\nl
\VUgras{El esas tre amata da omni,} e ne : \VUgras{el esas da omni tre amata.}\nl
--- \VUgras{Ta soldato esis grave vundata.}

\textit{La participo} o \textit{l'adjektivo} devas sempre \textit{sequesar} da sua\nl
komplementi, direta o ne direta, e ne separesar de sua\nl
substantivo : \VUgras{La homo estimata da omni,} e ne : \textit{la da omni}\nl
\textit{estimata homo,} o : \textit{la homo da omni estimata.} --- \VUgras{La fervoyi}\nl
\VUgras{formacas reto kovranta la mondo,} e ne : ... \textit{reto la mondo}\nl
\textit{kovranta,} e mem mine : \textit{la mondo kovranta reto.}

\VUgras{Me recevis letro skribita en linguo ne konocata da me,}\nl
e ne : \textit{me recevis letro en da me ne konocata linguo,} e mem\nl
mine : \textit{me recevis en da me ne konocata skribita letro.} ---\nl
\VUgras{El esas laudata ed amata da ti omna qui prizas la beleso}\nl
\VUgras{ed afableso,} e ne : \textit{el esas da ti omna qui prizas la beleso}\nl
\textit{ed afableso amata è laudata.} --- \VUgras{Ta urbo esas abunde provi}\cc
\VUgras{zita ye omna kozi utila por la vivo,} e ne : \textit{ta urbo esas ye}\nl
\textit{omna por la vivo utila kozi abunde provizita.} --- \VUgras{La mastro,}\nl
\VUgras{ebria de furio, bastonagis l'asno ne volanta avancar,} e ne :\nl
\textit{la de furio ebria mastro bastonagis la avancar ne volanta}\nl
\textit{asno.} --- \VUgras{Ni refutis omna objecioni prizentita da nekompe}\cc
\VUgras{tenta personi ed ofte genitita da nesuficanta konoco di nia}\nl
\VUgras{linguo,} e ne : \textit{ni refutis omna da nekompetenta personi}\nl
\textit{prizentita ed ofte da nesuficanta konoco di nia linguo geni}\cc
\textit{tita objecioni.} --- \VUgras{Tablo kovrita per ancien reda tapiso ma}\cc
\VUgras{kulizita,} e ne : \textit{tablo per ancien reda makulizita tapiso}\nl
\textit{kovrita.} --- \VUgras{Me propagas la linguo helpanta Ido, konstruktita}\nl
\VUgras{sur logikal principi e sencese developata per ciencal dis}\cc
\VUgras{kutado,} e ne : \textit{me propagas la helpanta, sur logikal principi}\nl
\textit{konstruktita e per ciencal diskutado sencese developata linguo}\nl
\textit{Ido.} --- \VUgras{La linguo helpanta esas l'inventuro maxim utila por}\nl
\VUgras{la developo di l'internaciona relati e maxim fekunda de}\nl
\VUgras{grava konsequi por la profito di la homaro futura,} e ne :\nl
\textit{la helpanta linguo esas la por la developo di l'internaciona}\nl
\end{VUpage}

% -------------------------------------------------------------------
%  FEUILLET 221 — OUVERTURE DE L'APENDICO 8, « PUNTIZADO ». Le filet
%  court de section echappe encore au detecteur : 100 px de large,
%  5 px d'epais, a y 830-834, x 590-689 — meme abscisse de depart qu'au
%  feuillet 187, ou l'homologue mesure 92 px.
%  Le filet de NOTE, lui, a ete verifie a la main : y=1191, 213 px,
%  cale sur la marge (x=101 pour une marge a 103). L'ordonnee de
%  l'outil est confirmee.
%  COQUILLES conservees : point final absent apres « quan ol referas »
%  (le titre suivant ouvre pourtant sur une capitale) ; titre
%  « PUNTIZADO » sans point final, seul du volume dans ce cas ; et
%  italique interrompue au milieu de la citation allemande de la
%  note (2).
% -------------------------------------------------------------------
\begin{VUpage}[221]{217}
\VUnotes{108.12mm}{%
(1) Multi ek ta exempli prenesis de la XXVI exerco di « Exer\cc
caro ». Ica finas per konsilo quan S\textsuperscript{ro} \textsc{Couturat} judikis, quale me,\nl
kom importanta e \textit{konform a l'evoluciono} di nia lingui : preferar la\nl
voco aktiva a la pasiva. Ex. : « Dicez per la formo aktiva : « On\nl
(\textit{o} ni) sequis la voyo, acensis la kolino, vizitis la kastelo, ed admiris\nl
la bela panoramo, quan on vidas del somito », prefere kam per la\nl
formo pasiva : « La voyo sequesis, la kolino acensesis, la kastelo\nl
vizitesis, e la panoramo admiresis, qua videsas del somito. »

(2) Okazione ni volas refutar un foyo pluse ti qui asertas, ke\nl
l'akuzativo esas necesa por posibligar l'inversigi, quin ula lingui\nl
tante prizas. Yen exemplo di D. frazo hazarde prenita : « \textit{Das}\nl
\textit{Wissen eignet sich das Kind} bereits in allerfrühester Jugend an... »\nl
Ube esas l'akuzativo? \textit{Das Wissen} o \textit{das Kind?} Nur la senco e la\nl
komuna raciono indikas ol, e posibligas komprenar ta frazo. E to\nl
eventas tre ofte en D., nam l'akuzativo esas identa a la nominativo\nl
en la pluralo, ed en la feminal e neutra singularo, do mezvalore\nl
5 foyi ek 6! Esas do multe plu simpla e plu sekura sequar sempre la\nl
normala vortordino, e dicar ex. en Ido : « L'infanto aquiras la savo\nl
ja en la maxim frua yuneso. »%
}

\VUcontinue
\textit{relati maxim utila e por la profito di la futura homaro de}\nl
\textit{grava konsequi maxim fekunda inventuro.} --- \VUgras{Vu esas ne}\cc
\VUgras{justa segun mea opiniono,} e ne : \textit{vu esas segun mea opiniono}\nl
\textit{nejusta} (1).

Ek \textit{Progreso}, VII, 161 : « \textit{Quo esas verajo decidar nur}\nl
\textit{povas intelekto.} » Yen bela exemplo di kompleta inversigo\nl
di la natural ordino di la pensi. Unlatere, oportis pozar\nl
\textit{decidar} ante lua komplemento « quo esas verajo ». Altralatere,\nl
oportis pozar \textit{nur} nemediate ante \textit{intelekto}, quan ol referas\nl
On povis do uzar un de la sequanta frazordini : « Decidar\nl
quo esas verajo, nur intelekto povas. » --- « Quo esas verajo,\nl
nur intelekto povas decidar. » --- « Nur intelekto povas deci\cc
dar, quo esas verajo. » E la triesma formo semblas a ni\nl
la maxim klara por omni, do preferinda. » (2).

\VUsaut{4.61mm}
\VUfilet{8.47mm}
\VUsaut{3.20mm}
\VUtitre{10.27pt}{38}{\VUetroit{0.857}{PUNTIZADO}}
\VUsaut{0.87mm}
\VUsoustitre{5.67pt}{285}{\textit{(Apendico 8-ma.)}}
\VUsaut{3.60mm}

\textit{Punto} (.) uzesas por separar la \textit{frazi}, t. e. expresuri di\nl
integra e kompleta pensi tote nedependanta. Bona precepto\nl
di stilo esas ta, quan S\textsuperscript{ro} \textsc{Peus} recevis de ula sua maestro :\nl
\end{VUpage}

% -------------------------------------------------------------------
%  FEUILLET 222 — LE FILET DE NOTE EST ATYPIQUEMENT LONG : 226 px la ou
%  huit feuillets temoins en donnent 210 a 214. Verifie a l'oeil et a
%  la main (y 1828-1830, x 260-505, seul sur sa rangee, cale sur la
%  marge) : il est bien reel, seulement plus long de 6 %. Ni l'echelle
%  du scan ni le pas des lignes n'expliquent l'ecart. On garde la
%  largeur reglementaire du volume, faute de savoir si le compositeur
%  a change de fer ou si le trait a bave.
%  Le blanc filet -> premiere note ne vaut ici que 18 px : la note
%  s'ouvre sur une parenthese, qui monte plus haut qu'une capitale.
% -------------------------------------------------------------------
\begin{VUpage}[222]{218}
\VUnotes{159.09mm}{%
(1) Pro to en Germania on separas per lakuni : 9 365 724.%
}

\VUcontinue
\textit{Uzez multa punti}, t. e. facez kurta frazi. En matematiko, la\nl
punto uzesas : 1\textsuperscript{e} en longa nombri por separar la mili, ex. :\nl
9.365.724 plu klara kam 9365724; 2\textsuperscript{e} en algebro lu havas\nl
la sama signifiko kam × (1).

\textit{Mayuskuli}. --- Okazione la punto, on devas regulizar l'uzo\nl
di la mayuskuli. On uzas komencal litero mayuskula :\parplein

1\textsuperscript{e} por la propra nomi, inkluzante en ici la nomi di landi,\nl
monti, fluvii, riveri, la nomi di populi, societi od institu\cc
curi, qui nature kontenas la religii e lia praktikanti. On\nl
lasas al derivaji di ta nomi la mayuskulo di la radiko.\nl
Ex. : \textit{la Franco, la Franci, la Franca} (\textit{linguo}); \textit{la Kristanismo}\nl
(de \textit{Kristo}), \textit{la Kristani; Alpala, Alpano} (de \textit{Alpi}). Co evitas\nl
la dicerno a qua obligas ula lingui pri la kazo en qua on\nl
konservas la mayuskulo od on remplasigas lu per minuskulo.\parplein

2\textsuperscript{e} por l'unesma vorto di singla frazo, e konseque, pos\nl
omna punto.

3\textsuperscript{e} pro politeso, en ula kazi, en tituli. Ex. : \textit{Sioro Pro}\cc
\textit{fesoro. En la nomo dil Rejo.}

On atencez la skribo di kelka mayuskuli :

1\textsuperscript{e} On evitez skribar I quale J, quale ofte eventas en\nl
Germania. On skribez Ido, ne Jdo.

2\textsuperscript{e} On distingez, skribante, la mayuskuli K e R.

3\textsuperscript{e} On distingez F e T : l'infra parto di T flexesas ad\nl
dextre; ta di F ad sinistre.

\textit{Komo} (,) uzesas por separar la \textit{propozicioni} qui kompozas\nl
frazo. Pos la relativa pronomi, on adoptas la franca uzo di\nl
la komo (D. 1062) : On distingas du uzi tre diversa di la\nl
relativa pronomi, qui konstitucas preske du diversa senci.\parplein

La frazo : « Me ne amas la infanti, qui facas bruiso »\nl
povas havar du senci : 1\textsuperscript{e} Me ne amas ti de la infanti, qui\nl
facas bruiso; 2\textsuperscript{e} Me ne amas la infanti, pro ke li (konocate\nl
e generale) facas bruiso. En l'unesma kazo, la relativo esas\nl
\textit{determinanta} : il \textit{determinas} la speco de infanti, quin me ne\nl
amas. En la duesma ol esas \textit{qualifikanta} od \textit{explikanta} : ol\nl
qualifikas generale omna infanti. On uzas la komo (avan\nl
\textit{qua}) en la duesma kazo, ne en l'unesma.

Altra exempli : « La homo qua lektas la jurnali devas ne\nl
kredar omno quon lu lektas. --- Regardez ta homo, qua\nl
\end{VUpage}

% -------------------------------------------------------------------
%  FEUILLET 223 — cinquieme page du volume sans note. Les trois
%  segments du balayage grossier (y 364, 1525, 1774) ont ete lus a
%  l'agrandissement : ce sont des suites de lettres du texte courant.
%  Aucun run horizontal continu de 60 px sur toute la page, hors bord
%  de scan. COQUILLE conservee : « separer » pour « separar », alors
%  que la meme page donne « separesas » et « separesar ».
% -------------------------------------------------------------------
\begin{VUpage}[223]{219}
\VUcontinue
lektas jurnalo : lu audas tamen omno quon ni dicas. » « La\nl
amiko qua ne audacas dicar la verajo ne esas vera amiko. ---\nl
Mea olda amiko, qua kustumis dicar la verajo, ne audacis\nl
takaze dicar ol a me. » « La unesma soldato qua vidis l'ene\cc
miki hastis avertar la generalo. --- Napoléon, qua vidis l'ene\cc
miki avancar, sendis kontre li sua olda guardo. » Ica exempli\nl
sugestas utila remarko. Omnafoye kande la substantivo esas\nl
determinita, la \textit{qua} esas nur explikanta, do postulas komo.\nl
Or ico eventas : 1\textsuperscript{e} kande la substantivo esas propra nomo :\nl
evidente ne existas plura Napoléon, do nulo \textit{qua} povus deter\cc
minar lu plu komplete; 2\textsuperscript{e} kande ol esas akompanata da\nl
\textit{determinanta} adjektivo, sive demonstrativa, sive posedala :\nl
« ta homo » esas ja suficante e komplete determinita per\nl
la demonstrativo (o per la gesto, qua akompanas ol); « mea\nl
olda amiko » anke esas determinita per ica epiteti.

Komprenende, la regulo valoras anke, se \textit{qua} esas preirita\nl
da prepoziciono : « La amiko de qua me demandis konsilo\nl
aprobis mea penso-maniero. --- Mea patro, de qua me de\cc
mandis konsilo, aprobis mea penso-maniero. »

En matematiko, en la nombri decimala, la \textit{integri} separ\cc
esas de la \textit{decimali} per \textit{komo} (ne per punto, qua havas altra\nl
uzo, V. p. 217). Ex. : 3, 14.

\textit{Punto-komo} (;) uzesas por separer propozicioni o frazi\nl
gramatikale nedependanta, ma ligita per la senco.

\textit{Bi-punto} (:) uzesas por anuncar expliko : ol dikas, ke\nl
la sequanta frazo explikas la preiranta. Ol anuncas anke\nl
citajo, ma lore sequesas da \textit{cito-hoketi}.

\textit{Puntaro} (...) indikas interrupto di la frazo, sive da altra\nl
parolanto, sive da la parolanto ipsa, qua haltas por retenar\nl
o chanjar l'expreso di sua penso. La Germani uzas vice to\nl
la \textit{Gedankenstrich} (vortope : \textit{penso-streko}) ma to esas mala\nl
uzado, nam la \textit{streko} havas altra senco e necesa uzado. (Vid.\nl
p. 221).

\textit{Parentezi} () uzesas por inkluzar frazo, propoziciono o\nl
vorto, qua esas aparta, e devas separesar de la cetera texto.\nl
To esas generale remarko laterala, quan on insertas, e qua\nl
interruptas la rekta ordino di la penso. De to konsequas\nl
konsilo : on ne trouzez la parentezo, pro ke ol riskas trublar\nl
la kompreno.

Pro ke la parentezi esas esence destinata ad \textit{inkluzar}, on\nl
\end{VUpage}

% -------------------------------------------------------------------
%  FEUILLET 224 — DEUX HORS-TEXTE, et TROIS RANGEES DE POINTS que le
%  detecteur a manquees. La geometrie n'annonce que 34 lignes la ou la
%  page en porte 37 : les trois rangees de points qui accompagnent
%  l'accolade verticale ne portent chacune que six points de 5 px.
%  La preuve est arithmetique : lignes de base 828,2 — 911 — 952 — 993
%  — 1075,8, soit des ecarts de 2, 1, 1, 2 pas exactement.
%
%  L'ACCOLADE VERTICALE est OUVRANTE, pointe a gauche vers l'unique
%  rangee : c'est ce que la page enonce trois lignes plus bas, « la
%  pinto devas turnesar ad l'unika lineo ». Encre x 786-796, y 877-995.
%
%  L'ACCOLADE HORIZONTALE de l'arbre genealogique est la seule du
%  volume ; d'ou \VUaccoladeH, ajoutee au preambule. Encre x 524-1063,
%  y 1339-1355, pointe au milieu exact.
%
%  COQUILLES conservees : « (e ne inverse!) Ni konsilas » sans point
%  apres la parenthese ; « 1) ) » avec son blanc, qui est la
%  demonstration meme de l'auteur ; et les accolades soudees « {} » de
%  la ligne 10, que la ligne 11 separe pourtant.
% -------------------------------------------------------------------
\begin{VUpage}[224]{220}
\VUcontinue
darfas nultempe uzar una sen l'altra, quale on facas ofte\nl
kun la nombri. Vice « 1) », on devas skribar, o « (1) », o\nl
simple « 1\textsuperscript{e} » o 1-e (abreviuro di \textit{unesme}). La uzado « 1) »\nl
havas grava detrimento : kande on vidas tala parentezo\nl
\textit{klozanta}, on serchas spontane la parentezo \textit{apertanta}, ed on\nl
ne trovas ol. Pluse, eventas ofte ke ta parentezo « vidva »\nl
trovesas en od apud altra vera (duopla) parentezo, e lore\nl
on havas, sive : « 1) ) », sive « )1) », du dispozuri egale\nl
absurda. Do : \textit{nultempe uzez un parentezo sola!}

\textit{Kramponi} [] ed \textit{Embracili} \{\} uzesas, en matematiko, kom\nl
parentezi di duesma o triesma grado, ex. : \{ ...[...(...)...]...\}.\nl
En la « prozo », on povas uzar la \textit{kramponi} en simila kazo,\nl
o kom aparta parentezi. On uzas l'\textit{embracilo} (unika) por\nl
korespondigar un lineo (unlatere) a plura linei (altra\cc
latere), tale :

\VUblanc{3.50mm}
\VUrang{\VUcase{48.77mm}{.}\VUcase{50.63mm}{.}\VUcase{52.49mm}{.}}
\VUrang{\VUcase{38.10mm}{.}\VUcase{39.96mm}{.}\VUcase{41.83mm}{.}%
\VUcase{45.13mm}{\VUaccolade{28.33pt}{1.47pt}}%
\VUcase{48.86mm}{.}\VUcase{50.72mm}{.}\VUcase{52.49mm}{.}}
\VUrang{\VUcase{48.68mm}{.}\VUcase{50.55mm}{.}\VUcase{52.41mm}{.}}
\VUblanc{3.46mm}

La \textit{pinto} devas turnesar ad l'unika lineo, e la \textit{branchi} (la\nl
konkaveso) ad la plura linei (e ne inverse!) Ni konsilas\nl
uzar minim ofte l'embracilo, pro ke ol komplikas e des\cc
faciligas la kompostado. Ol esas utila nur en sinoptikala\nl
tabeli. Ol povas uzesar anke horizontale, ex. en genealogikal\nl
arbori :

\VUrang{\VUcase{38.10mm}{Ludovikus}}
\VUblanc{2.58mm}
\VUrang{\VUcase{22.95mm}{\VUdecale{9.98pt}{\VUaccoladeH{45.63mm}}}%
\VUcase{23.20mm}{Petrus}\VUcase{34.55mm}{Paulus}%
\VUcase{46.74mm}{Ioannes}\VUcase{59.86mm}{Maria}}
\VUblanc{3.48mm}

\textit{Streketo} (-) esas ortografiala signo, e funcionas kom ligilo.\nl
Ol unionas la parti di vorto kompozita (V. ye \textit{Kompozado}).\nl
Ol indikas anke la seko di vorto de un lineo ad altra, pri\nl
qua ni quik parolos.

\textit{Seko di la Vorti}. [D. 485]. On admisas kompleta libereso\nl
en la seko di la vorti de lineo a lineo, ecepte ke singla parto\nl
devas kontenar vokalo, e ke la digrami o diftongi devas ne\nl
dividesar. Ex. on darfas sekar tale la vorto \textit{mustar} : \textit{mu-star},\nl
\textit{mus-tar}, o \textit{must-ar}. Ma \textit{neutro}, \textit{mashino} sekesas \textit{neu-tro},\nl
\textit{ma-shino}, e ne \textit{ne-utro}, \textit{mas-hino}. \textit{On nultempe uzez duopla}\nl
\end{VUpage}

% -------------------------------------------------------------------
%  FEUILLET 225 — septieme page du volume sans note. Les cinq segments
%  du balayage grossier tombent tous a trois ou quatre pixels sous une
%  ligne de base annoncee : ce sont des empattements soudes.
%  COQUILLE conservee : la virgule apres « questionanta » la ou le sens
%  appelle un point (ligne 37).
%  Le fac-simile porte ici la CROIX de renvoi de note, sans precedent
%  dans le volume.
% -------------------------------------------------------------------
\begin{VUpage}[225]{221}
\VUcontinue
\textit{streko} (=) rezervenda a la signo di egaleso, nek iterez la\nl
streketo ye komenco di la lineo sequanta.

\textit{Streko} (---), plu longa kam streketo, esas, tote kontree,\nl
\textit{separilo}, e la maxim grava. Ol indikas, sive chanjo di paro\cc
lanto (t. e. separas la dici, respondi, di diversa parolanti),\nl
sive chanjo di temo. Ol separas do multe plu kam la punto.\nl
Ma en la lasta kazo esas preferinda uzar la \textit{alineo}, qua in\cc
dikas plu klare ta separo; ol havas l'avantajo insertar en la\nl
texto « blankaji », vakuaji, qui repozas la okulo e furnisas\nl
ad ol halto-punto.

\textit{Cito-hoketi} (« »). La cito-hoketi indikas ed inkluzas la\nl
paroli o vorti, quin on citas. Li devas turnar sua konkaveso\nl
ad interne, t. e. ad la texto inkluzata, same kam la parentezi.\nl
Omna altra uzado esas mala, pro ke ol esas konfuziganta.\parplein

\textit{Noto-referi}. Por indikar la noti (ped-noti), nula moyeno\nl
esas tam simpla e komoda kam la \textit{numeri} : nam li esas la\nl
maxim bona e klara signo, e li esas en senlimita provizuro.\nl
Se on uzas steleti (*) o kruci (†), quale ula populi, e se on\nl
havas multa noti en un sama pagino, on ne plus savas quale\nl
helpar su : la procedo, iterar la steleti o la kruci, esas vere\nl
« sovaja », t. e. apartenas a la primitiva stando di civilizeso,\nl
en qua on ne posedis la cifri e reprezentis la nombri per\nl
streki; od on esas obligata rekursar anke ad altra signi,\nl
qui havas altra senci od uzi (ex. §, qua signifikas para\cc
grafo). Pluse, on bezonas ofte la steleto e la kruco por altra\nl
signifiki : la steleto indikas en linguistiko formo konjektita,\nl
ne atestita; en Ido, formo ne oficala, o teknikala; la kruco\nl
indikas generale mortinto, o dato di morto, ed en filologio,\nl
arkaika formo. To omna pruvas, ke ol esas nur remediacho,\nl
e ke nur la \textit{numeri} esas uzenda por la noto-referi.

\textit{Klamo-punto} (!) uzesas pos klamo o frazo klamanta o kla\cc
mata. Kelka populi uzas ol pos omna voko od interpelo,\nl
mem pos \textit{Sioro} en komenco di letri, unvorte, pos omna voka\cc
tivo. Ico esas forsan exajero, nam kande on parolas od ulu,\nl
on ne klamas : « Sioro! »

\textit{Question-punto} (?) uzesas pos frazo \textit{direte} questionanta,\nl
ne pos subordinala propoziciono questionanta, Ex. : « Qua\nl
venas? », ma : « Me questionas, qua venas. »

Pri ca du lasta signi, on propozis sequar l'exemplo di la\nl
Hispana, t. e. pozar li, renversita, ante la frazo klamanta\nl
\end{VUpage}

% -------------------------------------------------------------------
%  FEUILLET 226 — FIN DE L'APENDICO 8, OUVERTURE DE L'APENDICO 9.
%  Le filet court de fin de section echappe au detecteur pour la
%  troisieme fois du volume : 96 px de large, 5 px d'epais, a y
%  1288-1293, x 754-849, centre a 801,5 pour un milieu de bloc a 804,5.
%  Il tombe entre les indices 24 et 25, dans un intervalle que la liste
%  d'exces ne signalait pas du tout (188,8 px, soit 4,5 pas).
%  Huitieme page du volume sans note : les deux segments de y 1764-1765
%  sont les empattements de « komencas per la no... », soudes par le
%  seuil grossier trois pixels au-dessus de leur ligne de base.
% -------------------------------------------------------------------
\begin{VUpage}[226]{222}
\VUcontinue
o questionanta. La motivo esas : avertar la lektanto pri la\nl
karaktero di tala frazo, por ke lu konforme modifikez sua\nl
pronunco. Esas vera ke, pro manko di tala averto, on kelka\cc
foye eroras pri la karaktero di la frazo e pronuncas ol male,\nl
o mem devas repetar ol kun la justa pronunco. Ta propozo\nl
esas do konsiderinda; la unika kontrea argumento esas, ke\nl
l'altra uzado esas plu internaciona. Se on adoptus ol, la\nl
signi ! e ? divenus \textit{duopla} ed inkluzanta, same kam la paren\cc
tezi e cito-hoketi. Ed on vidas ke, pro la sama motivo di\nl
klareso, li devas prizentar du formi inversa, quale la \textit{ques}\cc
\textit{tion-punti} en la Hispana.

\textit{Apostrofo} (') uzesas por indikar eliziono (V. § 6).

\VUblanc{1.96mm}
\textit{Generala remarko}. --- On devas ne mis-evaluar l'importo\nl
di la puntizado, e desprizar la koncernanta reguli kom super\cc
flua o minucioza. La puntizado esas necesa por la klareso,\nl
do por la perfekta e sekura interkomprenado; ol havas signi\cc
fiko por la pronuncado. La punto reprezentas halto pasable\nl
granda; la bi-punto e la punto-komo, halto min granda;\nl
la komo, halto plu mikra, ma sentebla. Egardar ta signi e\nl
facar la korespondanta pauzi, esas l'unesma regulo di la\nl
diciono, e la maxim importanta.

[Segun \textit{Progr.}, IV, 531 (L. \textsc{Couturat}) e VI, 380]. Decido,\nl
866 : On adoptas la reguli propozita por la puntizado (IV,\nl
531). Decido 1062 pri la komo (VI, 211).

\VUsaut{6.86mm}
\VUfilet{8.13mm}
\VUsaut{3.20mm}
\VUtitre{10.45pt}{-49}{NOMI. ADRESI.}
\VUsaut{1.56mm}
\VUsoustitre{5.85pt}{260}{\textit{(Apendico 9-ma.)}}
\VUsaut{4.24mm}

[D. 1144] ... l'Akademio rekomendas enuncar la personal\nl
nomi komencante per la prenomo. (Ica decido ne koncernas\nl
la nomi en listi, adresari, e. c., ube on pozas la familial\nl
nomo unesme, por l'alfabetal ordino.)

On rekomendas por l'adresi la sistemo (Franca), qua\nl
komencas per la nomo di persono e finas per l'urbo e la\nl
lando, ed on repulsas la sistemo inversa (Rusa) [\textit{Progr.}, IV,\nl
470; VI, 52; VII, 162].
\end{VUpage}

% -------------------------------------------------------------------
%  FEUILLET 227 — QUATRIEME FILET COURT MANQUE PAR LE DETECTEUR, et un
%  cinquieme filet, plus long, qui clot l'apendico 10.
%  Le court : x 575-673 (99 px), y 766-768, centre 624 pour un milieu de
%  bloc a 626. Il se loge dans un intervalle de 178,4 px que la liste
%  d'exces ne signalait pas du tout, comme au folio 222.
%  Le long : x 525-726 (202 px), y 1793-1797, centre 625,5, a 64 px sous
%  la derniere ligne. Il fait le DOUBLE des filets de section ; c'est une
%  division plus forte, non un filet de note (qui se pose au fer a
%  gauche, et sous lequel il y aurait du texte).
%  Le bloc en escalier des indices 2 a 5 est un modele d'adresse : les
%  bords droits ne s'alignent pas (838, 853, 809, 767), donc ce n'est pas
%  un bloc centre mais quatre lignes a renfoncement croissant.
%  Dixieme page du volume sans note.
% -------------------------------------------------------------------
\begin{VUpage}[227]{223}
\VUornement[17.10mm]{159.09mm}

Exemplo di korekta adreso :

\VUblanc{1.94mm}
\VUrang{\VUcase{26.75mm}{Monsieur Louis Legrand}}
\VUblanc{0.60mm}
\VUrang{\VUcase{34.97mm}{64, rue Notre-Dame}}
\VUblanc{0.58mm}
\VUrang{\VUcase{41.23mm}{\textit{Sens (Yonne)}}}
\VUrang{\VUcase{46.82mm}{\textsc{France}}}
\VUblanc{1.96mm}

Se on volas indikar, dop nomo di persono, ta di lua urbo\nl
o lando (ex. : nomo di delegito, en raporto pri kongreso), ni\nl
konsilas pozar ta lasta nomo \textit{inter parentezi}, ex. : \textit{S}\textsuperscript{ro} \textit{Martin}\nl
(\textit{Paris}), e ne, segun la Germana kustumo, separar la du nomi\nl
per streketo : \textit{S}\textsuperscript{ro} \textit{Martin-Paris}, qua kredigas, ke la nomo di\nl
la persono esas \textit{Martin-Paris}.

\VUsaut{5.99mm}
\VUfilet{8.38mm}
\VUsaut{3.20mm}
\VUtitre{10.27pt}{42}{FORMULI DI POLITESO EN LETRI.}
\VUsaut{1.21mm}
\VUsoustitre{5.67pt}{222}{\textit{(Apendico 10-ma.)}}
\VUsaut{4.24mm}

Ta formuli esas afero di nacionala kustumo e stilo, e la\nl
simpla traduko di tala nacionala formuli genitus ne nur\nl
senfina diverseso, ma frazi stranja, nekomprenebla o mis\cc
komprenebla. Semblis do necesa fixigar, per konvenciono,\nl
to quo devas konsideresar kom \textit{polita} formuli.

Ye la komenco di letro, ni uzez nur \textit{Sioro}, e se la kores\cc
pondanto havas ula titulo, funciono o profesiono, qua kon\cc
sideresas en nia korespondado, ni skribez : \textit{Sioro Prezidero,}\nl
\textit{Sioro Profesoro,} e. c. Por iti, qui havas funciono o situeso,\nl
por qua la simpla \textit{Sioro} ne semblas suficanta, ni havas la\nl
vorto \textit{Sinioro} : \textit{Sinioro Episkopo, Sinioro Ministro}.

Ye la fino di letro, ni generale uzez : \textit{Kun sincera saluto}.\nl
Por siorini e la personi, quin ni qualifikas « sinioro » ni\nl
dicez : \textit{Kun respektoza saluto}.

Komprenende ta reguli ne koncernas la korespondado kun\nl
amiki, kamaradi, parenti, qua admisas tre granda diverseso\nl
en ta formuli.

[Segun \textit{Progr.}, II, 679; IV, 470].
\end{VUpage}

% -------------------------------------------------------------------
%  FEUILLET 228 — DERNIERE PAGE DES APENDICI. Sept lignes, puis
%  l'ORNEMENT DE CLOTURE : 296 px = 25,06 mm, le plus long filet du
%  volume, a y 1122-1125, centre a 806,5 pour un milieu de bloc a 809,5
%  — le meme decalage de 3 px que celui du folio 222. Sa longueur est
%  celle de \VUornementLargeur (25,8 mm), l'ornement de fin de partie.
%  Ordonnee : 24,30 + (1123-210) x 25,4/300 = 101,60 mm.
%  Onzieme page du volume sans note ; balayage complet a trois seuils,
%  rien d'autre sur la page que le texte et cet ornement.
% -------------------------------------------------------------------
\begin{VUpage}[228]{224}
\VUornement[25.06mm]{101.60mm}

\textit{Hike ni pozas la du prepozicioni} \VUgras{tra e trans} \textit{omisita} pos\nl
til p. 82.

\VUblanc{1.99mm}
\VUgras{Tra} = de un latero od extremajo a l'altra, interne. Ex. :\nl
\VUgras{Li pasis tra la foresto. Il sinkis sua espado tra ilua pektoro.}\parplein

\VUblanc{1.99mm}
\VUgras{Trans} = adsur l'altra latero. Ex. : \VUgras{On pasas tra la rivero}\nl
\VUgras{per vadeyo e trans olu per ponto. Ne irez trans la lago, nam}\nl
\VUgras{ibe vu jenesus dal turisti.}
\end{VUpage}

% -------------------------------------------------------------------
%  FEUILLET 229 — OUVERTURE DU TABELO, l'index analytique. Mise en page
%  nouvelle : corps plus petit (hauteur d'x 17 px contre 21), points de
%  conduite au pas de 17,25 px, numeros de page au fer droit. D'ou
%  \VUpageIndex, \VUindex et \VUindexNu, ajoutees au preambule.
%  LA JUSTIFICATION EST CELLE DU VOLUME : les 976 a 980 px rendus par la
%  geometrie sont un artefact de text_region, qui coupe la colonne des
%  numeros, separee des points par un blanc de 73 px.
%  LA PAGE NE PORTE PAS DE FOLIO — d'ou l'argument vide.
%  LIGNE INVISIBLE RETROUVEE : « Pagini. », 88 px de large, sous le seuil
%  du detecteur. La page compte 30 lignes, non 29.
%  FILET CENTRE de 155 px a y 510-513 : une cinquieme longueur, apres
%  90-100, 202 et 296 px.
%  L'ORDONNEE DE NOTE RENDUE PAR L'OUTIL (117,35 mm) EST FAUSSE : elle
%  tombe au milieu d'une ligne de texte. Cette page n'a aucune note.
% -------------------------------------------------------------------
\begin{VUpage}[229]{}
\VUsignature{168.70mm}{8}

\VUcentre{15.5pt}{352}{TABELO}
\VUsaut{3.20mm}
\VUfilet{13.12mm}
\VUsaut{5.50mm}
\VUcentre{8.6pt}{190}{UNESMA PARTO}
\VUsaut{4.40mm}
\VUcentre{10.1pt}{152}{\VUgras{MORFOLOGIO E SINTAXO}}
\VUsaut{0.54mm}

\VUpageIndex
\VUindexNu{}{Pagini.}
\VUblanc{1.10mm}
\VUindex{\VUgras{Alfabeto}}{11}
\VUindex{\VUgras{Pronunco dil vokali}}{11}
\VUindex{\VUgras{Pronunco dil konsonanti e digrami}}{13}
\VUindex{\VUgras{Acento tonika}}{15}
\VUindex{\VUgras{Artiklo}}{18}
\VUindex{\VUgras{Substantivo}}{21}
\VUindex{\VUgras{Propra nomi}}{24}
\VUindex{\VUgras{Adjektivo qualifikanta e lua plaso}}{27}
\VUindex{\VUgras{Gradi komparala}}{31}
\VUindex{\VUgras{Personal pronomi}}{32}
\VUindex{\VUgras{Posedal adjektivi e pronomi}}{34}
\VUindex{\VUgras{Demonstativ adjektivi-pronomi}}{36}
\VUindex{\VUgras{Relativa e questionala adjektivi-pronomi}}{38}
\VUindex{\VUgras{Pronomo « lo »}}{39}
\VUindex{\VUgras{Adjektivi-pronomi nedefinita}}{40}
\VUindexNu{\VUgras{Verbo :}}{}
\VUindex[\VUindexRenfoncement]{transitiva}{46 e 51}
\VUindex[\VUindexRenfoncement]{netransitiva}{51}
\VUindex[\VUindexRenfoncement]{pasiva}{48}
\VUindex[\VUindexRenfoncement]{mixita}{52}
\VUindex[\VUindexRenfoncement]{reflektiva}{53 e 54}
\VUindex[\VUindexRenfoncement]{reciproka}{54}
\VUindex[\VUindexRenfoncement]{unpersona}{54}
\VUindexNu{\VUgras{Adverbi :}}{56}
\VUindex[\VUindexRenfoncement]{di quanteso}{57}
\VUindex[\VUindexRenfoncement]{di tempo}{59}
\end{VUpage}

% -------------------------------------------------------------------
%  FEUILLET 230 — suite du TABELO. La page PORTE un folio, contrairement
%  a la precedente. Le detecteur de filet se trompe une seconde fois de
%  suite : les 45,81 mm qu'il annonce tombent sur le MENEUR DE POINTS de
%  l'entree « an », une ligne d'encre horizontale de x 300 a 1233 qu'il
%  prend pour un filet. Aucune note sur cette page non plus.
%  « Prepozicioni : » est une tete de rubrique, avec son propre numero
%  et SANS points de conduite : ce n'est pas un titre, il ne passe pas a
%  \VUtitre.
% -------------------------------------------------------------------
\begin{VUpage}[230]{226}
\VUpageIndex
\VUindex[\VUindexRenfoncement]{di loko}{62}
\VUindex[\VUindexRenfoncement]{di maniero}{63}
\VUindex[\VUindexRenfoncement]{di afirmo, nego, dubo}{66}
\VUblanc{2.91mm}
\VUindexNu{\VUgras{Prepozicioni :}}{67}
\VUblanc{0.43mm}
\VUindex[\VUindexRenfoncement]{ad (a)}{68}
\VUindex[\VUindexRenfoncement]{alonge}{68}
\VUindex[\VUindexRenfoncement]{an}{68}
\VUindex[\VUindexRenfoncement]{ante}{69}
\VUindex[\VUindexRenfoncement]{apud}{70}
\VUindex[\VUindexRenfoncement]{avan}{70}
\VUindex[\VUindexRenfoncement]{che}{71}
\VUindex[\VUindexRenfoncement]{cirkum}{71}
\VUindex[\VUindexRenfoncement]{cis}{71}
\VUindex[\VUindexRenfoncement]{da}{71}
\VUindex[\VUindexRenfoncement]{de}{71}
\VUindex[\VUindexRenfoncement]{di}{73}
\VUindex[\VUindexRenfoncement]{dop}{74}
\VUindex[\VUindexRenfoncement]{dum}{74}
\VUindex[\VUindexRenfoncement]{ek}{75}
\VUindex[\VUindexRenfoncement]{en}{76}
\VUindex[\VUindexRenfoncement]{erste}{76}
\VUindex[\VUindexRenfoncement]{exter}{76}
\VUindex[\VUindexRenfoncement]{for}{76}
\VUindex[\VUindexRenfoncement]{inter}{77}
\VUindex[\VUindexRenfoncement]{kontre}{77}
\VUindex[\VUindexRenfoncement]{koram}{77}
\VUindex[\VUindexRenfoncement]{kun}{78}
\VUindex[\VUindexRenfoncement]{lor}{78}
\VUindex[\VUindexRenfoncement]{malgre}{79}
\VUindex[\VUindexRenfoncement]{per}{79}
\VUindex[\VUindexRenfoncement]{po}{79}
\VUindex[\VUindexRenfoncement]{por}{80}
\VUindex[\VUindexRenfoncement]{pos}{80}
\VUindex[\VUindexRenfoncement]{preter}{80}
\VUindex[\VUindexRenfoncement]{pri}{80}
\VUindex[\VUindexRenfoncement]{pro}{81}
\VUindex[\VUindexRenfoncement]{proxim}{81}
\VUindex[\VUindexRenfoncement]{segun}{81}
\VUindex[\VUindexRenfoncement]{sen}{81}
\end{VUpage}

% -------------------------------------------------------------------
%  FEUILLET 231 — TROISIEME MENSONGE DU DETECTEUR DE FILET SUR LE
%  TABELO. Il rendait 74,08 mm ; a cette hauteur se trouve la LIGNE DE
%  BASE de la vedette demi-grasse « Konjuncioni : », dont la fermeture
%  morphologique soude les empattements en une plage continue de 185 px
%  partant de la marge. Deux indices suffisaient a la recuser : 185 px
%  au lieu des 201-246 reglementaires, et une hauteur d'empattement.
%  outils/filet.py resserre desormais sa fenetre de longueur.
%  LA POSITION N'EMPORTE PAS LA GRAISSE : six lignes au fer sur dix sont
%  ici en romain maigre. On ne deduit donc jamais l'une de l'autre.
%  Douzieme page du volume sans note.
% -------------------------------------------------------------------
\begin{VUpage}[231]{227}
\VUpageIndex
\VUindex[\VUindexRenfoncement]{sub}{82}
\VUindex[\VUindexRenfoncement]{super}{82}
\VUindex[\VUindexRenfoncement]{sur}{82}
\VUindex[\VUindexRenfoncement]{til}{82}
\VUindex[\VUindexRenfoncement]{tra}{224}
\VUindex[\VUindexRenfoncement]{trans}{224}
\VUindex[\VUindexRenfoncement]{ultre}{82}
\VUindex[\VUindexRenfoncement]{vice}{83}
\VUindex[\VUindexRenfoncement]{ye}{83}
\VUblanc{0.92mm}
\VUindex{Komplemento di adverbo prepoziciona}{85}
\VUblanc{0.64mm}
\VUindex{Plaso dil komplemento di prepoziciono}{86}
\VUblanc{0.72mm}
\VUindex{Prepozicioni avan infinitivo}{86}
\VUblanc{0.68mm}
\VUindex{Prepozicioni kun verbi}{87}
\VUblanc{1.44mm}
\VUindexNu{\VUgras{Konjuncioni :}}{88}
\VUblanc{0.25mm}
\VUindexNu{Koordinala :}{88}
\VUindex[\VUindexRenfoncement]{do}{88}
\VUindex[\VUindexRenfoncement]{e ed}{88}
\VUindex[\VUindexRenfoncement]{ma}{88}
\VUindex[\VUindexRenfoncement]{nam}{88}
\VUindex[\VUindexRenfoncement]{or}{89}
\VUindex[\VUindexRenfoncement]{od od}{89}
\VUindex[\VUindexRenfoncement]{sive}{89}
\VUindex[\VUindexRenfoncement]{tamen}{89}
\VUindex[\VUindexRenfoncement]{lore}{90}
\VUindex[\VUindexRenfoncement]{nek}{90}
\VUindex[\VUindexRenfoncement]{yen}{90}
\VUblanc{0.30mm}
\VUindexNu{Subordinala :}{90}
\VUindex[\VUindexRenfoncement]{ke}{90}
\VUindex[\VUindexRenfoncement]{se}{90}
\VUindex[\VUindexRenfoncement]{se nur}{90}
\VUindex[\VUindexRenfoncement]{quale se}{90}
\VUindex[\VUindexRenfoncement]{quankam}{91}
\VUblanc{0.34mm}
\VUindex{Konjuncioni kompozita}{91}
\VUblanc{0.72mm}
\VUindex{Konjunciono questionala}{92}
\VUblanc{1.06mm}
\VUindexNu{\VUgras{Interjecioni :}}{93}
\VUindex[\VUindexRenfoncement]{aye}{93}
\VUindex[\VUindexRenfoncement]{ba}{93}
\VUindex[\VUindexRenfoncement]{fi}{93}
\end{VUpage}

% -------------------------------------------------------------------
%  FEUILLET 232 — page du TABELO. « Filet non trouve » est ici exact :
%  le seul candidat est la soudure d'empattements de la vedette
%  « shamo », 65 px, calee sur le renfoncement des sous-entrees et non
%  sur le fer. Treizieme page du volume sans note.
%  UN SECOND EXCES ECHAPPAIT A LA GEOMETRIE : le blanc de +15,2 px qui
%  precede « Nombri : », manque parce que le detecteur n'a pas su lire
%  la ligne de base de « yen » (bas = -1,0 dans la geometrie). Les deux
%  blancs de la page sont de meme nature et de meme calibre.
% -------------------------------------------------------------------
\begin{VUpage}[232]{228}
\VUpageIndex
\VUindex[\VUindexRenfoncement]{ha}{93}
\VUindex[\VUindexRenfoncement]{he}{93}
\VUindex[\VUindexRenfoncement]{hem}{93}
\VUindex[\VUindexRenfoncement]{ho}{93}
\VUindex[\VUindexRenfoncement]{hola}{93}
\VUindex[\VUindexRenfoncement]{hop}{93}
\VUindex[\VUindexRenfoncement]{hu-hu}{93}
\VUindex[\VUindexRenfoncement]{hura}{93}
\VUindex[\VUindexRenfoncement]{krak}{93}
\VUindex[\VUindexRenfoncement]{krik}{93}
\VUindex[\VUindexRenfoncement]{krik-krak}{93}
\VUindex[\VUindexRenfoncement]{nu}{93}
\VUindex[\VUindexRenfoncement]{paf}{93}
\VUindex[\VUindexRenfoncement]{plump}{93}
\VUindex[\VUindexRenfoncement]{psit}{93}
\VUindex[\VUindexRenfoncement]{shut}{93}
\VUindex[\VUindexRenfoncement]{sus}{93}
\VUblanc{1.37mm}
\VUindexNu{\VUgras{Vorti interjeciona :}}{94}
\VUindex[\VUindexRenfoncement]{abase}{94}
\VUindex[\VUindexRenfoncement]{adavane}{94}
\VUindex[\VUindexRenfoncement]{addope}{94}
\VUindex[\VUindexRenfoncement]{adio}{94}
\VUindex[\VUindexRenfoncement]{bone}{94}
\VUindex[\VUindexRenfoncement]{brave}{94}
\VUindex[\VUindexRenfoncement]{certe}{94}
\VUindex[\VUindexRenfoncement]{fore}{94}
\VUindex[\VUindexRenfoncement]{haltez}{94}
\VUindex[\VUindexRenfoncement]{helpo}{94}
\VUindex[\VUindexRenfoncement]{silencez}{94}
\VUindex[\VUindexRenfoncement]{shamo}{94}
\VUindex[\VUindexRenfoncement]{sokurso}{94}
\VUindex[\VUindexRenfoncement]{tacez}{94}
\VUindex[\VUindexRenfoncement]{vere}{94}
\VUindex[\VUindexRenfoncement]{ya}{94}
\VUindex[\VUindexRenfoncement]{yen}{94}
\VUblanc{1.29mm}
\VUindexNu{\VUgras{Nombri :}}{94}
\VUindex[\VUindexRenfoncement]{kardinala}{94}
\VUindex[\VUindexRenfoncement]{kompozita}{95}
\VUindex[\VUindexRenfoncement]{ordinala}{97}
\end{VUpage}

% -------------------------------------------------------------------
%  FEUILLET 233 — LE TABELO PASSE A LA DEUXIEME PARTIE. Le filet centre
%  qui l'annonce mesure 178 px : une SIXIEME longueur, apres 90-100,
%  155, 202 et 296. Il est a y 1069-1071, x 533-710, centre a 621,5
%  pour un centre de justification de 618 a 628. Le detecteur l'a
%  manque comme les cinq autres, et il ne figurait dans AUCUN exces :
%  l'intervalle qui le porte vaut 186,5 px, et la liste des exces
%  n'annoncait que les deux blancs des titres, plus bas.
%  Les deux gros exces de 47 et 44 px sont les blancs qui encadrent
%  « VORTIFADO » ; rien n'y est cache, verifie a trois seuils.
%  Le numero de « indikativo » s'imprime « 1(6 » — le zero n'a pris
%  l'encre que sur sa moitie gauche. Compose 106, ce qu'impose la suite.
% -------------------------------------------------------------------
\begin{VUpage}[233]{229}
\VUpageIndex
\VUindex[\VUindexRenfoncement]{fracionala}{97}
\VUindex[\VUindexRenfoncement]{multiplikera}{98}
\VUindex[\VUindexRenfoncement]{nombro de foyi}{98}
\VUindex[\VUindexRenfoncement]{Expresuri distributiva}{98}
\VUblanc{0.74mm}
\VUindex{\VUgras{La hori}}{99}
\VUindex{En matematiko}{100}
\VUblanc{0.70mm}
\VUindex{\VUgras{Sintaxo}}{100}
\VUindex{Vortordino}{100}
\VUindex{Inversigo e \textit{n} inversigala (akuzativo)}{102}
\VUblanc{0.75mm}
\VUindexNu{\VUgras{Tempi e modi :}}{106}
\VUindex[\VUindexRenfoncement]{indikativo (prezenta, pasinta, futura)}{106}
\VUindex[\VUindexRenfoncement]{kondicionalo}{107}
\VUindex[\VUindexRenfoncement]{la konjuncioni ne influas la modo}{109}
\VUindex[\VUindexRenfoncement]{imperativo o volitivo}{107 e 110}
\VUindex[\VUindexRenfoncement]{expresuro \textit{por ke} e volitivo}{110}
\VUindex[\VUindexRenfoncement]{infinitivi}{112}
\VUindex[\VUindexRenfoncement]{participi}{112}
\VUindex[\VUindexRenfoncement]{absoluta participo}{112}
\VUindex[\VUindexRenfoncement]{participi substantiva}{113}

\VUsaut{4.30mm}
\VUfilet{15.07mm}
\VUsaut{4.90mm}
\VUcentre{8.5pt}{175}{DUESMA PARTO}
\VUsaut{3.60mm}
\VUcentre{10.27pt}{166}{\VUgras{VORTIFADO}}
\VUsaut{3.70mm}

\VUpageIndex
\VUindex{\VUgras{Elementi di vorto}}{117}
\VUblanc{0.69mm}
\VUindex{\VUgras{Procedi di vortifado}}{117}
\VUindex{Radiki}{117}
\VUindex{Dezinenci}{118}
\VUblanc{0.83mm}
\VUindex{\VUgras{Nemediata derivado}}{118}
\VUindex{Substantivo nemediate formacita de verbo}{118}
\VUindex{Adjektivo nemediate formacita de substantivo}{119}
\VUindex{Adverbo nemediate formacita de adjektivo}{120}
\VUindex{Adjektivo de adverbo}{120}
\VUindex{Verbo de radiko neverbala}{121}
\VUblanc{0.80mm}
\VUindex{\VUgras{Mediata derivado}}{121}
\VUindex{Afixi}{122}
\end{VUpage}

% -------------------------------------------------------------------
%  FEUILLET 234 — LA GEOMETRIE N'EN VOYAIT QUE VINGT-SIX LIGNES SUR
%  QUARANTE. Le bloc detecte s'arretait a y = 1300 : les treize
%  dernieres entrees, de « -ad- » a « -em- », tombaient hors de lui, et
%  les bandes scan/H234_*.png ne les montrent meme pas. Une quatorzieme
%  ligne, « ri- », manquait au milieu du bloc.
%  C'est la panne la plus grave du detecteur a ce jour : non plus une
%  ligne courte perdue, mais UN TIERS DE LA PAGE. L'exces de 52,4 px
%  qu'elle annoncait a l'indice 21 etait la somme d'une ligne entiere
%  sautee (40 px) et du blanc de rubrique (12 px).
%  Les prefixes et suffixes cites sont en ITALIQUE, non en gras.
%  Coquilles conservees : « anti » et « auto » sans tiret final quand
%  tous leurs voisins en portent un ; « sen- » entre « ne- » et « par- »,
%  hors de l'ordre alphabetique ; « -ebl.- » avec un point avant le
%  tiret fermant.
% -------------------------------------------------------------------
\begin{VUpage}[234]{230}
\VUpageIndex
\VUindexNu{\VUgras{Prefixi :}}{122}
\VUindex[\VUindexRenfoncement]{\textit{anti}}{122}
\VUindex[\VUindexRenfoncement]{\textit{arki-}}{123}
\VUindex[\VUindexRenfoncement]{\textit{auto}}{123}
\VUindex[\VUindexRenfoncement]{\textit{bi-}}{123}
\VUindex[\VUindexRenfoncement]{\textit{bo-}}{123}
\VUindex[\VUindexRenfoncement]{\textit{des-}}{124}
\VUindex[\VUindexRenfoncement]{\textit{dis-}}{126}
\VUindex[\VUindexRenfoncement]{\textit{ex-}}{126}
\VUindex[\VUindexRenfoncement]{\textit{ge-}}{126}
\VUindex[\VUindexRenfoncement]{\textit{mi-}}{127}
\VUindex[\VUindexRenfoncement]{\textit{mis-}}{127}
\VUindex[\VUindexRenfoncement]{\textit{ne-}}{127}
\VUindex[\VUindexRenfoncement]{\textit{sen-}}{128}
\VUindex[\VUindexRenfoncement]{\textit{par-}}{128}
\VUindex[\VUindexRenfoncement]{\textit{para-}}{128}
\VUindex[\VUindexRenfoncement]{\textit{pre-}}{128}
\VUindex[\VUindexRenfoncement]{\textit{pseudo-}}{129}
\VUindex[\VUindexRenfoncement]{\textit{quadri-}}{129}
\VUindex[\VUindexRenfoncement]{\textit{retro-}}{129}
\VUindex[\VUindexRenfoncement]{\textit{ri-}}{129}
\VUblanc{1.02mm}
\VUindex{\VUgras{Prefixi teknikala (\textit{equi-, ko-, mono-})}}{130}
\VUblanc{0.59mm}
\VUindex{Prepozicioni prefixa}{132}
\VUblanc{1.10mm}
\VUindexNu{\VUgras{Sufixi :}}{132}
\VUindex[\VUindexRenfoncement]{\textit{-ab-}}{132}
\VUindex[\VUindexRenfoncement]{\textit{-ach-}}{133}
\VUindex[\VUindexRenfoncement]{\textit{-ad-}}{134}
\VUindex[\VUindexRenfoncement]{\textit{-ag-}}{135}
\VUindex[\VUindexRenfoncement]{\textit{-aj-}}{135}
\VUindex[\VUindexRenfoncement]{\textit{-al-}}{137}
\VUindex[\VUindexRenfoncement]{\textit{-an-}}{139}
\VUindex[\VUindexRenfoncement]{\textit{-ar-}}{140}
\VUindex[\VUindexRenfoncement]{\textit{-ari-}}{141}
\VUindex[\VUindexRenfoncement]{\textit{-atr-}}{142}
\VUindex[\VUindexRenfoncement]{\textit{-e-}}{143}
\VUindex[\VUindexRenfoncement]{\textit{-ebl.-}}{143}
\VUindex[\VUindexRenfoncement]{\textit{-ed-}}{144}
\VUindex[\VUindexRenfoncement]{\textit{-eg-}}{144}
\VUindex[\VUindexRenfoncement]{\textit{-em-}}{145}
\end{VUpage}

% -------------------------------------------------------------------
%  FEUILLET 235 — QUATRE LIGNES MANQUAIENT A LA GEOMETRIE : « -i- »,
%  « -if- », « -il- » et « -iz- ». Les quatre exces d'un interligne
%  entier qu'elle annoncait leur correspondent un pour un.
%  QUATRIEME MENSONGE DU DETECTEUR DE FILET SUR LE TABELO : il rendait
%  122,43 mm, ordonnee qui tombe A L'INTERIEUR de la vedette grasse
%  « Kompozado : », dont la fermeture soude les empattements sur 181 px.
%  La fenetre de longueur resserree l'ecarte desormais, et le detecteur
%  rend enfin « filet non trouve », qui est la bonne reponse.
%  UN FILET CENTRE DE 143 px clot la page (x 531-673, y 1835-1837) :
%  septieme longueur du volume, apres 90-100, 155, 178, 202 et 296.
% -------------------------------------------------------------------
\begin{VUpage}[235]{231}
\VUpageIndex
\VUindex[\VUindexRenfoncement]{\textit{-end-}}{147}
\VUindex[\VUindexRenfoncement]{\textit{-er-}}{147}
\VUindex[\VUindexRenfoncement]{\textit{-eri-}}{149}
\VUindex[\VUindexRenfoncement]{\textit{-es-}}{149}
\VUindex[\VUindexRenfoncement]{\textit{-esk-}}{150}
\VUindex[\VUindexRenfoncement]{\textit{-estr-}}{151}
\VUindex[\VUindexRenfoncement]{\textit{-et-}}{152}
\VUindex[\VUindexRenfoncement]{\textit{-ey-}}{152}
\VUindex[\VUindexRenfoncement]{\textit{-i-}}{153}
\VUindex[\VUindexRenfoncement]{\textit{-id-}}{153}
\VUindex[\VUindexRenfoncement]{\textit{-ier-}}{153}
\VUindex[\VUindexRenfoncement]{\textit{-if-}}{154}
\VUindex[\VUindexRenfoncement]{\textit{-ig-}}{154}
\VUindex[\VUindexRenfoncement]{\textit{-ik-}}{157}
\VUindex[\VUindexRenfoncement]{\textit{-il-}}{157}
\VUindex[\VUindexRenfoncement]{\textit{-in-}}{158}
\VUindex[\VUindexRenfoncement]{\textit{-ind-}}{159}
\VUindex[\VUindexRenfoncement]{\textit{-ism-}}{160}
\VUindex[\VUindexRenfoncement]{\textit{-ist-}}{160}
\VUindex[\VUindexRenfoncement]{\textit{-iv-}}{161}
\VUindex[\VUindexRenfoncement]{\textit{-iz-}}{162}
\VUindex[\VUindexRenfoncement]{\textit{-oz-}}{163}
\VUindex[\VUindexRenfoncement]{\textit{-ul-}}{163}
\VUindex[\VUindexRenfoncement]{\textit{-um-}}{164}
\VUindex[\VUindexRenfoncement]{\textit{-un-}}{165}
\VUindex[\VUindexRenfoncement]{\textit{-ur-}}{165}
\VUindex[\VUindexRenfoncement]{\textit{-uy-}}{167}
\VUindex[\VUindexRenfoncement]{\textit{-yun-}}{167}
\VUblanc{0.55mm}
\VUindexNu{\VUgras{Kompozado :}}{168}
\VUindex{1\textsuperscript{e} Substantivo kun substantivo}{168}
\VUindex{2\textsuperscript{e} Substantivo kun verbo}{168}
\VUindex{3\textsuperscript{e} Substantivo kun adjektivo (o radiko adjektivigita)}{168}
\VUindex{4\textsuperscript{e} Prepoziciono kun verbo}{168}
\VUindex{5\textsuperscript{e} Prepoziciono o nombro-nomo kun adjektivo o substantivo}{168}
\VUindex{6\textsuperscript{e} Irga adjektivo kun radiko igita adjektivo o adverbo}{168}
\VUblanc{0.72mm}
\VUindex{\VUgras{Regulo di analizo o deskompozo}}{169}
\VUblanc{0.72mm}
\VUindex{\VUgras{Kompozado per prepozicioni}}{172}
\VUblanc{0.72mm}
\VUindex{\VUgras{La kompozaji e la sufixi}}{173}
\VUsaut{2.20mm}
\VUfilet{12.10mm}
\end{VUpage}

% -------------------------------------------------------------------
%  FEUILLET 236 — LE FLEURON. Seul ornement FIGURE du volume : une
%  vignette vegetale, tige ondulee portant feuilles, baies et un trefle
%  terminal (encre x 734-937, y 870-892). Il ne se compose avec aucune
%  fonte du projet ; il est DECOUPE DU SCAN et pose en image. C'est le
%  seul endroit ou la transcription cesse d'etre une composition et
%  devient un fac-simile — il faut le dire.
%  Le detecteur avait manque TROIS objets : le folio (qu'il prend
%  pourtant partout ailleurs), le fleuron, et le filet de PLEINE
%  JUSTIFICATION (1084 px) — trop long pour son test, qui cherche un
%  trait de 195 a 260 px cale sur la marge.
%  L'intervalle de 349 px qui porte le fleuron ne figurait dans AUCUN
%  exces : le seul blanc de la page que la liste ne signalait pas est
%  celui qui portait quelque chose.
%  DEUX PAS DE POINTS DE CONDUITE sur la meme page : 17,0 px pour la
%  liste des apendici, 43,3 px pour les deux catalogues.
% -------------------------------------------------------------------
\begin{VUpage}[236]{232}
\VUtitre{10.63pt}{60}{APENDICI}
\VUsaut{3.00mm}

\VUpageIndex
\VUindex{1. L'acentizo en Ido}{179}
\VUindex{2. La pluralo per \textit{-i}}{183}
\VUindex{3. Genro e maskulismo}{188}
\VUindex{4. Substantivigo dil adjektivo}{195}
\VUindex{5. Ca, ta e qua}{205}
\VUindex{6. La Konjugo-sistemo di Ido}{206}
\VUindex{7. Vortordino}{215}
\VUindex{8. Puntizado}{217}
\VUindex{9. Nomi. Adresi}{222}
\VUindex{10. Formuli di politeso en letri}{223}

\VUsaut{9.00mm}
\VUfleuron
\VUsaut{11.60mm}
\VUtitre{9.60pt}{40}{MONDO-BIBLIOTEKO}
\VUsaut{2.30mm}

\VUpageIndex
\VUindexLarge{1. D'Mumm se's, Komedio da \textsc{Dicks}. Trad. da H. \textsc{Meier}.}{Kr. \ 0.40}
\VUindexLarge{2. \textit{Nova Varianti di Antiqua Temo}. Originalo da A. \textsc{Kofman}.}{0.25}
\VUindexNu{3. \textit{La Karaktero Usana}, da Prof. \textsc{Brander Matthews}. Trad.}{}
\VUindexLarge[8.85mm]{da Dr. J. \textsc{Marelius}}{0.25}
\VUindexLarge{4. \textit{La Nimfo}, da \textsc{Elin-Pelin}. Trad. da \textsc{Th. Kaneff}}{0.25}
\VUindexNu{5. \textit{Kelka Maestroverki dil Moderna Liriko Araba}, da 26 autori.}{}
\VUindexLarge[8.85mm]{Trad. da \textsc{Rafael Nakhla} S. J.}{0.75}
\VUindexNu{6. \textit{La Mento Internaciona}, da Prof. R. D. \textsc{Carmichael} \& \textsc{Cecil}}{}
\VUindexLarge[8.85mm]{\textsc{Delisle Burns}. Trad. da D\textsuperscript{ro} J. \textsc{Marelius}}{0.25}

\VUsaut{2.90mm}
\VUfilet{\VUtexteLargeur}
\VUsaut{2.60mm}
\VUtitre{9.60pt}{40}{RADIO-LEXIKI}
\VUsaut{2.20mm}

\VUpageIndex
\VUindexNu{1. \textit{Internaciona Radio-Lexiko} en Ido e Germana, Angla,}{}
\VUindexNu[8.85mm]{Franca, Italiana e Hispana, da K. \textsc{Feder} \& J.}{}
\VUindexLarge[8.85mm]{\textsc{Nordin}}{Kr. \ 3.25}
\VUindexNu{2. \textit{Svenskt-Internationellt Radio-Lexikon}, da \textsc{Alfr. Liljes}-}{}
\VUindexLarge[8.85mm]{\textsc{tröm} \& J. \textsc{Nordin}}{0.25}

\VUsaut{2.10mm}
\VUcentre{10.2pt}{60}{EDITERIO P. AHLBERG}
\VUcentre{10.2pt}{0}{Surbrunnsgatan 37, Stockholm, Suedia. --- Postgirokonto 32.}
\end{VUpage}

% -------------------------------------------------------------------
%  FEUILLET 237 — MA CONJECTURE DE STRUCTURE ETAIT FAUSSE. J'avais
%  annonce au releveur une annonce d'editeur, d'apres le releve de
%  structure du § 3 ; la page porte en realite la fin des « Opinioni
%  pri Ido », puis le titre « Difuzeso di Ido. », puis le bloc U. L. I.
%  Le releveur a mesure au lieu de me croire, et l'a dit.
%  L'INTERLIGNE EST PLUS SERRE QUE PARTOUT AILLEURS : 33 px = 7,95 pt
%  entre lignes d'un meme alinea, contre 41,45 px au corps ordinaire.
%  Le CORPS, lui, ne change pas : hauteur d'x de 16-17 px, celle du
%  texte courant. La page est composee serree, non en petit oeil.
%  Ni filet, ni ornement, ni folio imprime.
%  Coquilles conservees : « esor » pour esar ; « Kobenhavn » avec trema
%  la ou le danois ecrit Kobenhavn ; « Schrijen » pour Schrijnen ;
%  « College » avec accent grave sur un nom anglais.
% -------------------------------------------------------------------
\begin{VUpage}[237]{}
% LE BLOC DE CETTE PAGE MONTE PLUS HAUT QUE \VUmargeSup. Ses 46 lignes
% au pas de 33 px, plus les blancs des deux titres, demandent 1614 px
% quand \VUtexteHauteur n'en offre que 1564 : la page deborde de 9 mm
% par le bas si on la cale comme les autres. Les deux dernieres pages
% du volume sont composees hors du cadre ordinaire — elles n'ont pas
% de folio non plus, ce qui va de pair.
% Faute de folio pour caler la page, la valeur est prise sur le
% debordement mesure et non sur une lecture : c'est un ajustement, et
% il faut le dire.
\vspace*{-9.0mm}
\VUpageSerre{8.26pt}{7.95pt}

\VUcontinue
plu granda. » (La sam autoro en \textit{Les Langues dans l'Europe}\nl
\textit{Nouvelle}, p. 323, editita en 1918.)

« Kom filologial produkturo artificala, Ido superesas Esperanto.\nl
Ido esas \textit{plu simpla, plu transparanta, sociale plu utiligebla kam}\nl
\textit{Esperanto}. » (Profesoro \textsc{Schrijen}, Utrecht, en la revuo \VUgras{Studien}.)\parplein

« Ido ne destronizos la lingui vivanta, ma ol dispensos la cien\cc
cisti esar samtempe poligloti. La simpleso dil morfologio e dil\nl
sintaxo, same kam la logiko rigoroza qua direktis lua kompozado,\nl
renversas l'objecioni quin on povus expresar. \VUgras{On devas dezirar tot}\cc
\VUgras{kordie la suceso di ta interesanta entraprezo.} » (La \textit{Revue Augusti}\cc
\textit{nienne}, 15 di agosto 1910.)

« La rezultajo esas linguo quan singlu povas lernar tre facile; ol\nl
havas super la cetera lingui artificala ica avantajo esar fondita sur\nl
ciencal e teknikal principi racionoza, e konseque ol devas ne timar\nl
esor uldie remplasata da linguo plu bona ed esence diferanta, qua\nl
fine konquestus la mondo. » (D\textsuperscript{ro} Profesoro \textsc{Otto Jespersen} en\nl
\textit{La Langue internationale et la Science}, pag. 30.)

« ... Ne penvalorus uzar tala linguo internaciona, se ol esus tam\nl
neperfekta kam esis exemple Volapük, o mem Esperanto. On devas\nl
formacar ol segun principi vere ciencala, por ke ol esez tote prakti\cc
kala. Tarelate la laborado facita inter 1908 e 1914 dal Idisti, en\nl
lia revuo monatala \textit{Progreso}, indikas tre importanta pazo adavane,\nl
e la linguo \VUgras{Ido}, per ta esforci di multalanda kunlaboreri atingis\nl
grado di perfekteso tante alta, ke me ne hezitus konsilar lua adopto\nl
kom linguo oficala di la Nacionar-ligo. » (\textsc{Otto Jespersen} en serio\nl
de diskursi che \textit{University Collège}, Londono, junio 1920.)

La granda linguisto skribas e parolas Ido; mem il facis kursi pri\nl
ca linguo a studenti dil Universitato en \textit{Köbenhavn}, ube il esis\nl
profesoro pri linguistiko. Altraparte il enuncis pri Ido ica atesto\nl
tre preciza :

« Por omna nacioni ed aparte por mea Dana sampatriani e por\nl
la cetera nacioni Skandinaviana, la linguo Ido certe esas plu facile\nl
lernebla e praktikebla kam Esperanto. Pro ta motivo ed anke pro\nl
sua ciencala konstrukteso e sua granda flexebleso, Ido semblas a\nl
me plu kapabla kam irg altra linguo plear la rolo di helpolinguo\nl
por omna nacioni civilizita. »

\VUsaut{2.00mm}
\VUcentre{11.0pt}{0}{\VUgras{Difuzeso di Ido.}}
\VUsaut{2.30mm}

Atinginte nur en la mezo di 1913 sua plena developeso, Ido vidis\nl
quik la milito desorganizar lua societi e grupi. Tamen nun ol esas\nl
difuzita en 55 landi di la kin parti dil mondo. Ol havis kongreso\nl
en \textit{Wien} (1921), \textit{Dessau} (1922), \textit{Cassel} (1923), \textit{Luxembourg} (1924),\nl
e \textit{Torino} (1925).

\VUsaut{1.70mm}
\VUcentre{11.0pt}{190}{\VUgras{U. L. I.}}
\VUsaut{1.20mm}
\VUcentre{11.0pt}{0}{\VUgras{Uniono por la linguo internaciona.}}
\VUsaut{0.90mm}
\VUcentre{10.2pt}{0}{Statuti recevebla che la sekretario Red. P. \textsc{Ahlberg}}
\VUcentre{10.2pt}{0}{Surbrunnsgatan, 37, Stockholm, Suedia.}
\end{VUpage}

% -------------------------------------------------------------------
%  FEUILLET 238 — CATALOGUE DE PRESSE ET DE LIBRAIRIE IDISTE. Ni points
%  de conduite, ni filet, ni ornement, ni folio imprime.
%  L'INTERLIGNE EST SERRE : 32,8 px = 7,90 pt a l'interieur d'une
%  notice, contre 41,45 px au corps ordinaire ; entre deux notices,
%  15 a 17 px de blanc en plus. Le « pas 41,28 px » de la geometrie
%  n'est l'ecart que d'un seul intervalle, celui des deux lignes
%  centrees du haut : le vrai pas est le contraire.
%  Coquilles conservees : « Ofical » a un seul f ; « Chekoslovakian »
%  sans -a ; et « L. de Beaufront » en romain a une notice, « L. DE
%  BEAUFRONT » en petites capitales a une autre.
% -------------------------------------------------------------------
% LA COUVERTURE N'EST PAS COMPOSEE PLUS LARGE QUE LE VOLUME, contre ce
% que j'avais conclu. Les 1172 a 1210 px que j'avais mesures venaient du
% GRAIN DU PAPIER VERT, qu'un simple seuillage compte pour de l'encre :
% le bord des lignes s'en trouvait gonfle de cent pixels. Repris avec
% PG.lines_of, qui travaille sur un masque de glyphes, ce feuillet donne
% une largeur maximale de 1108 px et une mediane de 1047, pour une
% justification du volume de 1083. Les deux lignes qui suivent le titre
% n'ont donc aucune raison de deborder : leur interlettrage est cale sur
% les 1079 et 1046 px mesures.
\begin{VUpage}[238]{}
\VUpageSerre{8.26pt}{7.90pt}
\VUcentre{11.0pt}{80}{\VUgras{MONDO}}
\VUsaut{2.10mm}
\VUcentre{10.8pt}{0}{\VUgras{Revuo por la Mondo-Linguo.}}
\VUsaut{2.10mm}
\VUcentre{10.2pt}{-60}{Ofical organo di Uniono por la Linguo Internaciona e di lua Akademio.}
\VUcentre{10.2pt}{-52}{Redaktero : P. \textsc{Ahlberg}, Surbrunnsgatan 37, Stockholm, Suedia.}
\VUsaut{1.35mm}

\textit{La Belga Idisto}, monatala revuo, preco : 6.20 fr. che J. J. V. Eug.\nl
Mathys, rue Léopold Vanderkelen, Louvain (Belgia).

\textit{Ido}, qua aparas 18 foye en yaro. Abonpreco : 5 fr. Suisa valoro.\nl
Redaktero : Albert Noetzli, Idokontoro, Oerlikon-Zürich (Suisia).

\textit{Bulletin Français-Ido de la Langue Auxiliaire}, organo di lingual\nl
expozo e defenso, buletino di la \textit{Société Idiste Française}. Direktero\nl
redaktanta, L. de Beaufront; abonpreco : 5 fr. kun suskripto a la\nl
societo, sive che la sekretario di ica : S\textsuperscript{ro} L. de Guesnet, 83, rue de\nl
Rochechouart, Paris (IX\textsuperscript{e}), sive che Ed. Bréon, Mussy-sur-Seine\nl
(Aube), Francia.

\textit{Die Weltsprache}, Zeitschrift für die internationale Hilfssprache\nl
\textit{Ido}, jährlich 2.--- Mk. Organo dil German Ido-Federuro, adreso :\nl
Bockenheimer Anlage, 45, Frankfurt am Main, Germania.

\textit{Linguo Internaciona}, monatal organo di Chekoslovakian Ido\cc
federuro. Abonpreco : 25 K. Ch. Administrero : V. Markvart, Praha\cc
Zizkov, Premyslovska 1481 (Chekoslovakia).

\textit{L'Idisto Katolika}, abonpreco : 6 fr. che Idokontoro, Thaon-les\cc
Vosges, Francia.

\textit{Revuo Internaciona}, organo dil laboristal \textit{Ido-Uniono}. Abonpreco :\nl
3 Germ. marki, che M. Grümmer, Sophienstr., 7, Leipzig, Germania.

\textit{Lexiko Ido-Angla}, da L. H. \textsc{Dyer}, preco : 10 shilingi, che J. W.\nl
Baxter, 57 Limes Grove, Lewisham, London S.E., Anglia.

\textit{Lexique-Manuel Ido-Français}, par J. \textsc{Guignon}, Vulaines-sur-Seine\nl
(S.-et-M.), Francia. Preco : fr. 7.50.

\textit{Vocabolario Moderno Ido-Italiano ed Italiano-Ido} da prof. \textsc{Paolo}\nl
\textsc{Lusana}. Preco : 10 Lire che profesoro G. Meazzini, San Giovanni\nl
Valdarno, Firenze (Italia).

\textit{Vollständiges Wurzelwörterbuch Ido-Deutsch}, da K. \textsc{Feder} e\nl
\textsc{Fr. Schneeberger}, bindita fr. Suisa 4.50 plus afranko; broshita\nl
fr. Suisa 5.50.

\textit{Taschenwörterbuch Ido und Deutsch}, brosh. fr. S. 2.50; bind.\nl
fr. S. 3.00. Che « Ido-Weltsprache Verlag », Weissenbühlweg 13,\nl
Bern (Suisia).

\textit{Dictionnaire Français-Ido}, 600 pag., da L. \textsc{de Beaufront} et\nl
L. \textsc{Couturat}. --- 12.00 fr. Franca.

\textit{Opiniono di linguisto} (prof. \textsc{Meillet}). --- 1.35 fr. che H. Vinez,\nl
35, rue Charlot, Paris 3. (Francia).

\textit{Relief-mapo} « \textit{La Nova Europa} », editita dal Acionsocieto Choko\cc
lado Tobler en Bern (Suisia).
\end{VUpage}
